% =====================================================================
%  Harald Høffding, PHILOSOPHIEN I TYDSKLAND EFTER HEGEL.  Kjøbenhavn: Forlagt af den
%  Gyldendalske Boghandel (F. Hegel), 1872 (I. Cohens Bogtrykkeri).  [IV], 325 pp.
%
%  TRANSCRIPTION -- GENERATED; DO NOT EDIT BY HAND.  Made in steps, all rerunnable
%  (ebook-method/tydskland-efter-hegel/, see its scripts' docstrings):
%    python3 ebook-method/tydskland-efter-hegel/draft.py --emit --force
%      this template + the Internet Archive's ABBYY OCR of Google's scan of the New York
%      Public Library copy (~/bibliotek/Høffding, Harald/1872-philosophien-i-tydskland-
%      efter-hegel.pdf, bound by abbyy_scan.py), then patches.py and spaced.tsv;
%    python3 ebook-method/collate/check.py tydskland-efter-hegel --apply
%      the decisions in decisions.tsv, taken against an INDEPENDENT reading of the same
%      images: tesseract (dan, tessdata_best).  SCANS.tsv has the scan's sha256.
%  No e-book exists.
%
%  STATUS (2026-10-10): TRANSCRIPTION COMPLETE, pp. 1-325 and the front matter collated
%  (OPEN 0; see RESUME-NOTES.md for how).  Translated 2026-10-09/10 (translation.tex).  pp. 236-237 (a leaf missing
%  from the NYPL copy) and pp. 22, 105, 109, 113, 139, 142, 143, 223, 229, 279, 281, 309 (degraded
%  or clipped in the NYPL scan) are taken from the Minnesota copy (HathiTrust).
%  WORD-ACCURACY.  The draft (ABBYY) was collated word by word, with punctuation, against
%  tesseract; ~4300 disagreements, over 1500 settled by rule (mostly by the Minnesota copy's layer
%  as a third witness), the rest read at the image in blind crops with hidden control crops
%  (45 of 45 caught) or set from whole-page readings.  Afterwards: footnotes read whole at the
%  image; letterspacing read at the image; unattested words swept; the whole text compared
%  word by word with the Minnesota layer (umnsweep.py: ~40 faults both NYPL OCRs shared);
%  all Greek read at 600 dpi.  Blind whole-page reads of 9 random pages found 3 faults
%  (tra/fra, O/0, forskjel lige), each mended.  TRANSLATION PASS (2026-10-09/10, TRANSLATION-PLAYBOOK
%  §6c): before the translators, 10 false paragraph breaks inside a sentence mended and the run-over of
%  p. 149's note removed from the foot of p. 150; translator E found a second copy of the p. 139 reading
%  inside p. 143 (34 lines, from the page-set reading), removed at its source (patches.PAGESET[143]); the
%  translators' 200 odd spots read blind at the image (two readers, controls 5/5; 14 re-cropped): 102
%  transcription faults mended (unclosed quotation marks, specks read as points, words split or glued at a
%  line end, o: for ɔ:, letterspacing missed: Bruno Bauer, Den yngre Fichte, er, ét) and 60 misprints marked
%  PRINTED AS IS (patches.POSTMIS), the doubled words at line turns (og|og, i|i, at|at) among them.
%  Not established: a misprint both copies' OCRs
%  read as a real word, where no reader looked.
%  GREEK.  The print sets all its Greek in the period's cursive Greek face; that is the
%  face, not an emphasis, so the Greek is set upright here.
%
%  CONVENTIONS.  "% --- p. N ---" + \opage{N} at the first word of printed page N
%  (a word broken over the turn is split with %).  Footnotes as printed, marked
%  *), **), ***) and numbered per page (\footnote[k]).  Italics \textit; letterspacing
%  \emph.  A note running over a page is set whole at its mark.  Folios and signature
%  marks are not transcribed.  The print governs; misprints are kept and marked
%  PRINTED AS IS.
% =====================================================================
\documentclass[12pt,a4paper,openany]{book}
\usepackage[utf8]{inputenc}
\usepackage[T1]{fontenc}
\usepackage{textalpha}   % Greek (p. 54: πρῶτον ψεῦδος)
\usepackage{libertinus}
\usepackage{libertinust1math}
\usepackage[danish]{babel}
\usepackage[protrusion=true,expansion=true]{microtype}
\usepackage{setspace}
\usepackage[margin=1.2in]{geometry}
\usepackage{fancyhdr}
\setlength{\headheight}{28pt}
\usepackage{hyperref}
\hypersetup{pdftitle={Philosophien i Tydskland efter Hegel (1872)}, pdfauthor={Harald Høffding}, hidelinks}
\pagestyle{fancy}
\fancyhf{}
\fancyhead[LE,RO]{\thepage}
\fancyhead[RE]{\textit{Harald Høffding: Philosophien i Tydskland efter Hegel}}
\fancyhead[LO]{\textit{1872}}
\fancypagestyle{plain}{\fancyhf{}\fancyhead[LE,RO]{\thepage}}
\setstretch{1.3}
\usepackage{marginnote}
\usepackage{graphicx}
\DeclareUnicodeCharacter{0254}{\reflectbox{c}}   % ɔ (id est)
\renewcommand{\thefootnote}{%
  \ifcase\value{footnote}\or *)\or **)\or ***)\else \arabic{footnote})\fi}
\newcommand{\tocline}[2]{\noindent #1\dotfill\ #2\par}
\newcommand{\opage}[1]{\marginnote{\footnotesize\textit{[#1]}}}
% a chapter of the print: new page, centred title, a bookmark and a contents line
\newcommand{\chap}[1]{\clearpage\addcontentsline{toc}{chapter}{#1}\vspace*{2em}\begin{center}{\large #1}\end{center}\medskip}
% a chapter that begins in mid-page in the print (Chapter 2 on p. 17): no page break
\newcommand{\chaphere}[1]{\addcontentsline{toc}{chapter}{#1}\begin{center}{\large\bfseries #1}\end{center}}
\pdfstringdefDisableCommands{\renewcommand{\footnote}[2][]{}\renewcommand{\opage}[1]{}}


\begin{document}

% --- front matter: read at the image 2026-10-09 (r-l-front.txt), compared with the ABBYY layer; library stamps
% and pen marks on the NYPL copy not transcribed.  Title PDF 9, motto PDF 10, Indhold PDF 11. ---
\begin{titlepage}
\begin{center}
\vspace*{3em}
{\LARGE\bfseries Philosophien i Tydskland}\\[0.8em]
{\large \emph{efter Hegel.}}\\[2em]
Af\\[1em]
{\large\bfseries Dr. H. Høffding.}\\[2em]
\vfill
{\small\bfseries Kjøbenhavn.}\\[0.4em]
\emph{Forlagt af den Gyldendalske Boghandel (F. Hegel).}\\[0.4em]
{\small \emph{I. Cohens Bogtrykkeri.}}\\[0.4em]
{\small 1872.}
\end{center}
\end{titlepage}
\thispagestyle{empty}
\vspace*{4em}
\begin{center}\small
„Auch wir können noch an \emph{wirkende}, aber wir können\\
nicht mehr an \emph{hexende} Ideen glauben.“
\end{center}
\begin{flushright}\footnotesize
H. Lotze: Geschichte der Aesthetik in Deutschland S. 415.
\end{flushright}
\clearpage
\thispagestyle{empty}
\begin{center}{\large\bfseries \emph{Indhold.}}\end{center}
\begin{center}\rule{3em}{0.4pt}\end{center}
\hfill{\small Side}\par
\tocline{\textbf{Første Kapitel:} \emph{Indledning}}{1}
{\small 1) Historisk Orientering p. 1. 2) Karakteristik af Hegels System p. 7. 3) Hovedproblemer i den efterhegelske Philosophi p. 14.}\par\smallskip
\tocline{\textbf{Andet Kapitel:} \emph{Den kritiske Bevidsthed}}{17}
{\small 1) Striden om den individuelle Udødelighed p. 17. 2) Strauss p. 29. 3) Bruno Bauer p. 40. 4) Ruge p. 44. 5) Feuerbach p. 49. 6) Max Stirner p. 88. 7) Das Verstandesthum und das Individuum p. 92.}\par\smallskip
\tocline{\textbf{Tredie Kapitel:} \emph{Den moderne Materialisme}}{94}
{\small 1) Den empiriske Materialisme (Vogt, Moleschott, Büchner) p. 99. 2) Den ethisk - spekulative Materialisme (Czolbe) p. 109.}\par\smallskip
\tocline{\textbf{Fjerde Kapitel:} \emph{Den spekulative Theisme}}{117}
{\small 1) Schellings senere Lære p. 121. 2) Den yngre Fichte p. 149. 3) Weisse p. 181.}\par\smallskip
\tocline{\textbf{Femte Kapitel:} \emph{Virkelighedsidealismen}}{220}
{\small 1) Trendelenburg p. 222. 2) Fechner p. 246. 3) Hartmann p. 268. 4) Lotze p. 286.}\par\smallskip
\tocline{\textbf{Sjette Kapitel:} \emph{Karakteristik og Resultater}}{319}
\begin{center}\rule{3em}{0.4pt}\end{center}
\clearpage
% CHECK: the Indhold's section numbers may be small raised figures with a parenthesis (reader); "ethisk - spekulative"
% is spaced around the hyphen in the print.  To confirm at the image before the book is called done.


% --- p. 1 ---
\setcounter{footnote}{0}%
\chap{FØRSTE KAPITEL}
\begin{center}\rule{2em}{0.4pt}\end{center}
\opage{1}


1. Det kunde synes at være et ligesaa utaknemligt
som vanskeligt Arbeide at give en Fremstilling af den
efterhegelske Philosophi i Tydskland. Thi der knytter
sig til denne Periode ikke den Interesse, som fremkaldes
ved store epochegjørende Fremskridt og Gjennembrud;
det er væsentlig en Epigonernes Tid efter det forudgaaende
halve Aarhundredes storartede philosophiske Produktion.
Og paa den anden Side ligger den efterhegelske
Philosopheren os saa nær, at en virkelig historisk og
objektiv Fremstilling maa være vanskelig, om ikke umulig.

Disse Betænkeligheder forsvinde dog ved den Betragtning,
at det er nødvendigt for Enhver, der stræber
efter en philosophisk Verdensanskuelse, at komme til
Klarhed over, hvilke afgjørende Problemer der røre sig
i Tidsalderens Tænkning, og i hvilken Retning deres Behandling
og Besvarelse gaaer. Saa dyb og inderlig en
Berigelse, der vindes ved bestandig at vende tilbage til
de store Mesteres Skrifter, saa vigtigt og uundværligt
er det derfor paa den anden Side at orientere sig med
Hensyn til Samtidens Tænkere, selv om de ikke skulde
kunne regnes til Mestrenes Klasse. Og naar det er
denne Interesse, der ansporer Forskningen, da falder
ogsaa det rent objektive og historiske Synspunkt bort
% --- p. 2 ---
\setcounter{footnote}{0}%
\opage{2}og giver Plads for en mere egentlig philosophisk og kritisk
Betragtningsmaade.

Udgangspunktet betegnes ved den hegelske Philosophi,
der ved sin universelle Karakter greb saa dybt ind
i hele Tidens aandelige Liv, og tillige selv var Resultatet
af en idealistisk Retning, der gaaer igjennem hele
den nyere Philosophi. Dog er Hegelianismen kun vort
relative Udgangspunkt. Den danner ikke, som man ofte
har ment, et Knudepunkt i den nyere Philosophis Historie:
det ligger allerede i, at den er den mest afgjorte
Idealisme. Knudepunktet maa derimod søges hos den
Philosoph, hos hvem den idealistiske og den realistiske
Udviklingsrække skjære hinanden og for en Tid gaa i Et
paa Basis af et dybere liggende Princip. Denne Philosoph,
den nyere Tids centrale Tænker, er Immanuel
Kant. Dette lader sig i Korthed vise baade med Hensyn
til hans Udviklingsgang, hans Princip og hans Methode.

Efterat Kant længe havde syslet dels med naturvidenskabelige
Studier dels med Studiet af den Leibnitz-Wolfske,
idealistiske Philosophi, var det efter hans egen
Erklæring David Hume, den gjennemførte Empirismes Repræsentant,
der vakte ham af hans „dogmatiske Slummer“.
Og at han nu saa sin Opgave i at finde Enhedspunktet for de
to modsatte Retninger, ses blandt Andet deraf, at han i
sit berømte Værk, „Kritik der reinen Vernunft“, efter
ethvert Hovedafsnit søger at vise, hvorledes hans Lære
baade lader Leibnitz og Locke komme til deres Ret. ---
Hvilket er da det philosophiske Princip, hvorved Kant
naaer til dette overlegne Standpunkt? Dette Princip
maa søges i selve den Maade, hvorpaa han opfatter den
philosophiske Erkjendelses Væsen. Han tager ikke som
Realismen sit Standpunkt i det Ydre, i Gjenstandene,
som skulde paavirke Mennesket udenfra og præge Erkjendelsen
ind i det, og ligesaalidt lader han idealistisk
Menneskeaanden udvikle Erkjendelsen af sit eget Indre.
Tanken er, for at bruge Bacos Billede, hverken Myren,
% --- p. 3 ---
\setcounter{footnote}{0}%
\opage{3}der blot sammenbringer og ophober Stoffet, eller Edderkoppen,
der udspinder sin Væv af sig selv, men Bien, der
uddrager Saften af Blomsterne og paa eiendommelig Maade
fordøier og omdanner den. Det er Tankens indre Vexelforhold
til Tingene i Erkjendelsen, Kant retter sit Øie imod;
han søger at bestemme de to Faktorers Væsen og Lovene
for deres Vexelvirkning. Erkjendelsen af disse Love
kalder Kant den transscendentale Erkjendelse. Denne
ligger til Grund for enhver anden Erkjendelse, hvad
enten denne er af metaphysisk eller empirisk Beskaffenhed.
--- Af dette Princip følger en tilsvarende Methode.
Denne er paa engang induktiv og deduktiv. Den gaaer
ud fra Erkjendelsens Kjendsgjerning og de i den indeholdte
Elementer, og forsaavidt bliver Undersøgelsen induktiv
og analytisk; men da Erkjendelsen ikke er en
ydre Kjendsgjerning, men sig selv iboende, sit eget Princip,
hvoraf alle dens Former og alt dens Indhold afhænger,
bliver Udviklingen synthetisk og deduktiv. Denne
Enhed lægger sig for Dagen i Forbindelsen af det Kritiske
og Dialektiske i Kants Undersøgelser. Kant distingverer
kritisk, men han gjør det for at lade Modsætningerne
træde frem og komme til deres Ret, for saa
herigjennem atter at naa den til Grund liggende Enhed.
Det er dette, man har kaldt det Drillende hos Kant, at
man ikke faaer Lov til at slaa sig til Ro ved de skarpe
Distinktioner, men stedse paa Ny vises ud over dem.
Dialektiken bryder frem netop ved Kritikens Gjennemførelse.


Dog efterlod Kant Opgaver nok for sine Efterfølgere.
Han havde i Principet angivet det Maal, der skulde
naaes, men i Virkeligheden havde han ikke betvunget
de stridende Magter. Navnlig havde han ved Sondringen
mellem Erkjendelsens subjektive Former og Tingen i
og for sig aabnet Veien for en Dualisme, der truede
selve det store Maal, han saa energisk havde vist hen
til, den indre Enhed af Idealisme og Realisme. Man
% --- p. 4 ---
\setcounter{footnote}{0}%
\opage{4}saa desuden snart, at denne Sondring beroede paa en
Selvmodsigelse. Thi under Forudsætning af vore Erkjendelsesformers
blot subjektive Betydning er Antagelsen
af en Ting i og for sig umulig: ingen Tankebestemmelser
--- og altsaa heller ikke Begreberne Væren og
Aarsag --- kunne jo saa konsekvent anvendes paa den;
den kan hverken \emph{være} eller være \emph{Aarsag} til Fornemmelsen
af en ydre Gjenstand. Uden Begrebet om „Ding
an sich“ kunde man, som Jakobi rigtig bemærkede, ikke
komme ind i Kants System, og med det kunde man
ikke blive der. Det var saaledes en i Sagens Natur
liggende Nødvendighed, at man først udskilte det realistiske
Element i Kants Lære. For at overvinde Dualismen
maatte man holde sig til det i sig Klare og Sikre,
Tanken selv og den sidste Kilde for dens Former, Selvbevidsthedens
oprindelige Enhed, hvad Kant kaldte
„Selvbevidsthedens transscendentale Apperception“. Overfor
denne rige og lyse Formverden stod „das Ding an
sich“ som en ubestemt, taageagtig Masse, hvis Sammenhæng
med hin ikke kunde paavises. Fichte konstruerede
Verden ud af Subjektiviteten og satte det som dennes
Opgave bestandig paa Ny at overvinde den saaledes oprundne
Tilværelse. Det Objektive staaer altsaa hos ham
endnu tilbage som et negativt Princip, der i evig Uro
sættes for atter at ophæves. Denne subjektive Idealismes
Skranker overskrider Schelling. Al Videnskab beroer,
siger han, paa den Forudsætning, at det absolut
Ideale ogsaa i sig selv er det absolut Reale. Det Absolute
er den rene Identitet, hvor ingen Forskjelle og
Modsætninger længere gjælde. Dette er den abolute % PRINTED AS IS (for absolute; read at the image, sweep s01)
Idealisme: det Ideale er Alt, fordi Alt er Ide. Den
Schellingske Naturphilosophi vil bevise Naturens Identitet
med Ideverdenen og sætter som sit Maal „den fuldkomne
Fremstilling af Intellektualverdenen i den phænomenale
Verdens Love og Former, og altsaa omvendt
Begriben af disse Love og Former ud fra Intellektual%
% --- p. 5 ---
\setcounter{footnote}{0}%
\opage{5}verdenen“. Schellings Naturerkjendelse gaaer altsaa ene
ud paa Phænomenernes ideale Betydning, hvorimod den
Opgave at forklare Tingenes virkelige Sammenhæng ikke
existerede for ham. Han saa en ny Periode ifærd med
at begynde, hvor Menneskeheden fra den tidligere ufuldkomne
Erkjendelse skulde føres til Skuen af det Absolute,
den ene sande Realitet. Denne prophetiske Skuen
bragte Hegel under Begrebets Form, idet han ved sin
dialektiske Methode udviklede Virkelighedens Indhold i
dets Enhed med Tankens (Ideens) Væsen. Herved forsvandt
endog det Skjær af Subjektivisme, som endnu
klæbede ved Schellings Konstruktioner, og de idealistiske
Konsekventser af Kants Lære udvikledes til deres
Høidepunkt.

Medens saaledes her det realistiske Element i Kants
Lære efterhaanden ganske elimineredes, tog dog andre
Tænkere deres Udgangspunkt netop i dette Element.
Allerede som ung Student og Fichtes Tilhører i Jena
udarbeidede Herbart en Kritik af Schellings første
Skrifter, hvori han bestrider Berettigelsen af at overføre
Tankens, Videns Form, den sig udviklende, aktive Enhed,
paa Videns Indhold og Gjenstand. Thi Videns Indhold
er den absolute Væren, men denne er uforenelig
med Begrebet om at sætte og frembringe sig selv; den
er Et med den absolute Ro og Stilhed, medens Jeget
er en evig frem- og tilbagestrømmende Hvirvel, derfor
maa begge holdes vel ude fra hinanden. I denne Kritik
er Herbarts senere Lære allerede indeholdt. Viden skal
trække sig tilbage i sig selv, for at Indholdet frit kan
fremstille sig. Indholdet stammer fra Erfaringen, og
Tanken har kun at undersøge, om Erfaringsbegreberne
ere frie for Modsigelser. Da nu dette langt fra er Tilfældet,
maa det Givne behandles saaledes, at man derigjennem
kan trænge ind til Sagens sande Væsen: „so
viel Schein, so viel Hindeutung af Sein“. Tanken faaer % PRINTED AS IS (af Sein: the image, af for auf; translation odd spot O003)
altsaa her kun formel Betydning; den skal ved sin Kri%
% --- p. 6 ---
\setcounter{footnote}{0}%
\opage{6}tik udskille den absolute Position, den sande Væren,
som ubevidst indeholdes i den sandselige Fornemmelse.
--- Schopenhauer er forsaavidt enig med Herbart, som
han i sin Kritik af den kantiske Philosophi erklærer, at
Kants Feil laa i den Maade, hvorpaa han udledte „das
Ding an sich“, ikke i, at han overhovedet anerkjendte
den. Herved kom han i en saa afgjort Opposition
til den spekulative Idealisme, at den endog udartede
til formelig Fanatisme, idet han kun mente at kunne
udlede, hvad han kaldte den spekulative Philosophis
Misforstaaelse af Kant, fra en Tilsidesættelse af Sandheden
for personlige og egoistiske Øiemed. Tingen i og
for sig bliver dog hos Schopenhauer ikke staaende som
et ubestemt Noget; tvertimod, dette Grundlag for alle
Phænomener, for den hele Natur er for ham Et med,
hvad vi finde i vort eget Indre som Villie. Men Villien
er forskjellig fra og absolut uafhængig af Erkjendelsen,
som kun er et Hjernephænomen, en flygtig Illusion,
medens Villien er det eneste Oprindelige og Reale.
„Das Ding an sich“ er saaledes her blevet Et og Alt,
ligesom omvendt hos Hegel Tanken, Begrebet var Et og
Alt. Naar Schopenhauer siger, at han har gjennemført
hvad Kant havde ladet ufuldendt, da kunde Hegel fra
sit Standpunkt sige det Samme.

Vi se altsaa, at efter Kant skilles Veiene atter,
efterat de en kort Tid havde forenet sig i ét Knudepunkt.
Men den idealistiske Retning blev baade paa
Grund af dens eget Væsen og Tidsalderens Karakter
den Alt beherskende i den paa Kant følgende Menneskealder.
Aarhundredets Begyndelse var idealistisk og romantisk;
man svælgede i de store Ideer, der vare opgaaede
for Bevidstheden under den nyere Tids Kamp
mod Fortidens Magter. Og den spekulative Idealisme
raadede netop over Former, i hvilke de saaledes vakte
Interesser kunde finde deres Udtryk og komme til Klarhed
over sig selv. Hegelianismens store kulturhistoriske
% --- p. 7 ---
\setcounter{footnote}{0}%
\opage{7}Betydning beroede netop paa dens Stræben efter at føre
alle væsentlige Modsætninger tilbage til en omfattende,
Alt forsonende Enhed. Først efter det hegelske Systems
Opløsning faa Herbart og Schopenhauer en mere indgribende
Betydning, saaledes at der navnlig om den
første af disse Mænd danner sig en fast sluttet Skole
(hvis Organ er „Zeitschrift für exakte Philosophie“).
Men dette hindrer ikke, at deres Plads i Philosophiens
Historie maa anvises dem i en Række, der parallelt med
den spekulative Idealisme udvikler sig af Kants Lære.

De philosophiske Phænomener, som vi i dette Skrift
ville sysselsætte os med, staa i et særligt Forhold til
begge disse Rækker. Udviklingen indenfor den Periode,
vi ville behandle (1830---1870), betinges fra først af ved
Forholdet til den hegelske Spekulation; men under dennes
Kritik og Opløsning gjøre da naturligen Spekulationens
Antipoder (Herbart og Schopenhauer) deres Indflydelse gjældende,
ligesaavel som de andre Elementer, den havde trængt
tilbage, og som netop nu udvikle sig med forbausende
Kraft (den exakte Erfaringsvidenskab). Og under alt
Dette søger man stedse tilbage til Kant, som den, fra
hvem den hele Bevægelse er udgaaet, for hos ham at
orientere sig med Hensyn til alle vigtige Hovedproblemer.

2. Efter denne almindelige historiske Overveielse,
hvorved vi have bestemt vor Periodes Stilling indenfor
det større Afsnit i Philosophiens Historie, til hvilket
den hører, maa vi nu nøiere karakterisere vort særlige
Udgangspunkt, den hegelske Philosophi, hvis Herredømme
over Aanderne netop ved Stifterens Død (1831)
var paa sit Høidepunkt.

Hegel satte fra første Færd af sin Anskuelse i Modsætning
til tvende ved Aarhundredets Begyndelse herskende
Retninger. Ifølge den ene kunde Mennesket
kun gjennem Anelsen eller den umiddelbare Følelse
komme i Forhold til det Uendelige. Men herved kastes
en Taage over Tankens og Tilværelsens Liv og Mang%
% --- p. 8 ---
\setcounter{footnote}{0}%
\opage{8}foldighed, og man nøies med „den ubestemte Nydelse
af den ubestemte Guddom“. Den anden Retning hævder
vel Tankens Betydning, men opfatter det Absolute som
umiddelbar Identitet, i hvilken alle Modsætninger ere
forsvundne, og som derfor kun kan gribes i en intellektuel
Anskuelse med kunstnerisk Genialitet. Men ogsaa dette
er uvidenskabeligt. Det Absolute kan ikke fremstilles
eller gribes umiddelbart („wie aus der Pistole geschossen“);
det vilde som det forskjelsløse Mørke ikke indeholde
nogen Forklaring af Tilværelsen. Den umiddelbare
Følelse og Anskuelse maa hæves til virkelig, spekulativ
Viden ved en videnskabelig Methode og i Kraft
af et videnskabeligt Princip.

I selve Fordringen om en Methode ligger, at det
Umiddelbare ikke lades i Ro, men underkastes en Prøvelse,
underlægges Reflexionen. Reflexionens negerende
Kraft trænger sondrende ind i den umiddelbare Enhed,
Momenterne stilles imod hverandre, der bliver Strid og
Modsætning istedetfor Fred og Hvile. Men herigjennem
kommer netop Enhedens sande Væsen for Dagen. Negationen
er ikke blot den opløsende, men ogsaa den forenende
Magt. De enkelte løsladte Momenter faa ikke umiddelbar
Selvstændighed; derved vilde der kun være sat en
umiddelbar Mangfoldighed istedetfor en umiddelbar Enhed.
Men de enkelte Momenter udelukke og negere hverandre,
og gjennem dette System af Negationer naaes en høiere,
gjennem Reflexionens Fuldendelse vunden Form for den
oprindelige Enhed. Methoden er altsaa dialektisk ɔ: den
lader Indholdet selv udvikle sig i hele sin Fylde gjennem
Modsigelsen, ved Negationens fremaddrivende Kraft.

Ved at anbringe Negationen føier Tanken ikke Noget
til af sit Eget. Negationen er immanent, indeholdes i
selve Væsenet. Det Sande er ikke, som Spinoza vilde,
den absolut positive Substants. Spinoza har Ret i, at
enhver Bestemmelse er en Negation, men han overser,
at enhver Negation tillige er en Bestemmelse, og at
% --- p. 9 ---
\setcounter{footnote}{0}%
\opage{9}kun gjennem Modsigelsen naaes det konkrete, virkelige
Indhold. Den substantielle Enhed maa, for at være
virkelig, artikuleres, og Artikulationen sker i Kraft af
den immanente Negation, der lader Forskjellene træde
frem af Enheden. Herved udvikles Substantsen til Frihed,
til Subjektivitet.

Heraf følger, at i den dialektiske Methode falder
Tankens Bevægelse sammen med Indholdets. Den absolute
Subjektivitet udvikler sig selv, kommer til sig
selv gjennem den spekulative Begrebsudvikling. I den
dialektiske Spekulation har Aanden da gjenvunden den
tabte Hvile i det Uendelige, som Skepticismen og Kriticismen
havde gjort Ende paa, og som Følelses- og Anskuelsesphilosophiens
Vinger ikke vare kraftige nok til
at naa igjen. Tanken griber nu med Frihed og Bevidsthed,
hvad Fædrene med umiddelbar Tro hengav sig til;
dette Fremskridt er det nu „lykkedes Verdensaanden“
at gjøre.

Den konkrete Tilværelse er som saadan den logiske
Tilværelse. Thi det er Bestemtheden, Kvaliteten, der
betegner den konkrete Tilværelse; Bestemtheden er Forholdet
til, Ligheden med sig selv, men denne Identitet
er Abstraktion, altsaa Tanke. Omvendt er det rene Begreb,
der kun forholder sig til sig selv, umiddelbar Identitet
ɔ: Væren. Væren er det aldeles abstrakte, umiddelbare
Forhold til sig selv, Begrebets abstrakte Moment,
der netop ved sin absolut abstrakte Karakter ligger
under for Negationen og derved aabner den dialektiske
Bevægelse. Ved denne Bevægelse vindes den videre
Udvikling af, hvad der i sig (an sich) allerede laa i Udgangspunktet,
den oprindelige Enhed af Tænken og Væren.
For Hegel er derfor Logikens Indhold den sande
Virkelighed. „Kun i sit Begreb,“ hedder det i Indledningen
til Logiken, „har Noget Virkelighed; forsaavidt
det er forskjelligt fra sit Begreb, hører det op at være
virkeligt og er et Ikkeværende (ein Nichtiges); til denne
% --- p. 10 ---
\setcounter{footnote}{0}%
\opage{10}Ikkeværen hører Haandgribeligheden og den sandselige
Udenforsigværen“. Hvad vi altsaa ikke støde paa under
den logisk-dialektiske Deduktion, det er kun Slakker,
der ikke høre med til Virkelighedens rige Erts. --- Den
umiddelbare Væren slaaer ved sin absolut abstrakte Karakter
over i sin Modsætning, og saaledes føres vi bestandig
fra et Moment til et andet. Men Udvortesheden
mellem disse Momenter forsvinder, naar det indses, at
Værens Bestemmelse ved sit Andet egentlig er dens Bestemmelse
ved sig selv, at Negationen væsentlig er en
indre. Væsenet er den paa engang ophævede og bekræftede
Væren; det har sin Modsætning i Skinnet, men
sætter selv Skinnet for derigjennem at hævde sig som
Væsen, som Grunden i Existentsen, som det Indre i det
Ydre. Den absolute Virkelighed er Enheden af disse
Modsætninger. Og denne Virkelighed er det Absolute
selv, der i sin indre Nødvendighed, som bliver gjennemsigtig
for Tanken, aabner sig i et Frihedens og Inderlighedens
Rige og derved godtgjør sig som Subjektivitet,
som Aand. Fra dette høieste Synspunkt gjælder det
altsaa: kun det Aandige et det Virkelige. Videnskaben % PRINTED AS IS (Aandige et: the image, et for er; O005)
er Aandens Virkelighed, det Rige, den bygger sig i sit
eget Element.

Logiken maatte ifølge dette egentlig være den eneste
Videnskab. Hvad den ikke kunde udvikle, kunde
som ot „Nichtiges“ slet ikke blive Videnskabens Gjenstand. % PRINTED AS IS: „ot“
Dog er Hegels Erkjendelsestrang her stærkere
end hans Konsekvents; den kunde ikke lade sig nøie
med den logiske Skyggeverden. Men hvad Andet kan
der da gives foruden den absolute Virkelighed, og hvorledes
kan Overgangen gjøres dertil fra Logiken?

Vi staa her ved det hegelske Systems Achilleshæl,
Overgangen fra Logiken til Naturphilosophien. Vanskeligheden
lægger sig for Dagen i den dobbelte Vending,
Hegel selv giver Sagen. Ved Logikens Slutning hedder
det, at Ideen, der er absolut sikker paa sig selv og
% --- p. 11 ---
\setcounter{footnote}{0}%
\opage{11}hviler i sig selv, frit udlader sig og omsætter sig i Udvorteshedens Form, som „äusserliche Idee“, for saa gjennem Naturens og Aandens stigende Former atter at
hæve sig til Begrebets Inderlighed. Men i Naturphilosophien betegnes Overgangen som et Ideens Affald fra
sig selv. Naturen er den rene Udvortesheds Sphære og
som saadan uskikket til at fastholde Begrebet, der kun
aldeles abstrakt fuldbyrdes i Naturen. Begrebet er i
Naturen, men er ikke Naturen selv, medens i Logiken
Begreb og Virkelighed vare Et. De Metamorphoser,
gjennem hvilke Ideen Skridt for Skridt kæmper sig frem
til Frihed gjennem denne sin Udlændighedstilstand, vedkomme
derfor ikke Naturen, men kun Begrebet i sig
selv, der er Naturtingenes sande Indre. Begrebets dialektiske
Proces og den reale Naturbevægelse have Intet
med hinanden at gjøre.

Et lignende Forhold finde vi i Aandsphilosophien:
en Modsætning mellem Aandsbegrebets dialektiske Udvikling
og det virkelige Aandsliv. Det er en hegelsk
Hovedsætning, at hvad der er fornuftigt er ogsaa virkeligt,
hvad der er virkeligt ogsaa fornuftigt. Ideen er
ikke saa afmægtig, at den kun skal være, og ikke virkelig
er. Hegel klager selv over, at man har misforstaaet
disse Sætninger og opfordrer til at søge deres Forklaring
i hans Logik. Ifølge Grundtanken i denne er nu,
som vi have vist, Ideen selv, det fuldt udviklede Begreb
den sande, den eneste Virkelighed, medens det Tilfældige,
det blotte Phænomen ikke har Krav paa Virkelighedens
Navn. Det er altsaa klart, at „Virkelighed“
her betegner den ideale, den fornuftige Virkelighed; men
det er da ogsaa klart, at Hegel overser det væsentlige
Problem, der her maa fremstille sig, nemlig Spørgsmaalet
om, i hvilket Forhold denne ideale Virkelighed
staaer til den ydre, faktiske Virkelighed. Sætningen om
Fornuftens og Virkelighedens Enhed er derfor tvetydig,
men netop i sin Tvetydighed karakteristisk for Hegel og
% --- p. 12 ---
\setcounter{footnote}{0}%
\opage{12}hans Tidsalder. Navnlig indeholder den Forklaringen af,
hvorledes han, der i sin Ungdom bekjendte sig til
fuldstændig radikale Anskuelser saavel i religiøs som i
politisk Henseende, i sin modne Manddom uden at fornegte
sit formelle Princip kunde optræde som det orthodoxe
Dogmes og den konservative Politiks Forkæmper,
--- og hvorledes saa atter den saakaldte venstre Side af
hans Skole ved konsekvent Fastholden af dette samme
Princip kunde komme til aldeles negative og revolutionære
Resultater baade paa det religiøse og det sociale
Omraade.

Det Godes Ide bestemmes i Logiken som en blot
subjektiv Fordring eller Skullen, der bestandig har Virkeligheden
udenfor sig og derved stedse hindrer sin egen
absolute Virkeliggjørelse, hvorfor den maa finde sin Udfyldning
i det Sandes Ide, i Kraft af hvilken Begrebet
og Virkeligheden ere Et, saa at Ideen altsaa ikke trænger
til den subjektive Virksomhed for at opnaa Realitet.
Det enkelte Subjekt forsvinder da som saadant
overfor den absolute Ides evige Virkelighed. I Konsekvents
heraf lærer saa Retsphilosophien, at den sædelige
Aand i og for sig har Virkelighed og ikke afhænger af
Individerne. Det Sædelige er i sig det Almene, der vel
behøver Individets Bevidsthed derom og den derved
vakte subjektive Lidenskab, men saaledes at det Individuelle
kun er et af de forsvindende Momenter, Ideen i
sin Udvikling sætter og atter ophæver som enkelte Former
for sin Virkelighed. Den ethiske Stræben, det frie
personlige Liv , den historiske Udvikling faaer saaledes
ikke nogen sand Betydning. Ideens Væsen i dens absolute
Autarki er og bliver den eneste Virkelighed. Den
subjektive Aand gaaer aldeles op i den objektive Aand,
d. v. s. i Staten og de af den omsluttede Samfund, der
ere den sædelige Ides Udtryk.

Dog hæver sig endnu ud over den objektive Aand
en høieste Sphære, den absolute Aand, der i sig sam%
% --- p. 13 ---
\setcounter{footnote}{0}%
\opage{13}menfatter alle tidligere Stadier. Her finder den høieste
Forsoning Sted gjennem Kunst, Religion og Videnskab,
og vi træffe her i Forholdet mellem de to sidste Magter
en særegen Anvendelse af Sætningen om det Fornuftiges
og det Virkeliges Enhed. Det religiøse Standpunkt betegner
her Virkeligheden. Det er for den speculative
Viden vel kun et endeligt Standpunkt, da det bevæger
sig i Forestillingens og Følelsens umiddelbare og med
det Individuelle behæftede Former; men dog skal der
i Sagen selv ingen Modsætning være. Hvad der tidligere
blev troet som Mysterium, skal nu aabenbares for
Aanden i Tankens almene og nødvendige Form. Og
under denne Formforandring bevares Indholdet uforvansket.
Thi baade Philosophi og Religion have det Uendelige
til Indhold, kun at det i Religionen, der netop herved
viser ud over sig selv, fremtræder under endelige
Former. Dette er Principerne for den spekulative Forsoning
af Videnskab og Kirkelære. ---

Med indre Nødvendighed blev Fremstillingen af det
hegelske Systems Grundtanker tillige en Kritik, idet
Modsigelsen mellem det Maal, den spekulative Idealisme
satte sig, og det Resultat, den virkelig naaede, paa de
afgjørende Punkter lagde sig for Dagen i uvilkaarlige
Tilstaaelser. Modsigelsen angaaer selve Grundproblemet
om Forholdet mellem Fornuft og Virkelighed. Men dog
maa det ikke overses, hvilke store Fortjenester den spekulative
Idealisme har indlagt sig af dette samme Problem.
Kant var ved sin skarpe Kritik bleven ført til
en Sondren mellem de to Elementer, saaledes at Tanken
med sine Former stod paa den ene Side, Tingene med
deres virkelige Væsen paa den anden Side; Tanken
blev derved rent formel og subjektiv; Tingenes Væsen
det Fjerne, Formløse og Ubestemte. Den spekulative
Idealisme har ved med Begeistring og Energi at fordybe
sig i Tankens Væsen godtgjort, at dette idealt indbefatter
Gjenstandenes Væsen i sig; fra den ideelle Side set
% --- p. 14 ---
\setcounter{footnote}{0}%
\opage{14}er Modsætningen mellem det Subjektive og Objektive
hævet. Tanken bærer i sig Tilværelsens hele Rigdom
som sit mulige Indhold. Den spekulative Idealismes
store Vildfarelse var nu den at forvexle Mulighed med
Virkelighed, den logiske Virkelighed med den reale.
Derfor mente den ad rent logisk Vei at kunne naa den
konkrete Tilværelses Mangfoldighed. Naar derfor ved
den spekulative Philosophis Opløsning Problemet om
Forholdet mellem Fornuft og Virkelighed fremstiller sig
paa Ny, saa betegner Fornuften her ikke blot Tankens
subjektive Form, men Selvbevidsthedens og Ideens hele
mangfoldige og indholdsrige Verden, Virkeligheden ikke
den formløse „Ting i og for sig“, men den i sine konkrete
og individuelle Former aabenbarede Tilværelse. Spekulationen
havde reduceret Forskjellem mellem disse Elementer % PRINTED AS IS (Forskjellem: the image, for Forskjellen; O006)
til en blot Formforskjel, der dialektisk kan sættes
for atter at ophæves; den reale Virkelighed blev
kun den ufuldkomne Form for den ideale Virkelighed.
Hos en Tidsalder, hvis hele Retning var idealistisk,
maatte dette storartede Forsøg paa at gjennemføre det
Ideale som den eneste sande Virkelighed vække Begeistring
og Tilslutning; men naar Virkelighedens Interesse
vaagner, maa Resultatet vise sig at bero paa en Illusion.

3. Spørgsmaalet om \emph{Forholdet mellem Fornuft
og Virkelighed} kan opfattes fra to forskjellige
Sider. Fornuften betegner i Modsætning til Virkeligheden
det Almene, Loven for Enkeltvæsenernes Væren
og Virken, Formaalet for og Værdien af deres reale
Udvikling. Saaledes bliver Spørgsmaalet da om \emph{Forholdet
mellem det Almene og det Individuelle}.
Men som Udtryk for den almene Lov og Formaalet staaer
Fornuften for det Andet som det Ideelle overfor det
Virkelige som det i materiel Realitet Existerende. Den
anden særlige Form for det almindelige Problem bliver
da Spørgsmaalet om \emph{Forholdet mellem det Aandelige
og det Materielle}.

% --- p. 15 ---
\setcounter{footnote}{0}%
\opage{15}Den spekulative Idealisme betragtede, som vi have
set, Individerne som forsvindende, i sidste Instants betydningsløse Momenter i den almene Ides Udvikling.

Saasnart Kritiken vaagnede, kastede den sig over dette
Punkt. Først gjøres der Ende paa al Tvetydighed ved,
at man \emph{pantheistisk} drager de sidste Konsekventser
af Begrebet om Ideen som eneste Virkelighed. Men
snart trænger saa Bevidstheden om det Enkeltes og Individuelles
Betydning sig atter frem, og man søger at
hævde det en Realitet, der gaaer videre end at være et
blot dialektisk Gjennemgangspunkt. Denne Bestræbelse
fører igjen i en dobbelt Retning. Dels ser man Enkeltheden
kun i den endelige Individualitet, og idet dette
Synspunkt gjennemføres, forsvinder Skridt for Skridt alt
Alment og nedsættes tilsidst til at være det Individuelles
Egenskab eller Prædikat. Dette Standpunkt er
det \emph{anthropologiske}, i sit Extrem det \emph{egoistiske}.
Dels ses Enkeltheden ogsaa som den Koncentration, ved
hvilken det Almene fatter sig selv, og som derfor er en
nødvendig Betingelse for det Almenes Realitet. Man
søger da at udvikle det spekulative Begreb om det Absolute
til en fyldigere og konkretere Gudside, i hvilken
Personlighedens Moment kommer til sin Ret. Denne
Bestræbelse gaaer altsaa i \emph{theistisk} Retning.

Ved sin Forflygtigelse af det Individuelle havde den
spekulative Idealisme krænket det Reale i det Hele. De
virkende Aarsager blive tilsidesatte for den høiere, ofte
mystiske „Idenødvendighed“, og det materielle Stof opløste
sig for den spekulative Konstruktions Kraft ligesom
Isen for Solstraalernes. Men den følgende Tidsalder
udviklede sig mere og mere i realistisk Retning,
og Reaktionen mod Idealismen steg til en saadan Grad,
at man i Ideen kun saa et Hjernespind og i det Materielle
den eneste Virkelighed. Overfor denne \emph{Materialisme}
træder dog en anden Retning, der ligesaameget som
den, men med mere Besindighed og Videnskabelighed,
% --- p. 16 ---
\setcounter{footnote}{0}%
\opage{16}anerkjender Nødvendigheden af det reale Grundlag og
de mechaniske Virkeformer, men som i alt Dette kun
ser Former og Midler for høiere, betydningsfulde Formaals
og Ideers Virkeliggjørelse, en Retning, som vi
derfor med Rette kunne kalde \emph{Virkelighedsidealismen}.


Kampen om det første Problem føres især i de to
første Decennier efter Hegels Død (Aarene 30 og 40),
Kampen om det andet Problem i de to følgende Decennier
(Aarene 50 og 60). Dog vil det være naturligst
ikke at følge Tidsordenen, men at sammenstille de beslægtede
Besvarelser af begge Problemer. Vi ville da
faa fire Hovedgrupper: 1) den kritiske Bevidsthed; 2) den
moderne Materialisme; 3) den spekulative Theisme; 4)
Virkelighedsidealismen, af hvilke de to første kunne betegnes
som negative i Modsætning til de to sidste som
positive. Thi den kritiske Bevidsthed udvikler sig gjennem
den spekulative Idealismes Opløsning og fører kun
indirekte og under Fordringens Form til et virkeligt,
nyt Indhold; og Materialismen, der tager sit Princip
udenfor Tanken og Erkjendelsen, maa konsekvent komme
til at staa i et skeptisk og fjendligt Forhold til Philosophien.
De to sidste Grupper gjøre derimod virkelige Forsøg
paa at udvikle en philosophisk Verdensanskuelse.

Overfor den storartede Produktivitet, som den tydske
Philosophi udfoldede, efterat Kant først med en genial
Forening af Alsidighed og Dybsindighed havde brudt de
nye Baner for Tænkningen, kan der synes at være noget
Ubetydeligt og blot Foreløbigt, noget Uselvstændigt og
Ufuldendt i den efterhegelske Philosophis Karakter. Vi
møde her, som allerede bemærket, ingen Tænkere af
første Rang, ingen epochegjørende, verdenshistoriske Opdagelser
i Tankens Verden. Men man maa betænke
den Opgave, der er stillet. Den spekulative Idealisme
kunde hurtig fuldende sine storartede Systemer, naar
den først ad Intuitionens Vei havde fundet Grundtanken;
% --- p. 17 ---
\setcounter{footnote}{0}%
\opage{17}Methoden var den abstrakt dialektiske, og Bevægelsen
foregik i Tankens rene Ætherregion, som saa forvexledes
med Virkeligheden selv. Den efterhegelske Philosophi er
neppe saa enig i Noget som i den skarpe Kritik af
denne abstrakte Idealisme. Den vil have at gjøre med
Virkeligheden som Virkelighed, men denne gribes ikke
paa én Dag af Tanken. Her maa gaaes fremad Skridt
for Skridt. Hverken Tankens eller Virkelighedens Ret
maa krænkes; de Punkter maa findes, hvor begge mødes,
Loven for denne Overensstemmelse findes, og det saaledes
Vundne efterhaanden samles til en omfattende Totalanskuelse.
Dette Arbeide er tilsyneladende mindre glimrende
end Spekulationens; men det fordrer ikke mindre
Anstrengelse og langt mere Resignation. Den spekulative
Philosophi hengav sig trygt til sin Grundtankes Konsekvents,
som Digteren til Phantasiens Magt; Spekulationen
er Philosophiens Lyrik. Hvor det derimod gjælder
at gribe Virkeligheden som Virkelighed og dog hævde
det sande Ideelles Ret, der gaaer Flugten ikke saa let,
og derfor heller ikke altid saa høit. Dog har heller
ikke her den evige Drift mod Frihed og Sandhed, som
er Menneskeaanden indplantet, ladet sig uden Vidnesbyrd,
men har lagt sig for Dagen i klare og dybe Tanker,
baarne af ædle Personligheder.

\begin{center}\rule{3em}{0.4pt}\end{center}
\chaphere{ANDET KAPITEL.}
\begin{center}\rule{2em}{0.4pt}\end{center}


I. \emph{Striden om den individuelle Udødelighed}.
--- Naar den hegelske Philosophi virkelig havde
„forsonet Idealet med Virkeligheden, vore Fordringer
med det, som vi besidde, vore Ønsker med det, som er
% --- p. 18 ---
\setcounter{footnote}{0}%
\opage{18}opnaaet“\footnote[1]{J. L. \emph{Heiberg}: Om Philosophiens Betydning for den nuværende Tid. S. 48.}, da kunde dens Grundtanke om Fornuftens
og Virkelighedens Enhed ikke staa hen i ren Almindelighed
--- den vilde jo saa selv være en af de Fordringer,
som nu skulde være tilfredsstillede. Den maatte idetmindste
være istand til at afgive et bestemt Svar paa
de afgjørende Spørgsmaal, det virkelige Liv leder til at
opkaste. Systemet maatte derfor først bestaa den Prøve
at stilles ligeoverfor de religiøse og sociale Problemer.
Men allerede med Hensyn til det første Punkt, som
blev gjort til Gjenstand for Debat, Spørgsmaalet om,
hvorvidt den hegelske Philosophi lærte Sjælens Udødelighed,
viste Systemets Ubestemthed overfor det Konkrete
og Individuelle sig klart. Hegel havde med sin
hele Intensitet fordybet sig i sin Grundtanke i og for
sig uden at gjøre de konkrete Problemer til Gjenstand
for særskilte Undersøgelser. Derfor kunde allerede denne
første Diskussion føre til en Spaltning indenfor hans
Skole i Retninger, af hvilke enhver mente at fortolke
Hegel paa den rette Maade.

Allerede før Hegels Død havde man bebreidet hans
Lære, at den nedbrød Udødelighedstroen. Og paa den
anden Side havde Ludwig Feuerbach i sit anonyme Skrift:
„Gedanken über Tod und Unsterblichkeit“ (1830) ud fra
selve det hegelske System pantheistisk negtet den individuelle
Udødelighed. Men den egentlige Strid udbrød
først, da Fr. Richter udgav sit Skrift: „Die Lehre von
den letzten Dingen“ (1833), hvori han med samvittighedsfuld
Nøiagtighed og Grundighed viste, at Konsekventsen
af Hegels Lære var Ophævelsen af den gamle Lære om
et Hinsides og et andet Liv. To Mænd, der allerede
vare optraadte med Kritik mod det hegelske System,
C. H. Weisse og I. H. Fichte (Søn af den berømte
I. G. Fichte), erklærede sig forsaavidt enige med Richter,
som de negtede Muligheden af at hævde den 
indivi%
% --- p. 19 ---
\setcounter{footnote}{0}%
\opage{19}duelle Udødelighed ud fra Hegels Forudsætninger, hvorimod
de fastholdt Muligheden af ad anden Vei at godtgjøre
den. Imod disse Fortolkninger af Hegel optraadte
saa den Mand, som paa denne Tid havde den største
Anseelse blandt Skolens Tilhængere, Carl Friederich Gö-
schel, for at redde Systemets Orthodoxi i dette Punkt,
med sin Afhandling: „Von den Beweisen für die Unsterblichkeit
der menschlichen Seele im Lichte der spekulativen
Philosophie“ (1835)\footnote[1]{Angaaende den yngre Fichtes, Weisses og Göschels Skrifter i denne Sag kan jevnføres \emph{Poul Møllers} interessante Afhandling: „Om Muligheden af Beviser for Menneskets Udødelighed“. Efterladte Skrifter 5te Bind.}. Med Göschels Opfattelse
begynde vi vor Fremstilling for derpaa at gaa over til
de forskjellige Dissenters.

Göschels Afhandling har sit Forbillede i det Skrift,
med hvilket Hegel syslede lige før sin Død, om Beviserne
for Guds Tilværelse. Hegel søgte at vise, hvorledes
de tre Hovedformer for disse Beviser vise tilbage
til én gjennemgaaende Tanke, der er Et med den spekulative
Philosophis Grundtanke om Forholdet mellem
det Uendelige og det Endelige, --- naar man kun gjennembryder
den udvortes og forstandsmæssige Skal, hvoraf
de i den populære Fremstilling dækkes. Noget Tilsvarende
forsøger nu Göschel med Beviserne for Sjælens
Udødelighed. Det saakaldte metaphysiske Bevis, der
svarer til det kosmologiske for Guds Tilværelse, gaaer
ud fra Sjælen som absolut enkelt og immaterielt Væsen
og slutter af denne Forudsætning, at den ikke kan
undergaa nogen Opløsning. Den slutter fra den ydre
Tilfældighed til den indre væsentlige Væren. Dette
Bevis viser ud over sig selv. Det er i Kraft af Bevidsthedens
Enhed, at Individet er det absolut enkelte
Væsen. Som Bevidsthed staaer Mennesket overfor den
øvrige Tilværelse. Det bevidste Sjæleliv er det Formaal,
til hvilket Livet sigter; Bevidstheden skal udvikle sig
% --- p. 20 ---
\setcounter{footnote}{0}%
\opage{20}saaledes, at den optager den hele Tilværelse i sin Inderlighed.
Som Bærer af et evigt Formaal maa det bevidste
Sjæleliv selv være evigt. Dog fører dette teleologiske
Bevis os ikke ud over Modsætningen mellem Subjekt og
Objekt, Selvhed og Andethed. Først naar Subjektet har
udviklet sig saaledes, at det ikke blot har overvundet
Legemligheden og tilegnet sig Tilværelsens objektive
Indhold, men og tilsidst har fundet og gjort sig til Et
med sin og Alts evige Grund og Kilde i den absolute
Subjektivitet, først da er den individuelle Bevidsthed
udviklet til Aand, og Kredsen er sluttet, idet den sidste
Begrundelse er funden for den tilfældige Individualitet,
med hvilken der begyndtes. Paa dette Høidepunkt vender
Sjælen tilbage til sit Udspring med et rigt udviklet personligt
og evigt Indhold, og dens Udødelighed kan derfor
udledes af dens eget Væsen, saa at vi her faa et
Bevis, vi kunne kalde det ontologiske.

Göschel lader sig ikke nøie med denne spekulative
Omsmeltning af de vulgære Beviser, men søger at godtgjøre,
at Hegels Lære ifølge hele sin Karakter maa føre
til at hævde Individets evige Realitet\footnote[1]{Jahrbücher für wissenschaftliche Kritik. 1834.}. Han henviser
da dels til den logiske Bestemmelse af Forholdet mellem
det Almene og det Individuelle, dels til Forholdet mellem
Aand og Natur, dels endelig til Opfattelsen af Gudsideen.
--- Den hegelske Logik sætter et indre Forhold
mellem det Almene og det Enkelte. Det Almene i
og for sig er goldt og abstrakt; for at være det sande
Almene maa det udvikle sig til Konkretion, sætte sine
indre Forskjelle som sine Momenter og saaledes blive
Totalitet, Enhed, Enkelthed. Omvendt er det Enkelte
i og for sig atomistisk og abstrakt; det ses først i sin
Sandhed ved at ses som det, der i sig indeslutter og
udtrykker almene Bestemmelser: det Enkeltes sande
Væren er at være det Almene. Paa disse dialektiske
% --- p. 21 ---
\setcounter{footnote}{0}%
\opage{21}Bestemmelser mener nu Göschel at kunne bygge et Bevis
for Personlighedens evige Gyldighed. Individet, siger
han, godtgjør sig som det Høieste, som det Almenes
Sandhed, fordi det har optaget det Almene i sig; derfor
er det Almene ikke længere en fremmed eller fjendtlig
Magt, overfor hvilken Individet skulde gaa til Grunde.
--- I Naturen gaaer det ganske vist saaledes, at Slægten,
det Almene, opretholdes gjennem Individernes Opstaaen
og Forgaaen; men for den Gjæld, Individet her
maa betale Slægten ved sin Død, er Mennesket befriet,
idet han som Aand er kommen ud over Modsætningen
af Slægt og Individ, det Almene og det Enkelte. I
Aandens Sphære er Slægtens Væsen frit optaget og
derved overvundet af Personligheden. --- Den vulgære
Pantheisme maatte ganske vist lade Individet gaa til
Grunde, fordi den kun opfatter det Almene som objektivt
og upersonligt; men den spekulative Logik udvikler
Gudsideen som absolut Subjektivitet, hvori Individet
kun kan forynges, ikke forgaa. Gud er den høieste Bevidsthed,
den evige Tanke, der opholder og opbevarer
de Forskjelle, de individuelle Bestemmelser, den tænker.
Den menneskelige Aand er en guddommelige Tanke, og % PRINTED AS IS (en guddommelige: the image, for guddommelig; O009)
Guds Tanke kan ikke forgaa, men maa være evig, som
Guds Væsen er det. Naar Individerne ikke bevaredes,
vilde den absolute Aands Erindring være om Død og Tilintetgjørelse,
og hvad blev der saa af dens eget Liv og
udødelige Kraft? ---

Den yngre \emph{Fichte} („Die Idee der Persönlichkeit
und der individuellen Fortdauer“) anerkjender det som
Hegels store Fortjeneste at have vist, hvor vidt den
apriorisk-dialektiske Deduktions Vei kan føre; men denne
Fortjeneste vilde man kun fordunkle ved at drage den
over paa et Omraade, hvor den ikke gjælder. Det Personlige
og Individuelle kan ikke erkjendes ad den dialektiske
Konstructions Vei; for at naa til det maa Erfaringen
gjenindsættes i sin Ret overfor den abstrakte
% --- p. 22 ---
\setcounter{footnote}{0}%
\opage{22}Spekulation. Göschel feiler derfor, naar han finder Guds
Personlighed indeholdt i Hegels Lære. Det er det hegelske
Udtryk „overgribende Subjektivitet“, der har
givet Anledning til denne Misforstaaelse, som bortfalder
ved nærmere Undersøgelse af den Maade, hvorpaa Hegel
udvikler Subjektivitetens Begreb. Gud er for Hegel ikke
blot Substants eller død Identitet, men den levende
Proces, der evig sætter og ophæver Modsætningerne,
stedse er et uendeligt Andet, men derigjennem stedse
vedbliver at være sig selv. Den udtømmer sig aldrig i
nogen enkelt Selvudformning, men griber ud over enhver
særskilt Form. Og just fordi Gud efter Hegels Opfattelse
er den absolute Proces i dens evige Uro, mangler
han ifølge I. H. Fichte den absolute Personligheds
Kjerne og Midtpunkt. Hegels Fortjeneste er begrebsmæssig
at have paavist Guddommens uendelige Liv og
Bevægelse; men den anden Hovedside lader han til
Gjengjæld træde tilbage, nemlig den evige Ro, Identitetens
stille Klarhed, den sig selv fattende Enheds
Baand, hvorved den Alvirksomme for sig selv bliver den
absolute Ene og Salige. Koncentrationen, „das unverrückbare
Selbst“, mangler. Derfor maa Gud i Hegels
System arbeide sig frem gjennem de lavere Stadier for
at naa til den absolute Aands Bestemmelse. Gud bliver
Resultat, ikke Princip.

Forholder det sig saaledes med den guddommelige
Personlighed, da kan det heller ikke gaa anderledes
med den menneskelige. Naar det Absolute er den evige
Proces, hvis Momenter ikke have blivende Betydning,
men sættes og hæves i det Uendelige, saa kan den konsekvente
hegelske Lære med Hensyn til det Individuelle
kun være, at det Individuelle er nødvendigt som dialektisk
Moment, men at det enkelte Individ som saadant
ikke kan være noget Substantielt og Vedvarende. Problemet
om den individuelle Udødelighed kan egentlig
slet ikke opstaa under Hegels Forudsætninger; det
% --- p. 23 ---
\setcounter{footnote}{0}%
\opage{23}fremstiller sig først,  og kan først finde sin Besvarelse,
naar man stiller sig paa Erfaringens, den physiologiske,
den psychologiske, den ethisk-religiøse Erfarings Standpunkt.
Hvorledes den yngre Fichte mente, at dette
Bevis kunde føres,  ville vi her ikke dvæle ved\footnote[1]{Hos Poul Møller (anf. Sted) vil man finde en udførligere Fremstilling.}, da
det staaer i nøie Sammenhæng med hans hele øvrige
Lære, som vi senere ville komme til at beskjæftige os
med. Det maa kun fremhæves, at han ved sin skarpe
Karakteristik af Hegels System bidrog meget til at klare
Synspunkterne, selv om hans egne Forsøg paa at betræde
„Erfaringens“ Vei endnu gaa vel meget i Retning
af det Phantastiske.

Ogsaa C. H. Weisse negtede, at Hegels Philosophi
kunde hævde Personlighedens Gyldighed. Han anerkjendte
den hegelske Logiks store og epochegjørende
Betydning, men tilskrev den kun formal Gyldighed, da den
ikke naaer selve Naturens og Aandens Virkelighed. Ja,
selv som Formalvidenskab er den ufuldkommen, da den
abstraherer fra Rum og Tid, disse Grundformer for al
Virkelighed. Derved er det sket, at Naturen kom til
at staa som den rene Udvorteshed, da Rummet og Tiden
ikke kunne ses i deres væsentlige Betydning for Ideen.
Anerkjendes denne Betydning, da træder Naturen Tanken
imøde som den frie, fyldige Virkelighed, der ikke
kan udtømmes ved logisk Konstruktion, men maa gjøres
til Gjenstand for en philosophisk Empiri. Hvad Aandsphilosophien
angaaer, da hævder \emph{Weisse}, hvis første
Studier især angik den klassiske Oldtid, hvorved der
vaktes hos ham en levende og klar Bevidsthed om det
Skjønnes Værd og Betydning ogsaa i det personlige
Liv, at Hegel med Urette har sat Videnskaben som
det høieste Trin. Sandhedsideen, der viser hen til den
almene Fornuft, maa indgaa en inderlig Enhed med
% --- p. 24 ---
\setcounter{footnote}{0}%
\opage{24}Skjønhedsideen, der viser os det individuelle Stofs Forklaring
ved det Uendelige, for saa endelig gjennem den
ethisk-religiøse Erfaring at gaa op i det Godes Ide.
Ad denne Vei alene naaes det høieste Udtryk for Gudsideen
som absolut Personlighed. Derfor hævder Weisse
ogsaa („Die philosophische Geheimlehre über die Unsterblichkeit
des menschlichen Individuums“), at Mennesket
kun ad den personlige, ethisk-religiøse Vei kan vinde
Overbevisningen om sin Udødelighed. Ifølge det hegelske
System har det enkelte Individ udspilt sin Rolle,
naar Verdensaanden er gaaet igjennem det, og denne
iler saa fremad til nye, rigere og fuldkomnere Former.
Men naar det Sande, Skjønne og Gode optages og indarbeides
i det personlige Liv, giver det dette en evig
og uendelig Betydning og optager det som gjenfødt i
det guddommelige Væsen. Hvad der saaledes er blevet
forklaret, kan ikke atter gaa til Grunde. Men paa den
anden Side følger det ogsaa af sig selv, at kun de Individer,
der paa denne Maade gjenfødes, leve evig; dem,
der ikke have naaet dette Aandstrin, tilkommer der
ingen Udødelighed. ---

Det viser sig saaledes som Resultatet af Weisses
og den yngre Fichtes Undersøgelser, at det hegelske System
egentlig talt slet ikke naaer det Punkt, hvor de
konkrete, individuelle Problemer opstaa. Begge vise de
ogsaa, skjøndt endnu paa ubestemt Maade hen til den
Magt, som Spekulationen havde krænket eller dog overset,
nemlig Erfaringen. Alligevel mente disse Forskere
at kunne fastholde den hegelske Methode idetmindste
indenfor Logiken og Metaphysiken; ved Overgangen til
de konkrete Sphærer skulde saa Erfaringen træde til
som Supplement. Med Rette indvendte Hegelianerne
herimod, at dette er en Umulighed, da den dialektiske
Methode netop er Indholdets Selvbevægelse. Göschel
hævdede i sit Skrift: „Der Monismus des Gedankens“,
at den Paastand, at Virkeligheden var mere end logisk,
% --- p. 25 ---
\setcounter{footnote}{0}%
\opage{25}mere end fornuftig og havde noget Eiendommeligt for
sig udenfor den logiske Form, førte til Dualisme og
Materialisme, forsaavidt som det Materielle postuleres
udenfra som Supplement til Logiken\footnote[1]{Smlgn. \emph{Göschels} Udtalelse i „Jahrbücher für wissenschaftliche Kritik“ 1835. p. 714: „Es ist die feinste Potenz des Materialismus, wenn sie meinen, dass im Leben mehr enthalten sei, als das Denken mit seinen Momenten, und zur Fülle des Lebens noch etwas Anderes hinzutrete als das Denken; vielmehr ist das, was mehr, was voller, dichter zu sein scheint, ein Mangel etlicher Kategorien.“}. --- Først senere
rykkede Kritiken frem til at tage fat paa selve den dialektiske
Methode, idet nemlig Trendelenburg i sine „Logische Untersuchungen“ (1840), som vi siden komme
til, gjennemførte Polemiken saaledes, at han trængte
ind til Systemets inderste Marv. Foreløbig er det, som
allerede bemærket, mere praktiske og konkrete Udgangspunkter,
Kritiken vælger sig. ---

Göschel var enig med sine to omtalte Modstandere
i at hævde den individuelle Udødelighed, men han paastod,
hvad de negtede, at den kunde udledes af Hegels
Principer. Mod begge Parter staa da Feuerbach og
Richter, der negte Udødeligheden i Kraft af disse samme
hegelske Principer.

Ifølge \emph{Feuerbach} lider den nyere Tid af en fordærvelig
Hulhed og Tomhed, fordi den, især i Protestantismen,
har sat Personligheden paa Thronen og lader
Alt dreie sig om den. Deraf den moderne Udødelighedslære.
De Gamle fordybede sig med barnlig Inderlighed
og Livsglæde i den virkelige Verden, og det faldt
dem ikke ind at attraa nogen Vedvaren ud over den;
forsaavidt de havde Forestillingen om et andet Liv, var
det som et Skyggeliv. Den gamle Christenhed og den
middelalderlige Kirke havde egentlig heller ingen Udødelighedstro
i moderne Forstand; de troede paa en Verdensdom,
paa Himmel og Helvede, men den enkelte Person%
% --- p. 26 ---
\setcounter{footnote}{0}%
\opage{26}lighed gik for dem saa aldeles op i de store Formaal og
Forbilleder, at de ikke havde nogen særlig Interesse for
dens rent individuelle Vedvaren. Det er først den moderne,
selviske Reflexion, der ikke kan slippe sit kjære
Jeg men endog forvandler Døden til en Skygge, et Forsvindende
overfor Individualiteten, i hvis Nydelse man
svælger. I den moderne Religion dreier det sig ikke
længer om Gud, men om Individet; Gud er Peripherien,
Individet selv Midtpunktet.

Døden maa igjen indsættes i sin Ret, anerkjendes
for at være mere end Skindød. Døden har nemlig for
det Første en stor ethisk Betydning. Den sande Handlen
er Kjærlighedens Selvhengivelse i Modsætning til det
naturlige Menneskes Egoisme. Det i Sandhed sædelige
og religiøse Menneske sætter sit Væsen i det, der staaer
over ham; hans sande Jeg er den uendelige Gjenstand,
til hvilken han hengiver sit endelige Jeg. Denne frie
-Kjærlighedsgjerning existerer i Naturen som Dødens % PRINTED AS IS (-Kjærlighedsgjerning: the image, a stray hyphen at the line start; recrop rc1)
Nødvendighed; der stilles det Spørgsmaal til Mennesket,
om han vil gjøre denne Nødvendighed til den
høieste Frihed gjennem Kjærlighedens Selvopgivelse;
kun saaledes, kun før Døden kan Døden overvindes.

I metaphysisk Henseende er Døden et Udtryk for
det Uendeliges Realitet: i Døden gjennembryder det
uendelige Væsen de endelige Skranker, indenfor hvilke
den individuelle Existents holdt det. Døden kommer altsaa
ikke af Mangel og Nød, men af Rigdom og Fylde. ---
I psychologisk Henseende gjælder det, at der i Selvbevidstheden
maa skjelnes mellem Bevidstheden selv og
dens Gjenstand. Det Bevidste, Gjenstanden er det enkelte
Individ; Bevidstheden selv er absolut almen, er
i alle Mennesker sig selv lig. De forskjellige Personer
vide sig selv i Bevidstheden, ligesom Farverne ses i
Lyset. Saa lidt som Lyset selv er Et med en enkelt
Farve, saa lidt er Bevidstheden i sig Et med den enkelte
Person. Døden er nu den Kjendsgjerning, i hvilken
% --- p. 27 ---
\setcounter{footnote}{0}%
\opage{27}Aanden, Bevidstheden godtgjør sin Frihed overfor de
enkelte Objekter; som Farverne vexle, fordi de ikke ere
Et med Lyset, saaledes vexle Personerne, fordi de ikke
ere Bevidstheden selv men kun dens Gjenstande, eller
fordi de kun ere Subjekter, ikke Subjektiviteten. Derfor
ere de dog ikke betydningsløse. Bevidstheden, den absolute
Subjektivitet er en evig Erindring, hvori Alt opbevares,
der har været Led i det store Bevidsthedshele.
At Mennesket er udødeligt, vil sige, at han er et Væsen
af Værd og Betydning.

Hvad Feuerbach havde fremsat med flammende
Ord og i epigrammatisk Korthed, udviklede Fr. Richter
med al mulig Vidtløftighed. Hans Skrift har størst
Interesse ved den troskyldige Ligefremhed, hvormed
han er sig bevidst at tale det Godes Sag, idet han tilintetgjør
den i hans Øine saa fordærvelige Lære om et
Hinsides. Han bebreider Hegel og hans Disciple, at de
ikke have draget de fulde Konsekventser af deres Lære.
Ved det hegelske System med dets gjennemførte Immanents
er der efter Richters Mening fremstaaet en ny
Verden; men dog har man ladet den bestaaende Tingenes
Orden staa urørt hen. Dette ligger maaske i Hegelianismens
ensidig theoretiske Karakter; for Hegel er
jo som for Aristoteles Theorien det Høieste. Men man
skal ikke blot tænke absolut, men ogsaa handle absolut,
leve spekulativt. Religionen er ikke noget særligt Bevidsthedstrin,
men indbefatter som Aandens inderligste
Fordybelse alle Bevidsthedstrin, ogsaa Videnskaben, i
sig. I Videnskaben naaer det religiøse Sindelag Erkjendelsen
af sig selv og af Virkeligheden og gaaer derigjennem
over til den praktiske Handlen. --- Fra dette
religiøse Standpunkt gjennemfører Fr. Richter Kritiken
af Forestillingerne om et Hinsides. Han opfatter den
nyere Tids Brøst paa samme Maade som Feuerbach.
Betingelsen for en virkelig ethisk Stræben er, at man
bliver sig bevidst, at denne Verden er en virkelig, i sig
% --- p. 28 ---
\setcounter{footnote}{0}%
\opage{28}afsluttet Verden, og at man hengiver sig til det almene
Formaal som ubetinget Tjener. Heri bestaaer det sande
evige Liv. Kritiken rører derfor ikke ved Udødelighedsdogmets
sande Indhold; den ser kun som nærværende,
hvad der i Dogmet skydes ud i en tilkommende,
hinsides Tilværelse. Motivet til at antage en saadan er
enten ørkesløse, ubegrændsede Ønsker uden Berettigelse
eller en ligesaa uethisk Fordring paa Løn for det Gode
i et andet Liv. Den gode Handlen har sin Løn og sit
Formaal i sig selv. Ved Døden erindres den Enkelte
om, at hans individuelle Liv ikke er Verdens sidste
Øiemed, at Slægten overlever ham, og saaledes lever
han ikke alene for sig selv men for Slægten. I Visheden
om sin Død har Mennesket den første Anledning til at
adskille sig fra Dyret. --- I anthropologisk Henseende
grunder Udødelighedstroen sig paa den gamle dualistiske
Opfattelse af Forholdet mellem Sjæl og Legeme. Sjælen
betragtes som boende i Legemet og forladende dette ved
Døden, der er begges Adskillelse. Men den sande Anthropologi
ser Mennesket som et Hele, som en harmonisk
Enhed fra Vuggen til Graven. ---

Det ene Spørgsmaal om den individuelle Udødelighed
maatte nødvendigen drage en Mængde andre dermed
sammenhængende efter sig; saa dybt griber det ind i
den hele Verdensanskuelse. Saaledes have vi da set,
hvorledes det førte til at drøfte Opfattelsen af Gudsideen,
Erkjendelsens Væsen (Spekulation --- Erfaring)
og Problemet om Immanents og Transscendents. Allerede
denne første Debat grundlagde derfor en Spaltning
indenfor den hegelske Skole, hvis modsatte Partier
Strauss senere (Streitschriften. 3die Hefte. p. 95) betegnede
som Høire og Venstre, efter Partierne i de franske
Forsamlinger. Spørger man, om det konservative Høire
eller det radikale Venstre nærmest udtrykker Skolens
Stifters egen Mening, da kan der ikke gives andet Svar,
end at Hegel vistnok selv personlig troede sig i Over%
% --- p. 29 ---
\setcounter{footnote}{0}%
\opage{29}ensstemmelse med Kirkelæren, men at den philosophiske
Konsekvents af hans System gaaer i modsat Retning.
Den fulde Konsekvents arbeider sig Skridt for Skridt
frem og fattes sjelden eller aldrig i sin hele Udstræk-
ning af Tankens første Ophavsmand. Hegel vilde, som
Ruge træffende bemærker (Deutsche Jahrbücher. 1841.
p. 142), neppe have betragtet den nyhegelske Udvikling
af hans Lære med gunstigere Øine end Kant Fichtes
Udvikling af Grundtankerne i „Kritik der reinen Vernunft.“
--- Vi gaa nu over til et endnu mere brændende
religionsphilosophisk Problem, hvorved hin Modsætning
mellem de to Retninger især træder frem.

2. \emph{David Friederich Strauss} (født 1808). ---
Den spekulative Philosophi havde stillet sig i et lignende
Forhold til den positive Religion, som det Mystiken,
ledet af sin praktisk-personlige Interesse, indtager. Ligesom
Mystikeren troer at anerkjende den hellige Historie,
naar han ser den som en indre, der foregaaer i hans
egen Sjæl under dens Kæmpen og Arbeiden for at føde
Gud i sig og ad Lidelsens Vei hendø i ham, saaledes
troer Spekulationen at have forklaret den historiske
Christendom, naar den har paavist det indre Forhold
mellem det Uendelige og det Endelige og Nødvendigheden
af dettes Opgaaen i hint. Hegel havde ogsaa
modtaget Spiren til sin Religionsphilosophi gjennem Studiet
af Mester Eckardt. Vi have allerede set Grundtanken
i den, men skulle nu betragte dennes Anvendelse
med Hensyn til Kirkelærens Midtpunkt, Christologien.

Den endelige Verden er efter Hegel en Fremtræden
i ydre Virkelighed af det Absolutes indre Forskjelle og
Momenter. Den endelige Aand er vel i sig selv Et med
Guddommen, og oprindelig hviler den ogsaa trygt i
denne Enhed, som umiddelbart fylder Bevidstheden. Men
Skridt for Skridt voxer den endelige Aands Selvbevidsthed
og dermed ogsaa dens Isolation. Det absolute Indhold
svinder mere og mere bort fra den og bliver til et
% --- p. 30 ---
\setcounter{footnote}{0}%
\opage{30}fremmed. Den antike Oldtids Opløsningsperiode er det
Tidspunkt, da det uendelige Indhold fuldstændig er forsvundet,
saa at den tomme Selvbevidsthed staaer tilbage
med stolt Resignation overfor den sammenstyrtede Verden
eller med negativ Glæde over, at Alt vises i sin Intethed.
Dette Standpunkt betegner Hegel som den ulykkelige
Selvbevidsthed: Selvbevidsthedens fuldkomne Endelighed
midt i dens tilsyneladende Uendelighed. Endeligheden tynger med hele sin Vægt, med al sin Smerte
og Jammer paa Aanderne, der have mistet de evige
Kræfter og ikke vide, hvorledes de skulle kalde dem tillive
igjen. Den eneste Redning er, at Menneskeaanden
erkjender, at den Endelighed, under hvilken den stønner,
ikke er fremmed for Guddommens Væsen, men er en
af dettes Former. Gud selv er ikke udenfor Endeligheden,
men er gaaet ind i den. Nøden og Smerten er
da kun Udtryk for Aandens Uendelighed, for dens Inkommensurabilitet
med den endelige Tilværelse, hvilken
den dog maa gaa igjennem for at udvikle sit eget Væsen.
Den endelige Aand ser sig nu ikke længer som
isoleret, Selvbevidstheden forenes atter med den uendelige
Substants. Dette er Gudmenneskets Ide, Ideen om
den guddommelige og den menneskelige Naturs Enhed,
der opgaaer for Bevidstheden.

Denne Konstruktion, hvori det Methaphysiske og % PRINTED AS IS (Methaphysiske: the image, O014)
det Historiske smelte sammen, behøver endnu et vigtigt
psychologisk Supplement for at føre til Maalet\footnote[1]{Den metaphysiske Konstruktion udvikles særlig i Religionsphilosophien og viser tilbage til Logiken, den historiske især i Phænomenologien, Philosophiens Historie og Historiens Philosophi. --- Allerede hos Schelling forekomme iøvrig Grundtrækkene til disse Konstruktioner. („Vorlesungen über das akademische Studium“. 8de og 9de Forelæsning).}. Ikke
blot spekulativt, for den tænkende Bevidsthed, maa Ideen
om det Guddommeliges og det Menneskeliges Enhed
aabenbare sig; den maa ogsaa træde frem for den umiddel%
% --- p. 31 ---
\setcounter{footnote}{0}%
\opage{31}bare, den folkelige Bevidsthed. Fra dette Synspunkt
viser hint Verdensaandens Fremskridt sig i Phænomenet
derved, at det bliver Verdens Tro, at Guddommen er
bleven til som en Selvbevidsthed, som et enkelt Menneske,
at denne Gudommelighed ses, føles, høres af den % PRINTED AS IS (Gudommelighed: for Guddommelighed)
troende Bevidsthed. Hvad der i sig selv er Menneskeaandens
Eiendom, træder da frem som et Ydre og Hinsides,
der staaer overfor den. Vi komme altsaa her tilbage til den
allerede tidligere berørte hegelske Lære, at Religionens
og Philosophiens Indhold er det samme, kun Formen
forskjellig. Enheden af den guddommelige og den
menneskelige Natur fremtræder spekulativt betragtet som
Ide, religiøst betragtet som den historiske Christus.
Denne Formforskjel skal ikke antaste Indholdet, uagtet
Hegel fastholder, at Begrebets Form er den høiere, til
hvilken der maa skrides frem fra den historiske eller
Forestillingsformen.

I denne spekulative Christologi saa en stor Del af
Hegels Disciple en Forsoning af Tro og Viden, og Theologien
feirede med Begeistring „den nye Æra“.

Da udkom i Aaret 1835 Strauss's „Leben Jesu, kritisch
bearbeitet“. Aaret 1835 var, som man med Rette
har sagt, i Theologien, hvad 1848 var i Politiken. ---
Strauss kom her til det Resultat, at den evangeliske
Historie er en mythisk Historie, der har udviklet sig i
den ældste Menigheds Bevidsthed som Udtryk for den store
Tanke, der var opgaaet for den om Enheden af det Guddommelige
og det Menneskelige. Ved Begrebet Mythe
hæves det Uklare og Tvetydige i den hegelske Christologi;
det afgiver det nødvendige Mellemled til at forklare,
hvorledes Ideen for Bevidstheden kan antage et
enkelt Individs Skikkelse. Denne Skikkelse er kun en
forsvindende Form, da den almene Ide ikke absolut
kan aabenbare sig i det individuelle Phænomen, men
behøver den hele Slægt til sit Medium. Alle Prædikater, som Kirkelæren tillægger Christus, maa derfor
% --- p. 32 ---
\setcounter{footnote}{0}%
\opage{32}overføres paa hele Menneskeheden; de modsige ogsaa
hverandre, hvis de virkelig skulle tilhøre et enkelt Individ.
Men Menneskeheden er den menneskeblevne Gud,
Søn af en synlig Moder og en usynlig Fader, af Naturen
og Aanden; den er Undergjøreren, forsaavidt Aanden
i Historiens Løb stedse fuldstændigere bemægtiger sig
Naturen; den er den Døende, Opstandne og Himmelfarne,
forsaavidt der af Ophævelsen af dens umiddelbare,
naturlige Tilstand stedse fremgaaer et høiere aandeligt Liv.

Naar den spekulative Christologi deducerede det
positive Dogme, da gik denne Deduktion egentlig kun
ud paa den psychologiske og historiske Nødvendighed af
den Illusion, hvori det Almene fremtræder som et individuelt
Enkelt. Hegels orthodoxe Disciple opfattede
naturligvis Sagen anderledes. Göschel fremhævede mod
Strauss, at naar man kun lader Ideen have Virkelighed
i den hele Slægt, synker man tilbage til Kants Standpunkt,
hvor Ideen kun har subjektiv Gyldighed: er Ideen
nemlig intetsteds fuldstændig virkeliggjort, saa existerer
den kun i den subjektive Sammenfatten af dens over
hele Slægten spredte Momenter, og den hegelske Grundtanke
om Tænkens og Værens Enhed er ophævet. ---
Strauss svarer hertil (Streitschriften. 3die Hefte. p. 99 f.),
at denne Indvending ogsaa kunde rettes mod Antagelsen
af Ideens Virkeliggjørelse i et enkelt Individ: thi da
denne Virkeliggjørelse ikke kan være absolut i noget
enkelt Moment af hans Liv, saa vilde den kun være til
for den tænkende og erindrende Sammenfatten af alle
Livsmomenterne. Det er Göschel selv, der falder tilbage
til det kantiske Standpunkt, da han ikke tiltroer
Ideen Kraft til at gjennemtrænge Mangfoldigheden uden
at tabe sin Enhed. Göschel fordrer Almenbegrebet
virkeliggjort udenfor Indbegrebet af de Individer, der
høre ind derunder, ligesom den Scholastiker, der forlangte
Frugt at spise, men Frugt \emph{overhovedet}, hverken
Æbler, Pærer eller nogen bestemt Frugt.

% --- p. 33 ---
\setcounter{footnote}{0}%
\opage{33}Man angreb ikke Strauss blot med dialektiske Vaaben.
Hans Bog havde vakt Røre i den theologiske Leir
og en forvirret Skare trængte sig frem for at svare eller
idetmindste udskjelde Forargeren. Orthodoxiens bekjendte
Forkæmper Hengstenberg med sin „Evangelische Kirchenzeitung“
blev ikke den sidste. Strauss har efter Hengstenberg
„et Hjerte som Leviathan, saa haardt som den
underste Møllesten“, og hører til dem, der ere begeistrede
for Afgrundens Aand, for det store Dyr, hvem
det blev givet at tale store Ting og Bespottelser o. s. v.
Herpaa svarede Strauss, der iøvrig satte mere Pris paa
Hengstenbergs orthodoxe Zelotisme end paa den sædvanlige
Apologetiks famlende og halve Fremgangsmaade,
med en værdig og i al sin Ro begeistret Udtalelse, hvor
han hævder den frie Videnskabs Ret og afviser alle
uvedkommende, fra det Moralske hentede Indvendinger
(Streitschriften. 3die Hefte. p. 22 ff.).

Vi maa særlig fremhæve, at Strauss ikke blot forsvarer
den philosophiske Tænknings, men ogsaa den historiske
Forskens Frihed. Ogsaa herved betegner Strauss
et Fremskridt i Forhold til Hegel. Denne yndede i det
Hele ikke historisk Kritik, ligesaalidt som Goethe. Begge
frygtede, at Oldtidens store Skikkelser skulde miste deres
Glands, naar Kritiken tog fat paa dem med sine pietetløse
Hænder. Hegel saa i Kritiken et Tilbagefald til
Kants Subjektivisme og fordrede, at man skulde fordybe
sig i det historiske Indhold uden at bryde sig om de
kritiske Vanskeligheder. Han haaner Schleiermachers
kritiske Undersøgelser om Platons Dialoger og Niebuhrs
om den ældste romerske Historie. Dette hænger
nøie sammen med Hegels hele Stilling overfor Virkeligheden.
Han viser den paa engang for megen og for
liden Respekt. Den sande Ærbødighed for Virkeligheden
bestaaer i en alvorlig Prøvelse af, hvor vidt Forestillingen
om det Virkelige ogsaa svarer til det Virkelige
i sig selv. Ligesaalidt som Naturforskeren anerkjen%
% --- p. 34 ---
\setcounter{footnote}{0}%
\opage{34}der det umiddelbare Sandsebillede for den virkelige Natur,
ligesaa lidt kan den historiske Forsker standse ved det,
den umiddelbare Tro fastholder som historisk Sandhed.
Dette Princip, der udspringer af selve den videnskabelige
Forskens Væsen, har Strauss først anvendt paa et Omraade,
der tidligere holdtes lukket for det. Hvor forvirret
en Opfattelse man endnu havde af Kritikens Betydning,
kan ses af den Fordring, der blev stillet af en
ivrig Hegelianer, der dengang hørte til yderste Høire,
men senere gik over til yderste Venstre, Bruno Bauer.
Kritiken skal efter hans Mening ikke være blot skeptisk
og negativ, men skal gjennem Negationen gjenopreise
sin Gjenstand i dens hele Omfang og ikke lade Noget
tilbage, som ikke positivt retfærdiggjøres. Først saa er
den sand, positiv Kritik. Men --- svarer Strauss --- Kritikeren
maa dog idetmindste forudsætte Muligheden af,
at der i hans Gjenstand er Noget, som ikke kan retfærdiggjøres,
eftersom han jo ellers intet Motiv havde
til at stille sig paa det kritiske Standpunkt. Hin Mulighed
kan kun holdes aldeles borte, naar man gaaer
ud fra, at Gjenstandens absolut guddommelige Karakter
gjør enhver Forvanskning eller Tilsætning umulig; men
dette er i sig selv en rent vilkaarlig Antagelse.

Fra et andet Synspunkt derimod har man med Rette
kunnet bebreide Strauss, at hans Kritik ikke var positiv.
Han standser Undersøgelsen ved de evangeliske
Skrifters Indhold og mente at være ved Enden, naar
han havde vist, hvilken Tanke der maatte ligge til
Grund for den mythiske Indklædning i de enkelte Fortællinger.
Det næste Skridt var at undersøge selve Skrifternes
Oprindelse, deres Sammenhæng med den hele
Tidsalder og med de aandelige Bevægelser, der i den
krydsede hverandre. Først herved, ved at indordnes i
den hele historiske Sammenhæng finder hint Indhold
sin fulde Forklaring. Paa denne Vei slog Strauss' Lærer,
den berømte Stifter af Tübingerskolen Ferd. Chr. Baur
(† 1861) ind. Han staaer paa dette Omraade ved Siden
% --- p. 35 ---
\setcounter{footnote}{0}%
\opage{35}af F. A. Wolf og Niebuhr paa den klassiske Philologis
Omraade. Vel har man ogsaa hos Tübingerne, og sikkert
ikke uden Grund, fundet ukritiske Eftervirkninger
af Hegelianismen. Men det beviser kun, at et i sig selv
rigtigt Princip kan faa en ensidig Anvendelse paa Grund
af de Forhold, hvorunder det først fremtraadte; det beviser
Intet imod Principets Sandhed. ---

Sit almindelige Standpunkt betegner Strauss som
den absolute Immanents. Hans hele Verdensanskuelse
fremtræder især i hans Skrift: „Den christelige Troeslære
i dens historiske Udvikling og i dens Kamp med
den moderne Videnskab“ (1840—41). Det er her navnlig
Opfattelsen af Gudsideen, det kommer an paa. Hegels
Lære om det Absolutes Subjektivitet var af hans orthodoxe
Disciple forstaaet saaledes, at deri ligefrem var
begrundet det gamle Dogma om Guds Personlighed og
Transscendents. Paa dette Punkt, ligesom i Udødelighedslæren
og i Christologien, ansaa de derfor deres
Mester for fuldt overensstemmende med Kirkelæren.
Men Strauss viser, at Hegel ved sin Lære om den absolute
Subjektivitet kun har villet indskærpe, at det Absolute
ikke kan opfattes som død, hvilende Substants,
men som det uendelige Liv, der er sit eget Produkt og
udvikler sig gjennem en Uendelighed af Forskjelle og
Modsætninger. Spinoza begreb ikke det Absolute som
den evige Selvbestemmelse; han viste, hvorledes Alt
strømmede tilbage i Substantsen, men ikke hvorledes
det strømmer ud af denne igjen, --- kun Venesystemet,
ikke Arteriesystemet, ikke Livets hele uendelige Kredsløb,
hvorved Substantsen gjennem sit Andet slutter sig
sammen med sig selv og saaledes bliver Person og Aand.
Denne Opfattelse undgaaer ifølge Strauss to farlige
Skjær. Den stiller ikke Guddommen paa udvortes Vis
fuldfærdig overfor Verden, som den vulgære Theisme;
den binder heller ikke Guds Væsen saaledes til Verdensudviklingen,
at han skulde behøve den endelige Til%
% --- p. 36 ---
\setcounter{footnote}{0}%
\opage{36}værelse og den menneskelige Aand for at være sig selv.
Ikke ved Mennesket, der jo ikke er ved sig selv, men
gjennem og i Mennesket, som det selv har sat, kommer
det Absolute til Selvbevidsthed og Personlighed.
Guds Personlighed maa ikke tænkes som Enkeltpersonlighed,
men som Alpersonlighed; istedetfor at personificere
det Absolute skal man begribe det som det, der i det
Uendelige personificerer sig selv.

De enkelte Personifikationer af det Absolute ere
kun forbigaaende Momenter. Hvis de fæstnede sig, vilde
det være et Tegn paa den absolute Aands Afmagt. Det
Substantielle er ikke i Individerne, men hinsides dem i
den absolute Aand, til hvilken de forholde sig som vexlende
Accidentser. Dette gjælder med Hensyn til hele
den i Tiden sig udviklende Tilværelse. Tiden er en
endelig Forestillingsform, der behøver ydre Støttepunkter
for Anskuelsen for at kunne holde sig. Derved fæstnes
de endelige Forskjelle, som i sig selv ingen Realitet
besidde. Enheden i Tidsmomenternes Forskjel er Evigheden;
ved at tænke Tiden berøver man den Begyndelse
og Ende, dens forestillede Støttepunkter, og den synker
af sig selv sammen i Evighedens Forskjelsløshed som
dens Grund, af hvilken den ligesaavel fremgaaer i det
Uendelige. Evigheden og Tiden forholde sig som Substantsen
og dens Accidentser.

Det fremgaaer heraf, at Strauss's Pantheisme ikke
kan undgaa tilsidst at falde tilbage i den spinozistiske
Substantslære. Det individuelle Liv er et blot Phænomen,
bag hvilket Tanken søger det Sande, som er Substantsens
forskjelsløse Evighed. Det er kun Forestillingen,
der endnu kan tilskrive Endeligheden virkelig Realitet,
og Forestillingen er jo ifølge Strauss en uadækvat
Erkjendelsesform.

Men der viser sig en endnu misligere Omstændighed.
Der forbliver stedse en Modsætning tilbage mellem
Substantsen og den individuelle Bevidsthed, i hvilken
% --- p. 37 ---
\setcounter{footnote}{0}%
\opage{37}den kommer til sig selv. Substantsen skal jo endog
ligge \emph{hinsides} Enkeltvæsenerne --- og dog er Strauss's
sidste Ord i hans Dogmatik, at det Hinsidige \emph{under}
alle Skikkelser er den Fjende, som den spekulative
Kritik skal bekæmpe. Immanentsens Standpunkt er
altsaa ikke gjennemført. Der er endnu et Skridt at gjøre
--- men det gaaer i en ganske anden Retning end til
Spinoza. Selvbevidstheden maa ophæve ethvert transscendent
Forhold for ganske at blive sig selv klar og gjennemsigtig,
den maa fjerne enhver fremmed Magt som
Skyggespil. Ad denne Vei føres vi ogsaa ved den „kritiske
Bevidstheds“ videre Udvikling hos Bruno Bauer
og Feuerbach.

Strauss's skarpe kritiske Tanke har saaledes løsnet
Enheden mellem de Momenter, der i Hegels Lære vare
i inderlig Forbindelse. Hans hele kritiske Retning bragte
ham i Modsætning til den dialektiske Methode og den
Alt medierende Spekulation, en Modsætning, der kun
kunde skærpes ved den Maade, hvorpaa de orthodoxe
Hegelianere misbrugte Læren om den dialektiske „Modsigelse“
til at ophæve Forskjellen mellem „Widersinn“
og „Widerspruch“ (Streitschriften. 3die Hefte. p. 23. 113).

Modsætningen mellem Substantsen og Individerne
viste sig at staa i nøie Sammenhæng med den psychologiske
Modsætning mellem Tanke og Forestilling. Betragtningen
af denne Modsætning vil kaste et afgjørende
Lys paa den straussiske og derved paa den hele spekulative
Religionsphilosophi. --- Strauss sætter det som
Christendommens Fortjeneste at have bragt Enheden af
det Guddommelige og Menneskelige til Menneskeslægtens
Bevidsthed. Denne Ide er derfor det Blivende i
Christendommen. Den gamle Verden var væsentlig dualistisk;
Folkene følte Guddommen som en fjern og fremmed
Magt, der kun i den ældste Tid, i den mythiske
Guldalder venlig færdedes blandt Menneskene. Overfor
denne Dualisme traadte Christendommen med sin Lære,
% --- p. 38 ---
\setcounter{footnote}{0}%
\opage{38}at Gud var bleven Menneske. Saaledes vare begge Verdener
forenede. Men Foreningen var kun fuldbyrdet
paa et enkelt Punkt, som en Undtagelse, et Under. Den
mirakuløse, altsaa dualistiske Betragtningsmaade forblev
herskende. Christendommen er altsaa det Guddommeliges
Dennesidigvorden under Forudsætning af dets Hinsidighed,
Begyndelsen af det moderne Immanentsens System,
men endnu paa Transscendentens gamle Jordbund.
Det kommer nu an paa, om man vil anse det for det
Eiendommelige ved Christendommen, at det Guddommeliges
Immanents indskrænker sig til et enkelt Punkt, der saa
ved en absolut, uoverstigelig Kløft adskilles fra den øvrige
Menneskehed, eller om Christendommens sande Væsen ikke
skulde kunne bevares, naar man udviklede den Immanentsens
Spire, den indeholder, ved Hjælp af den humane
Udviklings Resultater. Fastholder man den snevrere,
strengere og derfor væsentlig dualistiske Opfattelse af
Christendommen, saa kommer man i en uforsonlig Strid
med den humane Udvikling. Gaaer man derimod ind
paa en mere forsonlig Opfattelse, saa kan man sige, at
Christendommen stedse udvikler sig videre fra sit oprindelige
Udgangspunkt. Vi ville da have en Analogi
i den politiske Friheds Historie. De græske Stater anerkjendte
Friheden, men kun som et vist Antal Borgeres
Eiendom i Modsætning til Slaver og Metoiker. De vilde
altsaa kun vide af Friheden at sige paa dette Punkt.
Men den historiske Udvikling har ført til at udstrække
den til Alle, ifølge den vaagnende Bevidsthed om Alles
Delagtighed i den samme menneskelige Natur. (Zwei
friedliche Blätter. 1839). --- Med denne Betragtningsmaade
stemmer Strauss's Behandling af Dogmerne. Han
griber hvert enkelt Dogme i dets Udspring og følger det
gjennem den hele historiske Udviklingsgang. Gjennem
denne Proces sønderbryde Dogmerne deres endelige Form
og lutres til Begrebets Renhed. Forestillingsformen,
hvori Religionen har indklædt det uendelige Indhold,
% --- p. 39 ---
\setcounter{footnote}{0}%
\opage{39}er ikke blot en ydre Skal, men foraarsager en gjennemgaaende
Dualisme i den hele Anskuelse. Hegel har
derfor Uret i at anse denne Formforskjel for uvæsentlig.
Fra det Ufuldkomne maa der skrides frem til det Fuldkomne,
fra Tro til Viden.

I den philosophiske Verdensanskuelse, der afløser
den positive Kirkelære, finder Bevidstheden den samme
Tilfredsstillelse som den Troende i sine Dogmer, og de
Vanskeligheder og Modsigelser, som disse medføre, falde
hist bort. Men vil da denne Frigjørelse kunne blive
Alle tildel eller vil den stedse være indskrænket til nogle
Enkelte? Dette Spørgsmaal lader Strauss staa ubesvaret;
han ser i Modsætningen mellem de Vidende og de
Ikke-Philosopherende en Kløft mellem to Klasser af det
menneskelige Samfund, som maaske aldrig vil fyldes.
Han fremhæver kun, at det ikke kan nytte at tildække
denne Kløft, da den engang er der. „Saa lade da den
Troende den Vidende, og denne hin roligen vandre sin
Vei; vi lade dem beholde deres Tro, saa maa de ogsaa
lade os beholde vor Philosophi; og hvis det skulde lykkes
de overvættes Fromme at udelukke os fra deres
Kirke, saa ville vi agte det for en Vinding; af falske
Foreningsforsøg er der sket nok, kun Skærpelse af Modsætningerne
kan føre videre.“ (Troeslære. Dansk Overs.
I. p. 309).

Religion er altsaa kun en ufuldkommen Form for
Philosophi. Det Religiøse faaer egentlig kun en negativ
Bestemmelse; det er det, der skal ophæves, oversættes
paa et andet Sprog, et Barns Lallen i Forhold til den
Voxnes fornuftige Tale. Dette viser tilbage til en psychologisk
Forudsætning, hvorfra Strauss gaaer ud ligesaavel
som hele den spekulative Religionsphilosophi,
hvis sidste Konsekvents, han drager, --- den Forudsætning
nemlig, at Religionen har sit Udspring fra Erkjendelsesdriften
og altsaa væsentlig er et intellektuelt Phænomen.
Er denne Forudsætning rigtig, da skal der vanske%
% --- p. 40 ---
\setcounter{footnote}{0}%
\opage{40}lig kunne reises nogen afgjørende Indsigelse mod Strauss;
saa falder al Religion bort med den positive Religion.
Men er denne Forudsætning rigtig, da bliver det ogsaa
vanskeligt at forstaa den Kløft, der skal vedblive at bestaa
mellem de Vidende og de Troende. Hvorledes kan
Strauss tvivle om denne Kløfts Ophør, naar de to Elementer,
den adskiller, ere væsentlig beslægtede, ere Momenter
i samme Udvikling? Kan man antage, at der
skulde være sat en uoverskuelig Skranke mellem to
Erkjendelsestrin, og maa man ikke snarere spørge, om
den religiøse Bevidstheds Vægring ved at opløse sig i
spekulativ Viden, ikke kunde være grundet i, at den
ikke finder sit hele Væsen udtrykt i Erkjendelse, men
savner væsentlige Elementer, Elementer, der netop udgjøre
det i Religionen Religiøse?

Altsaa paa den ene Side den Lethed, hvormed Strauss
kan paavise Dogmets Opløsning i spekulativt Begreb,
paa den anden Side den Modstand, han indrømmer, at
den religiøse Bevidsthed gjør imod at foretage denne
Overgang, --- denne Modsætning vækker Tvivl om, at
den psychologiske Forudsætning, hvorfra der gaaes ud,
er rigtig. Paa en egentlig Prøvelse af den har Strauss
ikke indladt sig. Han har indledet en ny Epoche med
Hensyn til Religionsphilosophiens historiske Grundlag,
men han har overset det psychologiske Grundlag, der
er endnu vigtigere, da det historiske uden det ikke er
forstaaeligt. Ogsaa paa dette Punkt henvises vi til
Feuerbach, der gjennembrød de Skranker, der endnu
holdt Strauss fangen.

3. \emph{Bruno Bauer} (f. 1809) er et levende Vidnesbyrd
om Tvetydigheden i det hegelske Princip om Fornuftens
og Virkelighedens Enhed, et Princip, der ligesaavel
kan være reaktionært som revolutionært, ligesaavel
konservativt som radikalt. Saaledes kunde Bruno Bauer
fra det yderste Høire, den videstgaaende Orthodoxi pludselig
kaste sig over i det modsatte Extrem. I sin Kritik
% --- p. 41 ---
\setcounter{footnote}{0}%
\opage{41}af Strauss's „Leben Jesu“ bebreidede han, som allerede
omtalt, Strauss, at han ikke fuldstændig gjennemførte
Forstaaelsen af det Objektive; Kritiken standser paa et
vist Punkt og godtgjøres ikke som „Selvbevidsthedens
uovervindelige Magt“, der ikke lader Noget blive staaende
som uigjennemtrængelig Gjenstand, men negerer
alt Objektivt for at erkjende sig selv deri. Efter Bauers
Mening retfærdiggjør den sande, „positive“ Kritik Gjenstanden
i dens Helhed, idet den begriber den. --- Men
faa Aar efter angreb han Strauss fra en ganske anden Side.

Strauss er --- siger Bruno Bauer --- i sin Kritik
bleven staaende paa et væsentlig theologisk Standpunkt.
Han troer sin Opgave naaet, naar han har paavist de
evangeliske Fortællingers mythiske Karakter, og mener
saa, at disse Myther have dannet sig i den ældste Menighed
under Traditionens Indflydelse. Men han viser her
tilbage til noget Mysteriøst. Ethvert individuelt Værk maa
dog være fremgaaet af en individuel Tænken og Virken
og kan ikke voxe ud af en Almenbevidsthed,  der med
mystisk Magt beaander de Enkelte, uden at Nogen af disse
for sig skulde kunne kaldes Fortællingens Ophavsmand.
Nei, Kritiken maa gjennemføres saaledes, at man faaer en
enkelt, bestemt Forfatter at holde sig til. Evangelisterne
ere at sammenligne med Homer og Hesiodos, der, som
Herodot siger, have dannet Grækerne deres Guder.

Ligesom Strauss i sin historiske Kritik lader den
mystiske Magt overvælde den virkelige Selvbevidsthed,
saaledes gaaer hans hele Anskuelse ud paa at hævde
det Absolute som objektiv Substants. Derved er det, at
han væsentlig staaer paa et religiøst eller theologisk
Standpunkt. Det religiøse Standpunkt er nemlig det,
hvor Selvbevidstheden endnu har sit Indhold adskilt fra
sig, hvorved Indholdet bliver endeliggjort. Paa sit Høidepunkt
er den religiøse Bevidsthed absolut fremmed for
sig selv. I Naturreligionen betragtes Guderne endnu
som særlige, enkelte Væsener, og ere som saadanne be%
% --- p. 42 ---
\setcounter{footnote}{0}%
\opage{42}slægtede med Menneskene; derfor gjennemføres Fremmedheden
ikke, men Stat, Familieliv og Kunst faa umiddelbar
religiøs Betydning. „De Lænker, den menneskelige
Aand bar i disse Religioners Tjeneste, vare omvundne med
Blomster; herlig og festlig smykket som et Offerdyr bragte
Mennesket sig selv som Offer for sine religiøse Magter;
Lænkerne selv skuffede ham med Hensyn til hans Tjenestes
Haardhed.“ Men Blomsterne visnede, den antike
Naturverden sank sammen, og „den aandelige Abstraktions
Vampyr“ fuldbyrdede sit Værk. Kun Jeget, Selvbevidstheden
i dens hele Uendelighed blev tilbage ved
den undergaaede Verdens Ruiner, som det Dyb, hvori Alt
var sunket. Men Jeget kunde endnu ikke gjennemskue
sig selv; dets Uendelighed var endnu et fremmed Væsen
for det, en absolut, almægtig Guddom, der stod truende
overfor det. Dette er Christendommens afslørede Mysterium.
Selvbevidstheden har i de evangeliske Fortællinger
at gjøre med sig selv, men det vil sige, med den frygteligste
Parodi af sig selv, en Molok, der kræver ethvert
naturligt og virkeligt Forhold som Offer. Men da den
dog er Selvbevidsthedens Billede, saa forstaaes, hvorledes
den har kunnet fængsle Menneskeheden saalænge, paalægge
den saa skrækkelige Byrder. Menneskeslægten er
bleven opdraget i Trældom under sit eget Afbillede, for
at den des grundigere kunde vinde Friheden og befæste
sig i den\footnote[1]{Kritik der evangelischen Geschichte der Synoptiker. 3ter Band.}.

Pantheismen, der sætter en objektiv Substants som
Bevidsthedens Gjenstand, er endnu ikke naaet til denne
Frihed, er endnu indenfor Trældommen. Modsætningen
maa føres til Enden, hvor kun Atheismen er det rette
Navn. Den religiøse Bevidsthed maa fuldstændig negeres.
Det er aldeles falsk, naar Strauss mener, at Philosophien
skal træde i Religionens Sted, og at den kan
det, fordi de begge gaa ud paa det samme absolute Ind%
% --- p. 43 ---
\setcounter{footnote}{0}%
\opage{43}hold. Den religiøse Bevidsthed er vel Spiren til den
fuldkomne Selvbevidsthed, men denne Spire er, netop
fordi den kun er Spire, en falsk og fordærvelig Magt.
De, hos hvem Fornuften ikke er udviklet, have vel Anlæg
til sand Menneskelighed, men netop fordi det kun er et
Anlæg, ere de ikke Mennesker, men Umennesker. Derfor
kan der ikke tænkes paa noget Forlig med den religiøse
Bevidsthed. Sandhederne høre op at være Sandhed,
naar de omsættes i Religionens Form. Der er ikke
Sandhed i Religionen. Da den religiøse Selvbevidsthed
ikke er fri og ikke gjennemskuer sit eget Væsen, forholder
den sig endnu paa vilkaarlig og tilfældig Maade
til sin Gjenstand; det er den dunkle Drift, der sætter
Religionens Gjenstand, en Drift, der paa den ene Side
er en Angst for den mørke, uforstaaede Magt, som tynger
paa Bevidstheden, paa den anden Side er det lidenskabelige
Ønske, der vil se sine Phantasier realiserede\footnote[1]{Bruno Bauer i „Deutsche Jahrbücher.“ 1843. p. 82---95.}.

Naar da Kritiken er gjennemført til det Punkt, at
den mørke Magt er aldeles tilintetgjort, enhver Modstand,
ethvert Objekt fjernet, hvad bliver saa tilbage? Bruno
Bauer svarer: den rene Kritik. Intet har Sandhed uden
den kritiske Bevidsthed selv. Enhver Gjenstand sættes
kun for atter at negeres; saafremt den anerkjendes, hører
den op at være sand. Den, der søger Tilfredsstillelse
ved Noget, den, der vil begeistre sig for Noget, maa gaa
tilbage til Religionens dæmoniske Magt for at lade sig
berøve sin Frihed. Udtrykket Atheisme passer ikke mere,
da det endnu indeslutter et Hensyn til den Gjenstand,
der negeres, og derfor ikke betegner den absolut frie
Selvbevidsthed. Stat, Familie, Kunst og Videnskab opløses
i denne titaniske Bevægelse, i hvilken Broderen
Edgar Bauer særlig paatog sig Kritiken af Staten\footnote[2]{Se om ham „Nyt dansk Maanedsskrift“. 1871---72. p. 223.}.

Bevægelsen er her med rivende Fart ført til sin
% --- p. 44 ---
\setcounter{footnote}{0}%
\opage{44}yderste Grændse. Det er en Begrebets Fanatisme, her
optræder, først med Fordring om absolut Forstaaen, om
at gjennemskue Tilværelsens Dybder, saa ikke en eneste
Dunkelhed bliver tilbage, --- og da dette mislykkes, med
ligesaa absolut Hengivelse til den modsatte Stræben, at
lade Alt blive tilintetgjort af den Selvbevidsthed, i hvilken
det ikke vilde gaa op. Den hegelske Panlogisme er
her gjennemført saaledes, at den er vendt tilbage til sit
Udgangspunkt, det usynlige Punkt, hvor Tænken og
Væren er Et, fordi kun Tænken er. Dette usynlige
Punkt er det, Bauer betegner som den rene Kritik; men
istedetfor at Hegel gik ud fra dette Punkt for at erobre
den hele Verden, er Bruno Bauer naaet dertil ved
at lade Alt synke sammen i Intet. Han ender i Intet,
fordi han begynder med Intet.

4. \emph{Arnold Ruge} (f. 1802). --- Ligesom Hegel i
Religionsphilosophien havde anvendt sin Grundsætning
om Enheden af Fornuft og Virkelighed saaledes, at den
betød Dogmets Forsoning med den spekulative Tanke,
idet, som Strauss udtrykte sig, Dogmerne gjennem den
dialektiske Skjærsild kom ind i Ideens Himmel med Hud
og Haar, saaledes anvendte han den samme Sætning i Retsphilosophien
paa en saadan Maade, at Restaurationsperiodens
Stat, dette Udtryk for brudte Løfter, derved blev
godtgjort som „fornuftig“. Ligesaalidt her som hist anvendte
Hegel nogen egentlig Kritik af det Foreliggende. Da han var
overbevist om, at Ideen ikke er saa afmægtig, at den kun
skulde være et subjektivt Ideal, der ikke har vundet sig
nogen Virkelighedens Form, men at vi i Staten maa se
den Verden, den sædelige Ide har skabt sig, gjaldt det
for ham kun om umiddelbart at paavise Fornuften i de
historisk udviklede Statsformer.

Dog oplevede Hegel endnu, Aaret før sin Død, en
Virkelighed, han ikke fandt fornuftig, nemlig Julirevolutionen,
og hans sidste Dage forbitredes af en Strid, han
i den Anledning kom i med en af sine Disciple, Rets%
% --- p. 45 ---
\setcounter{footnote}{0}%
\opage{45}philosophen Gans, der vilde uddrage liberale Konsekventser
af hans Lære\footnote[1]{Gans begyndte sine „Vorlesungen über die Geschichte der letzten fünfzig Jahre“ med at bebreide den tydske Historieskrivning, at den udelukkende vendte sig mod Fortiden og tilsidesatte Samtidens Historie som „Tagesgeschihte“. % PRINTED AS IS (Tagesgeschihte)
}. Men det var først en halv Snes
Aar efter, at den spekulative Betragtning af Historien
gjennemgik en lignende Krise som den, der i religionsphilosophisk
Henseende var frembrudt med Strauss og
Bruno Bauer.

Arnold Ruge udgav fra Aaret 1838 af et Tidsskrift
ved Navn „Hallische Jahrbücher“, der oprindelig hørte
til, hvad man kaldte det hegelske Centrum. Det optraadte
med Ærefrygt for det Positive og ansaa med
Hegel Forskjellen mellem Religion og Philosophi for en
blot Formforskjel. Men Aar for Aar blev det mere radikalt.
Først Strauss's, siden Bruno Bauers og Feuerbachs
religionsphilosophiske Anskuelser bleve de herskende
i Aarbøgerne. Ruge skiller sig fra den hegelske
Skoles Centrum og Høire, idet han bebreider dem deres
Stilling som uvirksomme Tilskuere ved den verdenshistoriske
Proces. „Der er kun én Atheisme, Tvivlen
om Historiens Aand.“ Hermed fulgte en stedse skarpere
Kritik af de preussiske Statsforhold, som tilsidst
nødte Ruge til at drage til Dresden (1841), hvorfra han
nu udgav Tidsskriftet som „Deutsche Jahrbücher“.

Disse aabnes med et Forord, hvori han lægger Hegelianerne
til Last, at de have et færdigt System, en
forudsat Dogmatik og uden Videre finde Ideen i den
politiske status qvo. Mod disse tre Magter, den philosophiske,
den theologiske og den politiske Orthodoxi, vil
han aabne en Frihedskrig. Og spørges der, under hvilket
Banner denne Kamp skal føres, da svarer Ruge:
„Der gives kun ét Himmeltegn, som oplyser og varmer
den aandelige Verden, og denne Stjernes Navn er Phi%
% --- p. 46 ---
\setcounter{footnote}{0}%
\opage{46}losophien.“ Men saa kan Philosophien ikke opfattes paa
Hegels Vis, der mente, at dens Kald kun var at begribe
den forudgaaende Virkelighed. Hegel, siger Ruge, anerkjender
dog herved selv, at Tanken er et nyt og høiere
Trin, idet den dømmer Virkeligheden. Og idet Tanken
saaledes er det Sidste, er den, hvad Hegel oversaa, ogsaa
det Nye, en Foregriben af den tilkommende Virkelighed.
I Logiken lærer Hegel, at Negationen af det Foregaaende
fører over til et nyt Stadium; men i Verdenshistorien
vil han ikke indrømme dette. Og dog beroer Bevægelsen
i Verdenshistorien paa, at den menneskelige Bevidsthed
inderliggjør sig sit vundne Indhold, hvilket netop
sker gjennem den philosophiske Kritik. Herved skrides
der frem til en ny Opgave, der sætter et „skal“ istedetfor
et „er“, og saaledes fører den rette Gjennemførelse
af Hegels Dialektik til ogsaa at indse Sandheden af
Kants Lære om Villiens Autonomi og Fichtes Lære om
det sig selv bestemmende Jeg. Philosophen er en Fremtidens
Apostel, der theoretisk fuldbyrder den Aandsbevægelse,
der i Statslivet skal fuldbyrdes i Praxis. Det
er ingen Bebreidelse mod Theorien, at den er revolutionær:
„nur umwälzende Gedanken sind Gedanken.“

Kritiken søger altsaa den foregaaende Virkeligheds
Begreb for af det at udlede den efterfølgende Virkeligheds
Begreb. Men Hegel har sammenblandet det metaphysiske
og det historiske Begreb, og derved miskjendt
Kritikens Betydning. Det er det metaphysiske Begreb,
der driver Kritiken frem, idet det viser sig
ikke at have sit fuldstændige Udtryk i det historiske.
Statens Ide er saaledes ikke virkeliggjort i en enkelt
bestemt Stat. Hegel derimod hævede umiddelbart det
historisk Existerende til metaphysisk Værdighed. Ved
denne umiddelbare Overgang overses Modsætningen, og
dermed Kritiken. Kun gjennem denne gaaer paa den
anden Side Veien igjen tilbage fra Begrebet til Virkeligheden.
„Den sande Forbindelse af Begrebet og Virkelig%
% --- p. 47 ---
\setcounter{footnote}{0}%
\opage{47}heden er ikke Existentsens Apotheose til Begreb, men
Begrebets Inkarnation til Existents.“

Spørges der, hvorledes denne Begrebets Inkarnation
gaaer for sig, da bliver Ruges Bestemmelse af Religionens
Væsen at fremhæve. Den hegelske Philosophi og
den tydske Aand i det Hele har, siger han, skjult Theoriens
praktiske Betydning. Theorien bliver praktisk, naar
den gaaer over til Begeistring og Beslutning, o: naar
den bliver Religion. Thi Religion er i sig selv ikke
Andet end den praktiske Pathos for det Ideale, for Sandheden,
Troskab mod Ideen og opofrende Udøvelse af
den. Religionen er Praxis: den fører overalt, paa Grund
af de endelige Existentsers Modstand mod Ideen, sin
eiendommelige Kamp og Opofrelsen af jordiske Interesser
med sig. Dette er den sande Religion, det Dennesidiges
Religion, i Modsætning til de dualistiske Religioner,
der ved at gjøre Mennesket fremmed for sig selv
og for Tilværelsen, ikkun hindre ham i Frihedens Udøvelse
istedetfor at fremme den. Aldrig har Andet
været Religion end Friheden. Men nu bliver dens Form
og Indhold rigere end før. Humanitetens hele Idealverden
skal i Gemyttets Smeltedigel for deraf igjen at
fremgaa som en glødende Strøm og danne en ny Verden.

Den blotte Theori, der ikke fordyber sig i sig selv
til religiøs Pathos, kommer derimod ikke ud over den
saakaldte Liberalisme, der aldrig gaaer løs paa selve
Bevidsthedens Rod. Den troer, at det kun kommer an
paa at reformere de politiske Former, medens den forresten
lader Enhver tænke som han vil og blive salig
ved sin Tro. Denne Indifferentisme viser Liberalismens
Overfladiskhed. Den ser ikke, at hvad det gjælder om
er en Bevidsthedens Reform, og kommer derfor aldrig
længer end til de blotte Ønsker om Frihed. Med Bevidsthedens
Reform vilde Verdens Reform være given.
Naar denne Reform iværksættes, bliver Liberalismen til
Demokratisme; enhver fremmed Magt gaaer op i den
% --- p. 48 ---
\setcounter{footnote}{0}%
\opage{48}frie Selvbevidsthed, og i Staten forsvinder enhver Modsætning,
idet Kirken gaaer op i Skolen for gjennem 
Folkeopdragelsen at absorbere al Pøbel, og idet Krigs- og Retsvæsen umiddelbart varetages af det sig selv
organiserende Folk. ---

Paa dette Punkt i sin Udvikling blev Tidsskriftet 
forbudt af den sachsiske Regjering (1843). Ruge høstede
senere som Medlem af Frankfurterparlamentet 1848 den
Erfaring, at der endnu var langt tilbage for Virkeliggjørelsen
af Bevidsthedens Reform. Den Fordring, Ruge
havde stillet, gjør ham til det, han vilde, Philosophen
skulde være, en Fremtidens Prophet. Thi hin Fordring
om Bevidsthedens Reform er i sin Almindelighed den
hele Histories Opgave, og det er en storartet Idealisme
umiddelbart at gjøre den gjældende uden at have positive
Udgangspunkter for dens Virkeliggjørelse. Derfor
er Ruge ikke praktisk Politiker, men en ædel Discipel
af Platons Stat\footnote[1]{En Del af Afsnittet om Ruge har, med Tilføjelser og Forkortelser, været trykt i „Nyt dansk Maanedsskrift“. 1871---1872.}% the note runs over from p. 48 to p. 49 (joined by draft.py)
.

I philosophisk Henseende er det det Interessante
ved Ruge, at han har hævdet Philosophiens praktiske
Betydning. Han har hævdet Nødvendigheden af, at Ideen
ikke blot blev Gjenstand for Spekulation, men tilegnedes
af den levende personlige Bevidsthed. Derved har han
overskredet den hegelske Psychologis og Religionsphilosophis
Grændser, indenfor hvilke Gemyt og Phantasi
ikke kunde faa en saadan positiv Betydning. Ja, han
forlader ved sin Bestemmelse af Tankens Forhold til
den historiske Virkelighed selve den hegelske Grundtanke
om Enhed af Tænken og Væren. Er Kritiken en
nødvendig Mellembestemmelse mellem Begrebet og Virkeligheden,
saa er hin Enhed ikke umiddelbar og absolut.
Gjennem den historiske og religionsphilosophiske Kritik
øine vi saaledes nye metaphysiske Konsekventser.

% --- p. 49 ---
\setcounter{footnote}{0}%
\opage{49}5. \emph{Ludwig Feuerbach} (f. 1804). --- Ethvert af
de Phænomener, vi hidtil have beskjæftiget os med, fremtraadte
strax med en eiendommelig, fast udpræget Karakter.
Og forsaavidt som vi hos Bruno Bauer bleve
Vidne til en Svingning fra det ene Extrem til det andet,
skete denne Svingning i et Nu, med urolig Hast. Anderledes
den Tænker, vi nu staa ved, og som er den egentlig
typiske Repræsentant for den spekulative Idealismes
Selvopløsningsproces, men saaledes at han tillige paaviser
de positive Elementer for en ny Udvikling. Feuerbach
gjennemløber den hele Skala fra yderste Høire til
yderste Venstre, fra den absolute Idealisme til Materialismens
Grændse, saaledes at hvert enkelt Skridt tydelig
og eftertrykkelig markeres, hvorved Overgangene blive
successive, ja for det overfladiske Blik næsten umærkelige.
Han er opfyldt af en gjærende Uro, der følger
ham fra hans første Ungdom, da han som orthodox Theolog
glædede sig ved den spekulative Forsoning af Tro
og Viden, lige til den sensualistiske Subjektivisme, i hvilken
han ender. Stedse utilfredssillet ved det Givne, det % PRINTED AS IS (utilfredssillet: for utilfredsstillet)
hidtil Opnaaede, søger han stedse nye Former og Udtryk.
Han fortæller selv, med hvor liden Lyst han i Anledning
af Udgivelsen af sine samlede Værker vendte tilbage
til Skrifter, der for ham betegnede et tilbagelagt
Stadium --- og det gjorde ethvert af hans Skrifter, saasnart
han havde lagt den sidste Haand derpaa\footnote[1]{„Jeder fertigen Schrift sagte ich für immer Adieu; jede hatte mir nur meine Fehler und Mängel zu Bewusstsein gebracht und daher nichts andres in mir zurückgelassen, als das dringende Verlangen, ihr Andenken durch eine neue Schrift auszulöschen. Und nun wurde mir auf einmal zugemuthet, meinen unzufriedenen, schriftwidrigen, unbiblischen Geist auf alle meine längst meinem Sinn entschwundenen Schriften zu richten. Welche Zumuthung! Wider den Strom des Lebens soll ich schwimmen? wider den Lauf der Natur statt vorwärts, rückwärts gehen?“ (Werke. 1ster Band. p. V---VI).}. Denne Uro,
denne Stræben efter altid fornyet Fremgang giver ham
% --- p. 50 ---
\setcounter{footnote}{0}%
\opage{50}Karakteren af en ægte Tænker, der lidenskabelig vil
kæmpe sig frem til Sandhedens Kilder; men der er noget
Tantalisk i denne hans Stræben, idet det aldrig lykkes
ham at naa et Standpunkt, der tilfredsstiller ham, og
som han derefter kan udvikle og gjennemføre fuldstændig.
Han har Ret, naar han karakteriserer sig selv som
en „utilfreds Aand“. --- Den urolige, aldrig tilfredsstillede
Higen giver hans Sprog en eiendommglig Skjønhed; % PRINTED AS IS (eiendommglig: for eiendommelig)
i en Række af Vendinger, Antitheser, Billeder og Sammenligninger
søger han at udtømme det, han ikke kan
fastholde i skarpt afsluttet Begrebsform. Heri finder
han Enheden i sin Forfattervirksomhed udtrykt. „Du
tænkte som Philosoph“, siger han til sig selv, „men du
skrev ikke som Philosoph; du forvandlede stedse Tankevæsenet,
naar du udtalte det, til et Væsen af Kjød og
Blod. Du vidste, at Philosophien som saadan, den blotte
Fornuft, den rene Tanke ikke er Noget for Mennesket,
at man kun da kan overbevise Mennesket om en Sandhed,
naar man gjør den fra et Fornuftvæsen til et Væsen,
der er Mennesket ligt, et sandseligt Væsen\dots{} Overalt
har du knyttet det Abstrakte til det Konkrete, det
Usandselige til det Sandselige, det Logiske til det Anthropologiske.“
(Werke I. p. VII.---IX.). Udviklingen bestaaer
saa i, at hvad der fra først af kun var en Form,
en Indklædning, tilsidst bliver selve Væsenet. Det sproglige
Udtryk indgaaer altsaa en stedse inderligere Enhed
med Indholdet: „Tidligere tænkte du og udtalte det
ogsaa, om ikke ordret, saa dog faktisk: det Sande maa
være nærværende, virkeligt, sandseligt, anskueligt, menneskeligt;
nu siger du konsekvent omvendt: kun det
Virkelige, Sandselige, Menneskelige er det Sande.“

Vi kunne i Feuerbachs philosophiske Udviklingsgang
skjelne mellem tre forskjellige Stadier, naar vi forbigaa,
hvad der ligger forud for hans offentlige Fremtræden
(hans orthodox-theologiske Standpunkt). Disse tre Sta%
% --- p. 51 ---
\setcounter{footnote}{0}%
\opage{51}dier karakterisere vi i Almindelighed, førend vi gaa over
til nøiere at beskjæftige os med ethvert enkelt af dem.

Det første Stadium er pantheistisk. Bevægelsen
gaaer her absolut i Retning af det Almene, det Objektive,
Ideen. Alt Personligt og Individuelt søges ophævet
som Ufuldkommenhed. Ved at forenes med Subjektiviteten
krænkes Ideens Renhed; den kommer ikke længer
til at gjælde i og for sig. Men om Ideens objektive Realitet
tvivles der ikke; Ideen er tvertimod det ene sande
Værende. Som den, der har Gyldighed i og for sig kan
den kun opfattes og gribes af den theoretiske Erkjendelse.
Theorien er Sandheden. Religionen og Theologien
gaa glip af Sandheden, fordi de stille sig paa det praktiske
Standpunkt, hvor de endelige, personlige Motiver
gjælde og ikke Sagen i og for sig. Selvkjærligheden,
Lyksalighedsdriften betegnes allerede her som Religionens
Udspring, og der anbringes derved et Korrektiv
mod den spekulative Religionsphilosophi, der kun betragtede
Religion som en ufuldkommen Form for Theori. ---
Til dette første Stadium (1830---1838) høre følgende
Skrifter: Gedanken über Tod und Unsterblichkeit; Geschichte
der neuern Philosophie von Baco bis Spinoza;
Darstellung und Kritik der Leibnitzischen Philosophie;
Pierre Bayle.

Det reale Moment, Feuerbach saaledes var bleven
opmærksom paa ved sine religionsphilosophiske Undersøgelser,
vandt forøget Betydning for ham ved det ivrige
Studium af Naturvidenskaberne, hvortil han nu hengav
sig. Han blev derved opmærksom paa den sandselige
Erfarings Nødvendighed og Myndighed og vandt en inderlig
Overbevisning om Sandselighedens Realitet, der
maatte omstøde Ideens Gyldighed som den ene sande
Væren. Fra en Væren som objektiv Substants blev
Ideen nu henvist til subjektiv Væren som Slægtsbegreb.
De almene Begreber udvikler Mennesket i Kraft af Bevidstheden
om Slægten, hvis Væsen de udtrykke.  Ideen
% --- p. 52 ---
\setcounter{footnote}{0}%
\opage{52}er ikke længer Fornuftvæsen, men Væsen for Fornuften,
Fornuften selv. Som saadan har det Almene, Ideen
endnu Gyldighed overfor den endelige Empiri. Fornuften
ses som Menneskets sande Jeg, der giver de Love, under
hvilke Drift og Følelse skulle bøie sig. Men ligesom det
Objektive overhovedet kun er en Gjenspeiling af det Subjektive,
saaledes er Theologien Anthropologi, de religiøse
Forestillinger Frembringelser af Menneskets overstrømmende,
umættelige Higen efter Gemyttets Tilfredsstillelse.
--- Dette Stadium er navnlig repræsenteret ved
Feuerbachs Hovedskrift „Wesen der Christenthums“ (1841). % PRINTED AS IS (Wesen der Christenthums: the image, der for des; O027)

Det tredie Stadium indledes ved den Udtalelse (Philosophie
der Zukunft. 1843), at medens han tidligere
endnu har sat det Sandselige og Naturlige paa en underordnet
Plads som det, der skulde underordne sig Fornuften,
Erkjendelsen af det Almene, saa erklærer han
nu: kun det Sandselige er det Virkelige. Ikke Fornuften,
men Legemet i dets Totalitet er det sande Jeg.
Det Almene er kun en Illusion; den individuelle Drift
i sin Sympathi saavel som i sin Egoisme er det sande
Motiv. --- I „Wesen des Glaubens im Sinne Luthers“
og „Wesen der Religion“ forklares Religionen som udsprungen
af den egoistiske Lyksalighedsdrift, en Forklaring,
hvis extreme Gjennemførelse er given i Skriftet
„Theogonie“. Endelig gaaer Feuerbach i sit sidste Skrift:
„Gottheit, Freiheit, Unsterblichkeit“ (1866) lige til Materialismens
Grændse, idet han med særligt Velbehag
dvæler ved den Sætning: „der Mensch ist, was er isst.“
Hvad der holder ham tilbage er Overbevisningen om det
menneskelige Væsens indre Aktivitet. Medens han tidligere
betragtede de psychologiske Motiver, af hvilke Religionen
efter hans Mening udsprang, som fordærvelige,
saa bygger han nu selv sin Ethik paa dem. Det Standpunkt,
han paa sit første Stadium betegnede som det
praktiske, har han nu selv indtaget efterat have forjaget
Religionen derfra.

% --- p. 53 ---
\setcounter{footnote}{0}%
\opage{53}Efter denne foreløbige Oversigt kunne vi som den
almindelige Bevægelseslov, der behersker det Hele, betegne
Subjektiveringens Lov, ɔ: den Lov, at enhver objektiv
Væren, der sættes som Aandens Gjenstand, netop derved
er Aandens eget Værk, som derfor, naar Aanden
bliver sig selv bevidst, tages tilbage igjen, idet den sættende
Kraft reflekterer sig i sig selv og erkjender sin
Gjenstand som udsprungen af dens eget Væsen. Denne
Lov ville vi nu finde anvendt i det Enkelte.

a. Feuerbachs første Standpunkt er rent metaphysisk.
Vi have allerede omtalt hans „Gedanken über Tod
und Unsterblichkeit“, hvor alt Personligt ophæves i det
Almene, overfor hvilket det kun egoistisk kunde ville
gjøre sig gjældende. Dette Skrift kan betegnes som en
Rehabilitation af Døden: Døden er Udtrykket for det
Individuelles Intethed, er den almene Værens absolute
Seir. --- Det samme spekulativt-objektive Standpunkt er
gjennemgaaende i hans historiske Skrifter fra denne Periode.
I sin „Geschichte der neuern Philosophie“ lovpriser
han Spinoza og viser ved Udviklingen af det ontologiske
Bevis, at Guds Existents ikke er at adskille fra
hans Væsen; som Guds Væsen saaledes falder ogsaa
hans Existents indenfor Fornuften, der selv er den høieste
Væren. Det er den hegelske Logik, der endnu
fastholdes.

Bruddet med denne panlogiske Anskuelse indledes
ved den indtrængende Beskjæftigelse med den reale
Virkelighed, og først paa det religionsphilosophiske Omraade.
Feuerbach har med Iver fordybet sig i de individuelle
religiøse Phænomener, i den religiøse Psychologi;
deraf den Rigdom af karakteristiske og individuelle
Træk, som findes i hans Skrifter. Det var dette Individuelle,
der interesserede ham, medens den spekulative
Religionsphilosophi lod sig nøie med et Overblik over
Dogmernes almindelige, objektive Indhold uden Hensyn
til den psychologiske Jordbund, i hvilken det var op%
% --- p. 54 ---
\setcounter{footnote}{0}%
\opage{54}staaet. Feuerbach saa snart, at det var mere end en
Formforskjel, der adskilte Religion og Philosophi. Han
erklærede det for den spekulative Philosophis πρῶτον ψεῦδος at anse Erkjendelsesdriften for Religionens psychologiske
Kilde. I sit Skrift om Leibnitz henviser han
i al Almindelighed Religionen til det praktiske Omraade,
fordi den skylder sin Oprindelse og sin Magt til praktiske
Motiver, Haab og Frygt. Philosophien skal ikke,
som Hegel mente, søge at udfinde det fornuftige Begreb
af Religionens Indhold; Dogmet er ubegribeligt, fordi
det er ufornuftigt; den eneste Opgave er at forklare,
hvorledes det er opstaaet. I Afhandlingen om Bayle
finde vi den videre Udvikling heraf.

Theologien, hedder det her, har til sit Princip Underet,
idet den gaaer ud fra en vilkaarlig, ubetinget Villie
som Tilværelsens Aarsag. Spørgsmaalet om Tingenes
Grund existerer derfor slet ikke paa det theologiske Omraade;
det er engang for alle affærdiget med et Svar,
som intet Svar er; Villien er det asylum ignorantiæ,
hvorved al Videnskab afvises. Underet ophæver al Fornuft
og Erfaring. Man ser ved Underet kun de to Yderpunkter,
man ser f. Ex. Vin, hvor man før saa Vand;
men Overgangen, selve den egentlige Begivenhed er usyndig % PRINTED AS IS (usyndig: the image, usyn-|dig, for usynlig; recrop rc1)
og kan kun erfares med Troens Øine. Det virkelige
Under har derfor intet Fortrin for det opdigtede; thi det
Væsentlige ved begge, den overnaturlige Kausalitet lægges
ind af det opfattende Subjekt og udspringer ene af
dettes personlige Trang. Troen paa Underet er Underets
Væsen. --- I Modsætning hertil gaaer Philosophien
som Repræsentant for Videnskabens Ide ud fra Sagen
i sig selv og søger af Sagens Natur at finde den rette
Begrundelse. Philosophiens Princip er Fornuften, Lovmæssighedens
og Nødvendighedens Moder.

Naar Spørgsmaalet er om Christendommens Oprindelse,
da udleder Theologien den af en overnaturlig
Akt, der med Forbigaaelse af alle Betingelser umiddel%
% --- p. 55 ---
\setcounter{footnote}{0}%
\opage{55}bart virker paa dette bestemte Punkt i Historien. Men
det vil med andre Ord sige, at Theologien slet ikke anerkjender
dette Spørgsmaal for et Problem. Philosophien
derimod søger først at bestemme Religionens Natur
i Almindelighed og de almindelige Love, under hvilke
den fremtræder i Menneskeaanden og Historien. Disse
Love finde vi i de af Religionens Natur udspringende
Forestillinger, paa hvis Forskjellighed Forskjellen mellem
de enkelte Religioner beroer. En Religions Indhold kan
godt være irreligiøst, ɔ: staa i Strid med Religionens sande
Væsen. Den anden Faktor, Philosophien maa tage Hensyn
til, er Tiden, da Christendommen opstod. Under de
ulykkelige Tidsforhold, hvor Alt gik til Grunde, hvad
der havde baaret den antike Oldtid, kunde Sindet gjennemskue
alle endelige Vilkaars Intethed og frit hæve sig
over dem; under den almindelige sædelige Fordærvelse
maatte et rent, stille, inderligt Gemyt netop ved Modsætningen
til sine Omgivelser stødes tilbage i sig selv
for i den indre Verden at finde, hvad den ydre negtede.
Kun i en saadan Tid kunde Sædelighedens Ide fattes i
sin Renhed --- og det er det eneste Positive i Christendommen,
naar man fatter dens sidste Formaal i Øie.

Hvad Feuerbach her kalder det Positive i Religionen,
er dog ikke det, der gjør den til „positiv Religion“.
Ved Forskjellen mellem sand og falsk Religion er Indholdet
det afgjørende. Dette Indhold er Gudsbegrebet,
Maaden, hvorpaa Gud tænkes og er Gjenstand for Bevidstheden.
Men Gudsbegrebet er væsenlig Philosophiens
og Moralens Gjenstand, som Religionen gjør til Gjenstand
for Kultus. Hvad der gjør Religionen til Religion,
dens Forskjellighed fra Moral og Philosophi, Religionen
i sig selv er derfor endnu ikke god, hellig og guddommelig;
det er den kun, naar dens Indhold er et guddommeligt.
Helligheden er den eiendommelige religiøse Kategori.
Men Hellighed er ikke noget primitivt Begreb;
det Hellige maa først anerkjendes som fornuftigt og godt,
% --- p. 56 ---
\setcounter{footnote}{0}%
\opage{56}før det kan have Gyldighed. Kun det Sande er helligt,
men det Hellige som saadant er endnu ikke sandt. Naar
det Hellige adskilles fra denne dets berettigende Grund
og udvikles for sig selv, bliver Religionen positiv, „en
til Kirke verdsliggjort og ved sine Troesformler begrændset
Religion“, „et fra Fornuft og Sædelighed forskjelligt
Anliggende“. Kirkens Sag bliver Religionens Sag; naar
Mennesket kun troer, hvad Kirken lærer, og forretter
de hellige Handlinger, som den paabyder, saa er han
frelst. Den katholske Kirke er den eneste lykkede Definition
paa en positiv, kirkelig Religion; den har mest
trolig, mest konsekvent virkeliggjort dens Væsen, men
netop den har ogsaa paa det Skarpeste udtrykt den religiøse
eller kirkelige Interesses Indifferents imod Sædelighedens
Interesser. Den positive Religion betragter
Gud som et enkelt Væsen, overfor hvilket man har særlige
Pligter, der fylde Gemyttet med Frygt og Angst og
lade de humane, væsentlige og almene Pligtforhold træde
tilbage. Den positive Religion bliver nødvendig Lidenskabens
Sag. Angst for Salighedens Tab driver til Troen.
Selvopholdelses- og Lyksalighedsdriftens Begjær overvælder
Fornuften og kuer Evnen til ren og uhildet Sandhedskjærlighed.
Troen forsager vel Fornuften og det Høiere
i Mennesket, men ikke det Lavere, de egoistiske Motiver.
Det positive Religiøse staaer derfor i et udvortes
Forhold til det Ethiske. Dyden gjøres ikke til Princip,
men til Bisag: Troen, hedder det, gjør salig, men
Troen kan ikke være uden Gjerninger, ligesom Epikuræerne
sagde: Nydelsen gjør lyksalig, men Nydelsen er
ikke uden Dyd.

Som Dyden underordnes Troen, ɔ: Salighedsdriften,
saaledes underordnes det Gode den guddommelige Villie,
forestillet som vilkaarlig, personlig Villie. Det Gode har
for Philosophien sin Grund i sig selv, men for Theologien
er det Gode kun godt, fordi Gud vil det. Indvender
man herimod, at Gud kun vil det Gode, og at hans
% --- p. 57 ---
\setcounter{footnote}{0}%
\opage{57}Villie derfor ikke er nogen vilkaarlig Villie, saa har
man jo netop derved anerkjendt, at det Godes Ide er
det evige a priori, der bestemmer Villien og først gjør
den god. Kun i Kraft af det Godes Ide kan man tænke
det Gode som Prædikat. Derfor kommer det ud paa Et,
om man gjør Gud som det Gode eller det Gode i sig
selv til det Ethiskes Princip; Gud er her kun et Navn,
Væsenet er det Godes Begreb. Dog bør dette Navn udelades,
for at det Gode ikke blot skal ses som et maaskee
tilfældigt Prædikat ved et forudsat Subjekt, men kan
ses i sin absolute Selvstændighed. Kants kategoriske
Imperativ var det Manifest, i hvilket Ethiken forkyndte
sin Frihed og Selvstændighed for Verden. Naaede Kant
endnu kun den nøgne Pligtnødvendighed og kom ikke
ud over Ethikens Elementarlære, saa hævede Fichte sig
til saa ophøiede Ideer, at den positive Religion ikke har
Noget at stille ved Siden deraf. Vel ere Fichtes Ideer
strenge, næsten overmenneskelige, men Ethiken er heller
ingen Pædagogik; den akkomoderer sig ikke efter de
menneskelige Svagheder og Fornødenheder. Den sædelige
Aand er forsvunden, hvor det Godes Ide ikke tænkes
fri for al Personlighed og elskes og virkeliggjøres
for sin egen Skyld. ---

Feuerbachs Hovedbestræbelse gaaer her ud paa at
hævde Ideens Gyldighed. Skjøndt han ikke blot opfatter
den som abstrakt metaphysisk Ide, men i konkretere
Form, som den ethiske Ide, der er Enhed af Ethik og
Metaphysik, af Villie og Erkjendelse, saa afviser han
dog ethvert subjektivt Elements Indblanding. Methoden
beroer paa den tidligere angivne Subjektiveringens Lov:
han søger at opdage Bevidsthedens oprindelige Virken,
hvorved alt Andet begrundes. Her finder han da, at det
er den ethiske Ide, i Kraft af hvilken al Villie er god,
hvad enten den tænkes som guddommelig eller menneskelig
Villie; følgelig, slutter han, er den ethiske Ide
det Oprindelige, den gode Villie det Afledede. Og for
% --- p. 58 ---
\setcounter{footnote}{0}%
\opage{58}det Absolutes Vedkommende bortfalder saa Motivet til
at tænke sig et særligt, fra Ideen forskjelligt Subjekt.
Det er kun Mennesket, der som sandseligt Væsen ikke
falder sammen med Ideen, men netop derfor fremtræder
denne for ham som den ethiske Lov. I denne indeholdes
den sande Religion, medens den positive Religion
har sin Grund i, at den menneskelige Lyksalighedsdrift
ikke vil underordne sig Ideens Betingelser.

Ved den Vægt, Feuerbach her har lagt paa det
reale Moment baade i ethisk og i religiøs Henseende,
maatte hans Dom over den spekulative Philosophi modificeres.
Han priser det endnu som det Store ved Hegel,
at han har undersøgt de logiske og metaphysiske Grundbegreber
i og for sig, uden Forudsætning af et Subjekt.
Men dog udtaler han, at den systematiske Spekulations
Tid er forbi; der behøves en kritisk-genetisk Philosophi.
Den immanente Spekulation formaaer ikke at udvikle
den reale Virkelighed; Tanken maa kritisk bearbeide
denne, for saa --- genetisk --- at paavise dens Oprindelse
og derigjennem dens Væsen.

Men at Feuerbach ikke slog sig til Ro ved den Anskuelse
han havde udviklet i „Pierre Bayle“, dertil maa
Grunden søges i den Ubestemthed, hvormed „Ideens“
Forhold til det Empiriske opfattes, særlig i religionsphilosophisk
Henseende. Enten er Ideen det Oprindelige,
og det maa Feuerbach her forudsætte, men hvorledes
vil han da forklare de empiriske Forvanskninger
af dens Væsen, som han finder i de positive Religioner?
Eller ogsaa er det Empiriske det Første, men saa maa
der foregaa en Forandring i Bestemmelsen af det Ideale,
naar dette ikke længere er selvgyldigt, men betinges
ved hint.

Disse Spørgsmaal afgive Motivet til Feuerbachs følgende
Arbeider, i hvilke han indtræder i et nyt Stadium.

b. Overgangen til dette dannes af nogle mindre
Afhandlinger, i hvilke de Grundtanker allerede komme
% --- p. 59 ---
\setcounter{footnote}{0}%
\opage{59}frem, som ere de herskende i Feuerbachs Hovedværk
„Wesen des Christenthums“. Af disse Afhandlinger fremhæve
vi især „Kritiken des modernen Afterchristenthums“
og „Philosophie und Christenthum“. Ingen af Menneskets
væsentlige Evner --- saaledes er her Tankegangen
--- kan gaa ud over sig selv og fatte Noget, som er den
aldeles fremmed. Tænkningen kan ikke fatte en Tanke,
som ikke er dens egen: netop ved at fatte den vilde
den jo bevise, at det i Virkeligheden var dens egen
Tanke. Paa samme Maade gaaer det med Følelsen. Hvis
Følelsen skulde kunne gaa ud over den menneskelige
Tanke, maatte den ogsaa gaa ud over Menneskets Væsen,
og altsaa ikke være nogen menneskelig Følelse. Religionsphilosophien
maa derfor erkjende og behandle Religionen
som Psychologi. Modsætningen mellem Philosophi
og positiv Religion beroer paa Modsætningen mellem
Tænkningen og Hjertet paa den ene Side og Phantasien
og Gemyttet paa den anden Side. Hjertet er det naturlige,
sunde Gemyt, der gjør sig til Et med Fornuften.
Gemyttet derimod er det syge, overnaturlige Hjerte, der
forsmaaer Bestemmelse og Begrændsning, og for hvem
Tænkningen kun er en endelig Virksomhed, hvorimod
det tager sin Tilflugt til Phantasien, som hengiver sig til
Uendeligheden uden at lade sig beherske af fornuftige
og naturlige Love. Gemyttet og Phantasien hæve i Troen
paa en hinsides, himmelsk Uendelighed Mennesket ud
over alle endelige Forhold, men opløse derved ogsaa alle
Kjærlighedens og Menneskelighedens Baand.

I „Wesen des Christenthums“ udvikles disse Tanker
videre i dyb, indre Enhed og Sammenhæng. --- Mennesket
drives ikke som Dyret blindt og umiddelbart af
Slægtens Væsen, men dette bliver i Bevidstheden Gjenstand
for Mennesket. Da Væsenet som saadant er uendeligt,
maa Bevidstheden ogsaa være det. Og dog vindes
Bevidstheden kun ved Forholdet til en Gjenstand, altsaa
ved en Indskrænkning. Men denne Gjenstand har ingen
% --- p. 60 ---
\setcounter{footnote}{0}%
\opage{60}Betydning i og for sig, afset fra den Kraft, der sætter
den. Da Væsenet ikke kan gaa ud over sig selv, kan
den Gjenstand, til hvilken Subjektet staaer i et væsentligt
og nødvendigt Forhold, ikke være Andet end Subjektets
eget, objektivt udtrykte Væsen. Gjenstanden er
Subjektets aabenbare Jeg; dens Magt over Mennesket er
hans eget Væsens Magt. Fornuftgjenstandens Magt er
Fornuftens egen Magt; den elskede Gjenstand hersker
med Kjærlighedens egen Magt. Den, som tænker det
Uendelige, tænker og stadfæster Tankens Uendelighed;
den, som føler det Uendelige, føler og stadfæster Følelsens
Uendelighed.

Feuerbach angiver nærmere disse Principers Forhold
til tidligere religionsphilosophiske Anskuelser, navnlig
Schleiermachers og Hegels. Schleiermacher, denne
„Christendommens sidste Theolog“, havde gjort Følelsen
til det religiøse Organ. Men er Religionen en Følelsessag,
saa bliver Religionens Indhold ligegyldigt; det bliver
kun \emph{religiøst} Indhold ved at være Følelsens Gjenstand
og selve Følelsens Væsen bliver saa egentlig det
Guddommelige. Det var kun theologisk Hildethed, naar
Schleiermacher mente at kunne finde en Vei fra Følelsen
til den objektive Troeslære. --- Hegel derimod bestemte
Religionen, som det Bevidsthedstrin, hvor Menneskets
Viden om Gud blev Et med Guds Viden af sig selv.
Men atter her gjælder det: Organet er Et med Gjenstanden,
Midlet er selve Væsenet. Erkjendelsen af Gud
er Menneskets Selverkjendelse; Menneskets Viden om
Gud er Menneskets Viden om sig selv. Det mystiske
Forhold er et Forhold til sig selv.

Resultatet er altsaa: Gud er Menneskets aabenbare
Indre, hans udtalte Selv, den offentlige Bekjendelse af
hans Kjærlighedshemmeligheder. Dog gjælder dette kun
for den philosophiske Bevidsthed. Den religiøse Bevidsthed
er endnu kun Menneskets \emph{indirekte} Selvbevidsthed,
d. v. s. Mennesket bliver sig i Religionen sit eget
% --- p. 61 ---
\setcounter{footnote}{0}%
\opage{61}Væsen bevidst som et fremmed Væsen, henlægger først
sit Væsen uden for sig, før han finder det i sig selv.
Det historiske Fremskridt i Religionens Historie bestaaer
i, at denne Fremmedhed ophæves, at Selvbevidstheden
mere og mere erkjender sig selv i sin Gjenstand. Ethvert
Fremskridt i religiøs Henseende er derfor en dybere
Selverkjendelse. Dog er det eiendommeligt for
enhver positiv Religion, selv om den som Christendommen
gjør Menneskets eget Væsen til sit Indhold, at dette
Indhold staaer som et ydre, et fra Selvbevidstheden selv
forskjelligt. Denne Ufuldkommenhed hæves først, naar
man fra Religionens Sphære gaaer over i Philosophiens.

Vel indrømmer man, at de guddommelige Egenskaber
paa Grund af den menneskelige Erkjendelses Ufuldkommenhed
ere laante fra det menneskelige Væsen og
overførte paa Guddommen; men man fastholder alligevel,
at selve det guddommelige Subjekt ikke herved menneskeliggjøres.
Men det er en Distinktion, den religiøse Bevidsthed
ikke gjør og ikke kan gjøre. For den er Gud
er virkeligt, levende Væsen, ikke det i sig Bestemmelsesløse, % PRINTED AS IS (er virkeligt,: the image, er for et; O032)
som ligger bag Egenskaberne. Og desuden ---
hvor er Grændsen? Ere Egenskaberne Anthropomorphismer,
saa er Subjektet ogsaa Anthropomorphisme. Er
Kjærlighed, Godhed og Personlighed Anthropomorphismer,
saa er ogsaa selve Existentsen en anthropomorphistisk
Bestemmelse. Snarere forholder det sig omvendt; det
er Egenskaberne, Prædikaterne, som have Gyldighed i
sig selv, og Subjektet er kun det existerende Prædikat.
Den egentlige Atheisme er den, for hvem de guddommelige
Prædikater, Kjærligheden, Visdommen, Retfærdigheden
ere Intet, ikke den, der blot ophæver disse Prædikaters
Subjekt. Subjektet deler Skjæbne med Prædikaterne,
ikke omvendt.

Det følger af sig selv, at Mennesket, der overfører
væsentlige Prædikater fra sig selv til Guddommen og
derved først egentlig frembringer denne, ikke kan ved%
% --- p. 62 ---
\setcounter{footnote}{0}%
\opage{62}blive at tillægge sig selv de samme Prædikater, i det
Mindste ikke i samme Grad og Betydning. Han maa
nødvendigvis benegte hos sig selv, hvad han sætter i sin
Gud. Heraf kommer det Mærkelige, at jo menneskeligere
et Væsen Guddommen bliver, des større bliver
Forskjellen mellem Gud og Menneske. Mennesket plyndrer
sig selv for at berige sin Gud: jo mere Gud kommer
til at ligne Mennesket, des mindre kommer Mennesket
til at ligne Gud. Mennesket benegter sin Viden
for i Gud at se den Alvidende, sin Lyksalighed for i
Gud at se den absolut Salige, sin gode Natur for i Gud
at se den absolut Hellige.

Dette er dog kun endnu den ene Side af den religiøse
Bevidstheds Væsen, hvilken man kunde kalde det
religiøse Arteriesystem eller den religiøse Repulsionskraft.
Mennesket støder her sit eget Væsen ud fra sig og gjør
det til et objektivt Væsen. Den anden Side er, at Repulsionskraften
ophæves ved den modsatte Kraft, Tiltrækningen,
at der overfor Arteriesystemet staaer et Venesystem,
hvor Blodet strømmer tilbage til Hjertet. Vel
hæver den religiøse Bevidsthed Gud uendelig høit over
sig selv, men den ser dog sig selv som den væsentlige
Gjenstand for den guddommelige Virksomhed. I Gjenstanden
ligger Væsenet udtrykt; som den, der har Mennesket
til sin Virkens Gjenstand og Øiemed, er Guddommen
altsaa alligevel et menneskeligt Væsen. Religionens
Sandhed bestaaer i, at den stiller Mennesket i Forhold
til sit eget Væsen; dens Usandhed og Skranke bestaaer
i, at dette Væsen opfattes som et fremmed, ja modsat
Væsen --- heri ligger Kilden til den religiøse Fanatisme,
„det metaphysiske Princip for de blodige Menneskeofre.“

Naar Religionen udspringer af Menneskets Væsen
og dog sætter en Modsætning mellem Gud og Mennesket,
saa maa Spiren til denne Modsætning, med hvilken
Religionen begynder, kunne paavises i Menneskets eget
Væsen. Denne Spise ligger i den psychologiske Mod% PRINTED AS IS (Spise: for Spire; the Minnesota copy prints it too)
% --- p. 63 ---
\setcounter{footnote}{0}%
\opage{63}sætning mellem Forstand og Hjerte. Forstanden er uberørt
af Hjertets Lidelser, har ingen Drifter og Lidenskaber,
ingen Fornødenheder og derfor heller ingen
Mangler og Svagheder, saaledes som Hjertet. Medens
dette vender sig mod det Særlige, det Individuelle, er
Forstanden Repræsentanten for det Almene, Intelligentsens
rene, affektløse Lys. Gud som det metaphysisk
og moralsk fuldkomne Væsen er derfor Forstandens objektiverede
Væsen. Men overfor dette føler Mennesket
sig som det endelige og syndige. Hjertets Følelser komme
ikke til deres Ret. Derfor forestiller Hjertet sig paa
Grundlag af hint abstrakte Forstandsvæsen, der kun er
Religionens mathematiske Punkt, Gud som den uendelige
Kjærlighed. Derved forsones Mennesket med Gud.
Den dybeste Anskuelse af den guddommelige Kjærlighed
lader endog Gud selv blive Menneske. Inkarnationen er
den sandselige Aabenbaring af Guds menneskelige Natur;
Gud følte den menneskelige Trang.

Men denne Selvfordobling er ikke en uskadelig Illusion.
Det viser sig deri, at uagtet Guddommen ses som
Kjærlighed, saa er dog Kjærligheden ikke dens Væsensbestemmelse,
ikke gyldig i og for sig, men knyttes til
et forudsat guddommeligt Subjekt, der ligesom lurer i
Kjærlighedens Baggrund og endnu er Noget for sig afset
fra Kjærligheden, --- „et dæmonisk Væsen, hvis fra
Kjærligheden adskillelige og virkelig adskilte Personlighed
fryder sig ved Kjætternes og de Vantros Blod, den religiøse
Fanatismes Phantom.“ Hvorfra hidrøre da disse fordærvelige
Elementer, dette Usande i Religionens Væsen?

Forholdet til Gud som til et andet Væsen er dels
et naturligt, uvilkaarligt og ubevidst Forhold, dels et
bevidst og reflekteret. Det første har sin Rod i selve
Religionens Oprindelse og beroer paa dens væsentlige
Standpunkt, det praktiske. Menneskets Forhold til Gud
er Et med Forholdet til hans egen Salighed: Gud er selv
den virkeliggjorte Salighed og den ubegrændsede Magt
% --- p. 64 ---
\setcounter{footnote}{0}%
\opage{64}til at virkeliggjøre Menneskets Salighed. Gud som Religionens
Gjenstand --- og kun som saadan er han Gud
--- er Gemyttets, ikke Fornuftens Gjenstand, Gjenstand
for Praxis, ikke for den ubekymrede Theori, for Hjertenøden,
ikke for Tankefriheden. Fornuften, det Almene
virker paa det praktiske Menneske kun gjennem Forestillingen
om et personligt Væsen. I Religionen er Aandens
Lys spaltet i Phantasiens og Gemyttets Medium,
hvorved det ene og samme Væsen anskueliggjøres som
fordoblet. Om det, der ligger bag denne praktiske Bevidsthed
bekymrer Religionen sig ikke. Den umiddelbare
religiøse Bevidsthed hviler med barnlig Tryghed i
denne Enhed af det guddommelige og menneskelige Væsen.
Men jo mere Reflexionen vaagner, des mere arbeides
der paa at ophæve Enheden, fordi man aner, at den
i sine Konsekventser kunde føre til at ophæve den religiøse
Gjenstand som saadan. Den Forskjel, der paa
det umiddelbare Standpunkt kun er Brydningen af et
og samme Lys, fæstnes nu til en dogmatisk Modsætning.
Denne Bestræbelse er gjennemgaaende i de saakaldte
Beviser for Guds Tilværelse. De ville godtgjøre, at
Religionen har Ret, naar den sætter Menneskets Væsen
som et fremmed og objektivt. Guds Existents opfattes
da her som en særegen Bestemmelse, der tillægges den
religiøse Gjenstand, og hvorved denne sættes udenfor Bevidstheden.
Gud skal ikke blot være i Troen, i Følelsen,
i Gemyttets Inderlighed, men have en derfra forskjellig
og uafhængig Væren. Men en real, fra Tanken
forskjellig og uafhængig Væren kan kun være en sandselig
Væren. Kun ved at tillægge Gud en saadan, kan
man hævde hans Existents „udenfor os“. Theologien
vil undgaa dette ved at erklære, at dette „udenfor“ ikke
skal tages i ligefrem Betydning. Men i hvilken Betydning
da? En aandelig Væren er kun en Væren i Tanken,
aldeles ikke udenfor os. Reflexionen fører derfor
nødvendig til Atheisme, i den Forstand, at Guds Existents
% --- p. 65 ---
\setcounter{footnote}{0}%
\opage{65}negeres som en ydre, empirisk, af Tanken uafhængig
Væren. Men det sande Guddommelige er, som vi allerede
 have set, ikke denne blotte og bare Existents, men
indeholdes i „de guddommelige Egenskaber“. Vil Theologien
ikke gjøre dette Skridt, saa er der ikke Andet
for den at gjøre end at gaa tilbage til det umiddelbare
religiøse Bevidsthedstrin, der ikke sondrer mellem Indre
og Ydre. Ved Indbildningskraften, som er den til Gemyttet
svarende intellektuelle Evne, sættes en ydre Existents,
som dog ikke er ydre, en sandselig Væren, som
dog ikke er sandselig. Denne Svæven mellem det Sandselige
og det Usandselige er netop Indbildningskraftens
Væsen. Kun Phantasien bevarer for Atheisme.

Den Angst og Lidenskab, hvormed Mennesket under
den vaagnende Reflexion fastholder den religiøse Gjenstands
objektive Existents, staaer i nøie Forbindelse
med, at det ikke blot er det Sande, Almene og Berettigede
i Menneskets Væsen, der gjøres til Gud, men
ogsaa det Individuelle, Vilkaarlige og Endelige. Dette
Religionens dobbeltsidige Væsen finder Feuerbach især
udtrykt i Christendommens Lære om Forholdet mellem
Tro og Kjærlighed. Kjærligheden gjør Mennesket til Et
med Gud og bestemmer Gud som en menneskelig Gud.
Troen adskiller Mennesket fra Gud, idet den sætter ham
som forskjellig fra og hævet over de humane Væsensbestemmelser.
Dette gjør, at Kjærligheden begrændses,
ikke frit kan røre sig; Gud er vel Kjærlighed, men Kjærligheden
er ikke Gud. Kjærligheden har Gud i sig, men
Troen har ham uden for sig, og Kjærligheden gjøres i
Christendommen afhængig af Troen. Idet Gud adskilles
fra Mennesket, adskilles det ene Menneske fra det andet,
thi Gud er kun Menneskets mystiske Slægtbegreb, hans
Adskillelse fra Mennesket altsaa Opløsningen af det fælles
Baand. Troen er i sit Væsen exklusiv: dette Bestemte
skal du tro, naar du ei vil miste din Salighed.
Den Troende har Gud for sig, den Ikketroende har ham
% --- p. 66 ---
\setcounter{footnote}{0}%
\opage{66}imod sig; Tro er Godhed, Vantro Ondskab og Gjenstridighed.
Det Bud, at man skal elske sine Fjender, gjælder
ikke overfor Guds Fjender. Derfor ligger der i Troen
et ondt Princip: „Helvedes Flammer ere kun Gnisterne
af det tilintetgjørende, af Vrede glødende Blik, som
Troen kaster paa de Vantro.“ I Kjærlighedens Væsen
ligger derimod Universaliteten. Den sande Kjærlighed
er sig selv nok og fordrer ingen særlig Adkomst eller
Autoritet. Den erkjender Dyden i selve Synden, Sandheden
i selve Vildfarelsen. I Kjærligheden virkeliggjøres
Slægtens Væsen paa subjektiv Vis, gjennem Sindelaget.
Den er derfor i sig Et med Fornuften, som er Slægtens
objektive Realitet. I Modsætning til Fornuften og det
med den væsensbeslægtede Hjerte staaer Gemyttet, Individualitetens
Selvfølelse som saadan. Christendommen
er Modsigelsen mellem Hjerte og Gemyt, fordi den er
Modsigelsen mellem Kjærlighed og Tro: Troen kommer
fra Gemyttet, Kjærligheden fra Hjertet. Christusfiguren
er det Billede, under hvilken Slægtens Enhed paatrængte % PRINTED AS IS (under hvilken: the image, hvilken for hvilket (Billede); O036)
sig Folkebevidstheden. Slægtens Væsen, der lægger sig
for Dagen i Kjærlighedens Energi, overvinder enhver
Partikularitet, enhver Skranke. Hvor derfor Bevidstheden
om Slægten som Slægt opstaaer, der forsvinder
Christus, uden at hans sande Væsen forgaaer.

I dette Forhold mellem Troen og Kjærligheden ser
Feuerbach det tydeligste Udtryk for Nødvendigheden af
at gaa ud over det Religiøse. Er det godtgjort, at Theologien
er Anthropologi, at Guds Væsen er Menneskets
eget Væsen, saa er det nødvendigt at lade det komme
til det Vendepunkt, hvor denne klare Selvbevidsthed afløser
den tidligere, for sig selv fremmede religiøse Bevidsthed.
Og det ikke blot af theoretiske Grunde. Thi
Selvfordoblingen i den religiøse Bevidsthed er ikke en
uskadelig Illusion. For det Første er Gemyttet, Religionens
egentlige psychologiske Faktor og Kilde, det
med Naturen splidagtige og derfor ud over Naturen higende
% --- p. 67 ---
\setcounter{footnote}{0}%
\opage{67}Sind, der kun paa mirakuløs Maade kan se sine Ønsker
tilfredsstillede. I Modsætning hertil maa det hævdes, at
Naturen hører med til Menneskets Væsen og har sand,
real Gyldighed og Betydning i sig selv. Dernæst maa
det Ethiske udløses af sine Skranker, ikke længere udledes
af Troen. Det Rette, Sande og Gode har stedse
Grunden til sin Hellighed i sig selv. De ethiske Forhold
ere i og for sig religiøse Forhold. Livet er overhovedet
af guddommelig Natur og behøver ikke nogen
særegen Indvielse. ---

Hvad der i dette Feuerbachs Hovedskrift især tildrager
sig Opmærksomheden, er den Maade, hvorpaa
Religionen opfattes. Idet den føres tilbage til Gemyttet,
de ubegrændsede Ønskers Kilde og bestemmes som det
rent praktiske Standpunkt, hvor Frygt og Haab ere de
ledende Motiver, træder Modsætningen til Hegels intellektualistiske
Opfattelse af Religionen klart frem. Feuerbach
har selv erklæret, at man med Urette betragter
hans Lære som en videre Udvikling af den hegelske;
den er netop opstaaet af Opposition mod denne. Hvad
der for Hegel var det Sekundære og blot Formelle, Gemyttet
og Følelsen, i hvilke Religionens objektive Indhold
sænker sig for at blive Menneskets Eiendom, det
er for Feuerbach netop det Væsentlige og Primitive.
Nærmere staaer han derimod Schleiermacher, som gjorde
Følelsen til det religiøse Organ; og medens Hegel polemiserede
mod Schleiermacher, fordi denne tog saa subjektivt
et Udgangspunkt, er det, som allerede anført,
Feuerbachs Indvending mod Schleiermacher, at han ikke
er subjektiv nok, ikke absolut gjennemførte sit subjektive
Udgangspunkt.

Fremskridtet i Forhold til Afhandlingen om Bayle
bestaaer netop i den bestemtere Anvendelse af det subjektive
Synspunkt paa den positive Religion. Det ses
f. Ex. ved Behandlingen af Underbegrebet. Medens han
i „Pierre Bayle“ godtgjør Underets Intethed ved  at efter%
% --- p. 68 ---
\setcounter{footnote}{0}%
\opage{68}vise dets Uoverensstemmelse med det guddommelige
Væsen, betragter han det i „Wesen des Christenthums“
som den supranaturalistiske Tilfredsstillelse af et menneskeligt
Ønske og anlægger en subjektiv-ethisk, ikke
længer en objektiv-metaphysisk Kritik deraf.

Det individuelle Gemytsliv, hvoraf den positive Religion
udspringer, er sygeligt og uberettiget, fordi det
vil hæve sig over Fornuftens almene Væsenslov. Det
Individuelle underordnes altsaa endnu bestemt det Almene.
Medens Gemyttet ofrer Slægten for Individet,
ofrer Fornuften Individet for Slægten. Men det Almene
opfattes ikke længere, som i „Pierre Bayle“, under
den objektive Ides Form. Ideen ses nu som Menneskets
eget Væsen, som Slægtens Lov i Mennesket, om end
over Mennesket. Det Almene bevarer sin Realitet, kun
at denne forlægges fra den objektive Væren til den subjektive.
Feuerbach erklærer udtrykkelig, at det ved Kritiken
af Religionen er nødvendigt at støtte sig til det
Almene: „kun Slægten er istand til paa engang at ophæve
og erstatte Religionen.“

Men hans rastløse Aand drev ham stedse videre. I
Kraft af Subjektiveringsmethoden maatte han nødvendigvis
opkaste det Spørgsmaal, om ikke ogsaa det Almenes
\emph{subjektive} Væren skulde være en Illusion, der
ligesom Theologien opløser sig i ren Anthropologi. Derved
vilde da Veien banes til Overvindelse af den Modsætning,
som er gjennemgaaende i hans tidligere Skrifter,
Modsætningen mellem det Almene og det Individuelle.
Dette er Opgaven, hvis Løsning Feuerbach søger
paa det tredie Stadium af hans Forfattervirksomhed.

c. Feuerbach har selv betegnet Stadierne i sin Udvikling
saaledes: „Gud var min første Tanke, Fornuften
min anden, Mennesket min tredie og sidste Tanke“.
(Sämmtl. Werke. II p. 410). Dette sidste Stadium har
han dog endnu ikke ganske naaet i „Wesen des Christenthums“,
skjøndt han der vil gjennemføre den Sætning,
% --- p. 69 ---
\setcounter{footnote}{0}%
\opage{69}at Theologien er Anthropologi. Han hævder nemlig der
endnu Fornuften, Bevidstheden om Slægtens almene
Væsen som det, der vel er i Mennesket, men dog \emph{over}
Mennesket. Der er endnu et Transscendentsens Moment
tilbage i Immanentsen. Feuerbach er forsaavidt endnu
Idealist, idet han ved Bestemmelsen af Menneskets Væsen
gaaer ud fra Bevidstheden om det Almene og deraf udleder
det ethiske Princip.

Dog findes der allerede i en Afhandling i „Deutsche
Jahrbücher“ for 1841 (optaget i 2det Bind af hans Værker):
„Ueber den Anfang der Philosophie“, en Proklamation
af Realismen. Den philosophiske og empiriske
Videns Begyndelse erklæres her for en og samme Akt.
Philosophien skal ikke blot tilegne sig Empiriens Resultater;
den kan ikke komme til Realiteten uden ved
at begynde med Realiteten. Den maa i den empiriske
Virksomhed se en philosophisk Virksomhed, i Synet en
Tænken, i Sandseredskaberne Philosophiens Organer.
Psychologien er ikke en speciel Videnskab, men den
første almindelige Videnskab, der af Jegets forskjellige
Forhold deducerer de særskilte Principer. Til Jeget
hører da ogsaa Legemet, der som Jegets Passivitet ikke
blot har naturhistorisk eller empirisk, men ogsaa spekulativ
og metaphysisk Betydning.

Det her Antydede udvikles videre i Skriftet: „Grundsätze
der Philosophie der Zukunft“ (1843) paa følgende
Maade. Den nyere Tids Opgave var den at udvikle
Theologien til Anthropologi. Denne Opgave løste Protestantismen
fra sit praktiske Standpunkt, idet den kun
tog Hensyn til, hvad Gud var som Menneske og for
Mennesket; paa rationel, theoretisk Vei er Opgaven løst
i den nyere spekulative Philosophi, hvis Fuldendelse er
den hegelske Philosophi. Guds Væsen blev her udviklet
saaledes, at det blev Et med Fornuftens eget Væsen,
idet de guddommelige Væsensbestemmelser („Egenskaber“)
viste sig at være identiske med Fornuftens, og særlig
% --- p. 70 ---
\setcounter{footnote}{0}%
\opage{70}den spekulative Fornufts Væsensbestemmelser. Uendelighed,
Nødvendighed, Uforanderlighed, Evighed tilkomme
de evige Sandheder og Fornuftlove, og som Følge deraf
Fornuften selv. I Fornuften er det Tænkende og det
Tænkte Et. Medens derfor den almindelige Bevidsthed
skjelner mellem Gud som tænkende og som tænkt af
Mennesket, saa falde disse to Sider nu sammen som
identiske. Den guddommelige Tænken er kun en forestillet
Tænken, der har sin Virkelighed i den spekulative
Philosophis aprioriske Viden; og naar man ogsaa
tænkte sig Gud som den, der vidste alt Sandseligt, saa
har denne guddommelige Viden fundet sin Virkeliggjørelse
i den nyere Tids empiriske Videnskaber.

Men idet den nyere Philosophi forvandlede Gud til
Fornuften, lod den Fornuften selv antage guddommelige
Egenskaber. Pantheismen er en theologisk Atheisme, en
Negation af Theologien, men paa Theologiens eget Omraade.
Dette viser sig navnlig i Spekulationens Abstraktion,
der er den theoretiske Eftervirkning af Middelalderens
Askese. Hegel har derfor kun bevist Forstandens
Guddommelighed i dens Modsætning til Sandseligheden,
Verden, Mennesket selv. Er Theologiens Hinsides
blevet et Dennesides, saa er til Gjengjæld selve
den virkelige Verden bleven et Hinsides.

Den nye Philosophi, Fremtidens Philosphi maa derfor % PRINTED AS IS (Philosphi: for Philosophi)
træde i skarp Modsætning til Spekulationen og
først og fremmest korrigere Værens Begreb. Væren er
intet almindeligt Begreb, der kan adskilles fra de existerende
Ting selv. Væren er Et med det, som er.
Naar Hegel i Phænomenologien viser, at da enhver Ting
er et „Dette her“, saa er det Enkelte umiddelbart det
Almene, da har han ikke godtgjort Andet end Modsigelsen
mellem Ordet, som er alment, og Sagen, som
altid er enkelt. Ud over denne Modsigelse kan Tanken,
der stedse maa støtte sig til Ordet, ikke komme. Men
Spørgsmaalet om Væren er ikke blot theoretisk; det er
% --- p. 71 ---
\setcounter{footnote}{0}%
\opage{71}et praktisk Spørgsmaale, „eine Frage auf Tod und Leben“. % PRINTED AS IS (Spørgsmaale: the image, for Spørgsmaal; O038)
Ved sin Modsætning til det almene Ord bliver Væren
uudsigelig. Livet, Værens Hemmelighed, begynder først,
hvor Ordene høre op. Naar derfor Uudsigelighed er det
Samme som Ufornuftighed, saa er al Existents Ufornuft,
netop fordi den er denne bestemte Existents; --- men
saaledes er det ikke: Existentsen har i og for sig selv
Betydning og Fornuft, skjøndt den er uudsigelig.

Den nye Philosophi gjør altsaa Alvor af det Konkret-Enkelte,
som den spekulative Philosophi mente at lade
komme til sin Ret, naar den optog Enkeltheden som
Moment i Begrebet. Ikke som Begrebsmoment, men
som et sandt og virkeligt Væsen, og det vil sige: som
et sandseligt Væsen, skal det Virkelige være Erkjendelsens
Gjenstand. Sandhed, Virkelighed og Sandselighed
ere Et. Kun ved Sandserne \emph{gives} en Gjenstand. Væren
kan kun opfattes af Væren, af Existents, af Sandsning,
Anskuelse, Følelse, Kjærlighed: „das Sein ist also ein
Geheimniss der Anschauung, der Empfindung, der Liebe.“
Kjærligheden er Lidenskab, og kun Lidenskaben er Existentsens
Kjendemærke. Den abstrakte, lidenskabsløse
Tænken ophæver Forskjellen mellem Væren og Ikkeværen.
Naar den tidligere Philosophi sagde: hvad der
ikke kan tænkes, er ikke, saa siger den nye Philosophi:
hvad der ikke kan elskes, er ikke. Kun hvad der kan
elskes og tilbedes og altsaa kan være Gjenstand for Religion,
kun det kan ogsaa være Philosophiens Gjenstand.

Alt kan fornemmes paa sandselig Vis, om ikke umiddelbart,
saa dog middelbart, om ikke med de raa og
pøbelagtige, saa med de uddannede Sandser, om ikke
med Chemikerens og Anatomens, saa dog med Philosophens
Øine. Heller ikke er det Sandselige det, der ligger
lige for Haanden, det Tankeløse, som forstaaer sig
af sig selv. Det er kun gjennem en lang og vanskelig
Uddannelse, at Mennesket opdrages til Anskuelsen af
det Sandselige; umiddelbart ligge Phantasiens, Forestil%
% --- p. 72 ---
\setcounter{footnote}{0}%
\opage{72}lingens og den abstrakte Forstands Væsener ham langt
nærmere. Anskuelsen og Tænkningen maa gjensidig bestemme
hinanden. Anskuelsen for sig har ingen Grundsætninger,
og Tænkningen i og for sig kommer ikke ud
over den logiske Enhed med sig selv. Tankens Bevægelse
maa afbrydes af den sandselige Anskuelse og kan
derfor ikke, som Hegel vilde, afslutte sig til „en Kreds
af Kredse“. Den nye Philosophis Symbol bliver ikke
en Cirkel, men en Ellipse. Dog er det ingen fremmed
Magt, der bestemmer Tankens Bevægelse, thi Virkelighedens
Love ere Tankens egne Love; denne danner
ikke en Stat i Staten.

Den nye Philosophi gjør Mennesket til sit Erkjendelsesprincip,
ikke Fornuften. Mennesket er Fornuftens
Realitet og Subjekt; det er Mennesket, der tænker, ikke
selve Fornuften. Naar den tidligere Philosophi sagde:
„kun det Fornuftige er det Sande og Virkelige“, siger
derimod den nye Philosophi: „kun det Menneskelige
er det Sande og Virkelige.“ Men ved det Menneskelige
maa der ikke blot tænkes paa Forstanden. Det er ikke
blot ved Tænkningen, at Mennesket adskiller sig fra
Dyret, men ved sit hele Væsen, af hvilket Tanken kun
er en enkelt Følge. Da det Menneskelige med Indbegreb
af Naturen som Menneskets Basis, nu bliver
Philophiens eneste Gjenstand, saa bliver Anthropologien % PRINTED AS IS (Philophiens: the image, for Philosophiens; O039)
Universalvidenskab.

Men det enkelte Menneske har ikke Menneskets hele
Væsen i sig. Først i det ene Menneskes Fællesskab og
Enhed med det andet Menneske aabenbarer Menneskets
Væsen sig. Kun Enheden af Jeg og Du er Gud. Det
er denne Sandhed, der er udtrykt i Treenighedsdogmet,
den Sandhed, at intet Væsen i og for sig er sandt og
fuldkomment, men at Sandheden og Fuldkommenheden
kun er Enheden og Forbindelsen af væsenslige Væsener. % PRINTED AS IS (væsenslige: the image, for væsentlige?; O040)
Philosophiens sidste og høieste Princip er derfor Menne%
% --- p. 73 ---
\setcounter{footnote}{0}%
\opage{73}skets Enhed med Mennesket. Af denne Enhed udspringe
alle væsentlige Forhold og Principerne for alle de forskjellige
Videnskaber.

Medens den tidligere Philosophi havde en dobbelt
Sandhed, en theoretisk og en praktisk, idet den stillede
sig som Sandheden i og for sig i Modsætning til Religionen,
Sandheden for Mennesket, saa overvinder den
nye Philosophi denne Modsætning: den har ikke blot
samme Gjenstand som Religionen, men ogsaa Religionens
Væsen i sig, er selv den sande Religion. ---

„Den nye Philosophi“ skal efter Feuerbach være
væsensforskjellig, toto genere forskjellig fra den tidligere,
og den afgjørende Modsætning skal bero paa, at den
med Bevidsthed anerkjender Sandseligheden, hvad Spekulationen
kun indirekte havde gjort ved i det høieste
Princip at se det Reelle indesluttet som Moment. Opløsningen
af den absolute Idealisme maatte med Nødvendighed
føre til et Punkt, hvor man blev opmærksom
paa de positive Elementer, som vare blevne miskjendte
og oversete af Spekulationen. Feuerbach repræsenterer
netop det Stadium i Udviklingen, hvor den negative Fart
bringer Kritikeren Ansigt til Ansigt med det Positive
og Reale. Men netop hans negative og kritiske Tankegang,
der selv i høi Grad er paavirket af de theologiske
og spekulative Modstandere, den har bekæmpet, gjør
ham det vanskeligt paa positiv Vis at udvikle de nye
Elementer. Han har mere skuet det Nye med begeistret,
næsten tumultuarisk Følelse end klart og stringent
angivet frugtbare Principer for dets Behandling,
mere stillet en stor Fordring end givet Bidrag til dens
Løsning. Dertil kommer, at hans urolige, „utilfredse
Aand ikke kunde undgaa Selvmodsigelser, da den ikke
kunde finde sig i at afklare logisk, hvad den med mystisk
impetus grebe gjennem Phantasien. „Hans Philosopheren
vedblev at være Divination. Hans „Sandse%
% --- p. 74 ---
\setcounter{footnote}{0}%
\opage{74}lighed“ er en ny Form for den absolute Tænken, der
ganske ser bort fra den praktiske Erfaring“\footnote[1]{\emph{Alb. Lange}: Geschichte des Materialismus. p. 284. 306.}.

Erkjendelsesprincipet skal være Mennesket selv efter
hans Væsens hele Omfang. Medens den tidligere Philosophi
søgte Videns sidste Begrundelse i Videns metaphysiske
Ide, der lægger sit Væsen for Dagen gjennem
Menneskets psychologisk betingede Viden, gjør Feuerbach
den psychologisk-menneskelige Viden til Et med
den metaphysiske. Slægten i sin Mangfoldighed skal
nn besidde, hvad man tidligere koncentrerede i Videns % PRINTED AS IS (nn: the image, for nu; O044)
Ide. Men hvilken Realitet kan Feuerbach tilskrive Slægten
som saadan? Naar den ikke kan samle sin Mangfoldighed
til en Enhed, bliver Viden kun et Aggregat,
hvis Brudstykker existere for hver Enkelts Forestilling,
uden at Helheden endog kun, som Kant vilde, kan
blive ledende Regel og Formaal. I Virkeligheden bliver
kun hvert enkelt Menneskes Viden tilbage. Og efter
Feuerbach véd det enkelte Menneske kun, hvad han
har opfattet gjennem Sandsningen. Den nye Philosophi,
siger han, er i Henseende til sin Basis ikke Andet end den
til Bevidsthed hævede Fornemmelse, Hjertet, som er
bragt til Forstand. Standpunktet bliver saaledes fuldstændig
Sensualisme. Feuerbach siger, at han vil føre
Tænkningen tilbage til Enhed med Erfaringen; men det
bliver i Virkeligheden en fuldkommen Afhængighed af
Erfaringen. Naar han siger, at Alt idetmindste middelbart
kan opfattes ved Sandsningen, om ikke ved den
raa, saa dog ved den udviklede Sands, saa antyder han
herved en Mellembestemmelse, der maatte have hindret
de sensualistiske Konsekventser; thi hvad der adskiller
den uddannede Sands fra den umiddelbare, det er netop
det aandelige Arbeide, Tankens Energi, hvis Resultater
umiddelbart lægge sig ind i selve den sandselige Anskuelse.

Det Sandselige skal være Erkjendelsens eneste Gjen%
% --- p. 75 ---
\setcounter{footnote}{0}%
\opage{75}stand. Men det Sandselige  er dette Enkelte og Individuelle,
der er uudsigeligt, fordi Ordet som det Almene
ikke kan gribe det i dets individuelle Bestemthed; ---
og dog skal det besidde Fornuft og Gyldighed trods
denne Uudsigelighed. --- Dette synes at maatte gjøre
den „nye Philosophi“ aldeles umulig. Thi den skulde
jo netop hæve Fornemmelsen til Bevidsthed, bringe
Hjertet til Forstand, altsaa omsætte det Enkelte i det
Almene. Naar Feuerbach hævder. at det Individuelle er % PRINTED AS IS: „hævder.“
 fornuftigt, skjøndt denne Fornuft ikke kan udsiges eller
 begribes, saa statuerer han Paradoxet. Hans Bestem-
melser danne her en fuldstændig Parallel til S. Kierke-
gaards Bestemmelser af „den Enkelte“ og „Gudmenne-
sket“. Ligesom Kierkegaard paa Grund af Umuligheden
af at begribe erklærer Troen for en Lidenskab, saaledes
er ogsaa for Feuerbach Lidenskab Existentsens
Kjendemærke: kun det er, som er Gjenstand for Lidenskab.
Sandheden kan kun gribes med Menneskets hele
Væsen, ikke med Tanken alene. I en Aphorisme fra
samme Tid som det Skrift, vi her beskjæftige os med,
hedder det: „Videnskaben, siger man, løser ikke Livets
Gaade. --- Indrømmet; men hvad følger deraf? At du
skal løbe over til Troen? Det vilde være at komme fra
Dynen i Halmen. Nei, men at du skal gaae over til Livet,
til Praxis. De Tvivl, som Theorien ikke kan løse,
løser Praxis for dig“. (Werke II p. 411). Den nye Philosophi
erkjender altsaa selv, at den ikke han gjøre Fyldest % PRINTED AS IS (ikke han: the image, han for kan; O046)
i theoretisk Henseende.

I praktisk Henseende har den nye Philosophi gjort
det Fremskridt at overvinde Modsætningen mellem det
videnskabelige og det religiøse Standpunkt; den er selv
Religion, idet den kun gaaer ud paa det, der er den
lidenskabelige Fornemmelses og Kjærligheds Gjenstand.
Dog opstaaer der her en lignende Vanskelighed som i
theoretisk Hesseende. Ikke det enkelte Menneske, men % PRINTED AS IS (Hesseende: for Henseende)
kun Foreningen af Jeg og Du skal udtrykke Menneskets
% --- p. 76 ---
\setcounter{footnote}{0}%
\opage{76}Væsen. Det Uklare er, om Mennesket søger denne Forening,
fordi Slægtens almene Væsen fører ham ud over
hans Individualitet, eller fordi denne lidenskabelig-egoistisk
finder sin Tilfredsstillelse i Du'et. „Mit Princip,“ siger
Feuerbach i en Aphorisme, „er Ego og Alter Ego, Egoisme
og Kommunisme; thi begge ere uadskillelige.“ Kjærligheden
maa nemlig forudsætte Andres Egoisme som
berettiget. Men heri ligger kun, at det Individuelle og
det Almene \emph{skulle} forenes, ikke, hvorledes de kunne
være Et. I „Wesen des Christenthums“ var det Almene
endnu ene herskende og Egoismen Rod til alt Ondt;
nu har det optaget denne som sideordnet, uden at Spændingen
ophører. I praktisk saavel som i theoretisk Henseende
bliver der bestandig en Modsigelsens Braad tilbage,
der driver den rastløse Tænker videre fremad.

Feuerbach modificerede snart sin Opfattelse af Religionen
i Overensstemmelse med „Grundsätze der Philosophie
der Zukunft“, og det saaledes, at derigjennem
selve de i denne Afhandling udtalte Principer udvikledes
videre og skarpere. „Das Wesen des Glaubens im Sinne
Luthers“ (1844) er, ifølge hans egen Erklæring i Fortalen
til hans samlede Værker, det første Skrift, i hvilket
det abstrakte Fornuftvæsen ikke længere spøger i
Hovedet paa ham, men hvor Naturens og Menneskehedens
sandselige, virkelige Væsen er kommet til sin Ret, og
og hvor derfor ogsaa det, han erkjender som det Væsentlige
i Religionen, tillige er det Væsentlige for ham
selv. Det Eiendommelige ved det nævnte Skrift er, at
det paa engang stiller det Religiøse lavere og høiere
end „Wesen des Christenthums“. Her saa han endnu
det Almene som en væsentlig Faktor i Religionen, uagtet
det blev overvældet og forvansket af Gemyttets individuelle
Trang. Men Afhandlingen om Luther bestemmer
Gud som Menneskets tilfredsstillede Lyksalighedsdrift,
som det Væsen, der virkeliggjør de menneskelige Ønsker,
hvis høieste Potents er Ønsket om Salighed. Det sættes
% --- p. 77 ---
\setcounter{footnote}{0}%
\opage{77}som Luthers Eiendommelighed, at han ikke med den
gamle Kirkelære lagde Hovedvægten paa Dogmets objektive
Indhold, men helt igjennem betonede det subjektive
Moment: „Alles was wir im Glauben erzählen, ist
für uns geschehen und kommet uns heim,“ siger Luther.
Og Feuerbach tilføier: „Luther erst hat das Geheimniss
des christlichen Glaubens ausgeplaudert“. Troen i Luthers
Forstand bestaaer i Troen paa Gud som et Væsen,
der væsentlig staaer i Forhold til Mennesket. Vel gjør
Luther Mennesket til Intet overfor Gud; hans Lære er
en Hymne til Gud, en Pasquill over Mennesket. Men
Gud kunde ikke være Alt for Mennesket, naar Mennesket
selv ikke var Intet. Luther er kun inhuman mod
Mennesket, fordi hans Gud er absolut human. Den Gud,
Luther fjerner absolut fra Mennesket, er en Gud, der
selv vil det Høieste, Mennesket kan ville for sig selv.
Troen er derfor i sidste Instants Kjærlighed til sig selv.
Men alligevel fordømmes Troen ikke, som i „Wesen
des Christenthums“, som den absolute Egoisme. Vel
fastholdes det endnu, at kun Slægten kan være Maalestok
for det Gode, naar der skal være en absolut, ikke
relativ og egoistisk Maalestok; men paa den anden Side
udtales, at naar Kjærligheden bygges paa Troen, hvilket
efter Feuerbachs Opfattelse vil sige det Samme som at
Kjærligheden til Andre bygges paa Kjærligheden til
sig selv, saa kan dette have den rigtige Betydning, at
jeg ikke kan gjøre Andre godt, naar jeg selv ikke har
noget Godt, at jeg derfor maa besidde for at kunne
meddele.

Heraf ses blandt Andet, at man har havt Uret i at
mene, at det skulde være Oppositionen mod Religionens
Lære, der senere førte Feuerbach til at hylde Egoismen
som eneste Moralprincip. Han ser mere og mere Selvkjærligheden
 som den eneste virkende Kraft i Religionen, og det er den vaagnende Sympathi med dette
pathologiske Element i Menneskenaturen, der
% --- p. 78 ---
\setcounter{footnote}{0}%
\opage{78}andrer hans ethiske Anskuelse. Man kan i Feuerbachs
Udvikling parallelt spore en Fremhæven af den pathologiske
Faktor som Religionens Kilde og af Egoismen
som Moralens Princip, idet som han selv udtrykker sig,
det Væsentlige i Religionen mere og mere bliver det
Væsentlige for ham selv.

De tidligere religionsphilosophiske Skrifter beskjæftigede
sig udelukkende med Christendommen, den høieste
Religionsform, hvor Troens Gjenstand var et over
Naturen absolut hævet Væsen. Der opstaaer nu det
Spørgsmaal, hvorledes den religiøse Bevidsthed har kunnet
udvikle sig til et saadant Høidepunkt, hvorledes den
har faaet Kraft og Midler til at sætte et supranaturalt
Væsen udenfor sig. Dette Spørgsmaal besvarer Feuerbach
i Afhandlingen „das Wesen der Religion“ (1845).
Han fører her Religionens Væsen tilbage til den Schleiermacherske
Bestemmelse, Afhængighedsfølelsen. Det,
hvoraf Mennesket er og føler sig afhængigt, er oprindelig
ikke Andet end Naturen, som derfor er Religionens
egentlige Gjenstand. Mennesket føler, at han ikke skylder
sig selv sin Existents og ærer saa som Gud, hvad
han føler sig afhængig af. Saa stor Pris han sætter
paa sin Existents, saa stor Pris sætter han paa Gud.
Og det er Naturen, der er Menneskets Livskilde; al
god Gave kommer ikke ovenfra, men nedenfra, fra Naturens
Dyb. Det Høiere forudsætter det Lavere, ikke
omvendt, og jo høiere et Væsen er, des mere forudsætter
det. Det mest sammensatte, mest afhængige Væsen er
det høieste. Det individualiserede Væsen er et mere
guddommeligt Væsen end det individualitetsløse; men
det er bundet til sine Betingelser, hører op med dem,
ligesom det begynder med dem. Menneskets Grundvold,
den evige Natur bestaaer under en uendelig Opstaaen
og Forgaaen; den alene kan ikke afledes af noget Høiere.
Gud betragtet som objektivt virkeligt Væsen er selv ikke
Andet end Naturens Væsen, som ved Abstraktionens
% --- p. 79 ---
\setcounter{footnote}{0}%
\opage{79}Hjælp er adskilt fra den virkelige, sandselige Natur, og
ved Phantasiens Hjælp forvandlet til et menneskeligt
eller menneskelignende Væsen. Gud forudsætter derfor
Naturen, som det Abstrakte forudsætter det Konkrete.
At ville udlede Naturen af Gud er at ville udlede Naturens
virkelige Væsen af dens abstrakte, kun i Tanken
existerende Væsen. Grunden til, at man stadig forlanger
og forsøger en saadan Deduktion er, at Tanken uvilkaarlig
gjør det Abstrakte og Almene, som staaer den
nærmere og for den er det Høiere, til Forudsætning for
det Konkrete og Virkelige, skjøndt Forholdet i Virkeligheden
er det omvendte.

Hvor høit den religiøse Tro end hæver sig over
Naturen, vedbliver denne dog at være dens egentlige
Gjenstand. Paa Grund af sine uopfyldte Ønsker føler
Mennesket sig afhængigt, og den nødvendige Naturmagt,
der overalt træder det imøde, snart med Imødekommen,
snart som Skranke, opfattes saa som den frie, vilkaarlige
Magt, der giver eller negter Ønskets Gjenstand.
Mennesket lægger sit eget Gemyts- og Villesliv ind i
Naturen. Dette er det første Trin, den egentlige Naturreligion.
Naar Mennesket begynder at adskille sit eget
Væsen fra Naturen, at se sit sande Liv ikke som et
umiddelbart Naturliv, men som et menneskeligt Samfundsliv,
naar han altsaa hæver sig over Naturen i sin Forstand
og sin Villie, saa overfører han ogsaa de høiere
intellektuelle og moralske Egenskaber paa Guddommen
og adskiller derved mere og mere denne fra Naturen.
Denne Adskillelse gjennemføres, hvor som i Christendommen
den religiøse Tro udspringer af Ønsket om en
ubegrændset Salighed, et Ønske, der fuldstændig overspringer
de naturlige Skranker, indenfor hvilke Hedningenes
Ønsker bevægede sig. Grækerne gjorde ikke det
guddommelige ɔ: det mulige Væsen til Forbillede og
Maalestok for det virkelige, men det virkelige Væsen til
Maalestok for det mulige. Men de Christnes Gud er
% --- p. 80 ---
\setcounter{footnote}{0}%
\opage{80}transscendent ligesom deres Ønske. Aandsreligionen foretager
den modsatte Bevægelse af Naturreligionen. Denne
gjør de virkelige Væsener til forestillede Væsener, til
Væsener for Indbildningskraften, idet den er Tro paa
Naturen som et menneskeligt Væsen; Aandsreligionen
gjør forestillede Væsener, Tankens og Phantasiens Gjenstande
til virkelige Væsener, idet den er Tro paa det
menneskelige Væsen som Naturens Væsen.

Her fremtræder tydelig Feuerbachs Sympathi med
de positive Religioners reale Grundlag. Han lægger nu
ikke længere det Ethiske og Almene til Grund ved Bedømmelsen
af dem, men gjør Naturen til Maalestok og
fordrer, at de i den satte Grændser skulle respekteres.
Han føres herved over i en Naturalisme, der i sin videre
Udvikling vilde kunne føre tilbage til hans ældre pantheistiske
Anskuelse, hvis han ikke ved den Maade,
hvorpaa han erklærer Aanden for et blot Afledet, en Abstraktion,
der kun har subjektiv Betydning, snarere nærmede
sig en fuldstændig Materialisme. Det Interessante
ved hans sidste Skrifter er nu Brydningen mellem denne
materialistiske Tendents (en naturlig Konsekvents af den
Vægt, han mere og mere lagde paa det Reale) og den
Energi, hvormed han fastholder sit anthropologiske Standpunkt
og dermed sit oprindelige ideelle Udgangspunkt.
Efter disse Skrifter\footnote[1]{Die Unsterblichkeitsfrage vom Standpunkte der Anthropologie (1846), Vorlesungen über das Wesen der Religion (1848), Theogonie (1857), Gottheit, Freiheit und Unsterblichkeit (1866).}%
 ville vi nu fremdrage Hovedtrækkene,
dels af hans almindelige philosophiske, dels særlig
af hans ethiske og religionsphilosophiske Anskuelse, saaledes
som den former sig ved hans Virksomheds Afslutning.


Det er først Opfattelsen af Forholdet mellem det
Almene og det Individuelle, der tildrager sig Opmærksomheden.
Det existerende Væsen er stedse kun Individet,
% --- p. 81 ---
\setcounter{footnote}{0}%
\opage{81}Arten er kun Prædikat. Men Tænkningen vender Forholdet
om, gjør Subjektet til Prædikat, Prædikatet til
Subjekt. Tanken fremstiller det Almene, der forekommer
hos alle Enkelte, som en Sag for sig selv og sætter
som et Kontinuum, hvad der i Virkeligheden er et Diskret,
Livets uendelige „Mangegange“ som et identisk
„Engang“. Det første Skridt til Visdom er at erkjende
denne væsentlige og uovervindelige Modsætning mellem
Tænkningen og Virkeligheden.

Da Feuerbach tidligere satte Selvbevidstheden i nøie
Forbindelse med Slægtens almene Væsen, der gjennem
den blev Gjenstand for Individet, maatte naturligvis nu
Opfattelsen af Bevidstheden modificeres ved den forandrede
Opfattelse af det Almene. Individualitetens Hemmelighed,
siger han, er Enheden af Jeg og Ikke-Jeg, og
denne Enhed er ogsaa Religionens Grund. Uden Ikke-Jeg
vilde Mennesket ingen Religion have, da han saa
selv vilde være Gud; uden Jeg eller som et Jeg, der
ikke adskilte sig fra Ikke-Jeget, vilde han kun være en
Plante eller et Dyr. Jo dybere Mennesket gaaer i sig
selv, des mere ser han Forskjellen mellem Natur og
Menneske forsvinde, des mere erkjender han, at han
kun er et bevidst Bevidstløst, et Ikke-Jeg, som er Jeg.
Derfor er Mennesket det allerdybeste og dybsindigste
Væsen. Men Mennesket kan ikke udholde eller begribe
sin egen Dybde og splitter derfor sit Væsen i et Jeg
uden Ikke-Jeg, som han kalder Gud, og et Ikke-Jeg
uden Jeg, som han kalder Natur.

Naar det Almene kun er som Prædikat for det Individuelle,
det Bevidste kun som Prædikat for det Bevidstløse,
kunde det synes, som om Feuerbach allerede
havde overskredet Materialismens Grændse. Men dog
protesterer han mod „den extravagante og transscendente
Materialisme“, der gaaer ud over Menneskets Væsen for
at forklare det af rent ydre Aarsager (Noget Feuerbach
unegtelig selv nærmer sig i sin Sætning: „der Mensch
% --- p. 82 ---
\setcounter{footnote}{0}%
\opage{82}ist, was er isst“). Han vil fastholde „et Archimedisk
Punkt mellem Materialisme og Spiritualisme“, en „immanent
Materialisme“, der holder sig indenfor det Anthropologiske,
og han hævder en umiddelbar Enhed af
Sjæl og Legeme, som ikke lader Noget ligge imellem og
ikke giver Plads for nogen Adskillelse eller Modsætning
mellem Materielt og Immaterielt. Denne Enhed finder
han koncentreret i Hjernen som punctum saliens, hvor
Legemet er Aand, Aanden Legeme. Denne uopløselige
Umiddelbarhed er Livets uantastelige Eiendom, ja
er Livets Væsen selv i Modsætning til Tænken og
Viden. Forgjeves søger Idealismen ud fra sit theoretiske
Standpunkt at løse Problemet om den objektive
Værens Realitet, medens den forkaster Fornemmelsen
som noget blot Subjektivt. Fornemmelsen er vel subjektiv,
men dens Grund er objektiv, og Identiteten af
Subjekt og Objekt aabenbarer sig for os i Følelsen af
Velvære. Forsaavidt har Religionen Ret, naar den udleder
Verdens Virkelighed ikke af Forstanden, men af
Villien.

Ved disse Bestemmelser har Feuerbach yderligere
udviklet, hvad vi ovenfor betegnede som hans Paradox:
Værens Uudsigelighed og Ubegribelighed. Væren opfattes
abstrakt individualistisk, Væsenet bliver et usynligt
Punkt, der ikke kan angives i Theorien, og Modsætningen
mellem Tanken og Livet lader kun en praktisk,
en religiøs Løsning aaben. Derfor ender Feuerbach
med at sige: „Meine Philosophie ist keine Philosophie.“

Men det er kun indirekte, at han kan betragtes som
Stamfaderen til den moderne Materialisme. I den energiske
Fastholden af det anthropologiske Synspunkt overfor
den transscendente Materialisme ligger der et idealistisk
Princip skjult, som han aldrig opgiver. Herpaa beroer
ogsaa den afgjørende Modsætning mellem Feuerbach
og hans Samtidige, Auguste Comte i Frankrig, med hvem
han ellers har Meget tilfælles, navnlig i Betragtningen
% --- p. 83 ---
\setcounter{footnote}{0}%
\opage{83}af Menneskeaandens Udviklingsstadier. Parallelt med Comtes
tre Stadier: Theologi --- Metaphysik --- positiv Viden
gaa Feuerbachs tre: Religion --- Spekulation (Pantheisme)
--- Anthropologi. Men Comtes „positive Viden“ er
netop,  hvad Feuerbach kalder transscendent, idet den
tager Udgangspunktet udenfor Mennesket og forklarer
dettes Væsen rent af ydre Aarsager\footnote[1]{Smlgn. om Comte „Nyt dansk Maanedsskrift“. 2det Bind. p. 200---208.}.

I „Theogonie“ beskjæftiger Feuerbach sig paa Ny
med Undersøgelsen af Religionens Oprindelse og gjennemfører
i en Mangfoldighed af Vendinger og Exempler,
hentede „fra den klassiske, hebraiske og christelige Oldtids
Kilder“, den allerede tidligere opstillede Lære, at
Menneskets Forhold til Gud er hans Forhold til sig selv,
ikke i den Forstand, hvori dette blev opfattet i „Wesen
des Christenthums“, som Menneskets Forhold til sit almene
Væsen, til Slægtens Ide, men som et Forhold til
sin egen Selvkjærlighed. Religionen er ikke mere den indirekte
Selvbevidsthed, men den indirekte Selvkjærlighed.
Ønsket er det theogoniske Princip: det er Udtryk for en
Mangel, en Skranke, en Lidelse, men et heftigt, et revolutionært
Udtryk, en Villen, at denne Mangel og Skranke
ikke skal være. I og med Ønsket opstaaer derfor Forestillingen
om en Gud, ɔ: et Væsen, der er hævet over denne
Skranke og kan hæve Mennesket over den. Hvis man
vil sige, at ikke Ønsket om Lyksalighed, men Ønsket
om at blive god, er Grunden til Religion, da svarer
Feuerbach, at Religionen netop med Rette grunder Dyden
paa Saligheden. Mennesket kan ikke være godt, naar
det ikke er lykkeligt. Rent theoretisk kan man gjøre
Lyksaligheden afhængig af Dyden, men i Livet gaaer
det anderledes til, hvor det er følende, trængende, attraaende
Væsener, der optræde. Den, som derfor vil
gjøre Menneskene bedre, maa først gjøre dem lykkeligere.
% --- p. 84 ---
\setcounter{footnote}{0}%
\opage{84}Religionen er selv den ældste Kulturform og har i sig
samme Formaal som Kulturen, men vil opnaa det uden
at bruge Midlerne.

Her viser det sig, at Feuerbach i Principet er enig
med Religionen, saaledes som han opfatter den. Ligesom
Religionen efter hans Mening, saaledes bygger han
nu selv fuldstændig det Ethiske paa Egoismen\footnote[1]{Feuerbachs Udviklingsgang udtaler sig i følgende, hverandre afløsende Sætninger: a) „Wesen des Christenthums“: Christendommens Fordærvelighed viser sig i, at den bygger Kjærligheden paa Troen. b) En Aphorisme (samtidig med „Phil. d. Zuk.“): Kjærligheden til Andre maa anerkjende deres Egoisme som berettiget. „Wesen des Glaubens“: Det kan have en rigtig Mening at udlede Kjærlighed af Tro (Kjærlighed til Andre af Kjærlighed til sig selv). c) „Theogonie“: Ingen Dyd uden Lyksalighed; derfor har Religionen Ret, og Egoismen er Moralprincip.}, og i
sin heftige Polemik mod enhver Ethik, der ikke gaaer
ud fra den, kan han da fra sit Standpunkt støtte sig
til Religionen som faktisk Vidnesbyrd for sin Lære,
ligesom han ovenfor gjorde det ved Spørgsmaalet om
den objektive Værens Realitet. Hvor forunderlig har
Stillingen ikke vexlet under Feuerbachs Udviklingsgang!
Han begyndte med at angribe Religionen ud fra Philosophiens
Standpunkt; derpaa bekæmpede han Philosophien,
fordi den kun var Religionens abstrakte Gjenbillede;
og tilsidst støtter han sig til Religionen i sin
Polemik mod Philosophien.

I sin „Pierre Bayle“ havde Feuerbach talt med Begeistring
om Kants og Fichtes Ethik. Nu forkaster han
den med Foragt. Pligten kan ikke gjøres til Grund
og Aarsag; den er en Følge, en Virkning, en Drift, der
gjør sin Ret gjældende mod en anden Drifts Herskesyge.
Den onde Samvittighed er den vrede Skygge af
en Drift, som en anden og stærkere Drift voldsomt har
bragt af Dage eller truet dermed. Den væsentlige Drift,
% --- p. 85 ---
\setcounter{footnote}{0}%
\opage{85}der ligger til Grund for alle andre Drifter, er Lyksalighedsdriften,
med hvilken Villien er Et. At ville, vil
sige, ikke at ville lide. Herved ophæves Retten og Moralen
ikke; de sættes kun paa deres naturlige Fødder.
I selve sin Egoisme har Mennesket Kriteriet for, hvad
der er Ret eller Uret. Tyven selv vil, at hans egen
Eiendom skal respekteres, Morderen, at hans Liv skal
være helligt. Modsigelsen mellem Forbryderens Villie
med Hensyn til ham selv og hans Handlemaade overfor
Andre er den indre Grund til Samvittigheden, til Bevidstheden
om Ret og Uret. Jeg kan ikke fordre min Eiendom,
min Egoisme anerkjendt, naar jeg ikke anerkjender
Andres. Og i hvert Tilfælde vil de Andres Egoisme
nok holde min Egoisme indenfor de rette Grændser; selv
om det er mig ligegyldigt, om jeg er god eller slet, er
det dog ikke de Andre ligegyldigt. Det Gode er, hvad
der svarer til alle Menneskers Egoisme, det Onde, hvad
der tilfredsstiller nogle Enkeltes Egoisme paa Andres
Bekostning.

I denne sidste Sætning antydes der dog et Princip,
som fører os igjen ud over Egoismens Moral. „Alle
Mennesker“ er nemlig kun et atomistisk Udtryk for,
hvad Feuerbach tidligere kaldte „Slægten“. Hermed
stemme ogsaa hans Udtalelser i praktisk-social Retning,
som ingenlunde kunne begrundes ved den blotte Egoismes
Princip. Leilighedsvis taler han f. Ex. Proletariernes
Sag. Han viser, at det er falsk, naar man siger,
at de Intet have og derfor søge at opnaa Alt; de have
tvertimod saare Meget, nemlig menneskelig Selvfølelse,
Dannelsesdrift og Arbeidslyst. Hvad de forlange, er, at
Arbeidets Værd skal anerkjendes, at Arbeideren ikke
skal betragtes som et blot Middel, men som et selvstændigt
Øiemed, der har Gyldighed i og for sig selv. Men
Feuerbachs Praxis kommer endnu mere i Strid med hans
Theori, naar han overfor ensidige politiske Anskuelser,
der mene, at Alt beroer paa politiske og sociale Refor%
% --- p. 86 ---
\setcounter{footnote}{0}%
\opage{86}mer og Forbedringer af Menneskets materielle Vilkaar,
eller, som Feuerbach udtrykker det, at Ondet ikke sidder
i Menneskehedens Hoved eller Hjerte, men i dens
Mave, naar han overfor saadanne Anskuelser hævder,
at der jo gives mange Mavesygdomme, der have
deres Grund i Hovedet, og erklærer, at han nu engung % PRINTED AS IS (engung: for engang)
særlig har gjort sig Udgrundelsen og Helbredelsen af
Menneskehedens Hoved- og Hjertesygdomme til Opgave.
Ligeledes erklærer han, at den, der endnu er Slave af
Fordomme, ikke har nogen Ret til at befri sig fra det
materielle og det politiske Tryk. Det var ogsaa Grunden
til, at han afholdt sig fra direkte Deltagelse i Revolutionen
1848; den udsprang nemlig af Troen paa
Ideer og Institutioner, fra hvilke Menneskeslægtens
Bevidsthed maa befries, førend der kan være Tale om
praktisk og politisk Befrielse. Feuerbach godtgjør her,
hvad der forøvrig er Motivet til hele hans Virksomhed
med alle de vexlende Standpunkter, en Interesse for
Sandhed og aandelig Frihed, som ingenlunde kan udledes
af Egoisme.

Saaledes finde vi os endnu paa dette sidste Punkt,
til hvilket vi kunne følge Feuerbach, stillede mellem
dobbeltsidige Indtryk. Denne forunderlige Skikkelse
lader sig ikke fange under nogen bestemt Form; altid
er han i Modsigelse med sig selv og derfor altid ifærd
med at udvikle sig, aldrig afsluttet. Han tumler
med Problemerne, men siger aldrig det sidste Ord og
ender egentlig med at erklære deres Løsning for umulig.
Oprindelig udgaaet fra den absolute Idealisme
har han faaet Øie paa Realismens Betydning og kan nu
ikke finde de nødvendige Mellembestemmelser, men hengiver
sig med umiddelbar Begeistring og uden ganske
at opgive Forudsætningerne fra sit forrige Standpunkt.
Naar han f. Ex. siger, at det Enkelte er uudsigeligt,
saa er det Hegels egen Sætning fra Phænomenologien,
han adopterer, kun at han erklærer det for den sande
% --- p. 87 ---
\setcounter{footnote}{0}%
\opage{87}Værens positive Udtryk, hvad Hegel netop saa som noget
Negativt, som et Udtryk for Ikke-Væren. Derfor er
Feuerbach kun naaet til at blive Realismens Prophet
og Mystiker.

Feuerbachs Hovedfortjeneste ligger paa det religionsphilosophiske og psychologiske Omraade. Her ere hans
Udviklinger af epochegjørende Betydning for Forklaringen
af Religionernes Oprindelse i Menneskeaanden. Men
der er dog ved Feuerbachs psychologiske Undersøgelser
den Mangel, at de ikke føre ud over det Konkrete og
Empiriske; hans Styrke er, hvad han kalder den dramatiske
Psychologi, ɔ: Psychologien i nøie Forbindelse
med de Gjenstande, i hvilke Menneskesjælen aabenbarer
sig i sin Totalitet. En dybere Begrundelse af disse
konkrete Bevægelser, som Feuerbach iøvrig skildrer med
saa stort Mesterskab, afskjærer han sig allerede derved
fra at give, at han gjør selve Psychologien eller Anthropologien
til Grundvidenskab. Det anthropologiske
Standpunkt hindrer ham vel i at standse i en ren Sensualisme
eller Materialisme, men kun ved at gjøre hans
Lære fuldstændig subjektivistisk i sidste Instants. Hertil
er han naaet gjennem Subjektiveringsmethoden: Fornuften
er Guds Subjekt, og Mennesket er Fornuftens
Subjekt. Men hvilken Betydning disse subjektive Former
have i sig selv og i Tilværelsen i det Hele, undersøges
ikke.

Denne Feuerbachs empiriske og subjektivistiske Individualisme
hindrer ham i at vinde en virkelig historisk
Anskuelse. Det Historiske beroer paa Vexelvirkningen
mellem det Almene, de objektive Formaal og Ideer og
de konkrete, individuelle Væsener, der søge at tilegne
sig dem og virkeliggjøre dem i sig og udenfor sig. Men
den ene af disse Faktorer mangler hos Feuerbach. Han
er tilfreds med det psychologiske Hvorledes og bekymrer
sig ikke om det metaphysiske og historiske Hvad.
Heri danner han en fuldstændig Modsætning til Hegel,
% --- p. 88 ---
\setcounter{footnote}{0}%
\opage{88}der i Historien kun tager Hensyn til de almene Ideer,
medens det handlende Individ og den hele psychologiske
Betragtning er ham ligegyldig og værdiløs. Dog
staaer Hegels Anskuelse forsaavidt Historien nærmere,
som den formaaer at tilegne sig dennes Indhold i dets
indre Rigdom og Forskjellighed, medens Feuerbach bindes
til en og samme Kreds af Begreber, da de psychologiske
Faktorer væsentlig ere de samme paa ethvert
Punkt af Historien. Saaledes bliver for ham alle positive
Religioner i Grunden et og samme Væsen, hvis
indre Forskjelligheder blive betydningsløse, fordi han
overfor det pathologiske Element i Religionens Væsen
overser de andre vigtige Elementer i dette, som vise ud
over Subjektets rent individuelle Drift og godtgjøre, at
der ogsaa gjennem disse Aandsformer arbeider sig et
Høiere frem. Saaledes lader han hverken det rationelle
Element ei heller Resignationens Element i Religionen
komme til sin Ret. Og det har ikke blot en mislig Indflydelse
paa hans Opfattelse af de historiske Religioner,
men ogsaa paa hans Opfattelse af selve Menneskets Væsen
saaledes som dette gjennem dem lægger sig for Dagen.

Feuerbach har selv godtgjort sit individualistiske
Princips Sandhed ved at paatrykke sin Anskuelse og
sine Skrifter sin egen Individualitets eiendommelige Præg.
For at være enig med ham maatte man derfor egentlig
ogsaa ligne ham i psychologisk Henseende, og det er
ikke forunderligt, om den objektive Kritik maa reise
mange Indvendinger overfor ham. Men dog skilles vi
fra ham med den Overbevisning, at vi have beskjæftiget
os med en af de faa virkelig originale Tænkere efter
Hegel. ---

6. \emph{Max Stirner}. --- Feuerbach havde drevet Oppositionen
mod den spekulative Philosophi saavidt, at han
tilsidst erklærede den sande Philosophi for ingen Philosophi
at være. Og dog er der mere Spekulation i ham
end han selv vil være ved. Hvad den spekulative Tanke
% --- p. 89 ---
\setcounter{footnote}{0}%
\opage{89}til enhver Tid søger, Enheden i Tilværelsens Modsætninger,
det er ogsaa Feuerbachs Opgave, skjøndt han
undertiden paa paradox Vis afskjærer Muligheden af
dens Løsning. Og selv naar han gjør dette, fastholder
han udtrykkelig Modsætningens Led, om end i ensidig
Form; det er kun Udtrykket for deres Enhed og Sammenhæng,
han mistvivler om at finde.

Hvis den formelle Konsekvents kunde være Maalestok
ved Bedømmelsen af philosophiske Anskuelser, indeholder
Feuerbach Modsigelser nok. Og disse Modsigelser
fremtræde des mere grelt, jo mere han af de virkelige
og formentlige Modsigelser, han opdager i de tidligere
Anskuelser, tager Anledning til lidenskabelige, ofte haanende
Angreb, og jo mere han selv mener at have aabnet
Indgangen til en hel ny Erkjendelse, „toto genere“
forskjellig fra den tidligere. Forsaavidt er det en retfærdig
Nemesis, naar det Vaaben, han selv saa lidenskabelig
rettede mod Andre, ogsaa vendes mod ham selv,
naar han, der mente at være den mest radikale Revolutionær,
nu, ligesom de hverandre afløsende Partier under
den franske Revolution, selv overfløies af endnu mere
yderliggaaende Anskuelser, af hvilke han ligefrem stilles
sammen med den gamle Verdensanskuelse, han mente
at have styrtet. Men alligevel standser den egentlige
philosophiske Bevægelse ved Feuerbach. Den formelle
Konsekvents, der tilsidst endog fremtræder blot i Experimentets
Form, kan ikke fortjene Navn af philosophisk
Tænkning, uagtet det har sin philosophiske Interesse at
se, hvorledes den absolute Idealisme tilsidst ender i den
rene Formalisme.

Skriftet „Der Einzige und sein Eigenthum“ (1845)
af Pseudonymen Max Stirner (den som socialøkonomisk
Forfatter bekjendte Dr. Schmidt, der døde i Berlin i
Begyndelsen af Aarene 60) var for det 19de Aarhundrede,
hvad Helvetius's Bog „de l’esprit“ var for det
18de: den mest hensynsløse og uforbeholdne Forkyndelse
% --- p. 90 ---
\setcounter{footnote}{0}%
\opage{90}af Egoismen som eneste virkelig og berettiget Bevæggrund
til Handlen. Men medens Helvetius støtter sig
til Materialismen, er Stirner forsaavidt Idealist, som han
gjennemfører den Feuerbachske Subjektiveringsmethode
ud over det Punkt, hvor Feuerbach selv standsede\footnote[1]{Vel retter Stirner nærmest sin Polemik mod „Wesen des Christenthums“, men hans Indsigelser ramme ogsaa de Skrifter af Feuerbach, der høre til hans tredie Stadium, og som ogsaa tildels allerede forelaa, da Stirners Bog udkom.}.

Barnet er Realist, Ynglingen Idealist og Manden
Egoist. Paa samme Maade gaaer det med Menneskehedens
Historie som Helhed. De Gamle saa Sandheden
i Verden og lod deres eget Jeg bøie sig for Verdensforholdene,
hvilke imidlertid deres Stræben gik ud paa
betvinge og opløse. Denne Opløsnings Fuldendelse er % PRINTED AS IS (betvinge: the image, at lost before betvinge; O055)
Christendommen, hvor Verden er et Intet, Aanden den
eneste Sandhed. I den nyere Tid gaaer Bevægelsen ud
paa at opløse Aandens Sandhed. Som Verden for de
Gamle, saaledes er Aanden for de Moderne „eine Wahrheit,
hinter deren Unwahrheit sie zu kommen suchen und
endlich wirklick kommen.“ Denne Opløsning har nærmet % PRINTED AS IS (wirklick: for wirklich)
sig sin Fuldendelse i den nyeste Philosophi, hvor
Gud er erkjendt som Menneskets eget Væsen. Men
ogsaa kun nærmet sig sin Fuldendelse. Thi selv Feuerbach
og Bruno Bauer staa endnu i Grunden paa det
theologiske Standpunkt. Hvad nytter det nemlig, at
Feuerbach siger, at det Hinsides kun er vort eget Væsen,
naar han saa gjør vort egen Væsen til Gud og sætter % PRINTED AS IS (vort egen: the image, egen for eget; O056)
det i en saadan Modsætning til os, at det bliver til et
nyt Hinsides, en ny Himmel indeni os? Vi ere ikke
det, som er i os, ligesaalidt som vi ere det, der er udenfor
os. Feuerbach spalter os i et væsentligt og et uvæsenligt
Jeg og støder os derved tilbage i Dualismen. Gud
vedbliver at være Gud, selv om man forjager ham fra
Himlen og indeslutter ham i det menneskelige Bryst.
% --- p. 91 ---
\setcounter{footnote}{0}%
\opage{91}Feuerbach repræsenterer derfor egentlig den sidste krampagtige
Anstrengelse for at fastholde Christendommen;
han er, hvad han selv sagde Schleiermacher at være,
Christendommens sidste Theolog.

Det sidste Skridt, Transscendentsens fuldkomne Ophævelse
er endnu ikke gjort, saalænge der endnu anerkjendes
et Alment, hvorover Individet ikke raader som
sin Eiendom. Alt Helligt er noget Fremmed. Dersom
du troer paa Sandheden, er du jo Sandhedens Tjener,
tilhører ikke dig selv. De almene Formaal lægge Beslag
paa mig for deres egen Skyld, ere altsaa Egoister;
men er min Egoisme ikke ligesaa berettiget? Lad os
da vende Forholdet om og gjøre det Almene til vor
Tjener! Den politiske Liberalisme arbeider for Frihed
og Ret, den sociale Liberalisme (Kommunismen) for
Samfundet; men Frihed, Ret og Samfund ere kun Chimærer.
Jeg gjør Alt for min egen Skyld. Kun jeg
selv bestemmer, hvad der skal tilhøre mig. Grib til!
Krig af Alle mod Alle er Løsenet. De, der tale om
Kjærlighed, ere besatte, have en Skrue løs. Thi Kjærligheden
gjør mig afhængig, jeg besiddes ligesaameget
som jeg besidder; Kjærlighedens Gjenstand kan ikke
absolut være Et med mig.

Naar enhver Myndighed baade i det Indre og det Ydre,
i Himlen og paa Jorden er fældet, og „den Eneste“ ved
at forvandle Alt til „sin Eiendom“ selv er bleven den
eneste Realitet, saa forstaaes det, at han kan udslynge
en Sætning som denne: „Ich begehe getrost die Sünde
wider den heiligen Geist!“ Bevægelsen er fuldbyrdet; „den
Eneste“ har gjennemskuet Aandens Usandhed.

Derfor synes en Kritik heller ikke vel mulig her;
thi naar Sandheden er hans Eiendom, kan han ikke
bindes ved nogen almen Grundsætning. Med Rette har
man kaldt Stirners Bog den mest radikale, som nogensinde
er skreven. Den er radikal ogsaa i den Forstand,
% --- p. 92 ---
\setcounter{footnote}{0}%
\opage{92}at den synes at ville oprykke Sandheden med Rode, for
at denne ikke skal bruges som Vaaben imod den.

Dog giver Max Stirner humoristisk sig selv en praktisk
Dementi ved foran Bogen at sætte følgende Tilegnelse:
„Meinem Liebchen, Marie Dähnhardt.“ Den
viser, at selv om det egoistiske Standpunkt formelt kan
gjennemføres i Theorien, er det dog umuligt i Praxis.

En theoretisk Blottelse giver Stirner sig dernæst
ved stedse at fremstille Egoismen som en ny Moral,
en ny Sandhed, der skal afløse den gamle. Det tager
sig forunderlig ud i samme Øieblik, som al Sandhed
skal være udryddet! Men derfor kunde ogsaa selv Stirner
endnu opnaa at blive regnet til Høire.

7. „\emph{Das Verstandesthum und das Individuum}“.
(1846). Dette anonyme Skrifs Forfatter var % PRINTED AS IS (Skrifs: for Skrifts?)
efter Nogle en Geistlig fra Köthen ved Navn Karl Schmidt,
der har villet vise, til hvilke Resultater den „kritiske
Bevidsthed“ tilsidst maatte komme, efterat den engang
har givet sig den formelle Konsekvents's dæmoniske Magt
i Vold. Efter ved Uddrag af de forskjellige Kritikeres
Skrifter at have vist, med hvilken rivende Fart Strauss,
Feuerbach, Brødrene Bauer og Max Stirner afløse og
gjensidig omstyrte hverandre, optages den kritiske Bevægelse
ved den sidstnævnte Forfatter, der selv opstiller
et Ideal i det Øieblik, han mener at have omstyrtet
alle Idealer. En Verden af Egoister skal nu
være den fuldkomne Verden. Men saa er Stirner ikke
længer „der Einzige“; han er En af Mange, ikke det
absolute Individ. Alt hvad der ender paa --- ist er
ikke individuelt.

Imod Max Stirners „Einzige“ stilles saa Individet,
det absolute Individ, der er frit for enhver Rest af Forstandsvæsen.
Men Individet finder sig existerende i en
dobbelt Verden, en physisk og en psychisk, en Naturens
og en Aandens Verden. Det første Spørgsmaal er da,
hvorledes Individet stiller sig til disse Verdener.

% --- p. 93 ---
\setcounter{footnote}{0}%
\opage{93}I Virkeligheden existerer kun det Enkelte. Naturen
som Helhed er et Tankefoster, som opstaaer ved, at
Naturforskerne se alt Enkelt i Sammenhæng og bringe
det ind under almene Love; derved blive de Dogmatikere.
„Individet“ derimod tænker ikke de Enkelte, Atomerne;
derved vilde det gjøre dem almene; --- det anskuer
dem, stirrer paa dem og lader ved denne Stirren
paa det Enkelte Verden, Helheden og Sammenhængen
forgaa.

Aanden er det afblegede, det blegsottige Individ,
det matte Graat i Graat istedetfor Virkelighedens brogede
Mangfoldighed. I „Aandens Rige“ gjælde kun de
almene Magter, ikke Individet. „Individet“ derimod ser alle Menneskeverdenens Indretninger og Former kun som
Menneskenes krystalliserede Forestillinger, deres fixe
Ideer; derfor have de ingen Gyldighed for det. Dog
gaaer det ind paa disse Former, det vil sige, leger med
dem, nyder dem, saalænge det behager. Aandens Opløsning
vil sige Opløsning af al Livsalvor, alle Lidenskaber,
al Begeistring, al Pathos.

Men hvad er da Individet selv? Det kan ikke udsiges,
da Sproget kun kan udtrykke det Almene. Individet
kan kun betegnes ved Gestus, ved at pege. Individet
er inkommensurabelt, har intet andet Maal end
sig selv. Individet er Intet og er Alt; det er en Proteus,
der stedse skifter Skikkelse. Det bryder stedse
med sin Fortid og gjør dog ingen Fremskridt. I saadanne
og utallige andre Vendinger vilde man kunne karakterisere
„Individet“, og dog vilde man ikke komme
Sagen nærmere. Karakteristiken maa bestandig tilbagekaldes
og korrigeres --- og tilsidst afbryder „Individet“
dog spottende den hele Bestræbelse med de Ord: „Ich
bin ich selbst allein!“ ---

Den tydske Spekulation, der begyndte med Jeget
som Monade, af hvilken Tilværelsen uendelige Indhold % PRINTED AS IS (Tilværelsen uendelige: the image, for Tilværelsens; O060)
udviklede sig, ender saaledes i Jeget som det absolute
% --- p. 94 ---
\setcounter{footnote}{0}%
\opage{94}Atom, der umuliggjør enhver Bestemmelse, ethvert Indhold.
Men de sidste Konsekventser ere her rigtignok
af den Natur, at de kun kunde naaes ad Tankeexperimentets
Vei.

\chaphere{TREDIE KAPITEL.}


Grundene til Materialismens Fremkomst i vor Tid
og den store Tilslutning, den har fundet, ere af forskjellig,
dels negativ dels positiv Beskaffenhed. Den storartede
philosophiske Begeistring, som havde ført Bevægelsen
ud over dens sande Grændser, maatte nødvendigvis
medføre dels en Slappelse dels en Reaktion, der
vendte sig skeptisk og polemisk mod ethvert egentlig
philosophisk Forsøg. Philosophien deler her kun Skjæbne
med de andre aandelige Magter, der i Aarhundredets
Midte synes at ligge i Dvale efter det kraftige Liv i
dets Begyndelse. Tidsalderen kastede sig med al sin
Energi ind i det ydre Livs Strømninger og vendte foreløbig
de theoretiske og æsthetiske Idealer Ryggen, af
hvilke den foregaaende Slægt havde været begeistret.
I sin Hensynsløshed synes den endog ogsaa at vende
de ethiske Idealer Ryggen. I ethvert Tilfælde tilfredsstiller
den ikke som sin Forgænger den religiøse Trang
ad den indre, frie Tænknings og Erfarings Vei; den vil
ogsaa her være „praktisk“, og heraf forstaaes det, at den
mest udvortes Bogstavtro deler Herredømmet i Mængdens
Bevidsthed med Materialismen.

Men ogsaa Philosophien selv var gjennem den absolute
Idealismes Selvopløsning kommen til et saadant
Extrem, at Materialismen kunde synes den eneste Mulighed,
der endnu stod aaben. Skridt for Skridt havde det
% --- p. 95 ---
\setcounter{footnote}{0}%
\opage{95}Almene, Ideen mistet Terrain og var det reale Enkelte
traadt i Forgrunden. Tilsidst stod dette som et uopløseligt
Atom, medens Tanken var opløst i et Intet. Hvad
var der da at gjøre uden at tage sit Udgangspunkt i
de reale Atomer og i deres Sammensætning finde en
Erstatning for den tabte Ide?

Materialisterne beraabe sig ofte paa Feuerbach. Han
havde villet vise, at Mennesket i Religion og Spekulation
kun lægger sit eget Væsens Former over i det Objektive,
danner dette efter sig. Altsaa, slutte Materialisterne,
kan det ikke være Tanken, Fornuften, der
raader i Tilværelsen. Tanken opstaaer kun i Mennesket
som materielle Aarsagers Produkt. Feuerbach fastholdt
den menneskelige Tanke som det eneste mulige
Standpunkt, der ikke igjen kunde udledes udenfra; men
Materialismen, som derfor af ham betegnes som transscendent
og extravagant, bryder denne Tryllekreds og. % PRINTED AS IS (Tryllekreds og.: the image, a point after og, for nothing; O061)
forudsætter Stoffet som faktisk Substrat, af hvilken Tanken % PRINTED AS IS (af hvilken Tanken: the image, hvilken for hvilket (Substrat); O062)
udledes. Ogsaa paa Strauss beraabe Materialisterne
sig, idet de af hans Pantheisme udskyde den ideelle
Faktor, saa der kun bliver den materielle Substants tilbage,
som bærer Alt og ligger til Grund for Alt.

Dog ville Materialisterne ikke finde sig i at udledes
af den tidligere philosophiske Bevægelse. Deres Lære
skal være det nødvendige Resultat af den exakte Erfaringsvidenskabs
storartede Opsving og Fremskridt i
vort Aarhundrede. I Tillid hertil gjøre de kort Proces
med hele den tidligere Philosophi. Saameget maa indrømmes,
at Erfaringsvidenskabens og de praktiske Interessers
Herredømme maaske er den vigtigste Aarsag til
Materialismens Fremkomt, ligesom ogsaa de fleste af % PRINTED AS IS (Fremkomt: for Fremkomst; the Minnesota copy prints it too)
dens Ordførere ere Naturforskere, om de end have tilegnet
sig en qvasi-philosophisk Dannelse. Men en anden
Sag er det, om Materialismen virkelig er den exakte
Erfaringsvidenskabs sande Konsekvents eller om ikke
Forbindelsen mellem begge kun er tilsyneladende og til%
% --- p. 96 ---
\setcounter{footnote}{0}%
\opage{96}fældig. Dette maa den følgende Fremstilling afgjøre. I
ethvert Tilfælde er det kun den mellemliggende videnskabelige
Udviklings store Resultater, ikke eiendommelige
og nye Tanker, der bevirke nogen Modsætning mellem det
18de og det 19de Aarhundredes Materialisme. „Hvis
Materialismen i vor Tid ikke vandt en eiendommelig
Interesse ved den overvættes falske Begeistring for en
i Sandhed stor Kreds af videnskabelig Dannelse, hvis
Resultat og Analogier den misbruger paa den letfærdigste
Maade, saa vilde vi ikke have nødig her at beskjæftige
os med den“\footnote[1]{H. \emph{Lotze}: Medicinische Psychologie. p. 30.}.

Alligevel har det sin store theoretiske og praktiske
Interesse at orientere sig med Hensyn til den moderne
Materialisme. Naar man fra alle Sider forener sig om
den Erkjendelse, at Philosophien maa tilegne sig det
reale Grundlag, saaledes som den exakte Erfaringsvidenskab
paaviser det, at den atter maa stige ned fra Ideens
Himmel til den virkelige Verden, saa udtaler den moderne
Materialisme kun den ensidige Konsekvents heraf,
idet den lader Philosophien gaa op i Fagvidenskabernes
formentlige Konsekventser. Materialismen er den ufornuftige,
uphilosophiske Restauration af Virkeligheden,
repræsenterer yderste Venstre i Bestræbelserne for en
Fremtidens Philosophi. --- Materialismen har tillige en
populær, demokratisk Tendents i Modsætning til den
tidligere Videnskabs Aandsaristokrati. Denne Stræben
efter det Populære gaaer saa vidt, at Büchner i Fortalen
til sit Skrift „Kraft und Stoff“ ligesom erklærer,
at philosophiske Udviklinger, der ikke kunne forstaaes
af enhver Dannet, ikke ere Bogtrykkersværten værd; Philosophien
skal være „geistiges Gemeingut“\footnote[2]{Rigtignok kan Büchner ikke selv have opfyldt denne Fordring, da han strax efter bebreider sine Modstandere, at de af Man%
\opage{97}gel paa videnskabelig Indsigt ikke havde forstaaet hans Anskuelses indre Kjerne. Konsekvent maatte han have beskyldt dem for Mangel paa Dannelse.}% the note runs over from p. 96 to p. 97 (joined by draft.py)
. Carl Vogt
% --- p. 97 ---
\setcounter{footnote}{0}%
\opage{97}reiser om som en Materialismens  Missionær og holder
folkelige Foredrag, idet han endog i Mangel af andet
Lokale optræder i Kirkerne. Og denne Materialismens
demokratiske Retning staaer i Forbindelse med dens
hele praktiske Tendents. Den betragter sig som Indehaver
af Løsningen af de sociale og politiske Problemer,
idet den vil gribe Sagen fat fra neden og derved udrette,
hvad de høieste ethiske og religiøse Ideer ikke
have formaaet. Til den philosophiske Interesse slutter
sig da her nøie den kulturhistoriske. ---

Det er med Hensyn til Materialismen i vort Aarhundrede
værdt at lægge Mærke til, at de to vigtigste
Sammenstød i den materialistiske Strid fandt Sted mellem
Parter, der alle stode paa Fagvidenskabens Grund,
men vare uenige om, hvorledes de sidste Konsekventser
skulde drages. Dette viser strax, hvorledes det hænger
sammen med Materialismens Paastand om at være Fagvidenskabens
sande Konsekvents. --- Den berømte Chemiker,
Justus Liebig, havde i sine „Chemiske Breve“
erklæret, at Naturvidenskaben langt fra at afskjære Veien
til Erkjendelsen af de høieste Ideer, tvertimod selv var
et Led i den Kjede, som fører til dem. Jakob Moleschott,
hvis bekjendte Sætning „Ohne Phosphor kein
Gedanke!“ Liebig særlig havde angrebet, svarede ham i
sin Bog: „Der Kreislauf des Lebens“ (1852), hvor han
hensynsløst fastholdt de materialistiske Lærdomme.
Samme Aar som denne Bog udkom, udbrød allerede en
anden Strid, mellem Rudolph Wagner i Göttingen og
Carl Vogt i Genf. Dog naaede denne først nogle Aar
efter sit Høidepunkt, da Wagner paa et Naturforskermøde
i Göttingen fremstillede som sin Troesbekjendelses
Indhold en ætherisk Sjælesubstants, der ikke døde med
% --- p. 98 ---
\setcounter{footnote}{0}%
\opage{98}Legemet, en Tro, i hvilken han saa den bibelske Sjælelære
udtrykt, men som han kun mente at kunne fastholde
ved absolut at sondre sin Tro fra sin Viden. Han
erklærer sig derfor for Tilhænger af en Kulsviertro, saasnart
man gaaer ud over den exakte Videnskabs Omraade:
„In Sachen des Glaubens liebe ich den schlichten, einfachen
Köhlerglauben am meisten.“ Herimod optraadte
Vogt med sit knusende Stridsskrift: „Köhlerglaube und
Wissenschaft“ (1855). Vogt er en Aand, som paa sit
Omraade har ikke lidet tilfælles med Feuerbach. Her
er samme impetus, samme hensynsløse Cynisme i Tanke
og Udtryk. Men hvad den egentlige Tænkning angaaer,
er Forskjellen rigtignok betydelig; ved Siden af Feuerbachs
fine, gjennemtrængende Analyser tage Vogts Udviklinger
sig temmelig bjørneagtige ud. Hvad der mest
tiltaler hos ham, er hans djærve Humor, som ikke har
forladt ham, selv efterat han som „Reichsregent“ 1848
led Skibbrud paa sine ideale Drømme. Før den Tid,
sagde han senere, havde han nok tvivlet om Menneskets
Fuldkommenhed, men dog anset det for det mest fuldkomne
Dyr; men i Frankfurt havde han ogsaa mistet
denne Tro.

Hvad der gjorde Moleschotts og Vogts Stilling forholdsvis
let, var, at deres Modstandere ikke vare klare
med Hensyn til den Form, i hvilken Idealitetens Sag
skal forsvares. Liebig taler paa en vag og forvirret
Maade om Erkjendelse og Aabenbaring, saa man ikke
ser, om det er den frie Erkjendelse af det Ideelle eller
den positive Kirkelære, han vil hævde. Naturligvis benytter
Moleschott sig strax heraf til at fremstille ham
som Apologet for en middelalderlig Dogmatik. Endnu
mere uklar er Rudolph Wagner. Hans „Kulsviertro“
viser sig ved nærmere Eftersyn at være et helt scholastisk
System, sammensat af Bibellæren, af hvilken han
navnlig hævder Læren om alle Menneskers Nedstammen
fra Adam og Eva som nødvendig til Salighed, og af
% --- p. 99 ---
\setcounter{footnote}{0}%
\opage{99}den phantastiske Antagelse af en udødelig Sjælesubstants,
der bevæger Hjernens Fibre som Musikeren Pianofortets,
og ved Deling forplanter sig fra Forældre til Børn. Det
er intet Under, at ikke blot en Materialist som Carl
Vogt, men ogsaa Naturforskere som Virchow og Philosopher
som Lotze, Zeller og Weisse maatte protestere
mod en saadan Lære. Men denne Uklarhed har været
til Skade for Stridens Udbytte.

Foruden Moleschott og Vogt maa vi endnu fremhæve
to materialistiske Hovedforfattere, Louis Büchner
og Heinrich Czolbe. Hvad der især giver dem Interesse,
er, at de i Materialismen --- saa underlig det lyder
--- virkelig søge en Tilfredsstillelse af en aandelig Trang.
De ledes af en philosophisk og ethisk Interesse. Navnlig
træder Czolbe i en Modsætning til de andre Materialister
derved, at han kommer til Materialismen ad Deduktionens
Vei, en Modsætning, der er vel skikket til
at orientere os med Hensyn til den moderne Materialismes
sande Eiendommelighed, hvorfor vi ville lade vor
Fremstilling falde i to Afdelinger, den empiriske og den
spekulative Materialisme.

I. \emph{Den empiriske Materialisme}. --- Ingen af
Materialisterne stiller sig saa resolut paa den blotte Fagvidenskabs
Standpunkt og fælder ud fra dette absolute
Domme angaaende de store Tilværelsesproblemer, som \emph{Carl}
\emph{Vogt}. Han angiver selv sit Standpunkt som et physiologisk-naturhistorisk.
Physiologien har løst sin Opgave, naar
den efterviser, hvilket Tankens Organ er, hvorledes det
træder i Funktion, og paa hvilke Betingelser denne Funktion
beroer. Menneskeslægtens Naturhistorie har ikke
at gjøre med Menneskets Væsen, saaledes som dette begrebsmæssig
kan bestemmes, men undersøger de Millioner
af Mennesker, der ere spredte over Jorden, og udforsker
deres Eiendommeligheder, og det tilbage til en
Tid, „hvor Mennesket knap har efterladt os flere Spor
af sin Existents end ethvert andet vildt Dyr, som be%
% --- p. 100 ---
\setcounter{footnote}{0}%
\opage{100}boede samme Egn“. Spørges der da om Bestemmelsen
af Menneskets Forhold til Dyrene, da komme ene de
anatomiske og physiologiske Karaktermærker i Betragtning;
alle andre Hensyn ere ligegyldige\footnote[1]{Se foruden „Köhlerglaube und Wissenschafft“ (1855) Forelæsningerne „über den Menschen, seine Stellung in der Schöpfung und in der Geschichte der Erde“ (1863).}. % PRINTED AS IS (Wissenschafft: the image, O064)

Paa et saadant Standpunkt maa Fagvidenskaben
nødvendigvis stille sig, naar den skal kunne gjennemføre
sin exakte Methode. Denne forudsætter altid en vis Abstraktion;
Phænomenet maa isoleres, drages ud af sin
levende, mangesidige Virkelighed for at kunne kontrolleres
fra en enkelt bestemt Side. Her ved Fagvidenskabens
Begyndelse støde vi derfor ogsaa strax paa dens Grændse.
Er det ene de anatomiske Karaktermærker, der komme
i Betragtning, saa er det ogsaa netop kun som anatomisk
Væsen, Mennesket er blevet betragtet, og man kan
ikke uden Videre gjøre denne Betragtning til den absolute.
Men en saadan Begrændsning protesterer nu Carl
Vogt paa det Bestemteste imod; han vil umiddelbart
gjøre Physiologien til den absolute Videnskab.

Geoffroy St. Hilaire og Quatrefages mente, at man
i den naturhistoriske Betragtning af Mennesket ogsaa
maatte tage Hensyn til det Intellektuelle, Moralske og
Religiøse, som netop udgjør det eiendommelig Menneskelige,
og støttende sig hertil opstillede de et særeget
Menneskerige i Modsætning til Dyre- og Planteriget.
Vogt lader sig ikke nøie med at afvise dette fra et rent
naturhistorisk Synspunkt. Han paatager sig at vise, at
der i de aandelige Elementer, som de franske Forskere
fremhæve, ikke ligger noget for Mennesket Eiendommeligt.
Naar man tilskriver Mennesket Religiøsitet, da skal
det kun sige, at Mennesket danner sig visse Forestillinger
om de Phænomener, hvis Grund det ikke kan fatte,
medens Dyret paa Grund af sin ringere aandelige Be%
% --- p. 101 ---
\setcounter{footnote}{0}%
\opage{101}gavelse slet ikke føler sig foranlediget til at tænke over
Grunden til saadanne Phænomener. Dog vilde Vogt herved
let komme til at sætte en større Modsætning mellem Menneske
og Dyr, end der stemmer med hans Mening; han
forklarer derfor det Religiøses egentlige Oprindelse af
Frygten for det Ubekjendte, hvilket betragtes som et
Overnaturligt, og denne Frygt, der hos Mennesket
bliver til Gudsfrygt og frembringer Religionen, skal ogsaa
Dyret være i Besiddelse af. Saaledes føle Hunden
og andre af vore „intelligente Husdyr“ Spøgelsefrygt
ligesaavel som Bretagneren eller Baskeren. Heller ikke
det Moralske skal være noget for Mennesket Eiendommeligt,
da Anlæg dertil allerede vise sig hos Dyrene,
navnlig dem, der leve i Samfund; tilmed ere Begreberne
om Godt og Ondt relative, afhænge af den forskjellige
Kultur, af Samfundets Fornødenheder og de Enkeltes
indbyrdes Forhold. Materialismens Psychologi reducerer
sig altsaa til det, som Mennesket har fælles med Dyret.

I Vogts physiologiske og anatomiske Undersøgelser
kunne vi her ikke fordybe os. Det er tilstrækkeligt at
fremdrage den Maade, hvorpaa han søger at præcisere
sine Resultater. Hjernen, hedder det, er alle forskjellige
saakaldte Sjælefunktioners Organ; disse Funktioner
ere bundne til visse Dele og Steder i Hjernen og kunne
kun udøves af dette Organ. Om han end, skjøndt nødig,
maa give Virchow Ret i, at det er umuligt naturvidenskabelig
at forklare Bevidsthedens Kjendsgjerning, mener
han dog, at man vil kunne opnaa med Bestemthed at
eftervise de Ganglieceller, ved hvis Incitering denne eller
hin specielle Fornemmelse vækkes. At Bevidstheden
ikke kan fremkomme uden Ganglieceller, kan bevises;
\emph{hvorledes} den opstaaer i dem, kan ikke siges; men
det staaer fast, at Funktion og Organ afhænge af hinanden
og ere uopløselig forbundne med hinanden. Funktionen
maa altsaa ogsaa være afhængig af Organets materielle
Tilstand og lide under dets Forandringer.

% --- p. 102 ---
\setcounter{footnote}{0}%
\opage{102}Det er klart, at her kun er paavist et Afhængigheds-,
et Betingelsesforhold. Fordi to Momenter betinge
hinanden, behøver det ene ikke at være frembragt af
det andet. Der kan godt være en kvalitativ Forskjellighed
tilstede (hvilken naturligvis ikke maa forvexles
med absolut Uensartethed). Forholdet mellem det Ensartede
og det Uensartede i de to Momenter vil man
finde ved at fremdrage ikke blot de Punkter, hvor Betingetheden
er absolut, men ogsaa de, hvor den viser sig
at være relativ, idet hvert Moment virker paa sin eiendommelige
Maade. Den kvalitative Forskjel, som bliver
tilbage, viser saa hen til et tredie, mere omfattende Element,
af hvilket begge hine Momenter udspringe. Her godtgjør
da, naar vi anvende dette paa det foreliggende Exempel,
den psychologiske Betragtnings Utilstrækkelighed
sig ligesaavel som den physiologiskes. Kun en metaphysisk % PRINTED AS IS (physiologiskes: both copies)
Undersøgelse kan fastsætte Forholdet mellem
Aand og Materie. Materialisterne kunne kun udlægge
Physiologiens Resultater i deres Faveur ved at forudsætte
det som givet, at Aanden kun er Materiens Produkt.
Det er deres tragiske Skjæbne, at de, der sky
Metaphysik som Pesten, kun opnaa at gaa ud fra den
rent krasse og raa Metaphysik. Det Metaphysiske lader
sig nu engang ikke udelukke.

Vogt og hans Meningsfællers Polemik er i høieste
Grad paavirket af selve de Modstandere, de bekæmpe.
Deres Arvefjende er Dualismen, der sætter en ydre Adskillelse,
et fremmed Forhold mellem Sjæl og Legeme,
og denne Fjende opdage de i enhver Opfattelse, der
ikke med dem vil udlede Aanden af Materien. Derfor
har Rudolph Wagner frembragt megen Konfusion ved
sin Lære om Sjælesubstantsen og om det specielle Troesorgan.
Han er en Modstander efter Materialisternes
Hoved, idet han tager ligesaa udvortes og umiddelbart
paa Tingene som de selv. Overfor saadanne Anskuelser
ere de i deres Ret; men de lade sig i Seirens Rus føre
% --- p. 103 ---
\setcounter{footnote}{0}%
\opage{103}videre og mene ved Beviset for Aandens Betingethed
ved det Materielle at have godtgjort dette som den fuldstændige
Aarsag til hin. ---

Medens Vogt gaaer ligefra Zoologien til Afgjørelsen
af Hovedproblemerne, træffe vi hos \emph{Moleschott} Forsøg
paa at give en dybere Begrundelse gjennem en Art Erkjendeslære % PRINTED AS IS (Erkjendeslære: the image, Er-|kjendeslære, for Erkjendelseslære; O065)
og Metaphysik. Hos ham mærker man Paavirkningen
af Spinoza, Hegel og Feuerbach fra hans
tidligere Studier. Philosophien er ifølge hans Opfattelse
det aandige Udtryk for den Sum af Iagttagelser, som
den sandselige Opfattelse til en vis Tid har vundet. At
philosophere er at tænke. Tænkning bestaaer i Bearbeidelsen
af den sandselige Iagttagelse. Philosophi og
Erfaring skulle gjensidig gaa op i hinanden, saa at Anskuelsen
er Tanke og Forstanden anskuer med Bevidsthed.
Man troer her at høre en Schellingianer. Men
snart viser det sig, at Enheden af det ideelle og det reelle
Element i Erkjendelsen er ensidig: Tankens Ret stammer,
som Moleschott udtrykker sig, kun fra Sandsernes
Naade. Tænkningen opstaaer af Iagttagelsen, og Iagttagelsen
er Opfattelse af det bevægede Stofs Indtryk
paa vore Nerver og gjennem dem paa Hjernen. Det
tænkende Menneske er derfor Summen af sine Sandser,
Bevidstheden en Egenskab ved Stoffet. Hvad der erkjendes,
er kun Tingenes Egenskaber, men Summen af
disse for Sandserne aabenbarede Egenskaber udgjør
Tingenes Væsen, ligesom Sandsernes Sum udgjorde Tanken.
Træet i og for sig er ikke forskjelligt fra Træet
saaledes som det opfattes; thi Træet er kun, forsaavidt
det opfattes. Uden et Forhold til Øiet, i hvilket det sender
Lysstraaler, er Træet ikke til. Al Væren er en Væren ved
Egenskaber, og enhver Egenskab bestaaer i et Forhold.

I disse Sætninger indeholdes adskillige Problemer,
som Moleschott dog umiddelbart gaaer hen over. Navnlig
mærker han ikke, hvorledes han selv har trykket en
Skepticismens Braad ind, som maa gjøre al Materialisme
% --- p. 104 ---
\setcounter{footnote}{0}%
\opage{104}umulig. Naar det Ydre kun er ved at speile sig i det
Indre, hvilken Forvisning kan man da have om den ydre
Virkeligheds Realitet, end sige, at man skulde kunne
sætte det Ydres Bevægelser som Aarsag til det Indre?
Moleschott svarer, at i Konsekvents af Feuerbachs Lære
maa Kraftens Begreb og dermed ethvert Begreb om Aktivitet
kun opfattes som en Anthropomorphisme uden objektiv
Betydning. Men saa maa det Samme gjælde om
Stoffets Begreb, det Begreb, hvorpaa Moleschott bygger
sin hele Lære.

Læren om Stoffets Kredsløb er nemlig ifølge Moleschott
den nye Verdensanskuelse, for hvilken Encyklopædisterne
i det 18de Aarhundrede kun vare de begeistrede
Forløbere; den er det Hængsel, hvorom Nutidens Philosophi
dreier sig. Den nyere Chemi har veiet Tidens
Tand. Den har opdaget, at Stoffet er udødeligt, aldrig
gaaer tilgrunde, men kun omsættes i andre og atter andre
Former. Mellem den organiske og uorganiske Verden,
mellem Planter og Dyr kredser Stoffet i uendelig Bevægelse.
F. Ex.: „Bjergmanden, som i Wetterau eller
Estramadura graver efter phosphorsur Kalk, søger Mere
end Guld; han graver efter Hvede, efter Mennesker.
Bjergmanden hæver den Aandens Skat, som Bonden sætter
i Omløb, idet han giver Tidens Hjul det første Stød.
Bjergmanden, der i sit Ansigts Sved og med Livsfare
tilkæmper sig sit Livs Ophold, véd ikke, om ikke maaske
Stoffet til det bedste Hoved glider gjennem hans
der. Han sætter maaske Aarhundreder i Bevægelse ved % PRINTED AS IS (der. Han: the image, hans|der.: Hæn- lost at the line turn (Hænder); recrop rc1)
sit skjulte Arbeide.“ Derfor, siger Moleschott, nævner
man ikke Stoffets Kredsløb uden en vis Ærefrygt. Thi
ligesom Handelen er Samfærdselens Sjæl, saaledes er
Stoffets evige Kredsen Verdens Sjæl. Paa Forholdet
til dette Kredsløb bero de forskjellige Væseners Eiendommelighed.
Stoffet ligger til Grund for Alt, gaaer
igjennem Alt. Hvilket er da Stoffets almindelige Begreb?
Det findes ved at abstrahere fra alle Egenskaber,
% --- p. 105 ---
\setcounter{footnote}{0}%
\opage{105}ved hvilke det ene Stof adskiller sig fra det andet. Der
bliver da tre fælles og afgjørende Egenskaber tilbage:
Tyngden, Rumopfyldelsen og Kraften eller Evnen til
Bevægelse. Kraften er altsaa kun en Egenskab ved
Stoffet. Dette gjælder, som vi allerede have set, ogsaa
om den høieste Kraft, Bevidstheden. Stoffet regerer Alt,
ogsaa Mennesket: befordre Kryderi ikke Fordøielsen,
vække Løg og Vanille ikke sandselige Drifter, beherske
Vin, Kaffe og The ikke „Hjernens Stemning“?

Moleschott har ved denne Begrebsbestemmelse ikke
gjennemført Abstraktionen. Det er nemlig klart, at af
de tre Egenskaber, i hvilke efter hans Lære Materiens
Væsen skal udtrykke sig, gaa de to, Tyngden og Rumopfyldelsen,
ind under den tredie, den almindelige Bestemmelse
af Kraft. Den nyere Physik fører ogsaa til
at opfatte Materien som den concentrerede Kraft. Men
dette viser os hen i en ganske anden Retning end Materialismen.

Af sin Lære om Stoffets evige Kredsløb udleder
Moleschott nu eiendommelige praktiske Konsekventser.
Ligesom Mennesket i det Hele fra først til sidst er bestemt
ved Stoffets Indflydelse, saaledes er ogsaa Villien
kun det nødvendige Udtryk for en ved ydre Indvirkninger
fremkaldt Tilstand i Hjernen. Begreberne Godt
og Ondt og Retsbegrebet fremgaa af Samfundets Fornødenheder.
Retten til Straf grunder sig paa Trangen
til Selvopholdelse. Af denne ubetingede Determinisme
udleder Moleschott en Mildhedens og Forsoningens Stemning
overfor Forbryderen, ifølge Mad. de Staëls Ord:
tout comprendre, ce serait tout pardonner, et Ord, der
ifølge hans Mening burde staa i Spidsen for den nyere
Tids Evangelium, ligesom Buddet om Kjærlighed til
Næsten er Kjernen i Christendommens Moral. Dog
glemmer Moleschott selv dette „gyldne Ord“ overfor det
hessiske Deputeretkammer, der til hans Indignation forkastede
et Forslag om Dødsstraffens Afskaffelse. En
% --- p. 106 ---
\setcounter{footnote}{0}%
\opage{106}anden materialistik Forfatter, H. Czolbe, har indset denne % PRINTED AS IS (materialistik: for materialistisk; both copies)
Inkonsekvents og erklærer, at Moleschotts „bløde Gemyt“
her er løbet af med ham. Dog bringer dette ikke Moleschott
ud af den optimistiske Glæde over Tilværelsen,
som han udleder af sin materialistiske Determinisme:
„Med Stoffet kredser Livet gjennem alle Verdens Dele,
med Livet Tanken, med Tanken den af Naturnødvendighed
gode Villie. Trods alle sine Onder er og bliver
Jorden dog et Paradis.“

Uagtet Enkeltvæsenerne kun dukke op i Stoffets
uendelige Kredsløb for atter at synke tilbage i dette,
faaer Menneskets Virken dog en særegen Betydning derved,
at den, ledet af den chemiske, physiske og physiologiske
Viden kan spare paa Stoffet ved at føre det ad
Baner og gjennem Forbindelser, hvor det ad den korteste
Vei kan øve den høieste Virkning. Ved den rigtige
Blanding af Kulsyre, Ammoniak o. s. v. arbeide vi paa
Menneskehedens høieste Udvikling. Naturforskeren er
vor Tids Prometheus, Chemien den centrale Videnskab.
Løsningen af vor Tids store sociale Spørgsmaal ligger i
Naturforskerens Haand; thi Fattigdom er umiddelbart
kun Mangel paa Stof, og da Læren om Stoffets Kredsløb
forvisser os om, at der i Naturen er Stof nok, maa
Videnskaben engang komme saavidt, at den lærer en
Fordeling af Stoffet, hvorved Fattigdom bliver umulig.

Hvad der interesserer os ved denne paa „ein stoffliches
Glaubensbekenntniss“ byggede Ethik, er at se,
hvorledes Materialismen her som paa saa mange andre
Punkter gaaer ud over sig selv, idet den vil drage sine
sidste Konsekventser. Moleschott kommer uvilkaarlig til
at stille „en hellig Opgave, et Formaal, der skal virkeliggjøres
ved Hjælp af de Love, efter hvilke Stoffets Bevægelser
retter sig. Men her bliver det da Tanken og
Loven, som ere det Styrende, og det absolute Afhængighedsforhold
til Materien brydes. Vel vil Moleschott
hertil svare, at Tankens Virken kun er en Sidebane,
% --- p. 107 ---
\setcounter{footnote}{0}%
\opage{107}paa hvilken Stofbevægelsen ligesom slaaer ind, for med
større Intensitet at vende tilbage igjen. Men allerede
heri ligger en Selvopgivelse af Materialismen. Stoffets
Begreb umuliggjør nemlig under materialistiske Forudsætninger
et Kredsløb i anden Forstand end den blotte
udvortes Succession af forskjellige indbyrdes afvexlende
Former. Saasnart der skal udvikles et Formaal gjennem
dette Kredsløb, er der sat en koncentreret Enhed, som
er forskjellig fra Stoffet som saadant, og til hvilken dette
forholder sig tjenende. ---

Den mest populære af de materialistiske Forfattere
er Louis Büchner, hvis Skrift: „Kraft und Stoff. Empirisch-naturphilosophische
Studien er blevet oversat paa
Engelsk, Fransk og Spansk og har oplevet en halv Snes
tydske og franske Oplag. Desuagtet er dette Skrift
meget overfladisk, men interesserer ved Büchners levende,
skjøndt ofte vel naive Begeistring for Sagen.

Vi finde her først en Stræben efter en skarp Opfattelse
af Erfaringens Væsen. Büchner lærer, at den
videnskabelige Forsken tilsidt støder paa en naturlig % PRINTED AS IS (tilsidt: the image, for tilsidst; recrop rc1)
Grændse, der ligger i selve Menneskets Erkjendelsesevne.
Vi kunne ikke gjøre os nogen Forestilling eller
have nogen Viden om det Absolute ɔ: om det, der gaaer
ud over den sandselige Verden. Al vor Viden og Forestillen
er kun relativ og beroer paa en gjensidig Sammenligning
af Tingene. Derfor naaer den heller ikke en
Forklaring af det indre Forhold mellem Aand og Materie.
Büchner erklærer gjentagne Gange, at hvad han vil paavise,
er kun det nødvendige og uopløselige Forhold mellem
Aand og Materie, Kraft og Stof, ikke dette Forholds
Hvorledes. Men skjøndt Forholdet mellem Aand
og Materie efter hans Mening fra Naturvidenskabens
Standpunkt ikke kan udvikles positivt, saa kan det dog
bestemmes negativt, det vil sige, Empirien afskjærer
Muligheden af et transscendent, dualistisk Forhold, idet
den paaviser Materien som Aandens nødvendige Betingelse.

% --- p. 108 ---
\setcounter{footnote}{0}%
\opage{108}Disse Sætninger karakterisere korrekt Erfaringsvidenskabens
Standpunkt. Ved at fastholde dem vilde Büchner
heller ikke være bleven Materialist. Men uheldigvis modsiger
han dem i samme Aandedræt, som han fremfører
dem. Allerede i Fortalen finder han i den spekulative
Naturphilosophis Skibbrud det tydeligste Bevis paa, at
Verden ikke er Virkeliggjørelsen af en Tanke, men et
Komplex af Ting og Kjendsgjerninger. Paa lignende
dogmatisk Vis kalder han Forholdet mellem Aand og
Natur (som han selv har erklæret for uforklarligt) et
kausalt Forhold og underforstaaer herved, at det er Materien,
der er Aarsagen, Aanden, der er Virkningen. Men
stærkest træder Uklarheden frem i Forholdet mellem Kraft
og Stof i Almindelighed. Han gaaer ud fra den Moleschottske
Sætning: Uden Stof ingen Kraft, uden Kraft intet
Stof, og ligesom Moleschott opfatter han nærmere denne
Sætning saaledes, at Kraften er en Egenskab ved Stoffet.
Denne Paastands Uholdbarhed have vi ovenfor set.
Da Büchner blev opmærksom paa Physikernes, navnlig
Helmholtz's Undersøgelser om Kraftens Opholdelse,
indskød han i 5te Oplag af sit Skrift et Kapitel om
„Kraftens Udødelighed“, hvori Kraftens Kredsløb stilles
som nødvendigt Korrelat ved Siden af Stoffets Kredsløb,
saaledes at „begge tilsammen fra Evighed og i Evighed
danne den Sum af Phænomener, som vi kalde Verden.“
Han tænker ikke paa, at han herved er kommen over
i en fuldstændig Dualisme mellem to paralelle Kredsløb, % PRINTED AS IS (paralelle: the image, O072)
men vender rolig tilbage til sin Forudsætning: „I Materien
bo alle Natur- og Aandskræfter; Materien er al
Værens Urgrund“. Og dog fastholder han, at al Kraft
er immateriel: „Vi have defineret Kraften som en Egenskab
ved Stoffet og set, at begge ere uadskillelige; dog
ligge begge \emph{efter Begrebet} langt borte fra hinanden,
ja negere i en vis Forstand hinanden. Idetmindste vide
vi ikke, hvorledes man vil definere Aand og Kraft anderledes
end som noget Immaterielt, der i sig selv udelukker
Materien og er den modsat.“ --- Men hvorledes kan
% --- p. 109 ---
\setcounter{footnote}{0}%
\opage{109}Noget, som i sig selv udelukker Materien, være en Egenskab
ved den?

Ligesom den immaterielle Kraft dog skal være en
Egenskab ved Materien, saaledes skal den immaterielle
Sjæl være et Produkt af Materien. Büchner henviser
her navnlig til den successive Udvikling. Ideerne opstaa
ved, at Mennesket opfatter det Fælles ved Tingene i
den ydre Verden, deraf danner sig en ideal Skikkelse
og endelig tillægger denne Prædikat af det Sande, Skjønne
eller Gode. Men herved forudsættes bestandig Hovedsagen:
selve den uendeliggjørende, idealiserende Evne
til at danne en saadan Skikkelse. Ved at forudsætte
den forudsætter man netop det, der skulde forklares. Fordi
Kunsten laaner sine enkelte Træk fra Naturen, beroer
Idealet ikke blot paa en Sammensætning af disse enkelte
Træk. Fordi f. Ex. Japaneserne anse sorte Tænder for en
Skjønhed, medens det forekommer dem afskyeligt at
have hvide Tænder „ligesom en Hund“, derfor er det
æsthetiske Ideal ikke noget rent Tilfældigt og Udvortes.

Hos Büchner brydes saaledes de forskjellige Retninger.
Han tilhører Materialismen af Princip, men uvilkaarlig
rører den sande Opfattelse sig hos ham, og han
kommer saaledes i Modsigelse med sig selv. Tænkes
denne Brydning og Selvmodsigelse nu gjennemført, medens
Materialismen alligevel fastholdes, da kommer man
til et Punkt, hvor man ikke længer gaaer ud fra Erfaringen,
men fra den indre Overbevisning om Materialismens
Sandhed, efter hvilken Overbevisning saa Erfaringens
Udsagn maa rette sig.

2. \emph{Den ethisk-spekulative Materialisme}. ---
Dette Navn synes at indeholde en Selvmodsigelse. Materialismen
synes at være saa nøie knyttet til Empirien,
at denne ikke kan overskrides uden at ogsaa Materialismen
opgives. Dog maa allerede selve Afsnittet om den
empiriske Materialisme have lært os at indse Urigtigheden
af denne almindelig udbredte Mening. Det viste
% --- p. 110 ---
\setcounter{footnote}{0}%
\opage{110}sig, at Materialisterne kun med Urette kunde paaberaabe
sig Erfaringsvidenksabens Autoritet. Dels indlader Erfaringsvidenskaben % PRINTED AS IS (Erfaringsvidenksabens: for Erfaringsvidenskabens; the Minnesota copy prints it too)
sig ikke paa dogmatiske Anskuelser,
og dels synes, naar Alt kommer til Alt, den exakte Erfaringsvidenskab
snarere at føre i idealistisk end i materialistisk
Retning. Dette vil især fremgaa af dens Behandling
af Materiens Begreb og af dens Lære om Sandseopfattelsen,
og vil i 5te Kapitel af dette Skrift vise sig
fra forskjellige Sider. Derfor er det et Vidnesbyrd om
en sjelden Uhildethed, naar den sidste af Materialismens
Forkæmpere, vi her ville fremdrage, \emph{Heinrich Czolbe}
(Læge, ligesom Büchner) erklærer det for et Bevis paa
Uvidenhed at paastaa, at den moderne Naturvidenskab
fører til Materialisme. Det er hans Overbevisning, at
naar man som Vogt og Moleschott antager den nyere
Physiologis Lære om Sandsningen, da har man ingen
anden Udvei end at kaste sig i Armene paa et eller
andet spekulativt System. Naar Materialismen hidtil har
bekæmpet Theologien og den spekulative Philosophi, da
staaer derfor nu den vigtigste Opgave tilbage, nemlig at
bekæmpe den moderne Naturvidenskab\footnote[1]{„Die Entstehung des Selbstbewusstseins“. Berlin 1856. --- „Die Elemente der Psychologie vom Standpunkte des Materialismus“ i 26de Bind af „Zeitschrift für Philosophie und philosophische Kritik.“}.

Den nyere Physik og Physiologi lærer, at Sandseorganerne
paavirkes udenfra paa rent mechanisk Vis, men
at disse physiske Paavirkninger i Organismen selv paa
en efter de enkelte Sandseorganers Eiendommelighed
forskjellig Maade omsættes til subjektiv Fornemmelse.
Saalænge man hylder denne Lære, har man ifølge Czolbe
ikke faaet Bugt med det Oversandselige. Overgangen
fra de physiske Virkninger af Stød og Tryk til Lys- og
og Lydphænomener staaer som noget Uforklarligt og fører % PRINTED AS IS (og Lydphænomener: the image, og doubled over the line turn (og|og); O073)
til Antagelsen af et eiendommeligt subjektivt Princip i
% --- p. 111 ---
\setcounter{footnote}{0}%
\opage{111}det fornemmende Væren. Sandsenervernes specifike Energi
er et fuldstændigt mystisk Væsen. Der kan i Kraft af
Materialismens Princip ikke være en saadan Kløft mellem
det physiske Agens og den psychiske Fornemmelse.
Sandsenerverne kunne kun være passive Substrater, og
Sandseorganerne have kun den Betydning at forstærke
Bevægelsen, førend den naaer Nerverne; og det maa
være væsentlig den samme Bevægelse, der udenfra forplanter
sig helt igjennem uden nogen kvalitativ Forandring.
Men hvorpaa beroer da Forskjellen mellem de
forskjellige Sandser? Dette Spørgsmaal, som man efter
Czolbes Mening pleier at overhugge istedetfor at besvare,
besvarer han ved at opstille Begrebet om intensiv Kvantitet
ɔ: en Kvantitet, der ikke maales efter Tal og Volumen
(extensivt) men efter Hurtighed og Styrke. Forskjellen
mellem de forskjellige Sandser beroer nu derpaa,
at ikke blot de forskjellige Bevægelser (Lyd, Lys
o. s. v.) forplante sig med forskjellig Hurtighed, men
ogsaa de forskjellige Sandsers Nerver (Høre- og Synsnerver
o. s. v.) have forskjellig Elasticitet, saa at Bevægelserne
udenfra forplantes paa forskjellig Maade. Ligesom
intensiv Kvantitet træder istedetfor specifik Kvalitet,
saaledes specifik Elasticitet istedetfor specifik Energi.
Og ligesom Bevægelsens forskjellige Hurtighed skal forklare
de forskjellige Sandser, saaledes fremkalder dens
forskjellige Styrke de forskjellige Følelser og Drifter:
for stor Styrke fremkalder Smerte, for liden Trang og
Drift, en Middelstyrke Følelsen af Lyst.

Alle indre Erfaringer (Fornemmelser, Følelser o. s.v.),
som paa denne Maade forklares rent mechanisk, omfattes
af en fælles Egenskab, Bevidstheden. Denne bestaaer
i alle indre Erfaringers i sig selv tilbageløbende
Retning, hvorved de danne en uopløselig Enhed. Bevidstheden
forudsætter da et Organ, som kan give de
Bevægelser, der forplante sig til det, en saadan i sig
selv tilbageløbende Retning. Et saadant Organ er Hjer%
% --- p. 112 ---
\setcounter{footnote}{0}%
\opage{112}nen, og Bevidsthedens Subjekt (det, som er bevidst) er
altsaa de Bevægelser, der ved Sandserne meddeles Hjernen.
Selvbevidsthedens Enhed udtrykkes ved, at disse
Bevægelser kredse ved Siden af hverandre; det er ikke
nødvendigt at antage en centripetal Retning af Bevægelserne
i Hjernen.

Man behøver ikke at være meget bekjendt med Physiologi
og Psychologi for at se, hvor dristig Czolbe træder
op mod alle hidtil gjældende Anskuelser. Czolbe er i
denne Henseende meget aabenhjertig. Han udtaler ligefrem,
at hans Lære strider mod den almindelige Physiologi
og er reaktionær, da den forsvarer en Anskuelse, der af
Naturforskerne betragtes som forældet. Men han erklærer
tillige, at naar man fra Materialismens Standpunkt,
som udelukker alt Oversandseligt, skal give en Forklaring
af Sandsefornemmelsen og Bevidstheden, maa man
komme til Resultater som de nu fremstillede. Det
dybeste Differentspunkt mellem den materialistiske og
den spiritualistiske Philosophi ligger efter hans Mening
i den forskjellige Afgjørelse af det Spørgsmaal, om de
ydre Agentia mechanisk forplante sig gjernem Nerverne % PRINTED AS IS (gjernem: for gjennem; both copies)
til Hjernen eller om de kun berøre Sandseorganet og i
dette undergaa en Omdannelse. Hvis han tillagde Erfaringsvidenskaben
afgjørende Betydning, maatte han
sætte hele Problemet om Materialisme og Spiritualisme
i Bero for foreløbig at komme til Overbevisning om, i
hvilken Retning den exakt opfattede Erfaring viser.
Men han gaaer ud fra Materialismens Princip og deducerer
deraf, hvorledes den faktiske Sammenhæng maa
være. Derfor er hans Naturphilosophi spekulativ trods
Schellings og Hegels. Principet gjennemføres ved en
Række Hypotheser, af hvilke den ene fører den anden
med sig. Og det Forunderligste er, at Czolbe netop
ved paa denne Maade at gjennemføre sit Princip føres
ud over det. Sandsningen forklares ved intensiv Kvantitet,
Bevidstheden som roterende Bevægelser; men saa
% --- p. 113 ---
\setcounter{footnote}{0}%
\opage{113}maa man ogsaa slutte omvendt: hvor der er intensiv
Kvantitet og roterende Bevægelse, der er ogsaa Sandsning
og Bevidsthed. Hele Naturen bliver altsaa besjælet.
Dette er, siger Czolbe, kun en real Opfattelse af
Pantheismen. Naar man opfatter Aanden som Natur,
maa man ogsaa opfatte Naturen som Aand. Skjøndt
han saaledes i Materialismens Omslag til Idealisme ser
en indre Nødvendighed, overraskes han dog selv ved
Konsekventserne af sit Princip: „Scheint es mir doch
selbst, dass ich durch die Konsekvenzen, zu denen das
Princip mich zwang, in eine märchenhafte Gedankenwelt
gerathen bin“. Sagen er, at den Illusion, hvorpaa
Materialismen grunder sig, her lægger sig for Dagen.
Czolbe forklarer ligesaalidt Fornemmelse og Bevidsthed
som den Physiologi, han forkaster. Det Afgjørende ved
Fornemmelse og Bevidsthed er den indre Opfattelse og
Fremstilling af det ydre Virkelige. Men i Czolbes Lære
se vi bestandig det ene ydre Virkelige virke paa det
andet; selv Bevidstheden skal kun være en vis Bevægelse
af ydre Dele i et ydre Rum. Her gjøres altsaa ingen
Overgang fra det Ydre til det Indre. Czolbe har selv
spærret sig denne ved at sætte et absolut Ydre. Der
gives intet absolut Ydre uden i Materialistens Tanke ---
og netop fordi det er i hans Tanke, er det alligevel ikke
absolut udvortes; han tænker uden at vide det Tanken
med, idet han tænker det Ydre. Saaledes er hos Czolbe
Cirkelbevægelsen i Grunden ikke Andet end Symbol paa
Aandens indre Enhed. Det er hans Trang til anskuelige
Begreber, der bringer ham til at lade dette Symbol
krystallisere til en fast endelig Form, hvilken han saa
betragter som Sagen selv.

Thi Materialismens Princip, hvilket han negativt bestemmer
som Udelukkelse af alt Oversandseligt i Tænkningen,
bestemmer han positivt som udspringende af
„das geistige Bedürfniss nach Anschaulichkeit des Denkens.“
Allerede heraf ses, at hans Udgangspunkt er
% --- p. 114 ---
\setcounter{footnote}{0}%
\opage{114}et subjektivt. Men en anskuelig Forklaring er for Czolbe
Et med en mechanisk Forklaring; han bruger ligefrem
„physikalisch“ og „anschaulich“ som Synonymer. Dog
viser det sig ved selve Materiens Begreb, at Physikens
Forklaringer ikke altid opfylde Anskuelighedens Fordringer,
idetmindste i Czolbes Forstand. Han fordrer
derfor udstrakte, begrændsede, uigjennemtrængelige Atomer
og afviser Physikens Kraftcentra som „et oversandseligt,
Materialismens Princip modsigende Begreb.“
Han vil fremdeles forklare Legemernes gjensidige Tiltrækning
ved Hjælp af „anskuelige Grundbetingelser“
uden at anvende Mystik --- paa følgende Maade: Legemerne
ere sondrede, men omsluttes dog alle af Rummet, og
Tiltrækningen er den nødvendige anskuelige Mellembestemmelse
mellem Sammenhæng og Adskillelse. Denne
Forklaring, hvis Fortjeneste Czolbe især sætter i, at
den udelukker Kraften, dette oversandselige Begreb („die
moderne Phrase: „kein Stoff ohne Kraft“, ist eine Unmöchlichkeit, % PRINTED AS IS (Unmöchlichkeit: for Unmöglichkeit (d:C180, zoomed: ch))
ein Absurdum“), er et Modstykke til den
spekulative Naturphilosophis Konstruktion af Materien.
Spekulationen konstruerer Materien ud af Kraften, Czolbe
Kraften (Tiltrækningen) ud af Materien. Her fremtræder
klart hans Materialismes spekulative Karakter.

Anskuelighedens Princip staaer hos Czolbe i Forbindelse
med en hel ethisk Verdensanskuelse, i hvilken,
som han selv klart er sig bevidst, det sidste Motiv
til hans Materialisme ligger. Han minder gjentagne
Gange om, at Anskuelighedsprincipet er beslægtet med
Græcitetens hele Aandsretning, og efter hans egen Udtalelse
var det Digteren Hölderlins af Begeistring for
Græciteten besjælede Naturalisme, der vakte de første
Spirer til hans Verdensanskuelse hos ham. Det Skadelige
i den nuværende Kulturtilstand bestaaer efter hans
Mening især i Levningerne af den supranaturalistiske
Verdensanskuelse. Transscendentsen, navnlig Læren om
den personlige Gud, anser han for en Rest af Heden%
% --- p. 115 ---
\setcounter{footnote}{0}%
\opage{115}skabets Egoisme og betragter derfor Materialismen som
Christendommens Fuldendelse. Den græske Tilfredshed
med det jordiske Liv strider mod Egoismen og forudsætter
megen Selvbeherskelse og Resignation; den fortjener
langt mere Navn af Fromhed end det, man i Almindelighed
kalder saaledes. Sin ethiske og kulturhistoriske
Grundanskuelse finder Czolbe udtrykt i Strauss's
bekjendte Ord (i Afhandlingen om Julian Apostata):
„Grækernes frie harmoniske Menneskelighed og Romernes
paa sig selv beroende Mandighed er det, hvortil vi nu
igjen ere i Begreb med at arbeide os frem af den lange
christelige Middelalder og berigede med det aandelige og
sædelige Udbytte fra denne.“ Utilfredsheden med Phænomenernes
Verden er en moralsk Svaghed. Men det
er ikke blot „der Blödsinn der Transscendenz“, som
Czolbe vil bekæmpe, det er ogsaa den abstraherende
Forstands ensidige Overvægt ikke blot i Philosophien
men ogsaa i de exakte Videnskaber. Han vil frembringe
Harmoni mellem Anskuelsen, Følelsen og Forstanden og
derved i den hele Personlighed.

I den Grad er den ethiske Opfattelse af Materialismens
Princip den til Grund liggende hos Czolbe, at han
i sidste Instants begrunder det ved en Appel til Følelsen.
Den umaadelige Kløft mellem materialistisk og
spiritualistisk Philosophi bevirkes af den ved forskjellige
Livserfaringer betingede forskjellige Gemytsretning saavel
i æsthetisk som i ethisk Henseende. I Virkeligheden
finder der altid et credo ut intelligam Sted; Følelsen
behersker Erkjendelsen, en rent objektiv Erkjendelse
gives ikke. Derfor nærer Czolbe megen Sympathi for
den klassiske theologiske Bekæmper af Materialismen,
Fabri; han ser ligesom denne i Erkjendelsens Resultater
Udtryk for Villiens Retning. Det er en ægte
theologisk Sætning, naar Czolbe udtaler, at det er „Daarskab
i Superlativ“ at ville philosophere over Materialismen
uden en inderlig Overbevisning om dens  Sandhed!

% --- p. 116 ---
\setcounter{footnote}{0}%
\opage{116}Czolbe er en Stoiker. Ogsaa for Stoikerne var det
ethisk-psychologiske Motiv Hovedsagen; de søge en Lære,
der kan tilfredsstille deres praktiske Tendents, og komme
derved paa lignende Maade som Czolbe til en materialistisk
Anskuelse. Begges Selvmodsigelse beroer paa,
at de postulere Materialismens Nødvendighed ud af et
idealt Princip. Men herved have de afsløret Materialismens
Hemmelighed. Selv Vogt og Moleschott gaa i
Grunden ud fra et Idealt: Erkjendelsens Sandhed og
Ret og den historiske Udviklings Formaal. Materialismen
slaaer ikke blot ved konsekvent Gjennemførelse
over i sin egen Modsætning, men forudsætter endog denne
som sit Motiv. Vi have dvælet længere ved Czolbe,
netop fordi han saa klart viser dette. Ikke blot er han
den ædleste og konsekventeste blandt Materialismens
Forkæmpere, men han har ogsaa den tydeligste Bevidsthed
om, at dens Betydning i Nutiden overveiende er
kulturhistorisk\footnote[1]{Czolbe optræder i det Hele med en hos Materialisterne sjelden Beskedenhed og Resignation. Han begynder ikke som disse med at bryde Staven over alle tidligere Lærdomme og Institutioner; ogsaa her er han konsekventere end hine, idet han anerkjender det historisk Givne som nødvendigt, medens hine vilde være Reformatorer. I inderlig Overbevisning om sin Sags Sandhed er han ogsaa forvisset om, at den vil seire, naar den naturlige Udvikling fører dertil. Derfor er han tolerant, hvilket heller ikke er nogen materialistisk Dyd, skjøndt den konsekvent maatte være det.}. Den moderne Materialisme betegner
den mest yderliggaaende Reaktion mod den absolute
Idealisme. Den er et radikalt Udtryk for Tidsalderens
realistiske Retning, et memento om theoretiske og praktiske
Problemer, der kræve den alvorligste Overveielse. Og
hvis det skulde forholde sig saaledes, som en pessimistisk
Betragtning kunde føre til at mene, at vi gaa en
ny Middelalder imøde, behersket ikke som den første af
% --- p. 117 ---
\setcounter{footnote}{0}%
\opage{117}Spiritualisme. men af Materialisme, da kunne vi i ethvert % PRINTED AS IS (the point after Spiritualisme, for a comma; the image, s01)
Tilfælde trøste os med, at denne allerede har erklæret
sin egen Usandhed baade i sit Princip og i sine Konsekventser.


\chaphere{FJERDE KAPITEL.}


Overfor en storartet Tankeudvikling som den absolute
Idealisme, hvis Ensidighed er bleven aabenbar, kan
der indtages forskjellige Standpunkter. Man kan aldeles
forkaste den og uden at indlade sig i nogensomhelst Forhandling
med den søge sig nye Udgangspunkter. Saaledes
Materialismen, der med barbarisk Ignorants ser ned
paa den foregaaende Tidsalders Aandsbevægelse. Man
kan ogsaa arbeide sig ud af den ved at opløse den Enhed,
til hvilken den havde ført Tilværelsens Elementer
tilbage, og derpaa ad de formelle Konsekventsers Vei
føre ethvert af disse Elementer ud i deres abstrakte
Extrem. Saaledes den kritiske Bevidsthed. Men i Modsætning
til disse Retninger er da den Bestræbelse naturlig
saameget som muligt at korrigere det overvundne
Standpunkt ved nye Elementer, der ikke opløse det, men
bringe det Sandheden og Virkeligheden nærmere. Denne
Bestræbelse redder den historiske Kontinuitet. Den ensidige
Kritik skjærer Traaden over og hindrer den virkelige
Tilegnelse af Udviklingens Resultater. For at denne
skal være mulig, forudsættes nødvendig en sympathetisk
Dvælen ved det Tidligere.

Vel kunde det synes, som om en saadan Bestræbelse
var umulig ved et Tankesystem som den absolute
Idealisme. Netop ved sin absolute Karakter synes dette
System ikke at være skikket til at indgaa som Element
% --- p. 118 ---
\setcounter{footnote}{0}%
\opage{118}i en mere omfattende Sammenhæng. Det vil mediere
Alt, men synes ikke selv at kunne medieres. Alt eller
Intet, synes det her at maatte hedde, som jo og den
kritiske Bevidsthed under sin Udvikling med rivende Fart
omsatte Hegelianismens Alt til et absolut Intet.

Men det var netop som allerede tidligere bemærket
Spekulationens Feilgreb, at den satte som absolut, hvad
der kun er relativt. Spekulationen havde Ret i, at Tankens
Væsen ikke er rent subjektivt og formelt, afsondret
fra al indre Sammenhæng med Virkeligheden; Tanken
er et Virkelighedselement ligesaavel som den ydre Virkelighed,
der i sin rene Objektivitet ogsaa kun er et Ele
ment. Ved at udvikle sit eget Væsen og Indhold efter
dets indre, nødvendige Love udvikler Tanken derfor forsaavidt
Virkelighedens Væsen. Men netop dette \emph{forsaavidt}
oversaa Spekulationen. Den betænkte ikke, at
hin Udvikling kun har ideel Gyldighed, men ikke udtømmer
hele Virkelighedens Væsen. Derved blev Spekulationen
Panlogisme; Virkelighedens Indhold saaes
fuldstændig udtrykt i Tankens Indhold. Ved en optisk
Illusion forvexledes den ideale og den reale Virkelighed.
Grunden til denne Illusion søge de Forskere, vi nu gaa
over til, deri, at Spekulationen under sin storartede Udvikling
glemte sit oprindelige Udgangspunkt, den kantiske
Fornuftkritik. Kant havde netop gjennemført og
fastholdt Sondringen mellem Tænken og Væren, det
Subjektive og Objektive, og det saa absolut, at der, som
vi have set i første Kapitel, behøvedes Spekulationens
hele Kraft for at overvinde den. Men denne Sondring
hvilede paa den uomstødelige Sandhed, at den menneskelige
Erkjendelse maa vende sig kritisk mod sig selv og
ikke umiddelbart betragte sine egne Former som virkelig
Væren. Kun da bliver der ogsaa Plads for en principiel
Indflydelse af Erfaringen paa den hele Verdensanskuelse.
Disse to Sider, Opgivelse af „Tankens Monisme“ i Betydning
af en umiddelbar Enhed af Tænken og Væren,
% --- p. 119 ---
\setcounter{footnote}{0}%
\opage{119}og Hævdelsen af det reale Element hænge nøie sammen
og vise tilbage til Kant\footnote[1]{Smlgn. med Hensyn til den Betydning, der nu igjen tilskrives Kant: C. H. \emph{Weisses} akademiske Tale: „In welchem Sinne die deutsche Philosophie jetzt wieder an Kant sich zu orientiren hat.“ Leipzig 1847. I. H. \emph{Fichte}: „Bericht über meine philosophische Selbstbildung“ (Vermischte Schriften 1ster Bd. 1869). --- I Berlin anbefalede \emph{Schelling} Studiet af Kants Kritik der reinen Vernunft som Grundlag til Forstaaelse af sine Forelæsninger.}.

Der reiser sig da det Spørgsmaal, hvilken Gyldighed
og Betydning den spekulative Erkjendelse nu kan
faa. Dette Spørgsmaal besvares forskjellig af de Forskere,
vi ville fremdrage som Repræsentanter for denne
Retning. Deres første Udvikling faldt under den spekulative
Philosophis Herredømme, og de søge saa hver paa
sin Vis at fastholde det Berettigede i den. Denne Bestræbelse
fremtræder navnlig som en konkretere Udvikling
af Gudsideen, hvorfor vi have betegnet denne Retning
som den spekulative Theisme. Skjøndt der paa
flere Punkter viser sig Indflydelse af positiv-dogmatiske
Forestillinger hos disse Tænkere, vilde man dog gjøre
dem Uret i at mene, at der til Grund for deres Bestræbelse
laa den Tendents at hævde Dogmet om en personlig
Gud for enhver Pris\footnote[2]{I Weisses Skrift: „Das philosophische Problem der Gegenwart. Sendschreiben an I. H. Fichte“ hedder det: „Ich weiss wohl, dass Sie sich dieses Ansdrucks [der göttlichen Persönlichkeit, der Persönlichkeit Gottes], den die Gegner freilich jetzt als das Schibolet Ihrer Philosophie, so wie nicht minder der meinigen zu betrachten lieben, in Ihren früheren Schriften nur anbequemungsweise, und nicht ohne den Vorbehalt, dass er auf Gott nur eine uneigentliche Anwendung leide, bedient habe. Auch ich meinerseits habe, obwohl aus andern Gründen und in anderer Weise als Sie, mich stets dagegen verwahrt, dass man mir das Bekenntniss des Begriffs der göttlichen Persönlichkeit ohne Weiteres und mit der Voraussetzung unterlege, als sei mir dieser Begriff ein für sich selbst klarer, und bedürfe es über die Bedingungen seiner Gültigkeit nicht \opage{120}erst einer nähern wissenschaftlichen Verständigung . . . als sei es mir im Allen nur um die Persönlichkeit Gottes quand même zu thun.“ (p. 10).}. Deres Bestræbelse har sin % PRINTED AS IS (Ansdrucks, in the note: for Ausdrucks; the note reader and the collation)
% --- p. 120 ---
\setcounter{footnote}{0}%
\opage{120}naturlige Grund i deres Opposition mod den spekulative
Panlogisme og Pantheisme\footnote[1]{Kun for den yngre Fichtes Vedkommende kunde der, som vi ville faa at se, være Grund til at tale om en ligefrem Tendents.}.

Ad den dialektiske Spekulations Vei kan man efter
deres Anskuelse kun naa den almindelige Horizont, hvorunder
den hele Virkelighed fremtræder, de høieste, mest
abstrakte Love for al Tilværelse. For at opnaa en Virkelighedserkjendelse
i egentlig Forstand maa man gjøre et
Skridt ud over hin Kreds af almene Love og Bestemmelser.
Dette Skridt bestaaer paa engang i en Tilegnelse
af den konkrete Erfaring og i den videre Udvikling af
Gudsideen fra den almene Fornuft til det bevidst villende
og virkende Væsen, som ene afgiver Begrundelsen
af den reale Virkelighed. Spekulationen gaaer paa dette
Punkt over i philosophisk Empiri og Religionsphilosophi.
Der er altsaa i denne Retnings Tankegang en dobbelt
Bevægelse, en opadstigende og en nedadstigende: først
søges der op til det høieste Princip, og naar dette er
vundet og sat som virkeligt, søges den givne Tilværelse
forklaret deraf.

Det følger af sig selv, at det Punkt, hvor den ene
af disse Bevægelser gaaer over i den anden, ogsaa maa
blive det vanskeligste Punkt i disse Anskuelser. Begrebsbevægelsen
skal føre til sin egen Ophøren, Aktiviteten
gaa over i Receptivitet. Dette Vendepunkt bliver en
Standsning, og Bevægelsen danner ikke en Ellipse, men
en Vinkel. Dette kommer, som vi allerede ovenfor (p. 24 f.)
have havt Leilighed til at bemærke, deraf, at man endnu
paa mange Maader staaer paa det spekulative Standpunkt,
man søger at komme ud over. Ikke uden Grund
har man fundet en Dualisme i saaledes at supplere den
% --- p. 121 ---
\setcounter{footnote}{0}%
\opage{121}spekulative Tænkning med Elementer, som ikke organisk
kunne udvikles deraf. Paa en saadan Anstødssten strande
jo i Regelen Forsøgene paa at forsone det Gamle med
det Nye. Dog bliver derfor denne Retnings Betydning
ikke ringe, og navnlig, som det vil vise sig, i religionsphilosophisk Henseende.

I. Schellings \emph{senere} Lære. --- Det er en af Heroerne
fra Spekulationens glimrende Periode vi først
komme til at beskjæftige os med. Schelling (f. 1775,
d. 1854), der i Aarhundredets Begyndelse optraadte med
sin Ungdoms geniale Værk, hvor den absolute Idealismes
Grundtanke henkastedes i kjække og store Træk, fremstod
henimod Aarhundredets Midte med en Lære, der
paa engang skulde fastholde den absolute Idealisme og
ud over dennes Omraade føre ind til en hel ny, hidtil
kun dunkelt anet Verden. Han er et Exempel paa ualmindelig
Livskraft. Tidlig moden (Magister i sit 17de,
philosophisk Forfatter i sit 19de Aar) lagde han ogsaa
tidlig Grunden til det System, der har sikkret ham en Plads
i Philosophiens Historie. Men trods denne tidlige Fuldendelse
i det Væsentlige, hvilede hans Tanke dog ikke;
den arbeidede paa en stedse videregaaende Udvikling af
de oprindelige Principer. Den senere Lære er opstaaet
ved en saadan organisk Udvikling, hvor de ledende Motiver
paa ethvert Sted kunne paavises. Og dog kommer
Resultatet af den senere Lære til at danne en fuldstændig
Modsætning til den ældre Lære. Hvad der tidligere
var det Eneste og Absolute, nedsættes senere til et Foreløbigt
og Relativt. Vi skulle nu i Korthed se, hvorledes
Schelling kom til den Anskuelse, vi her have med at gjøre.

Til enhver Tid --- siger Schelling i et af sine ældste
Skrifter --- staa to Hovedretninger i den philosophiske
Stræben overfor hinanden. Den ene tager sit Udgangspunkt
i det Værende og Reale og søger deraf at forklare
Tanken og det Ideale; den anden gaaer omvendt
ud fra det Ideale og Subjektive og søger at naa det Virke%
% --- p. 122 ---
\setcounter{footnote}{0}%
\opage{122}lige. Ingen af disse Retninger kan naa en absolut Fuldendelse,
da Modsætningen bestandig bliver staaende paa
Grund af det ensidige Udgangspunkt. I det Absolute
maa derimod begge Modsætningerne være Et, og den
sande Philosophi maa derfor indtage det absolute Standpunkt,
Fornuftens Standpunkt, der forholder sig ligegyldig,
indifferent til hine Modsætninger. Forsaavidt disse
alligevel træde frem i det Endelige, bliver deres Forhold
dog kun kvantitativ Differents, en Overvægt af det Ideale
over det Reale (Aanden) eller af det Reale over det
Ideale (Naturen). Men denne kvantitative Differents er
kun i Phænomenet; Fornuften, det Absolute, Guddommen
er i sig selv absolut Indifferents. Kun forsaavidt
dette absolute Standpunkt forlades og der sættes et Reflexionsforhold,
har den endelige Verden Realitet\footnote[1]{„Briefe über Dogmatismus und Kriticismus“ (1796). „Darstellung meines Systems“ (1801).}.

Men det Spørgsmaal kan ikke trænges tilbage, hvorfor
det da i det Hele er nødvendigt at gaa ind i dette
Reflexionsforhold istedetfor at fastholde det absolute Standpunkt.
I Skriftet „Philosophie und Religion“ (1804) søger
Schelling at besvare det og fremdrager derved det Punkt,
hvor Spiren til hans senere Anskuelse ligger. Han negter,
at der er en kontinuerlig Overgang fra det Absolute
til det Endelige. Sandseverdenens Oprindelse er kun
tænkelig ved et Brud, et Affald fra det Absolute, thi
der er i Gud ingen positiv Grund til de virkelige Tings
Opstaaen. I Jegheden ligger Principet for Affaldet og
dermed for Endeligheden. Affaldet, der er evigt som
det Absolute selv, bestaaer nemlig i en Væren i sig selv
istedetfor en Væren i det Absolute. Derved ophæves
den absolute Indifferents og Modsætningerne træde i Kraft.

Der er altsaa et irrationelt Element i det Endelige,
idet Væren ikke følger af Begrebet, det Særlige ikke af
det Almene, men noget af Begrebet Uafhængigt maa
% --- p. 123 ---
\setcounter{footnote}{0}%
\opage{123}komme til. Dette Irrationelle er et Lavere end det Absolute. Derimod er det umuligt at antage en høiere
Potents over det Absolute; derved vilde dette jo netop
nedsættes til et Relativt. Det Absolute er som saadant
Gud\footnote[1]{Allerede i et Brev til Hegel af 4de Februar 1795 siger Schelling: „Für uns sind die orthodoxen Begriffe von Gott nicht mehr. Wir reichen weiter noch als zum persönlichen Wesen.“}. ---

For en videre Udvikling stillede der sig nu en dobbelt
Mulighed. Dersom den absolute Identitet fastholdes,
vil det Endelige kun kunne blive et Negativt, en
ydre, usand Form for det Absolutes indre, positive Enhed.
Denne Tankegangs Gjennemførelse vilde, som Weisse
rigtig har bemærket, have ført til den hegelske Anskuelse.
Men den anden Mulighed er den, at hint irrationelle
Element, som er Principet for den endelige Frihedsudvikling,
accentueres stærkere, betragtes fra et mere
positivt Synspunkt. Da vil det blive nødvendigt ligesom
at lade den absolute Identitet blegne og at finde (hvad
Schelling i den nys anførte Afhandling udtrykkelig negtede
Muligheden af) et mere fyldigt og aktivt Udtryk
for det høieste Princip; thi den absolute Identitet tilsteder,
som vi have set, ikke nogen positiv Frihedsudvikling
i det Endelige. Denne sidste Vei slog Schelling
ind paa allerede i en Afhandling fra 1809: „Philosophische
Untersuchungen über das Wesen der menschlichen
Freiheit und die damit zusammenhängenden Gegenstände.“


Identitetssystemet, hedder det her, har overvundet
Modsætningen mellem Tanken og dens Gjenstand, mellem
Fornuften og Virkeligheden. Men der viser sig nu
en høiere Form for det samme Problem, idet der spørges
om Forholdet mellem Frihed og Nødvendighed.
Undersøgelsen heraf finder et Tilknytningspunkt i de
tidligere Undersøgelser ved den vigtige Adskillelse, Iden%
% --- p. 124 ---
\setcounter{footnote}{0}%
\opage{124}titetslæren gjorde mellem Væsenet som existerende og
Væsenet som Existentsens Grund. Denne Adskillelse
kom frem i Naturphilosophien, idet Schelling opfattede
Naturen som en Skala af kvantitative Differentser eller
Potentser, af hvilke den ene bestandig danner Basis for
den følgende. I denne Sondring mellem Grund og Existents
ser han den afgjørende Forskjel mellem sin
egen Lære og Spinozas. Grunden er ikke Gud absolut
betragtet ɔ: som existerende; den er det i Gud,
som ikke er Gud, Naturen i Gud, hans Prius (dog ikke
i Tid eller Væsen). Alle Ting ere i Gud, ikke i hans
egentlige, aktuelle Væsen, men i hans Væsens evige
Grund. Og Modsætningen mellem Grund og Existents
er ikke Dualisme, thi den forudsætter en Ur-Enhed, Urgrunden
eller, efter Jakob Böhmes Udtryk, Ugrunden,
som ligger forud for alle Modsætninger. Urgrunden er
den absolute Indifferents, som vi allerede kjende fra den
tidligere Lære, det Absolute som absolut (das schlechthin
betrachtete Absolute). Men denne Indifferents er altsaa
kun det absolute Udgangspunkt, den almene Forudsætning
for den konkrete Udvikling, som bestaaer i, at Modsætningerne
sættes og atter overvindes til Identitet ved
Kjærligheden, i den personlige Existents.

Tidligere kjendte Schelling ingen afgjørende Modsætning
mellem Indifferents og Identitet (Forskjelsløshed
og Enhed). I den authentiske Fremstilling fra 1801 bestemmes
udtrykkelig Fornuften først (§ 1) som Indifferents
og sættes dernæst som Et med Identiteten (§ 9).
I et noget senere Arbeide, som først blev trykt efter
hans Død (Werke 1ste Abth. 5ter Bd), skjelner han nøiere
saaledes, at det Absolute kaldes Gud, men Fornuften
Guds Reflex, der forholder sig til Gud som Indifferents
 til Identitet. Men ogsaa her fastholder han udtrykkelig
 det pantheistiske Standpunkt og behandler alle Problemer
ud derfra. Det Samme er Tilfældet i Skriftet „Philosophie
und Religion“. Naar Schelling altsaa i Afhand%
% --- p. 125 ---
\setcounter{footnote}{0}%
\opage{125}lingen om Friheden vil have hele sin tidligere Lære betragtet
som kun omfattende, hvad han der kalder Urgrunden
og Grunden (Naturen), da beroer dette paa
en Illusion; hans tidligere Standpunkt var absolut og
afsluttet i sig, men han søgte nu ud over det, søgte en
Existents, en Virkelighed, som han ikke fandt indenfor
dets Grændser.

Medens hans tidligere Anskuelse har en antik Karakter,
dels ved den Vægt, der lægges paa Naturen, dels
ved at Kunsten og Staten ses som Aandens høieste
Sphærer, kaster han sig nu over i en middelalderlig
Theosophi, hvor man dog trods meget rent Scholastisk
gjenkjender den geniale Opfinder af Identitetslæren i
mange overordentlig dybsindige Undersøgelser. Men hans
Sikkerhed og Seirsfølelse har forladt ham. Man har
træffende sammenlignet ham paa hans tidligere Standpunkt
med den store Napoleon, der med genial Sikkerhed
tvang Alt til at bøie sig for den overlegne Magt.
Og dette Billede har man gjennemført videre ved at
sammenligne Schelling i hans senere Lære med Napolens % PRINTED AS IS (Napolens: for Napoleons (Napo-|lens); the image, t01)
Forsøg paa konstitutionel Regering i de 100 Dage.
Her er den samme Vaklen og Ubestemthed i den hele
Optræden, Tegn paa, at Skueplads og Skuespiller ikke
længer passe for hinanden. „Istedetfor det Overmod,
hvormed Idenditetslæreren legede med Videnskabens kronede % PRINTED AS IS (Idenditetslæreren: for Identitetslæren, in a quotation; the image, s01)
Hoveder, er der traadt en usikker, ja ængstelig
Hensyntagen til det bedømmende Publikum“\footnote[1]{\emph{Erdmann}: Ueber Schelling, namentlich seine negative Philosophie. Halle 1857. p. 9.}. Schelling
gjorde mange Tilløb til at offentliggjøre sin nye
Lære, men trak sig stedse tilbage igjen. Endog trykte
Værker lod han tilintetgjøre, idet han betænkte sig i
det sidste Øieblik før Udgivelsen. Men paa samme Tid
offentliggjorde han Udtalelser fulde af Overmod og Selvtillid og af stor Bitterhed mod Hegels Lære, der havde
% --- p. 126 ---
\setcounter{footnote}{0}%
\opage{126}overfløiet hans og nu ved sin rige og storartede Systematik,
sin sindrige Gjennemførelse og sine store Forjættelser
i flere Decennier herskede over Aanderne. Først
nogle og tredive Aar efter Udgivelsen af Afhandlingen
om Friheden traadte han i Berlin op for en større Kreds
med sin nye Lære, og i authentiske trykte Skrifter udkom
den først efter hans Død.

I en af hans faa offentlige Udtalelser i hint Mellemrum
bestemmer han i almindelige Træk Modsætningen
mellem sin tidligere og senere Lære\footnote[1]{Fortale til den tydske Oversættelse af et Skrift af Cousin 1834 (Werke 1ster Abth. 10te Bd).}. Her fremkommer
for første Gang den berømte Adskillelse af negativ
og positiv Philosophi. Philosophien maa begynde med
det, der nødvendig maa tænkes, det, der ikke ikke-kan
tænkes. Dette, uden hvilket ingen Erkjendelse er mulig
(det Aprioriske), er det Negative i Erkjendelsen;
det, ved hvilket den virkelige Erkjendelse bringes istand,
er det Positive. Det absolute Prius, endog for Guddommen,
indeholdes i Fornuftvidenskaben, som udgrunder
det Værende selv, det høieste Værende (Subjekt-Objekt).
Dog maa dette ikke misforstaaes, som hos
Hegel („ein später Gekommener“), som om denne Fornuftvidenskab
skulde være rent abstrakt, en blot logisk Proces.
Denne Misforstaaelse har kun afgivet et Bevis mere
paa, at det er umuligt at naa Virkeligheden ud fra det
rent Rationelle. Ikke Væren, som Hegel vil, men det
Værende er Philosophiens Begyndelse; og naar Philosophien,
som negativ Philosophi, har udviklet dette efter
alle dets Væsens Muligheder, en Udvikling, der fører
os ud over Modsætningen mellem Rationalisme og Empirisme,
saa opstaaer Spørgsmaalet om det, der ikke
længere er negativ Betingelse, men som er positiv Aarsag.
Men her behøves da et nyt Princip, en ny Begyndelse:
„Jeg vil ikke det blot Værende; jeg vil det
% --- p. 127 ---
\setcounter{footnote}{0}%
\opage{127}Værende, som er eller \emph{existerer}“. I denne Betydning
forestaaer der Philosophien en stor og i Hovedsagen
sidste Forandring, som paa den ene Side vil yde den
positive Forklaring af Virkeligheden, medens paa den
anden Side Fornuften anerkjendes at være i Besiddelse
af det absolute Prius, selv for Guddommen.

Medens Polemiken mod Hegel i denne Udtalelse
vakte Forargelse blandt Hegelianerne, vandt de videre
Antydninger Bifald hos Mange, der allerede iforveien
vare traadte i Opposition mod det hegelske System. Dog
vare disse Antydninger endnu vel ubestemte; man kunde
endnu ikke se, hvilken Karakter Modsætningen mellem
negativ og positiv Philosophi vilde faa. Aldrig er derfor,
som man med Rette har sagt, Schelling bleven rost
fra saamange forskjellige Sider; Enhver tydede hans Udtalelse
efter sin Mening.

I Aaret 1841 blev han opfordret til at holde Forelæsninger
i Berlin --- for at modarbeide den negative
Kritik, som havde udviklet sig af Hegelianismen. Dette
System, der i sin Tid kunde gaa for det Bestaaendes
Philosophi, var nu i radikal Opposition, og Schellings
„positive Philosophi“ skulde stilles i Breschen. Fra et
theologisk Katheder i Berlin blev hans Optræden hilset
som en Forsynets Vending. Den første Forelæsning,
som han udgav, er ikke blot mærkelig ved den Selvtillid
og Myndighed, hvormed han udtaler sig\footnote[1]{En Modstander sagde om denne Forelæsning, at den i Tydskland blev læst som en Throntale: „og Ligheden var stor nok; Taleren udtrykte sig med megen Værdighed om sig selv, gjorde store Løfter og undgik Alt, hvad der kunde bringe ham i Forlegenhed.“ \emph{Noack}: Schelling und die Philosophie der Romantik. 2ter Bd. p. 466.}, men ogsaa
ved en tydeligere Bestemmelse af Forholdet mellem negativ
og positiv Philosophi. Han erklærer, at han ikke
er kommen for at nedrive, men for at opbygge. „Intet
skal ved mig gaa tabt af hvad der siden Kant er er%
% --- p. 128 ---
\setcounter{footnote}{0}%
\opage{128}hvervet for den ægte Videnskab; hvorledes skulde jeg
tilmed opgive den Philosophi, som jeg selv har grundlagt,
min Ungdoms Opfindelse? Ikke at sætte en anden
Philosophi i dens Sted, men til den at føie en ny Videnskab,
som man hidtil har anset for umulig, for derved
igjen at befæste den paa dens sande Grundlag, at gjengive
den den Holdning, som den havde tabt ved at gaa
ud over sine naturlige Grændser, idet man gjorde Noget
til det Hele, som kun kunde være et Brudstykke af et
større Hele, --- det er Opgaven og Hensigten.“ Philosophien
(ɔ: den hegelske) forsikrer at være religiøs i
sine Resultater, men dette bestrides fra Livets Side.
Dette Misforhold mellem Videnskab og Liv, som har
fremkaldt saa stærk en Reaktion, maa hæves; man maa
bringe Philosophien ud af dens vanskelige Stilling. Og
dette véd Schelling sig istand til, da han er i Besiddelse
„ikke af en Intet forklarende Philosophi, men af en saadan,
der yder Oplysninger, som man har ønsket med
Længsel og indstændig forlangt, --- en Philosophi, der
udvider den almindelige Bevidsthed ud over dens nuværende
Grændser“\footnote[1]{Erste Vorlesung in Berlin. Werke 2te Abth. 4ter Bd. p. 360 ff.}.

Modsætningen mellem negativ og positiv Philosophi
beroer altsaa paa, at den sidste i sit Væsen er religiøs
Philosophi. Og sammenholdes dette med Ytringerne i Fortalen
til Cousin, da bliver Forholdet det, at den positive
Philosophi gjennem det Religiøse finder Veien til en
Opfattelse og Forklaring af den virkelige Existents. I
de efterladte Skrifter (Werke, 2te Abth.) finde vi nu
den videre Udførelse deraf.

Der maa skjelnes mellem det Værende og det, som
er det Værende, det existerende Subjekt. Det Første
er det Virkelige, det Andet det Virkeliges Virkelighed.
Afset fra Virkeligheden er det Virkelige kun muligt, er
% --- p. 129 ---
\setcounter{footnote}{0}%
\opage{129}Virkelighedens mulige Indhold. Det Virkelige er et
Hvad (qvid), Virkeligheden, Existentsen et At (qvod). Til
den uendelige Mulighed af Væren svarer den uendelige
Mulighed af Erkjendelse, og gaaer man ud fra Fornuften
som det Princip, i hvilket begge Muligheder ere Et,
idet deres Modsætning er ophævet, saa bliver det Fornuftvidenskabens,
den første Videnskabs Opgave at gjennemgaa
det Virkeliges uendelige Potentser, udtømme
deres hele Fylde. Idet denne første Videnskab baner
Veien for Erkjendelsen af den sande, positive Virkelighed
ved at udskille det, der kun har Existentsens Mulighed,
er den kritisk, og svarer til Kants Kritik der
reinen Vernunft; den bliver herved tillige negativ, forsaavidt
den forholder sig afværgende mod det blot Værende
(det, som ikke er det Værende); og idet den kun
beskjæftiger sig med det, som nødvendigvis maa tænkes,
følger den den rene logiske Nødvendigheds Lov og er
rational og apriorisk Philosophi. Men dette Rationale
maa ikke forstaaes som Modsætning til det Empiriske.
Thi for at det Værendes Muligheder kunne udtømmes,
er det nødvendigt ogsaa at tænke dem som værende
og virkelige, hvorved de blive Gjenstand for Erfaringen.
Tanken gaaer ind i Erfaringen og følger den fra Potents
til Potents. Erfaringens Grændse er her, som Kant
lærte, ogsaa Tankens Grændse, fordi begges Væsen er
Et. Da man nu ikke fastholdt, at det virkelig Existerende
(det, som er det Værende) slet ikke har Noget
at gjøre med hele denne Udvikling, der drager sig gjennem
Tankens, Naturens og Historiens Riger for at udtømme
Mulighedernes uendelige Indhold, --- da man
(Hegel) sammenblandede Subjekt og Prædikat, Existents
og Potents, saa kom man til den Lære, at Gud maa
gaa gjennem Naturens og Historiens Sphærer for at
komme til Bevidsthed om sig selv, en Pantheisme, der
i sine Konsekventser nødvendigvis maatte føre til Atheisme.
Identitetslæren var ogsaa Pantheisme, og dette
% --- p. 130 ---
\setcounter{footnote}{0}%
\opage{130}er forsaavidt ogsaa nødvendigt for den negative Philosophi,
der stedse bevæger sig indenfor det Absolute,
den almene Ides Omraade og ikke kan gaa ud derover.
Men denne Pantheisme har kun logisk Betydning; thi
den negative Philosophi opkaster endnu ikke Spørgsmaalet
om det Existerende, men holder sig indenfor det
rene Begreb.

Hvilke Muligheder indeholdes der nu i dette rene
Begreb? Det Værende er for det Første (som primum
cogitabile) den rene Kunnen, potentia, det i sig selv
værende Subjekt; men medens dette mulig Værende (— A)
stedse gaaer ud over sig selv, har sin Væren udenfor
sig, er det andet, til dette svarende Moment det Værende
som det rent Værende (+ A), der er actus purus, rent
Objekt. Som disse Momenters Enhed, som den høieste
Ide (summum cogitabile) er det Værende endelig Subjekt-Objekt (± A), Enhed af potentia og actus, alle Muligheders
Indbegreb. Alle tre Momenter tilsammen er
det Absolute, men vel at mærke, bestandig kun som
Mulighed, som Ide. Her støder Tanken da paa sin
sidste Grændse, idet den opkaster Spørgsmaalet om,
hvad det er, som er Ideen, er det høieste Tænkelige.
Men dette Spørgsmaal kan endnu ikke besvares; det
kan end ikke i afgjørende Forstand komme frem paa
dette Stadium.

Ved at opsøge Potentserne i det Absolute og følge
dem i deres mulige Funktioner komme vi ikke udenfor
Ideen. Det mulig Værende vil, naar det sættes i
Funktion, brede sig i kvantitativ Uendelighed som det
græske ἄπειρον, uden Maal og Skranke. Idet nu det
andet Moment, det rent Værende, som actus og bestemt
Grændse (πέρας), reagerer herimod og bringer Negationen
ind i den uendelige Vorden, frembringer den faste Kvaliteter
og individuelle Former (τὸ περαῖνον, der tillige er
τὸ εἰδοποιοῦν). Fremtræder endelig det tredie, høiere Moment,
da have vi det aristoteliske οὗ ἕνεκα, Enheden af
% --- p. 131 ---
\setcounter{footnote}{0}%
\opage{131}Form og Materie. Den relative Enhed af de tre Momenter (den absolute Enhed kan der først spørges om
ved den negative Videnskabs Ende) er Sjælen, som
Verdenssjæl og individuel Sjæl. Igjennem Sjælen staaer
alt Værende i Forbindelse med Ideen og det, som er
Ideen. Med disse fire Principer er hele Ideverdenen
givet. Idet det ene Moment viser hen til det andet,
ligger deri er Anelse om, at der er en høiere Væren % PRINTED AS IS (deri er Anelse: the image, er for en; O080)
end den blotte Mulighed, den idebundne Væren. Men
uagtet al Overgang fra Potents til Aktus er en Villen,
og vi saaledes i Forholdet mellem Momenterne i det Absolute
have en Urvillie, som da al Væren kun er ved
en Villen, en Selvhævden (Gjenstand = Modstand), --- saa
er denne Villen endnu dog kun, som det allerede blev
udtrykt i Afhandlingen om Friheden, en Hunger og Tørst
efter Existents, ikke Existentsen selv.

Skulle vi naa det Punkt, hvor denne uendelige Trang
tilfredsstilles, da maa vi forlade det Absolute, det blot
Værendes Standpunkt. Veien bort fra dette finde vi
ved at fremkalde den sidste Mulighed, som endnu er tilbage:
Muligheden af en Væren udenfor Ideen. Denne
Mulighed kan ikke søges i selve Ideen. Det er I. G.
Fichtes udødelige Fortjeneste at have paavist, hvorledes
den virkelige Verden først fremstaaer ved, at Jeget fra
den umiddelbare Enhed med det Absolute træder ud i
Reflexionsforholdet, sætter sig som Selvbevidsthed og
derved i Modsætning til en objektiv Virkelighed (smlgn.
ovenfor p. 122). Muligheden heraf ligger i, at Sjælen,
der ifølge det Foregaaende er Potentsernes relative Enhed,
kun er dette, forsaavidt den ikke løsriver sig fra
det Absolute. Bevarer den ikke Enhedsforholdet til
dette, men forsøger at gjøre den relative Enhed til en
absolut --- og dette sker ved, at den opgiver Gud (det,
som er det Værende) og i Ideens Iboen mener at have
og være Alt --- saa opstaaer derved en Verden udenfor
Gud, hvis trinvise Udvikling Fornuftvidenskaben maa
% --- p. 132 ---
\setcounter{footnote}{0}%
\opage{132}følge som den sidste Række af Muligheder, der ligge
indenfor dens Omraade. Herved først opstaaer Materien
i physisk Forstand, som vis inertiæ, og Tid og Rum
som Former for Tingenes Udenforsigværen (Excentricitet,
Forrykkelse fra Centrum). Verden kommer nu til
at staa som fremmed overfor den bevidste Aand, som
ved en Akt, der ligger ud over Bevidsthedens Grændse,
og som den derfor paa dette Stadium ikke kan opfatte,
er sat i denne Modsætning. Den eneste Vei til ved egne
Kræfter at komme ud over denne er at gjennemtrænge og
overvinde Verden ved sin Erkjendelse og sin Handlen.

Allerede de Gamle have paa dybsindig Vis udtalt
dette. Prometheusskikkelsen, en af de Urtanker, der
selv arbeide sig frem“, betegner Menneskeaanden som
den, der ved selve sin Guddomsgnist vender sig bort fra
Guddommen og ved eget Snille søger at grundlægge en
ny Verden uden Gud. Og som Prometheus lænkes til
Klippen, hvor hans Lever stedse fortæres af Ørnen, og
først kan befries fra sin Kval, naar en Gud staaer frem
for at udløse ham, saaledes vandrer Menneskeaanden nu
den lange Vei for at søge den virkelige Gud, den har
tabt, men som er den sidste Aarsag, den personlige
Bærer af de Ideer, paa hvis Besiddelse den har grundet
sin Autonomi.

Jeget har ikke blot at kæmpe med den materielle
Verden, som omgiver det; det finder sig ogsaa omsluttet
af ideelle Skranker, som Led i en aandig Helhed, hvor
det Enkelte underordnes Totaliteten, det Forudgaaende
det Efterfølgende. Denne totale Orden kan ikke have
nogen vilkaarlig eller tilfældig Oprindelse; det er en
Lov, som udspringer af det Værendes Væsen, Kilden
til al Ret og Pligt. Det er Kants uforgjængelige Fortjeneste
at have fremdraget denne Lov og paavist dens
Uafhængighed og Selvstændighed. Kun forsaavidt kunde
hans Lære om Fornuftlovens Autonomi misforstaaes,
som man ved Fornuft her vilde tænke paa den subjek%
% --- p. 133 ---
\setcounter{footnote}{0}%
\opage{133}tive, menneskelige Fornuft. Loven har nemlig som sagt
sin Grund i selve det Værende. Men som Fornuftlov
tør den ikke umiddelbar føres tilbage til Gud; det vilde
være at negte den rationelle Ethiks og dermed den rationelle
Philosophis Realitet. --- Som faktisk Grundlag
for Virkeliggjørelsen af den rationelle Lov træder Staten
(den juridiske Lovgivning) op over for den Enkelte, dog
ikke som en Indskrænkning af hans Frihed, men tvertimod
for at gjøre dem mulig. Staten er i sit Væsen
kun tjenende, Middel for et frit Samfundsliv, i hvilket
de personlige Dyder kunne udøves. Ved Hjælp af dette
faste Grundlag skal den Enkelte vinde den indre Uafhængighed,
som tillige er det bedste Værn mod politisk
Undertrykkelse. Vel bestaaer Statens Udvikling i, at
stedse Mere af den almene Fornuftlov gaaer over i Statens
Virkelighed; men Historiens Maal kan dog ikke
søges i en fuldkommen Stat. En saadan hører kun
hjemme i Idealets, i Mulighedernes Verden\footnote[1]{Tidligere (f. Ex. i „Vorlesungen über die Methode des akademischen Studiums 1803“) lovpriste Schelling Platons Stat og stillede det antike Statsliv over det moderne, fordi det ikke som dette havde Privatlivet udenfor sig.}.

Lovens og Statens almene, upersonlige Magt er en
Tugt, Mennesket er undergiven, og fra hvis Tryk han
søger at befri sig. Mennesket er hildet i den Modsigelse,
at han har villet sig selv i Modsætning til det
Absolute, men at denne Villen dog netop fører til Afhængighed
af upersonlige Magter. Der kommer da et
Vendepunkt, hvor Jeget atter fra den udadvendte Stræben
bøier sig tilbage i sig selv, til det beskuende Liv. Det
resignerer nu, opgiver at naa Maalet ad Handlingens
Vei og søger ved Fordybelse i Ideen at trænge frem til
den Gud, det længes efter. Det første Trin er her den
mystiske Fromhed, hvor Sjælen stræber at opgive sig
selv, sin egen Villie (se désapproprier sa volonté. Fé%
% --- p. 134 ---
\setcounter{footnote}{0}%
\opage{134}nelon) for ganske at hensynke i Gud, med uinterresseret
Kjærlighed. Hertil slutter sig for det Andet Kunsten,
i hvilken Menneskeaanden i fuldstændig selvløs Produceren
gjør sig Guddommen lig. Endelig søger Aanden,
for det Tredie, sin reneste, fra det Potentielle mest befriede
Form i selve den rationelle Videnskab. Den skuer
nu Gud som den høieste Tanke, den rene νόησις τῆς νοήσεως, det absolute Subjekt-Objekt. Men det er dog
kun en rent negativ Bestemmelse, under hvilken Gud
her endnu opfattes. Gud ses vel i sin rene Selvværen,
men kun som det Sidste, hvortil man kommer, ikke som
den levende Aarsag, hvorfra man kan gaa ud. For Aristoteles
er Gud saaledes udtrykkelig Enden, det sidste
Formaal, der selv Intet bevæger, ikke gaaer ud over
sig selv, men er ἄπρακτος τὰς ἔξω πράξεις. Derfor er det
Misbrug her at tale om den absolute Aand. Vi ere
endnu ved den blot ideelle, idebundne Gud, der vel er
summum cogitabile, men dog hører ind under den negative
Philosophis Omraade.

Hvorledes kunne vi da overskride dette Omraades
Grændser? Ikke ved Tænkningen. Thi denne følger
kun den rene logiske Nødvendigheds Lov og kan derfor
ikke gaa ud over Ideen, men søger kun at udtømme
alle dennes Muligheder. Det maa være i Villien, at
Fordringen opstaaer om at sætte Gud ikke længer som
blot Ide, men over al Tanke og Fornuft. Den absolute
Forsoning naaes ikke i den rationelle Sphære. Udviklingen
indenfor denne førte os paa sit Høidepunkt til det
kontemplative Liv (Mystik, Kunst og Videnskab); men
fra dette drives Mennesket bestandig tilbage til det praktiske
Liv: der maa handles. Nu gaaer da den nuværende
Tilstands hele Forfængelighed op for ham, og han
søger ud over Ideen til den levende og handlende Gud,
der ikke længere er fortabt i det Værende (Ideen), men
er Værens Herre. Gud stødes ud af Ideen, søges og
gribes som Personlighed. Det er ikke Fornuften, det
% --- p. 135 ---
\setcounter{footnote}{0}%
\opage{135}Almene, men Individet selv, der længes efter Salighed
og finder denne i den virkelige Gud.

Den negative Philosophi blev, idet den gjennemløb
Mulighedernes Skala, tillige en Videnskabernes
Encyklopædi. Ethvert enkelt Stadium indenfor denne
Udvikling svarer til en af de særlige Videnskaber. Disse
ligge saaledes paa Veien fra den negative til den positive Philosophi. Den negative Philosophi er den almindelige
Videnskab, der ingen særlig Gjenstand har,
netop fordi den omfatter Erkjendelsens hele Potents.
Men idet den nu har fuldendt Udviklingen og udtømt
Mulighederne, og idet Overgangen til den positive Philosophi
er gjort ved, at Gud er befriet fra Ideen og ses
som det, der er det Værende, faaer Philosophien en
anden Karakter. Den faaer nu sin egen bestemte Gjenstand, bliver en Videnskab, ikke blot Videnskaben.
Thi som det absolut Idefrie er Gud ikke et Alment,
men et Enkelt, et Dette (ἕν τι, τόδε τι ὄν), ikke værende
Noget, men værende absolut. Der existerer overhovedet
intet Alment. Kun det Enkelte \emph{er}; det Almene er kun,
forsaavidt det absolute Enkeltvæsen er det Almene, har
det til sit Prædikat. Det er urigtigt at sige, at Ideen
er den sande, den eneste Realitet. Ideen er kun, forsaavidt
Idealet er. Idealet er Værens Aarsag (αἰτία τοῦ
εῖναι) for Ideen. Er det absolut immanente Begreb (den % PRINTED AS IS (εῖναι: without its breathing, for εἶναι; the image, 600 dpi)
absolute Ide) den negative Philosophis Endepunkt, saa
er det absolut transscendente Begreb (det existerende
absolute Enkeltvæsen) den positive Philosophis Begyndelsespunkt.
Den negative Philosophi var apriorisk Videnskab,
idet den fulgte det Værendes Potentser for at
udtømme deres Væsen; dens Omrids var de evige Muligheders,
de evige Sandheders Rige. Den var kun en
tænkende, ikke en egentlig vidende Videnskab, og bevarede
stedse en hypothetisk Karakteer: „naar det Værende
er, saa maa det være saaledes beskaffent.“ A priori ɔ: forud for Væren. Potentsernes Rige er nu gjennem%
% --- p. 136 ---
\setcounter{footnote}{0}%
\opage{136}løbet, og Tanken er standset ved den absolute Aktus,
som den ikke kan begribe, netop fordi alt Alment, al
Mulighed er udelukket. Tanken kan kun komme dertil,
maa forholde sig rent skuende.

Der er ingen immanent Overgang fra negativ til
positiv Philosophi. Det Immanentes Rige er gjennemvandret.
Det absolut Existerende, den idefrie Gud er
som uafhængig af al Ide ogsaa uafhængig af den negative
Philosophies høieste Ide. Derfor kommer den positive
Philosophi kun til at staa i tilfældig Sammenhæng
med den negative. Den, som ikke føler Trang til at
gaa over til den positive, vil slaa sig til Ro i den
negative, og omvendt vil man umiddelbart kunne gaa
over til den positive, uden først at have gjennemløbet
den negative\footnote[1]{Werke, 2te Abth. 1ster Bd. p. 564.}. Den negative Philosophi kommer
gjennem Immanentsen, naar den forstaaer sig selv, til
det transscendente Væsen. Den positive Philosophi begynder
med det blot Existerende, den rene Transscendents
og søger nu at gjøre den immanent, paavise det
qvid, der bæres af dette qvod. Med andre Ord, medens
den negative Philosophi udviklede Guddommens Væsen
indtil det Punkt, hvor Spørgsmaalet opstod om dette
Guddommeliges Existents, begynder den positive Philosophi
med det absolut Existerende og paaviser, at det
er Gud. Den positive Philosophi begynder saaledes med
en Væren, der ligger over og forud for al Tanke (das
unvordenkliche Sein), og ser denne Væren som Bærer
af Potentsernes hele Rige. Ved Udviklingen af Potentserne
paa denne Basis erkjendes Gud som Værens Herre.
Navnet Spekulation tilkommer særlig den positive Philosophi.
Vi ere her ikke længere paa den nødvendige
Tænknings Omraade, men det gjælder nu en fri
ɔ: spekulativ Tænken. At spekulere vil sige at se sig
om efter Muligheder, ved hvilke et vist Øiemed kan
% --- p. 137 ---
\setcounter{footnote}{0}%
\opage{137}naaes i Videnskaben. Disse Muligheder (Oversætninger
i en hypothetisk Slutning) paavises som virkelige a posteriori.
Thi Methoden optager paa ethvert Punkt den
virkelige Erfaring: „Naar det nødvendig Existerende er
Gud, saa ere a, b, c virkelig; nu existere ifølge Erfaringen
a, b, c virkelig; følgelig er det nødvendig Existerende
virkelig Gud.“ --- Det første Spørgsmaal bliver
nu det om den virkelige Verdens Mulighed.

Den første Potents, det mulig Værende, har en
stadig Tendents til at udbrede sig, ligesom strømme over.
Og da Potentsernes Udvikling ikke blot skal være en
rent nødvendig, som den, ved hvilken den anden Potents,
det rent Værende (Sønnen), fører den første tilbage
til (Aandens) Enhed i Gud selv, --- da Gud ikke
blot vil være Gud for sig selv, men ogsaa meddele sit
Væsen, ligesom Geniet stedse trænger til at producere,
--- saa opstaaer Muligheden af, at han sætter de Potentser,
hvis Enhed han er, i Spænding, i en Modsætning,
der skal overvindes gjennem en successiv Proces. Herved
vendes Enheden om, og Universet (unum versum)
opstaaer. Men endnu kun i Gud. Muligheden af en
Verden extra deum (ikke blot præter deum) ligger i, hvad
allerede den negative Philosophi paa sin Vis lærte, at
Bevidstheden (Mennesket), i hvilken de udstrømmende
Potentser igjen vende tilbage til Gud og som altsaa er
Baandet mellem Verden og Gud, atter kan løsne dette
Baand, fordi den vil være Potentsernes Herre, ligesom
Gud er det. Hvad Mennesket glemmer, er, at medens
Gud ved sit Væsens Nødvendighed er Potentsernes og
Værens Herre, ser Mennesket det kun i relativ Forstand,
under Forudsætning af sin umiddelbare Enhed
med Gud. Det Moment, der frister Mennesket til atter
at sætte Potentserne i Spænding, er igjen hin første
Potents, det uendelig fødende og frembringende Princip,
der i Imaginationen fremstiller en Verdens Mulighed
for ham. Heri virker tillige den Magt, den græske
% --- p. 138 ---
\setcounter{footnote}{0}%
\opage{138}Mythologi fremstiller som Nemesis, den Verdenslov, der
vil, at Alt skal være klart og uden Tvetydighed, at intet
tilfældigt Forhold maa bestaa, at den uprøvede Lykke
er utilstedelig. I Christendommen fremstilles dette Princip
som Satan, der synder ἀπ’ἀρχῆς ɔ: er det solliciterende
Princip, der rører sig i Menneskets Villie, for at denne
kan blive prøvet. Opfattet efter sin philosophiske Ide
er Satan derfor det egentlige primum movens i Historien,
uden hvilket Verdenslivet vilde stagnere. Mod
denne Magt optræder den anden Potents (Logos, Sønnen),
ikke for absolut at negere den, men for at underlægge
sig den, gjøre den til Stof og Basis og saaledes
atter føre den tilbage til et indre Forhold til Gud. Det
er denne Proces, som Mythologiens og Aabenbaringens
Philosophi forfølger, hvorved den (den vigtigste Del af
den positive Philosophi) bliver den høieste Historiens
Philosophi.

Den Forudsætning, hvorfra saavel Mythologi som
Aabenbaring gaa ud, og uden hvilken de ere aldeles
uforklarlige Phænomener, er, at der har været et oprindeligt
realt Forhold mellem Mennesket og Gud. En Psychologi,
der er hentet fra de nu sædvanlige Bevidsthedstilstande,
forslaaer her ligesaalidt som de mechaniske
Love, der gjælde under Universets nuværende Tilstand,
kunne overføres paa den oprindelige Vorden, den første
levende Opstaaen. Heri ligger Grunden til, at alle tidligere
Forsøg paa at forklare Mythologi og Aabenbaring
ere slaaede feil.

Da Mennesket ved en Akt, der ligger ud over al
Historie og al Bevidsthed, satte Potentserne i Spænding,
ophævedes det oprindelige, virkelige Gudsforhold\footnote[1]{En anden Følge af denne Katastrophe er særlig bleven betragtet i den negative Philosophi, nemlig Tilblivelsen af en ydre, materiel Verden. Derimod kunde den negative Philosophi ifølge sit Væsen ikke gaa ind paa, hvorledes „Værens \opage{139}Herre“ stiller sig til hin Katastrophe, og hvorledes ifølge denne Forholdet bliver mellem ham og Mennesket, eller rettere, den kunde kun opfatte dette fra den negative Side.}% the note runs over from p. 138 to p. 139 (joined by draft.py)
, og
den uendelige Magt virker nu som en Kraft, der truer med
% --- p. 139 ---
\setcounter{footnote}{0}%
\opage{139}med at søndersprænge Bevidstheden, naar den ikke trænges % PRINTED AS IS (med at søndersprænge: the image, med doubled over the page turn (p. 138 ends "truer med"); O087)
tilbage indenfor faste Grændser. Den begrændsende,
Skridt for Skridt seirende Magt, der i Menneskeaanden
tager Kampen op med hin dæmoniske Kraft, og som
derigjennem selv arbeider sig frem til Virkelighed, er,
som allerede anført, den anden guddommelige Potents
(Sønnen). Denne theogoniske Proces er foreløbig rent
mythologisk. Den gaaer ud paa at materialisere den
første, chaotiske Potents, nedsætte den til Tjener, til
Stof og Grundlag. Polytheismen er det første Resultat
af denne Proces. Thi da den sande Enhed er sat ud
af Kraft, idet den første Potents vil herske alene, medens
de andre Potentser stedse stræbe igjen at indtage
deres Plads, saa forme under denne Strid de forskjellige
Stadier sig for Bevidstheden til forskjellige Gudeskikkelser,
Knudepunkter paa den Vei, Bevidstheden maa
tilbagelægge for at finde igjen, hvad den har forskjertset.

Det første Stadium er Stjernedyrkelsen. Her er
forsaavidt det egentlig Mythologiske endnu ikke begyndt,
som den virkende Potents endnu staaer paa Overgangen
mellem at være Centrum og Peripheri, ligesom ogsaa i
Planeterne den oprindelige Tendents til at være Centrum
endnu er tilstede, og ved bestandig at trænges tilbage
bliver Grunden til den bestandige Bevægelse. Det næste
Stadium er det, hvor Potentsen nedsættes til Materie,
til Peripheri, hvilket mythologisk betegnes ved, at der
optræder kvindelige Guddomme ved Siden af de mandlige
(Urania ved Siden af Uranos: Herod. I, 131; den
persiske, babyloniske og arabiske Religion). Som Mythologiens
Philosophi i det Hele er en Naturphilosophi,
kun i en høire Sphære (begge fremstille, hver paa sin % PRINTED AS IS (høire Sphære: the image, høire for høiere; O088)
Vis, en Theogoni), saaledes svarer dette Punkt i den
% --- p. 140 ---
\setcounter{footnote}{0}%
\opage{140}mythologiske Udvikling særlig til, hvad der maa tænkes
som Naturens egentlige Begyndelse, hvor det, der først
selv var et Værende nedsættes til Grund og Stof (mater—
materia). Snart optræder ogsaa den høiere Potents (Dionysos:
Herod. III, 8) i sine første Bebudelser, endnu
dunkle og uforstaaede. Men Guden selv er ældre end
hans Navn, hans Nærværelse i Bevidstheden er ældre
end hans fuldkomne Virkeliggjørelse deri. Nu fyldes
Bevidstheden af Kampen mellem de to Magter, den
materielle og den ideale Gud (phønicisk, phrygisk, ægyptisk,
indisk Mythologi). Med den materielle Gud frygter
Bevidstheden at miste Gud i det Hele; kun gjennem
Smerte og Død sker Befrielsen og den ideale Guds Seir.
Deraf Modsætningerne i den mythologiske Bevidsthed:
Sorg og Jammer fordi Guden er død, Glæde, fordi han
atter er født. Den græske Mythologi danner endelig
paa engang Afslutningen og Rekapitulationen af den
mythologiske Proces. I den græske Gudeverdens Historie
have vi i Successionen af Uranos, Kronos og Zeus
en Afspeiling af Kampen mellem det ideale og materielle
Princip, indtil tilsidst den klare olympiske Harmoni
behersker Alt. Men den græske Mythologi har foruden
denne exoteriske Side ogsaa en esoterisk, i Mysterierne,
hvis Indhold er den opdukkende Bevidsthed om Mythologiens
eget Væsen. Demeter betegner den høiere, ideale,
mod Gud vendte Bevidsthed. Hun sørger ved Persephones
Død, fordi hun med hende (Bevidstheden som
den, der er underkastet det materielle, mythologiske
Princip) troer at have mistet Alt. Den mysteriøse Forsoning
sker saa ved, at hun istedetfor den undergaaede
Gud faaer den blivende, som aldrig kan gaa under. Men
heri ligger et Fremtidens Moment: den mystiske Dionysos
(Ἴακχος) er den stedse kommende (Eleusis --- Advent).
Her er da Mythologiens Ende, hvor den ikke blot har
forklaret sig selv, men ogsaa vist ud over sig selv.

Overgangen fra Mythologi til Aabenbaring sker ikke
% --- p. 141 ---
\setcounter{footnote}{0}%
\opage{141}ved, at der optages et nyt substantielt Indhold. Forskjellen
beroer paa det virkende Princip. Baade Mythologi
og Aabenbaring forudsætte, som allerede tidligere
bemærket, et ud over al bevidst Tænken og Digten liggende
realt Gudsforhold. Men dette Forhold kan enten være
blot naturligt og ufrit, Nødvendigheden underlagt, hvor
da paa Grund af Blindheden og Hildetheden Overtroen
snart opstaaer. Saaledes i de naturlige, vildtvoxende Religioner.
Eller ogsaa er det ikke blot en naturlig Potents,
der virker, men Mennesket befries fra den blinde,
ufri Religion derved, at det er en Personlighed, der fremtræder
som virkende. Derfor er Christi Person og Historie
Christendommens sande Indhold, hvorimod de
christelige Dogmer kun ere Produkter af en jammerlig
Philosophi\footnote[1]{Tidligere (Methode des akad. Studiums 1803) havde Schelling netop rost Kirkefædrene, fordi de ved Philosophiens Hjælp havde forstaaet at udvikle et saa dybsindigt System som Kirkelæren af de bibelske Bøgers tarvelige Indhold, der staaer langt tilbage for de indiske Religionsbøgers.}. I Hedenskabet foregik Processen rent
indenfor Bevidstheden. En subjektiv Forsoning var her
tilstrækkelig, forsaavidt kun Følgerne af Adskillelsen fra
Gud skulde overvindes. Anderledes naar det gjælder
at ophæve selve Aarsagen, fuldstændig at bryde den
guddommelige Uvillie (den første Potents); dette kunde
kun ske ved en virkelig historisk Handling. Istedetfor
den ekstatiske Historie i Mythologien træder nu, i Inkarnationen,
den virkelige, ydre Historie.

Men den sande Forstaaelse af denne historiske Kjendsgjerning
har man hidtil ikke naaet, fordi man ikke satte
den i Forbindelse med den hele religionshistoriske Fortid.
Man saa ikke, at Christus implicite ogsaa var det
i Naturreligionen Virkende. Gjennem Mythologiens theogoniske
Proces har den anden Potents underlagt sig den
førte og staaer som Seirherre i Menneskeaanden. Spørgsmaalet % PRINTED AS IS (førte og staaer: the image, førte for første; O090)
er nu, hvorledes den, Menneskets Søn, som den
% --- p. 142 ---
\setcounter{footnote}{0}%
\opage{142}med Rette kan kaldes, da den menneskelige Bevidsthed
har bevirket dens Vorden under denne Form, vil benytte
den erhvervede Magt og Myndighed (μόρφη θεοῦ): om den % PRINTED AS IS (μόρφη: accent on the ο and none on the η, for μορφῇ; the Minnesota image, 600 dpi)
vil beholde denne for sig selv (ἁρπαγμὸν ἡγεῖσθαι Phil. II,
6) eller benytte den til at føre Verden tilbage til Gud
(παραδιδόναι τὴν βασιλείαν τῷ θεῷ καὶ πατρί. 1 Kor. XV, 24).
(Smlgn. Fristelseshistorien). Menneskeblivelsen, hvor den
anden Potents istedetfor Gudligheden paatager sig en
Tjenerskikkelse for at lide Døden, er Svaret herpaa. I
Inkarnationen og de historiske Begivenheder, der slutte
sig til den, har den indre, transscendente, guddommelige
Historie gjennembrudt den ydre Histories Skranker, og
denne faaer nu kun Betydning, forsaavidt den kan godtgjøre
sin Sammenhæng med hin. Nu kommer ogsaa
den tredie Potents, den hellige Aand, til Gjennembrud,
idet den første Potents ikke længere chaotisk og satanisk
unddrager sig den højere besjælende Enhed.

Loven for Kirkens Udviklingshistorie er indeholdt i
Forholdet mellem de tre Hovedapostle: Petrus staaer
parallelt med Moses, er Lovgiveren, det Stabiles Repræsentant;
Paulus er en Elias, Frihedens og Bevægelsens
Repræsentant; Johannes endelig er for Christendommen,
hvad Johannes den Døber var for Jødedommen,
Fremtidens Apostel. Den petrinske Kirke er derfor Katholicismen
med dens faste, lovbundne Enhed; den paulinske
er Protestantismen med dens radikale Frihedsprincip;
den frie Enhed har sin Type i Johannes, der
forener Peters Enfold med Pauli dialektiske Skarphed.
Thi Maalet, som skal naaes, er hin først i Sandhed almindelige
Kirke (om Kirke her endnu er det rette Ord),
der alene opbygges i Aanden og kun bestaaer i den fuldkomne
Forstaaelse af Christendommen og dennes virkelige
Sammensmeltning med den almindelige Videnskab og
Erkjendelse. Den frie Religion er vel indledet men ikke
virkeliggjort ved Christendommen. Bevidstheden maa
for at naa den frie Religion frigjøre sig fra Christen%
% --- p. 143 ---
\setcounter{footnote}{0}%
\opage{143}dommen ligesaavel som fra Naturreligionen; begge ere i
sig selv kun Kilder til ufri Erkjendelse. Men denne
frie Religion naaes ved selve Forstaaelsen af Christendommen,
naar man nemlig ser, at denne tillige er den
høieste Videnskab, der fører til den sidste historiske
Sammenhæng af Alt. Den philosophiske Religion eller
Aandens Religion er altsaa Et med Religionphilosophien, % PRINTED AS IS (Religionphilosophien: for Religionsphilosophien; the reading of the Minnesota image and its layer)
og adskiller sig fra den saakaldte Fornuftreligion
derved, at den har Mythologien og Aabenbaringen
til sine Forudsætninger og er deres reale Fuldendelse. ---

Schellings senere Lære har en afgjort historisk Karakter.
Alt er her Udvikling til et høieste Maal. Vel
var heller ikke i Identitetslæren Tilværelsen en død Væren;
Alt var Liv, nemlig det Absolutes uendelige Selverkjendelse,
dets evige Sammenslutten af sig selv som
Objekt med sig selv som Subjekt. Men det var en Bevægelse,
der var Et med absolut Hvile; idet Synspunktet
kun toges i det Absolute, forsvandt ethvert hemmende
Moment; der var ingen Skranker, ingen Strid. Efterhaanden
blev dog Schelling opmærksom paa et irrationelt
Element, der sprængte denne immanente Enhed og
aabnede Blikket for en Sondringens og Uroens Verden.
Og jo mere han fæstede sin Opmærksomhed paa dette
Element, der for ham blev Et med Villiens Væsen, des
mere traadte det Absolute i Skygge og betragtedes nu
som det, der laa forud for den konkrete Udvikling, den
forskjelsløse Enhed af de Modsætninger, der først ved
Udviklingens Ende kunde forsones til virkelig Identitet.

Men den afgjørende Modsætning mellem den ældre og
den senere Lære beroer dog paa, at Bevægelsen nu foregaaer
om to Centra. Den negative Philosophi fører til
et Punkt, hvor Transscendentsen sættes i Kraft, og den
positive Philosophi begynder saa med det transscendente
Princip for igjen at føre tilbage til det Immanente. Nødvendigheden
af det dobbelte Centrum søges i Modsætningen
mellem det Almene og det Enkelte. Den negative Philosophi
bevæger sig indenfor det Almene og naaer derfor
% --- p. 144 ---
\setcounter{footnote}{0}%
\opage{144}ikke Virkeligheden, thi kun det Enkelte existerer. Modsætningen
bliver her saa stærk, at kun den efter personlig
Forløsning higende Villie skal kunne gjøre Overgangen.
Alt det Virkelige skal Tænkningen kunne erkjende,
men ikke, at det Virkelige er (quid, ikke
quod). Existentsens Bestemmelse er ved at udskilles („udstødes“)
fra det Almene ganske vist bleven irrationel;
men med hvilken Ret tilskriver Schelling en saa uhyre
Abstraktion, som Sondringen mellem Existentsen og det
Existerende, Realitet? Afset fra det existerende Væsen
er Existentsen kun et Phantasibillede. Det absolut Enkelte
i denne Forstand kan kun skues, ikke tænkes, og
forsaavidt Schelling her taler om en fri Tænkning, bliver
denne kun en phantastisk Spekulation.

Træt af den storartede Bevægelse i Tankens Verden
søger Tænkeren dreven af en dunkel Trang ud over
denne. Den absolute Idealisme bliver nu kun en Indledning
til de høieste Mysterier. Saaledes begyndte Nyplatonikerne
netop, hvor Platon og Aristoteles endte.
Det høieste Princip, den sidste Aarsag til Alt naaes
ifølge Plotin kun, naar Sjælen vækkes af et Lys, der er
over Fornuften, til at søge ud over denne og vende sig
mod et Høiere. Hvad Nyplatonismen er i Forhold til
den klassisk-græske Philosophi, det er Ny-Schellingianismen
i Forhold til den spekulative Idealisme.

Ogsaa hos andre Tænkere, hvis Udvikling har ført
dem til et Brud med den spekulative Idealisme, finde
vi et lignende Forhold. Ludwig Feuerbach forkastede
Spekulationen som en Apotheose af Fornuften, af det
Almene i Modsætning til det Enkelte og Sandselige, og
korrigerer ligesom Schellings Værens Begreb saaledes:
„Væren er Et med det, som er.“ Spørgsmaalet om
Væren er for ham ikke længer et logisk Spørgsmaal,
men gjælder Liv og Død, er en Anskuelsens og Kjærlighedens
Hemmelighed. Naar han siger, at Existentsen
har Betydning og Fornuft i sig, skjøndt den ei kan
% --- p. 145 ---
\setcounter{footnote}{0}%
\opage{145}udsiges, da minder dette ganske om Plotins Ord om Forholdet
til det Høieste: ἔχειν οὐ κωλυόμεθα κἂν μὴ λέγωμεν.
I Skrifter, der omtrent ere samtidige med Schellings
Forelæsninger i Berlin og med Feuerbachs „Philosophie
der Zukunft“, distingverer Søren Kierkegaard paa tilsvarende
Maade mellem Viden og Virkelighed („Philosophiske
Smuler“, „Uvidenskabelig Efterskrift“): Hvad
Virkelighed er, kan ikke tilegnes i Viden: enhver Viden
om Virkelighed er Mulighed; Existents er bestandig det
Enkelte. Ogsaa her korrigeres Værens Begreb i Modsætning
til den spekulative Idealismes Forflygtigelse. Hos
Kierkegaard faaer Existentsens Bestemmelse Betydningen
af den ethisk betingede Virkelighed; det Enkelte, som
er, er det existerende Individ, der staaer under den
ethisk-religiøse Lov. Hos Feuerbach er Existentsen givet
for Sandsningen; det existerende Enkelte er den sandselige
Gjenstand. Hos Schelling endelig er Virkeligheden
den spekulative Skuens Gjenstand, der ene kan gribes,
naar man søger ud over Bevidsthedens sædvanlige Grændse;
det existerende Enkelte er her, ligesom hos Plotin, det
Phantasibillede, der fremtræder for Tanken, naar denne,
træt af at gjennemarbeide Kategorierne, synker hen i en
ekstatisk Anskuelse, der fremkaldes af den ikke mættede
„Hunger efter Existents“\footnote[1]{Endnu en Parallel ligger nær: mellem Schelling og Schopenhauer. Nylig har saaledes \emph{Edv. v. Hartmann} fremhævet, at Schellings senere Lære som Enhed af negativ og positiv Philosophi har den Fortjeneste at sammenfatte Hegels Princip (Ideen) og Schopenhauers (Villien) som lige berettigede og lige uundværlige Sider af det ene Princip. Philosophie des Unbewussten. 3te Aufl. p. 766.}.

Men at denne Modsætning ikke kan fastholdes, ses
hos Schelling deraf, at den positive Philosophi ikke lader
sig nøie med det blotte „at“ (qvod), men bygger en
hel transscendent Udvikling paa Basis af „das unvordenkliche Sein“. Modsætningen gjælder altsaa kun i selve
% --- p. 146 ---
\setcounter{footnote}{0}%
\opage{146}Begyndelsespunket. Og selv her føres vi ud over den; % PRINTED AS IS (Begyndelsespunket: for Begyndelsespunktet; the image, s02)
thi Tanken maa dog kjende sit eget Udgangspunkt, maa
berøre terminus a qvo for sin egen Bevægelse; ogsaa i
selve Udgangspunktet gaa altsaa Tanken og Væren sammen,
og Væren bliver ikke „unvordenklich“.

Dette kunde nu forsaavidt ogsaa synes at være Schellings
Mening, som han jo sætter den positive Philosophis
Opgave i at gjøre det Transscendente immanent. Dette
sker gjennem den mythologiske Proces, der er en Bevidsthedsproces,
om Bevidstheden her end er i en dunkel,
fatalistisk Magts Hænder, er ude af sig selv (ekstatisk).
Men ved Slutningen af denne Proces statuerer Schelling
atter det Transscendente, endog i mirakuløs Form: den
ekstatiske Historie træder, forat Forsoningen kan gjennemføres,
i Inkarnationen frem som virkelig Historie.
Men Muligheden og Nødvendigheden heraf paavises ikke.
Da det religiøse Forhold fra først af er bestemt som et
Bevidsthedsforhold , hvorfor kan saa Forsoningen ikke
gjennemføres gjennem en Bevidsthedsproces? Tilmed da
allerede den mythologiske Proces ifølge Schelling maa
vise Muligheden heraf, idet den jo Skridt for Skridt har
trængt det dunkle, adsplittende og fristende Princip tilbage.
Man kunde da tænke sig Udviklingen ført videre
ved den Overgang a potentia ad actum, som nu er Schellings
Yndlingsbestemmelse, tænke sig „den anden Potents“
fuldende sin Vorden og derved ogsaa gjøre den
tredie Potents's (Aandens) Virkelighed mulig, saa at Enheden
og Forsoningen naaes uden ny Transscendents.
Men saa maatte man rigtignok opgive at deducere den
kirkelige Treenigheds- og Inkarnationslære. Det er for
at gjøre Fyldest i saa Henseende, at Schelling her har
ophævet det rationelle Forhold mellem potentia og actus.

Den hele absolute Modsætning mellem „hvad“ og
„at“ viser, for Schellings, ligesom for Feuerbachs og
Kierkegaards Vedkommende, at den spekulative Idealismes
Forudsætninger endnu fastholdes. Ellers kunde
% --- p. 147 ---
\setcounter{footnote}{0}%
\opage{147}Schelling ikke mene, at Tanken ved at udvikle sine egne
Muligheder tillige udviklede Tingenes Væsen (qvid), at
altsaa Erfaring og Tænkning dækkede hinanden. Den
negative Philosophi staaer i sig selv ganske paa Spekulationens
Standpunkt, og navnlig beholdes Naturphilosophien
ganske fra den tidligere Lære. Ligesaa ukritisk,
som Schelling i denne forholdt sig til den virkelige Erfaring,
forholder han sig i den positive Philosophi til
Mythologi og Aabenbaring. Hans Lære hviler her paa
et meget usikkert historisk Grundlag\footnote[1]{Dette vil man let overbevise sig om ved et Blik paa p. 139 f. ovenfor. Særlig skal fremhæves, at Schelling opfatter Successionen af Uranos, Kronos og Zeus som en Kamp mellem mythologiske Forestillinger, der fortrænge hverandre, medens Forestillingen om Gudekampen dog kun er udsprungen af Forsøg paa at forklare de stridende Magter i Naturen eller at udmale sig Gudernes (de olympiske Guders) Oprindelse og Historie.}. Karakteristik % PRINTED AS IS (Karakteristik: the image, for Karakteristisk; recrop rc1)
er i denne Henseende den Ytring, at da hans Aabenbaringsphilosophi
beroer paa en Udvidelse af Philosophien
og den philosophiske Bevidsthed, der kun var
mulig ved at gaa tilbage til Principerne, men den historiske
Kritik ikke kjender denne Udvidelse af Bevidstheden,
saa kan han aldeles ikke tage noget Hensyn til
denne Kritik\footnote[2]{Werke. 2te Abth. 4ter Bd. p. 231.}. Han behandler med andre Ord Kritiken
paa samme Maade, som Theologien gjør det. Selve
den mythologiske Udvikling ser han som en indre, metaphysisk,
der ikke betinges eller berøres af physiske
eller psychologiske Indvirkninger.

Alligevel har Schelling den store Fortjeneste at have
hævdet Nødvendigheden af en positiv og sympathetisk
Forstaaelse af de historiske Religioner. Hans Religionsphilosophi
danner derved en Modsætning til den spekulativ-idealistiske.
Ogsaa Feuerbach danner en saadan
Modsætning; men han betonede saa ensidig det lavere,
pathologiske Element i Religionen, at det blev umuligt
% --- p. 148 ---
\setcounter{footnote}{0}%
\opage{148}at forklare, hvorledes den ogsaa kunde indeholde en virkelig
Sandhedskjerne. Schelling fastholder Fornuftsandheden
som Grundlag (Prius) og kan derved ogsaa hævde
det religiøse Forholds Sandhed.

Det Berettigede i Schellings Begreb om en positiv
Philosophi er i det Hele dette, at det Existerende staaer
selvstændig overfor den begribende Tanke og maa anerkjendes
af denne, ikke opløses og forflygtiges. Men
dette Positive kan ikke naaes ad anden Vei end den
exakte Erfaringsvidenskabs og den historiske Kritiks.
Og naar Erfaringen faaer selvstændig Betydning, vil
Tingenes Existents (quod) ogsaa vise sig at være en
nødvendig Forudsætning for, at Tanken kan begribe deres
Væsen (qvid). Schellings Betoning af Existentsens Begreb
er mere en prophetisk Bebudelse end en virkelig
Begyndelse. Han har skuet i Phantasien, hvad der er
Opgaven for en møisommelig Arbeiden i Virkeligheden.

Schellings positive Philosophi er ikke orthodox. Den
udelukker ikke blot Theologi, men ogsaa Troen som
positiv Tro kan umulig faa virkelig Betydning, da den
positive Philosophi jo netop skal føre ud over Mythologi
og Aabenbaring til den frie, philosophiske Religion,
der er Et med Religionsphilosophien. „Aabenbaringsphilosophien
er mere radikal, end man i Almindelighed
troer. Fra theologisk Side mente man strax at kunne
benytte Schellings nye System som Modvægt mod den
negative Kritik, og han var, som vi have set, ikke fri
for at kokettere med Orthodoxien. Men den gamle Tænker
har ikke fornegtet sin Fortid i den Grad, som det
fra de mest modsatte Sider, den negative Kritiks saavelsom
Orthodoxiens, blev antaget. Theologien har derfor
her, som saa ofte ovenfor philosophiske Systemer,
indskrænket sig til at acceptere enkelte Præmisser eller
Slagord, medens den forsigtig holder sig tilbage fra Konsekventserne.
Selve disse Konsekventser udtales iøvrig
kun i korte Antydninger; men allerede ved disse fast%
% --- p. 149 ---
\setcounter{footnote}{0}%
\opage{149}slaaes der et Princip, som afgiver det bedste Korrektiv
mod den store Tænkers egne ukritiske Lærdomme.

2. \emph{Den yngre Fichte}. --- Ogsaa dette Navn
minder om en af de største Skikkelser i den nyere Philosophis Historie. Denne Erindring ligger saameget
nærmere, som Sønnen selv erklærer i Faderens Tænkning
og Personlighed at have fundet de første Impulser
til sin egen Forsken. Denne fik derved et Præg, som
den selv under de forskjellige Phaser, den gjennemgik,
aldrig mistede. --- Johann Gottlieb Fichte (1762—1814)
greb og udviklede, som allerede tidligere omtalt (p. 4),
med rastløs Energi et centralt Punkt i den kantiske
Philosophi og førte derved denne ud over dens oprindelige
Begrændsning og Uklarhed. Han blev derved ført
til en subjektiv Idealisme, i hvilken han tillige fandt et
Grundlag for sin praktiske Stræben. Heroisk Arbeiden
for Arbeidets egen Skyld og med asketisk Foragt for
enhver Hvilen i det Givne og Opnaaede var Fichtes
ethiske Princip. Men senere søgte han ud over denne
subjektive og praktiske Idealisme, denne evige og aldrig
beroligede Vorden og Virken at gribe en evig Væren, ud
over den strenge Pligtnødvendighed at finde det religiøse,
det salige Liv, hvor det Absolute ikke længere
staaer som et fjernt, uopnaaeligt Maal, men er iboende i
Mennesket. Fra sit rent subjektive Udgangspunkt mente
han at være naaet til dette Høidepunkt, idet han ikke
længere betragtede Bevidstheden som absolut, men som
phænomenal Aabenbaring af det Absolute.

Det var fra denne senere Form af Faderens Lære,
at Sønnen, Immanuel Hermann Fichte (f. 1797) tog sit
Udgangspunkt. I Studiet af Faderens efterladte Skrifter
fandt han det theoretiske Udtryk for, hvad der i hans
Barndom var traadt ham imøde praktisk i begge Forældrenes
personlige Tænkemaade og Liv\footnote[1]{Smlgn. „J. G. Fichtes Leben und literarischer Briefwechsel“ \opage{150}2te Aufl. 1862. Von I. H. Fichte, --- en Bog, som Feuerbach (i „Pierre Bayle“) med Rette satte meget høit.}. Thi Mode%
% --- p. 150 ---
\setcounter{footnote}{0}%
\opage{150}ren var paa den inderligste Maade gaaet ind i sin Ægtefælles
hele Aand, og det blev paa Grund af hans tidlige
Død hende, gjennem hvem denne Aand maatte virke paa
Sønnen. „Hendes Indflydelse,“ siger denne\footnote[1]{Bericht über meine philosophische Selbstbildung. Verm. Schr. 1869. 1ster Bd. p. 32.}, „skylder
jeg Alt, hvad der har bevaret sig hos mig mit hele Liv
igjennem af høiere Vækkelse, af urokkelig Forvisning,
om end ikke philosophisk udtydet, endnu mindre philosophisk
begrebet. Men Driften, den uslukkelige Trang
var derved nedlagt i mig til at se hint Ethisk-Religiøse
retfærdiggjort og bragt til Forstaaelse for mig ogsaa ad
Begrebets Vei.“ Denne Udtalelse bidrager til Forklaringen
saavel af de stærke som af de svage Sider hos
den yngre Fichte. I Kraft af det Motiv, han her har
angivet for sin philosophiske Forsken (det var det, der
fra philologiske Studier førte ham over til Philosophien),
kalder han sig en Personlighedsphilosoph; Overbevisningen
om Personlighedens og det personlige Livs
Realitet er det Gjennemgaaende i hele hans lange videnskabelige
Virksomhed. Men denne personlige Interesse,
som fra først af er tilstede og fordrer sin Fyldestgjørelse,
hindrer da ogsaa paa mange Maader det frie, uhildede
Blik, den rene, konsekvente Udvikling. Saamegen
Fortjeneste den yngre Fichte end har indlagt sig ved
sin Opposition mod den ensidige Pantheisme og Idealisme
og senere mod Materialismen, saa svækkes dog
paa den anden Side Virkningen heraf ved det Tendentiøses,
det bevidste Øiemeds umiddelbare Fremtræden,
og ogsaa i selve Resultaterne fremkommer derved meget
Uklart og ligefrem Phantastisk.

Som Trangen til at klare den ethisk-religiøse Overbevisning
var det positive Motiv for I. H. Fichtes Phi%
% --- p. 151 ---
\setcounter{footnote}{0}%
\opage{151}losopheren, saaledes har han fremdeles erklæret, at den
Modsætning, hvori denne Overbevisning maatte komme
til det herskende System, det hegelske, har været en
„bestandig fremaddrivende Braad for ham. Den absolute
Idealisme, saaledes som Schelling og Hegel havde udviklet
den, maatte i en dobbelt Henseende virke frastødende
paa ham. For det Første ved sit pantheistiske
Resultat, ifølge hvilket alt Enkelt og Individuelt var et
Usandt og et Forsvindende; --- Personlighedens Realitet
var jo for ham fra først af en Grundsandhed. Men dernæst
ogsaa ved Postulatet om en absolut Erkjendelse.
For Fichte var det en nødvendig Fordring, at Erkjendelsens
Gyldighed først maatte godtgjøres, før Tanken
kunde hæve sig til det Absolute. Det subjektive Standpunkt,
som Schelling og Hegel mente at have overvundet,
havde endnu sin Berettigelse for ham. --- „Jeg maa
bekjende,“ siger han (anf. Skr. p. 53), „at den uvilkaarlige
Protest derimod [mod Resultaterne af Hegels System],
men tillige Ønsket om ad spekulativ Vei at gjøre
Fyldest for de derved opstaaede nye Fordringer, siden
hin Tid ere blevne den egentlig tilskyndende Spore til
min Philosopheren.“

Studiet af Philosophiens Historie førte ham\footnote[1]{„Beiträge zur Charakteristik der neuern Philosophie“ (1829. 2te Aufl. 1841). „Ueber Gegensatz, Wendepunkt und Ziel heutiger Philosophie“ (1832).} til at
skjelne mellem en dobbelt Række af Anskuelser, der
stode i en afgjørende Modsætning til hinanden. Den
ene er den, der som Schelling og Hegel strax tager et
objektivt Udgangspunkt og herfra udvikler en Lære om
Væren som objektiv Substants. Denne Retning hviler,
saavel i sine spekulative som i sine mystiske Former,
paa Enheden af Tænken og Væren og fører i sidste
Konsekvents til Pantheisme. Den anden Retning er
den subjektivt-reflekterende, som navnlig repræsenteres
% --- p. 152 ---
\setcounter{footnote}{0}%
\opage{152}af Kant; dennes Ensidighed bestaaer i, at den ikke kan
hæve sig over den subjektive Bevidstheds Skranker, ikke
kan naa en objektiv Erkjendelse. Den sande Forening
og Mægling mellem begge Retninger finder han i, at det
philosophiske System begynder med Erkjendelseslæren
og først gjennem denne hæver sig til det objektive Standpunkt.
I hans methaphysiske Hovedværk „Grundzüge % PRINTED AS IS (methaphysiske: the image, O099)
zum Systeme der Philosophie“ er derfor første Del Erkjendelseslære:
„das Erkennen als Selbsterkennen“ (1833),
anden Del Ontologi (1836), „Videnskaben om Virkelighedens
Former og, da kun det Absolute er det Virkelige,
om det Absolutes Virkelighedsformer“; tredie Del:
„die spekulative Theologie oder allgemeine Religionslehre“
(1846) danner Afslutningen, idet der fra Verdensbegrebet,
saaledes som Ontologien udviklede det ved at
udvikle Virkelighedsformerne, søges tilbage til Gudsideen,
den sidste Forklaring og Begrundelse af den givne Verdenskjendsgjerning.


Bevidstheden maa, saaledes indledes Erkjendelseslæren,
begynde med sig selv, sit umiddelbart givne Væsen.
Først Skridt for Skridt kan den udvikle og opdrage
sig selv til Sandhed. Bevidstheden er selv Begyndelsen,
fordi den er det umiddelbart Visse, og Erkjendelseslæren
er Bevidsthedens Udviklingshistorie til
Tænkning og under Tænkning. Fra sit første Standpunkt,
hvor den endnu er fortabt i Anskuelsen, i umiddelbar
Enhed med det Objektive, arbeider Bevidstheden
sig frem til Forestilling, Tanke og Erkjendelse. Naar
den nu paa sit høieste Trin besinder sig paa sit eget
Væsen og med gjennemtrængende Reflexion og opløsende
Skepsis har tilintetgjort ethvert umiddelbart Indhold,
bliver der kun selve dens eget Væsen tilbage som det
eneste Visse. Den hele Objektivitet er forsvunden i Subjektivitet.
Viden er derfor fra den ydre Realitets brede
Mangfoldighed svunden sammen til en snever Kreds, til
den formelle Viden af sig selv. Dette er et nødvendigt
% --- p. 153 ---
\setcounter{footnote}{0}%
\opage{153}Gjennemgangspunkt til den frie Erkjendelse af Sandheden.
At det kun kan være et Gjennemgangspunkt,
--- at Bevidstheden ikke kan holde sig paa dette Punkt,
ligger i, at den ved at gjøre Alt til Intet ogsaa vilde ophæve
sig selv, idet den jo vilde blive Bevidsthed om
Intet. Bevidstheden kan kun have Realitet, forsaavidt
den har et absolut Indhold som sin af Reflexionen uopløselige
Basis. Dette Indhold er \emph{absolut}, fordi det
ikke, som det endelige Objektive, sættes af Bevidstheden,
men selv er dennes Rod og Grundlag; det er en
absolut Væren, fordi det ikke er Bevidstheden selv,
men er i Bevidstheden. Det Absolute er al Bevidstheds
evige Grundbetingelse. Bevidstheden er altsaa kun et
Relativt, Afbillede og Aabenbaring af et ubetinget Værende,
der maa være, saa vist som der i det Hele er
Bevidsthed til. Heraf følger tillige, at den spekulative
Tanke ikke staaer i Modsætning til Anskuelsen; som
Bevidstheden i sin sidste Grund er en Fornemmen af
det Absolute, saaledes maa Tanken, naar den skal fatte
alle dettes Aabenbaringsformer bestandig optage Anskuelseselementer
i sig. Den høieste Erkjendelsesform
er derfor den spekulativt-anskuende Erkjendelse og al
Spekulation i sidste Instants Theosophi.

Bestemmelsen af det Absolutes Væsen og Begreb
maa Erkjendelseslæren overlade til Ontologien. Erkjendelseslæren
kan kun godtgjøre, at det Absolutes Begreb
har Realitet; Ontologien har den Opgave at udvikle de
evige Former, i hvilke alt konkret Virkeligt fremtræder.
Men den Tænkning, ved hvilken denne Udvikling gjennemføres,
er ikke længer som hos Hegel, den abstrakte
og indholdsløse; den har tvertimod gjennem Erkjendelseslæren
vundet Visheden om det Absolute og søger nu
kun nærmere at udvikle, hvad der ligger i dets Væsen.
Fra de mere abstrakte og ensidige Former søger den ad
den immanente Dialektiks Vei til stedse mere konkrete
Bestemmelser. Paa sit Høidepunkt, naar den har ud%
% --- p. 154 ---
\setcounter{footnote}{0}%
\opage{154}tømt Kategorierne, de evige Verdensformers System,
bliver Ontologien Idelære eller spekulativ Theologi, idet
den ender i et Realt, som ikke længer er behæftet med
Modsigelsens Braad, men indeholder Virkelighedens evige
Forbillede, til hvilket alle hine Verdensformer forholde sig
tjenende. I Idelæren giver den negative Dialektik Plads
for den positive; thi det er nu reale Forhold, det gjælder
om. Bevægelsen indenfor Ideverdenen er et Forbillede for
den reale Verdensproces. Naar derfor Idelæren fører til
det Resultat, at Aandens, Personlighedens Ide er den
høieste Virkelighed, saa er dette en Sandhed, som ogsaa
den reale Verdensudvikling faktisk godtgjør, hvilket det
dog ikke længere er Methaphysikens men Realphilosophiens % PRINTED AS IS (Methaphysikens: so twice in the book (also p. 214); the image, s03)
(Naturphilosophiens og Historiens Philosophis)
Sag at godtgjøre. Den hegelske Philosophi maatte føre
til Pantheisme, netop fordi den ikke gjorde denne Adskillelse
mellem Formal- og Realphilosophi. Hos Hegel
er Kategoriernes dialektiske Udvikling det Absolutes egen
Selvudvikling; derfor er det Absolute først fuldendt,
først den absolute Aand ved Udviklingens Slutning.
Men det maa nøie fastholdes, at den ontologiske Begrebsudvikling
sondrer, hvad der i Gud evig er Et, hvad han
er i uforanderlig Lighed med sig selv, ophøiet over al
Vorden og Udvikling. Forskjellig herfra er den Maade,
hvorpaa Guds Væsen lægger sig for Dagen indenfor den
virkelige Verdensudvikling, af hvilken hin evige Væren
i sin absolute Enhed er uberørt. („Gegensatz von Ewigkeit
und zeiterfüllender Genesis“). Som vi kunne erkjende
Rummets Love afset fra det Rumopfyldende, saaledes
kunne vi ogsaa gjennem Udviklingen af de aprioriske
Former erkjende Guds Væsen. Men det fører
ikke, som Pantheismen mener, til en adækvat og absolut
Begriben af det Guddommelige, end sige, at vor Tænken
skulde være identisk med den guddommelige Selverkjendelse.
Thi denne Formerkjendelse naaer endnu
ikke Gud i sig selv, det guddommelige Væsens Fylde
% --- p. 155 ---
\setcounter{footnote}{0}%
\opage{155}og uendelige Frihed. Medens vor ontologiske Tænkning
er rent apriorisk og dialektisk, er den guddommelige
Tanke, der aabenbarer sig i Verdensvirkeligheden, absolut
positiv og konkret. Modsætningen mellem det theocentriske
og det anthropocentriske Standpunkt maa fastholdes;
det første er os absolut utilgængeligt. ---

Det Karakteristiske ved dette Forsøg paa at komme
ud over den absolute Idealisme er, at man ved at give
den objektive Ideerkjendelse en subjektiv Begrundelse
gjennem Reflexion over Bevidsthedens Væsen modificerer
dens Væsen og Resultater og gjør den fra absolut til
relativ. Spørgsmaalet er da, om Ideerkjendelsen ikke
taber sin spekulative og objektive Betydning ved at støttes
til en reflekterende Erkjendelseslære. Mod dette
Punkt vendte Kritiken sig ogsaa, ikke blot fra Hegelianernes
Leir, som i Fichtes Forsøg saa en Tilbagevenden
til Kants subjektive Standpunkt, men ogsaa fra en Mand,
som ellers i sin hele Retning syntes at staa Fichte nær
og i Almindelighed stilledes sammen med ham. C. H.
\emph{Weisse} paaviste i sit Stridsskrift mod Fichte: „Das
philosophische Problem der Gegenwart“ (1842), at det ved
det Udgangspunkt, denne havde taget i Faderens subjektive
Idealisme, blev ham umuligt at vurdere end sige
at optage i sin Anskuelse, hvad Schelling og Hegel havde
lært om Ideens Væsen. Det Absolutes Ide, saaledes
som disse Tænkere opfattede og udviklede den, er den
i en intellektuel Anskuelse eller dialektisk Begrebsudvikling
godtgjorte Forvisning, at ikke det blotte Subjekt
er det i Sandhed Værende og Virkelige, men at
dertil udfordres Subjektets Overgang i det Objektive og
Tilbagevenden til sig selv igjen, kort sagt det absolute
Subjekts uendelige Potensation ved at medieres
med sin Objektivitet. Men denne Ide kan saa langt fra
vindes gjennem en reflekterende Erkjendelseslære, at
denne tvertimod stedse maa forudsætte hin. Erkjendelseslæren
kan ikke godtgjøre Ideens Realitet, da den jo,
% --- p. 156 ---
\setcounter{footnote}{0}%
\opage{156}som Fichte ogsaa indrømmer, selv hviler paa denne. Man
kommer i hvert Tilfælde til at bevæge sig i en Kreds:
Erkjendelsens Realitet skal bevise det Absolutes Realitet,
hvorpaa den selv beroer. Af denne Kreds kommer
man kun ud ved at skjelne mellem den Forskole, Individet
maa gjennemgaa for at komme til Sandheden, og
Sandheden selv, og saa ved Udviklingen af denne at gaa
ud fra det, som i sig selv er det Oprindelige.

Dog var det ikke denne Betragtning, der senere
bevægede Fichte til at søge en anden Begrundelse af
det Absolute og en anden Overgang fra Erkjendelseslæren
til de øvrige Dele af Systemet. Han fastholdt
tvertimod Nødvendigheden af, at Erkjendelseslæren gik
forud for Metaphysiken (Idelæren, Læren om det Absolute).
Men medens Overgangen tidligere var immanent,
idet Bevidsthedens eget Væsen førte tilbage til det Absolute
som sin Grund, forkastede han senere denne rene
Immanents som pantheistisk. Gud blev ved den kun
set som det iboende Subjekt-Objekt, i hvilket det endelige
erkjendende Subjekt dialektisk ophævedes gjennem
den uendelige Erkjendelsesakt. Derimod ophæves det
pantheistiske Princip, naar man betragter Erkjendelseslærens
Resultat som et nyt Problem, ved hvis Løsning
vi forlade den negative, immanente Udvikling. Dette
Problem bliver da: hvad er Grunden til den Identitet
af Subjekt og Objekt, hvorpaa al Erkjendelse beroer?

I sin „spekulative Theologi“, i hvilken han paa
denne Maade korrigerer sin tidligere Lære, udtaler han,
at denne kun kunde føre til at opfatte Gud som den
endelige Verdens aandelige Enhed. Det er navnlig naturvidenskabelige
og anthropologiske Studier, der have ført
ham til Overbevisningen om det Individuelles Realitet,
at det ikke er et blot Endeligt, der ophæves i det Absolutes
uendelige Proces. Hermed følger da ogsaa en
ny Methode. Var denne tidligere dialektisk-ontologisk,
idet den ad den immanente Spekulations Vei søgte at
% --- p. 157 ---
\setcounter{footnote}{0}%
\opage{157}udvikle Gudsideen, saa lægges der nu en Tankegang til
Grund, der svarer til det gamle kosmologiske Bevis,
hvor der fra Tilværelsen som given Kjendsgjerning sluttes
tilbage til Gud som dens Grund og Aarsag.

Den Enhed af Subjekt og Objekt, som Erkjendelsen
forudsætter, bliver nu kun en enkelt Kjendsgjerning,
som for at ses i sit rette Lys maa indordnes i den hele
Verdenskjendsgjerning. Denne viser os nemlig som det
virkelig Værende ikke, som Hegel mente, en Række af
forsvindende Momenter, men en Uendelighed af reale
Positioner, af Urkvaliteter, som udgjøre det Oprindelige
og Blivende i al Vexel, det Substantielle i det Endelige.
Pantheismen kjender kun to Led, det Absolute og
det vordende (phænomenale) Endelige; men skal den
reale Verdensudvikling i sin Mangfoldighed ikke nedsættes
til et blot Skin, saa bliver et tredie Led, et
Mellemled nødvendigt, Begrebet om det substantielle
Endelige, den reale Urposition. Herved gjendrives Pantheismen
ligesom fra neden af, ud fra den umiddelbare
Verdenskjendsgjerning. Men idet Forklaringen af det
Givne saaledes fører til Antagelsen af en uendelig, real
Mangfoldighed --- en Antagelse, som indeholder det Berettigede
i Herbarts Lære og i hans Opposition mod den
spekulative Idealisme --- saa viser det sig dog snart, at
denne Mangfoldighed kun bestaaer og kun kan bestaa
i og formedelst en uendelig, dem alle gjennemtrængende
Enhed. Hver Position er ikke Noget for sig selv, men
staaer i indre Vexelforhold til alle andre. Paa samme
Maade derfor, som Erkjendelsens Kjendsgjerning fører
til Antagelsen af en Enhed, der omfatter og betinger
saavel Subjektet som Tingene, viser det sig nu, at alt
Existerende i sin Mangfoldighed og individuelle Eiendommelighed
maa bæres og grundes ved en ubetinget
Enhed, der indeholder den sidste Grund til den hele
Verdenssammenhæng saavel som til enhver relativ Enhed
% --- p. 158 ---
\setcounter{footnote}{0}%
\opage{158}og Helhed indenfor denne. Verdensenheden kan kun
være et Produkt af den skaberiske Ur-Enhed.

Det næste Spørgsmaal bliver da, hvorledes en saadan
absolut Enhed af det Uendelige er mulig. Ved at antage
den absolute Enhed lægger man jo egentlig kun Vanskeligheden
et Skridt tilbage, fra det Betingede til det Ubetingede.
Ogsaa ved dette Spørgsmaal fremhæver Fichte,
at dets Besvarelse maa søges paa Grundlag af og i
Analogi med det oprindelig Givne. Han opstiller i det
Hele som sit ledende Princip, at kun af Verdensvirkeligheden
kan Gudsideen erkjendes. --- Den faktiske Sammenhæng
mellem Enkeltvæsenerne viser sig nærmere at
forudsætte en ideel Forbindelse, i hvilken det enes Begreb
betinges ved det andets Begreb og kræver det som
sin nødvendige Forudsætning: for at det ene skal være,
maa ogsaa det andet være. Hvor adskilte Enkeltvæsenerne
end ere i deres ydre Virkelighed, ideelt set ere
de dog nærværende i og for hverandre. Denne ideelle
Sammenhæng er en teleologisk Sammenhæng. Forholdet
mellem Øiemed og Middel er jo netop dette, at hint
forudsætter dette som nødvendig Betingelse, medens
omvendt dette først ved hint finder sin fuldstændige
Forklaring og Begrundelse. Denne ideelle, teleologiske
Sammenhæng forudsætter da igjen en overskuende og
sammenfattende Tanke, der tillige i sit reale Væsen faktisk
sætter og omslutter det existerende Mangfoldige.
Man har tænkt sig det Absolute som en Verdenssjæl,
der virker og former i Verden med allestedsnærværende,
men bevidstløs Fornuft. Den ideelle Sammenfatten er
her en skuende Imagination, og den skabende Virksomhed
bestaaer i, at denne imaginative Virksomhed stiger
til udtrykkelig, individuelt producerende Anskuen. Saaledes
i Schellings Naturphilosophi, hvor alt Existerende
er særegne Former for det Absolutes Selvanskuen. Et
høiere Trin naaes, naar man som Hegel indser, at det
Absolute ikke kan være fortabt i den objektive Anskuen,
% --- p. 159 ---
\setcounter{footnote}{0}%
\opage{159}men maa være fri som almen Fornuft, absolut Aand. Dog
er det endnu bestandig kun Begrebet om en Verdensaand, derved naaes, og Afslutningen bliver kun relativ, % PRINTED AS IS (derved naaes: the image, der lost before derved; recrop rc3)
stiller et nyt Problem. Thi Begrebet om Verdensaanden
er kun et koncentreret Udtryk for Verdensproblemet,
for den Kjendsgjerning, at i Universet er den høieste
Enhed af Fornuft og Virkelighed realiseret; men det
Spørgsmaal forbliver ubesvaret, hvorledes dette er muligt,
hvilken den sidste Grund til denne Kjendsgjerning
er. Heri ligger Nødvendigheden af at gaa ud over den
hegelske Philosophi\footnote[1]{Jvnfor ovenfor p. 21 f. Fichtes Karakteristik af den hegelske Philosophi i et tidligere Skrift.}.

Dertil kommer endnu en anden Vanskelighed. Den
absolute Idealisme kan ikke udvikle Aandens Begreb
saaledes, at den kommer til at staa frit overfor det Ubevidste.
Det ubevidste Liv, Naturen i snevrere Forstand
forbliver derfor uforklaret. Derimod fører Erkjendelsen
af, at alle umiddelbare Tids- og Rumsforskjelle, al umiddelbar
Existeren i det Hele er indoptaget i den teleologiske
Sammenhæng, til den Overbevisning, at denne hele
Sammenhæng maa besiddes af det Absolute som Verdens-Urbillede
i evig Klarhed. Den indre Hensigtsmæssighed
fører saaledes med Nødvendighed ud over sig
selv. Verdens-Urbilledet, den evige Tanke-Kosmos, Albevidstheden
fører atter tilbage til en Selvbevidsthedens
Enhed som sin Grund. Og saaledes staa vi da ved det
sidste Led af det lange teleologiske Bevis, som i sin
Gang Skridt for Skridt har støttet sig til Verdenskjendsgjerningen.
Kun en guddommelig Villiesmagt kan indeholde
Aarsagen til Verden, da denne ikke er Et med
det rene Begreb. Og en saadan Magt besidder Gud som
den selvbevidste Enhed, der evig sætter og forener de
reale Urpositioner.

Baade i det høieste Princips Begreb og i Begrebet
% --- p. 160 ---
\setcounter{footnote}{0}%
\opage{160}om det Endelige er der saaledes paavist Korrektiver mod
den abstrakte Spekulations pantheistiske Resultater. Og
begge Korrektiver staa i indre Forbindelse. Det Absolutes
koncentrerede Enhed som Selvbevidsthed vilde ikke
være mulig, naar der ikke indeholdtes en virkelig Realitet
deri, og paa den anden Side vilde det Endelige
kun være et usandt Skin, naar ikke dets Realgrund indeholdtes
i det guddommelige Væsen. Urpositionerne, som
her indtræde som den nødvendige Mellembestemmelse
mellem det Uendelige og det Endelige, ere derfor evige,
„urskabte ved Guds evige Villen af sig selv“. --- Spørgsmaalet
er nu, hvorledes Urpositionerne nærmere forholde
sig paa den ene Side til Gud, af hvis Væsen de udspringe,
paa den anden Side til den endelige Verden,
hvis Substants de udgjøre.

Hvad den første Del af dette Spørgsmaal angaaer,
da ligger der heri en Vanskelighed, som allerede Weisse
i sit ovenfor nævnte Stridsskrift med stor Skarpsindighed
har paavist. Idet nemlig Urpositionerne ved
en Slutning fra Virkning til Aarsag (via causalitatis)
føres tilbage fra den virkelige Tilværelse, hvis individuelle
Grundlag de udgjøre, til Guds Væsen, gaaer denne
Bevægelse egentlig ud paa det Samme, som Fichte selv
lægger Schelling og Hegel til Last, nemlig at søge en
Forklaring af det Givne ved at lægge Vanskeligheden et
Skridt længere tilbage. Men derved skærpes den snarere
end mildnes. „Hvorledes ,“ spørger Weisse med
Rette, „kan det forenes med Guds Begreb som Aand,
som fri, absolut Aand, at tænke nogensomhelst særskilt
endelig Bestemmelse eller, hvad der vel kommer ud paa
Samme, en Uendelighed af saadanne Bestemmelser sat
i ham paa evig og uforanderlig Vis ?“ Fichtes Besvarelse
af dette Spørgsmaal er uklar og vaklende. Medens det
nogle Steder synes, som om han opfatter Urpositionerne
som evige i deres absolute, monadiske Enkelthed („urskabte“
altsaa = uskabte) --- en Forestilling, der vilde
% --- p. 161 ---
\setcounter{footnote}{0}%
\opage{161}føre til en eventyrlig Præexistentslære, --- saa hedder
det paa andre Steder, at „Urpositionerne existere i den
evige Natur vel paa urbilledlig, men dog ikke paa udtrykkelig
individuel Maade“, at de ikke i sidste Instants
kunne forblive i deres atomistiske Enkelthed, men maa
gaa tilbage i det guddommelige Væsens Enhed som dettes
evige Grundkræfter. Dog erklærer han det Spørgsmaal
for uløseligt, „paa hvilken udtrykkelig Maade Gud
forener disse for os substantielle Væsener med sig, om
han sætter dem som perennerende eller tager dem som
flydende Momenter tilbage i sin Enhed.“ Ganske vist
kan dette Problem presses saaledes, at det ikke blot bliver
uløseligt, men ogsaa meningsløst. Dog kan det ikke
negtes, at Fichte, idet han fra Herbart optog Urpositionernes
Begreb, ved hvilket det efter hans egen Erklæring
(Ontologie p. 194) først blev ham muligt at gjøre
det afgjørende Skridt ud over Hegel, --- derved har bragt
noget Uensartet ind i sin Anskuelse, som kun vanskelig
kan forenes med dennes øvrige, fra den spekulative Idealisme
stammende Elementer. Derfor har Erdmann ogsaa
betegnet hans Lære som Eklekticisme. Eklekticisme
bliver ethvert Standpunkt, der ikke formaaer at lade de
nye og fremmede Elementer, det optager, omgaa en
indre Omdannelse ved Tilegnelsen.

Urpositionernes System udgjør for Fichte den indre
Natur i Gud. Ved Hjælp af dette Begreb mener han
ikke blot at komme ud over den pantheistiske, men ogsaa
over den vulgære theistiske Anskuelse. Pantheismens
Berettigelse ligger i den Sandhed, at uden Verden, uden
Universum vilde Gud ikke være Gud; med Guds Virkelighed
er der sat et Universum, hvori Alt er indesluttet
og Intet kan opstaa eller forgaa. Men paa den anden
Side gjælder det: den \emph{umiddelbare} Virkelighed er
ikke den guddommelige; Gud vilde være Gud ogsaa uden
denne givne Verden. --- Den sædvanlige Theisme lærer
i Modsætning til Pantheismen en mirakuløs Skabelse af
% --- p. 162 ---
\setcounter{footnote}{0}%
\opage{162}Intet, en virkelig Frembringelse af Noget, som før ikke
var men nu bliver til som Almagtens Produkt. Overfor
denne Anskuelse holder Fichte sig til sin Faders bekjendte
Ord, at den er al falsk Metaphysiks Grundvildfarelse,
og at Ingen har sagt et forstaaeligt Ord om den.
I Sammenligning med en saadan Anskuelse vilde Pantheismen
nødvendigvis have Fortrinnet; den staaer Sandheden
et Trin nærmere, da den ikke krænker Gudsideen
ved at tillægge Gud en absolut Vilkaarlighed. Den sande
Theisme er den, der har optaget Pantheismen i sig,
forsaavidt denne er berettiget, nemlig forsaavidt den
hævder Væsenssammenhængen mellem Gud og Verden.
Den kan i Modsætning til den vulgære, dualistiske Anskuelse
betegnes som den konkrete Theisme eller, med
et fra Krause laant Udtryk, som Panentheisme\footnote[1]{Smlgn. „Vermischte Schriften“ I. p. 275: „Ved Theisme forstaa vi den ganske almindelige Tanke, at det absolute Verdensprincip i intet Tilfælde kan tænkes som blind, bevidstløs Magt, som almindelig Substants eller abstrakt, upersonlig Fornuft, men kun som i og for sig værende Væsen, for hvis Grundegenskab den menneskelige Tænkning ikke raader over anden Analogi eller andet Udtryk end den absolute Selvbevidsthed.“}.

Dette fører os til den anden Del af det ovenfor opkastede
Spørgsmaal, Forholdet mellem Idealverdenen
i Gud og den endelige Tilværelse, eller med andre Ord
til det egentlige Skabelsesproblem. At Fichte negter
en Skabelse i den vulgære Betydning, som her tillige er
den kirkelig-dogmatiske, have vi allerede set. Skabelsesproblemet
maa for ham antage følgende Form: hvorfor
forblive Urpositionerne ikke i deres evige, uomskiftelige
Ro, men gaa ind i den endelige Tilværelses Begrændsning
og stedse omskiftende Vorden? --- Den guddommelige
Villie er ikke absolut producerende; den er kun tilladende.
Dette Begreb om den tilladende Villie bliver
i Skabelseslæren Mellembestemmelsen mellem den pan%
% --- p. 163 ---
\setcounter{footnote}{0}%
\opage{163}theistiske og den dualistiske Opfattelse, ligesom Urpositionernes
Begreb var det ved Opfattelsen af det almindelige
Forhold mellem Gud og Verden. --- Istedetfor den
indre, forbilledlige Maade, hvorpaa Urpositionerne ere i
den evige Idealverden, udlades de nu til isoleret, hverandre
udelukkende Existents, sættes under Rummets,
Tidens og Legemlighedens Form. Skabelsesproblemet
er ikke et særskilt Problem; det er ikke mere ubegribeligt
end Problemet om et hvilketsomhelst enkelt Naturphænomens
Opstaaen. Men den tilladende Villie er kun
den ene Side af Sagen. Urpositionerne stræbe nemlig
selv efter Selvstændighed; der er i enhver af dem en
prometheisk Spire, i hvis Vækkelse Skabelsen netop bestaaer,
en Spire, som gjør ethvert endeligt Væsen til
et eiendommeligt, og som driver Enkeltvæsenet til at
frembringe sig selv, idet det frembringes ved den guddommelige Villie. Denne Villies Almagt er altsaa ikke
længere transscendent; men den er derfor ikke mindre
ubetinget. ---

Denne klare og rationelle Udvikling af Forholdet
mellem Gud og Verden kommer dog, afset fra den tidligere
berørte Vanskelighed ved Urpositionernes Begreb,
til at lide af en Uklarhed, som stammer fra den Maade,
hvorpaa Guds Personlighed opfattes. Fichte hævder med
Føie, at der i Begrebet om en absolut Selvbevidsthed
ikke indeholdes nogen Modsigelse, som ikke ogsaa vilde
ligge i Begrebet om en blot Substantsenhed eller om en
upersonlig Verdensaand. Problemet forbliver det samme,
nemlig om Forholdet mellem Uendelighed og Enhed,
men det løses fuldstændigst og klarest ved den absolute
Selvbevidstheds Begreb. Det er ingen utilstedelig Anthropomorphisme
at overføre Selvbevidsthedens, Personlighedens
Bestemmelse paa Gud; tvertimod, kun Gud er et
rent Jeg, en fuldkommen Enhed med sig selv, medens
i den endelige Aand Jeget kun er Reflexionens og Abstraktionens
Værk og Enheden derfor formel og  tom. Hvad
% --- p. 164 ---
\setcounter{footnote}{0}%
\opage{164}der trænger til Forklaring er derfor snarere, hvorledes
Jegets eller Selvbevidsthedens Begreb ogsaa finder sine
endelige Repræsentanter. --- Men Fichte finder i Personlighedens
Princip, saaledes som han anvender det paa
Gud, dog endnu „et specifisk Mere“ end den blotte Nødvendighed;
Begrebet om den guddommelige Villie udtømmes
efter hans Opfattelse ikke ved, at den ses som den
høieste Enhed af Frihed og Nødvendighed --- det vilde,
som han udtrykker sig, kun give en Virken, ikke en
Handlen. Den guddommelige Handlen skal derimod
positivt udelukke det Nødvendige, det „Naturagtige“.
Gud kunde derfor have undladt at skabe Verden, og
han kan i ethvert Øieblik vende tilbage til sin indre, afsluttede
Selvtilfredsstillelse, uden Verden; naar Verden
alligevel bestaaer, kan det kun være ved hans Villie.
Fichte sætter her ikke blot sin rationelle Skabelseslære
i Fare, men træder ogsaa i Modsætning til den nyere
spekulative Philosophis hele Standpunkt, som han dog
ikke vil fornegte, men tvertimod føre videre. Skal
Alt i sidste Instants bero paa et Valg, til hvis Udfald
Tanken ikke kan finde nogen nødvendig Aarsag, hvad
er det da Andet end en Tilbagevenden til den gamle
creatio ex nihilo? Et ubetinget og ubegrundet Valg var
netop det Intet, af hvilket Gud ifølge Kirkelæren skabte
Verden. Og Fichte kommer derved ind paa Forestillingen
om en engang for alle færdig Gud, hvis Natur kun
ses som Villiens Produkt, en naturløs Gud. Han polemiserer
mod Schellings Forsøg paa at eftervise et ubevidst
Naturliv indenfor selve Guddommen. Naturen er
altsaa for Fichte væsentlig et Afledet, ikke et Oprindeligt\footnote[1]{„Der Geist der Natur, weil er sich nur in unbewusster Weisheit des Wirkens zu zeigen vermag, ist eben darum der nicht-absolute, nicht der Geist Gottes.“ „Beiträge zur Charakteristik der neuern Philosophie“. 2. Aufl. p. 765.}.
% --- p. 165 ---
\setcounter{footnote}{0}%
\opage{165}Det store Resultat, som den spekulative Idealisme havde
naaet med Hensyn til Gudsideen, skjøndt endnu under
ensidig Form, var derimod, at det guddommelige Væsen
ikke kunde være et abstrakt Subjekt, en ren, naturløs
Aand, men at det som den sidste Grund til Alt ogsaa
maa være væsenbeslægtet med Alt. Den yngre Fichte
forskjertser dette Resultat ved sin Spiritualisme, der
ikke har Blik for Naturlivets Betydning med Hensyn til
Bestemmelsen af Gudsideen. ---

Vi følge Fichte videre til hans Antydninger om Historiens
Philosophi og Religionsphilosophien. Idet Enkeltvæsenerne
fra den evige Urposition gaa over i den endelige
Tilværelse, den timelige Udvikling, bliver det egentlig
Centrale i disse Enkeltexistentser det aandelige Uranlæg,
i hvis eiendommelige Form Individualiteten bestaaer.
Dette er, hvad Fichte kalder Genius. Gjennem dette
ideale Element staaer Enkeltexistentsen i indre Forbindelse
med det guddommelige Væsen. Grunden til
det Onde ligger ikke i selve Endeligheden, ei heller i
en ubegribelig Vilkaarlighed hos Mennesket, men i Muligheden
af, at det endelige Væsen ikke naaer det Maal,
der er sat for det. Da derfor hverken Endeligheden eller
Synden medfører en absolut Modsætning til Gud, eller
fører udenfor Gud, saa kan dennes Virken i Tilværelsen
heller ikke gaa for sig paa ydre, mirakuløs Vis. Gjennem
det ideale Uranlæg (Genius) i det Endelige arbeider
Guddommen styrkende, fremmende, potentserende henimod Verdensudviklingens Maal. Ved denne guddommelige
Virken ophæves Skridt for Skridt det Isolerede og
Dualistiske i den endelige Tilværelse. Alle egentlig historiske
Bedrifter ere frembragte ved den guddommelige
Vækkelse af Menneskets Genius, ved den Begeistring,
der griber den menneskelige Frihed, saa at den, idet den
øver sin egen Gjerning, faaer den Bevidsthed, at dens
Indhold ligger ud over den selv, maa være den indgydt.
Paa det høieste Trin er det den religiøs-ethiske Ide,
% --- p. 166 ---
\setcounter{footnote}{0}%
\opage{166}der begeistrer, og Individet bliver derved Mellemmand
mellem Guddommen og den øvrige Menneskehed, religiøs
Genius, Prophet. Det høieste Maal, som skal naaes gjennem
denne stedse af guddommelige Kræfter baarne Udvikling,
er den inderligste og fuldstændigste Enhed af
Selvtilfredsstillelse og Hengivelse. Al Salighed bestaaer
i den bevidste Fuldendelse af sit eget Væsen i og ved
Kjærligheden.

Den guddommelige Virken er netop guddommelig,
fordi den ikke træder frem i ydre, raa Haandgribelighed,
men ganske iklæder sig Enkeltvæsenernes Frihed og den
endelige Tilværelses Love. Det sande Under bestaaer
i den Maade, hvorpaa gjennem denne indre Udvikling
nye Ideer pludselig kunne træde frem og Genier vækkes
til at bære dem. Det bestandige Under er den guddommelige
Aands Indtræden i den endelige, hvorved
denne bliver Medskaber og Deltager i Verdensudviklingens
Proces. Derimod vil ingen philosophisk Tænkning
nedlade sig til at antage egentlige Mirakler, Afbrydelser
eller Ophævelser af Naturordenen; thi dette er kun det 
Hensigtsløse og Aandløse, den sandselige Parodi paa
hine aandelige Undere. Guds Virken er universel, ikke
bunden indenfor visse Tidsbegivenheders Skranke.

Dog indtræder der her hos Fichte, ligesom ved Udviklingen
af Gudsideen og af Skabelsesproblemet, et Brud
paa det Rationelle. Ligesom han hist fordrede „et
specifisk Mere“ foruden Enheden af Frihed og Nødvendighed,
saaledes er han heller ikke her tilfredsstillet ved
den vel noget ubestemte, men dog i sit Princip rationelle
Bestemmelse af Historiens Philosophi, men fordrer,
under aabenbar Paavirkning af den positive Religion, at
Verdensudviklingen maa naa sit Høidepunkt i et Gudmenneske,
idet Enheden af den guddommelige og den
menneskelige Aand oprindelig kun skal kunne tænkes i
et Subjekt. „Det kan ikke være nogen Mangfoldighed
af Jeger, i hvilken Gud oprindelig bliver Menneske, men
% --- p. 167 ---
\setcounter{footnote}{0}%
\opage{167}som i den hele Skabning Intet vakler i ubestemte Omrids,
men er tilspidset til det Individuelles afgjorte Eiendommelighed, saaledes maa ogsaa den høieste Skabelsesgjerning ganske afsluttes og fuldendes i et enkelt Faktum.
Gudmenneskeligheden, hvorved Guds Aand paa
universel Maade skal gjøre sig til Et med Menneskeslægten,
er en fuldstændig uklar Tanke, ligesaa abstrakt
og halv som det pantheistiske Gudsbegreb, til hvilken
den ganske svarer, idet begge bero paa, at man holder
det Abstrakte, blot Almene for det Høieste. Gudmennesket
ene løser Historiens forvirrede Gaade, idet hans
Tilværelse og dens historiske Følger yder os den virksomme
Garanti for, at i ham er den guddommelige Aand midt
iblandt os og saa vist vil fuldbyrde Menneskeslægtens
sædelige Forløsning, som den først har sendt Gudmennesket.“ --- Dog skal dette kun være et almindeligt,
med Nødvendighed fremkommende Postulat; deraf skal
ikke følge videre Konsekventser med Hensyn til den
Maade, hvorpaa det historisk virkeliggjøres, skjøndt
Fichte ikke undlader i Forbigaaende at bemærke, at der
deraf opstaaer „et indre Bevis af eiendommelig Art“ for
Troværdigheden af de historiske Beretninger, der handle
om Gudmennesket.

Det er let at se, at alt Dette staaer i Strid med
de Præmisser, hvoraf det udledes. Men dette er en Strid,
som nødvendig fremkommer ved den uklare Maade, hvorpaa
Fichte opfatter Forholdet mellem Immanents og
Transscendents. Det skal, ifølge en oftere gjenkommende
Vending, være Immanentsen selv, der fører til Transscendents;
men deri ligger da og, at Transscendentsen
paa den anden Side betinges ved og omsluttes af Immanentsen.
Det er særlig en høist overfladisk Bevisførelse,
der ligger til Grund for Postulatet om Gudmennesket.
Vel er Strauss's Opfattelse af Ideen pantheistisk;
men alligevel er hans Udsagn rigtigt, at Ideen
aldrig fuldstændig udtømmer sig i en Enkeltexistents.
% --- p. 168 ---
\setcounter{footnote}{0}%
\opage{168}Derimod ligger den Maade, hvorpaa den Straussiske Opfattelse
skal korrigeres, i Fichtes egen Lære om Individualitetens
Betydning. Ideen om Guds og Menneskets
Enhed er i sig selv ikke nok, men skal personlig tilegnes
og leve i de enkelte Individer. Hver Enkelt skal,
som Mystikerne lærte, blive en Christus. Det er dette
Eftertryk paa Ideens Iboen i Personlighederne, der fører
ud over det Pantheistiske, ikke den phantastiske Fordring
om en enkelt og absolut Virkeliggjørelse af Ideen.

Andre Udtalelser tyde iøvrig paa, at det ikke var
Fichtes Mening ved hint Postulat at hylde det positive
Dogme. Han nævner det som en stor Velgjerning ved
sin Opdragelse, at han var bleven aldeles religiøst opdragen,
men at dog de egentlige Troeslærdomme med
deres „Hemmeligheder“ og Ubegribeligheder forbleve
ham fjernt. „Denne Forskaanelse for det helt Overflødige
bevarede mig for den farlige Konflikt, ved Indtrædelsen
af en modnere Dannelse at maatte bryde med
Troesautoriteten. Det Evige, Almenmenneskelige i Troen
forlod mig aldrig. Den ydre, historiske Indfatning kunde
jeg trøstig overlade til den sigtende Kritik (Verm.Schr.
I. p. 327). Philosophien maa, det hævder han overfor
Schelling og Weisse, fuldstændig sige sig løs fra den
theologiske Forestillingskreds. Den christelige Theisme
skal udvikle sig til den universelle, humanistiske Theisme.
„Thi vi give at betænke, at ikke blot den fra
Grunden af forandrede og udvidede kosmologiske Verdensanskuelse,
men endnu mere det ulige friere Overblik over
de store verdenshistoriske Religioners Væsen og fælles
Oprindelse, ligesom endelig den dybere Erkjendelse af
det Guddommelige i enhver ægte Kulturudvikling, --- at
alt Dette ogsaa spekulativt gjør et mere udvidet og tillige
dybere udviklet Gudsbegreb nødvendigt“ (p. 119). ---

De Skrifter af I. H. Fichte, vi hidtil have beskjæftiget
os med, nemlig „Grundzüge zum Systeme der Philosophie“
i deres tre Dele: Erkjendelseslære, Ontologi 
% --- p. 169 ---
\setcounter{footnote}{0}%
\opage{169}og spekulativ Theologi hidrøre fra Aarene 30 og 40.
Det var da, som vi oftere have fremhævet, det pantheistiske
Problem, der stod paa Dagsordenen. Trods alle
sine Uklarheder har Fichte bidraget meget til Kritiken
af den ensidige Pantheisme. Ogsaa med Hensyn til det
Problem, der i Aarene 50 og 60 paatrængte sig. Materialismen,
Spørgsmaalet om Forholdet mellem det Aandelige
og det Materielle, har han Fortjenester. --- Var det ved
det første Stridsspørgsmaal dialektiske og metaphysiske
Undersøgelser, det især gjaldt om, saa blev nu Psychologien
den philosophiske Videnskab, der tiltrak sig den
høieste Interesse. Allerede i sine første Skrifter appellerede
Fichte til Erfaringen som Spekulationens Supplement,
og det var Studiet af den, der overbeviste ham
om det Individuelles Realitet. Han udtaler endog, at
han ene grunder sin Lære paa Erfaringen og de paa
den byggede nødvendige Slutninger. Men det har tildels
allerede vist sig i det Foregaaende og vil under den
nærmere Betragtning af Fichtes anthropologiske og psychologiske
Skrifter tydeligere vise sig, at han kun vanskelig
formaaer at stille sig i et objektivt og uhildet Forhold
til det Givne. Ligesom han ved sit Forsøg paa at overvinde
den Schelling-Hegelske Spekulation blev Gjenstand
for en berettiget Kritik fra Weisse, som dog ellers i
mange Maader var enig med ham, saaledes mødte ogsaa
den Maade, hvorpaa han senere søgte at bearbeide Erfaringsvidenskabens
Resultater i anthropologisk og psychologisk
Henseende, en skarp Indsigelse fra \emph{Hermann}
\emph{Lotze} (Streitschriften. 1stes Heft. 1857), der ligeledes
hvad den hele almindelige Retning angaaer udtalte sin
Samstemmen med ham. Paa begge Punkter synes Forklaringen
af dette eiendommelige Forhold at ligge i, at
Fichte efter hans egen, ovenfor omtalte Erklæring er
Personlighedsphilosoph og stedse umiddelbart søger at
fyldestgjøre den religiøse Interesse. Han har ikke Resignation
nok til at bøie sig for Spekulationens og Er%
% --- p. 170 ---
\setcounter{footnote}{0}%
\opage{170}faringens objektive Magter for ved dem at naa den sande
Erkjendelse, som saa ogsaa maatte indeslutte den sande
Tilfredsstillelse af den religiøse Interesse; men denne
Interesse griber ubevidst ind i Undersøgelsernes Gang og
hindrer deres objektive, videnskabelige Gjennemførelse.

Dette ses blandt Andet af hans Polemik mod Atomlæren.
Den moderne Physik har ført til Nødvendigheden
af at antage Materien for i sidste Instants at bestaa af
diskrete Dele eller Atomer. Der kræves nemlig til den
exakte Forklaring af det Vexlende i Phænomenerne et
et Uforanderligt, et Konstant. „Det gjælder,“ siger W. % PRINTED AS IS (et Uforanderligt: the image, et doubled over the line turn (et|et); recrop rc2)
Weber\footnote[1]{I et Brev til Fechner, anført i dennes „Atomenlehre“ 2te Aufl. p. 88.}, „ved Bevægelsens Aarsager at udsondre en
saadan konstant Del, at Resten vel er foranderlig, men
at dens Forandringer kunne tænkes afhængige ene af
maalelige Rums- og Tidsforhold. Ad denne Vei kommer
man til et Begreb om Masse, hvortil Forestillingen om
rumlig Udstrækning slet ikke nødvendig knytter sig. Konsekvent
maales da ogsaa Atomernes Størrelse i den atomistiske
Anskuelse ingenlunde efter rumlig Udstrækning,
men efter deres Masse o: efter det ved ethvert Atom
konstante Forhold, i hvilket Kraften staaer til Hastighedens
Forøgelse.“ Atomlæren er i denne Form en
Hypothese, til hvilken Videnskaben nødvendigvis maa
komme. Men det er klart, at der herved ikke læres
nogen dogmatisk Anskuelse; Atomlæren er ingen Afslutning,
kun det for Physiken, altsaa relativt Sidste.
Men den Fordring maa stilles til den, som vil optage
den moderne Physiks Resultater i sin Verdensanskuelse,
at han ikke omgaaer, men virkelig gjennemtænker denne
nødvendige Forudsætning. Ligesaalidt her, som tidligere
overfor den spekulative Opfattelse af det Absolute, er
det nok med partielle Ændringer og Omtydninger.

Alligevel betragter Fichte i sin Anthropologi (1856)
% --- p. 171 ---
\setcounter{footnote}{0}%
\opage{171}Atomlæren mere som en Fiktion end som en nødvendig
videnskabelig Hypothese. Han søger derfor at bekæmpe
den med alle de Vaaben, der staa til hans Raadighed,
uden at betænke, at han netop mangler det Vaaben,
hvorpaa det her ene vilde komme an, nemlig den exakte
Methode. Og dog vil han angribe Atomlæren „i dens
egentlige Midtpunkt, paa selve Naturvidenskabens Omraade“!
Han bliver her det idealistiske Modbillede til
Czolbe, der i Kraft af Materialismens Princip polemiserede
mod den moderne Physiologi. Men medens Czolbe
vender sig mod Videnskabens klare, positive Lære, sammenblander
Fichte nogle enkelte materialistiske Atomikeres
Mening med Physikens Atomlære. Deraf kommer det,
at han betragter denne som en Lære, der kan gjendrives
med spekulative Vaaben, og han overser ganske den skeptiske
Karakter, enhver saadan fagvidenskabelig Hypothese
besidder i metaphysisk Henseende. Det er derfor
som det vil vise sig, mere Mangel paa skarp Opfattelse
af Erfaringsvidenskabens Princip og Methode end virkelig
Uoverensstemmelse med dens Resultater, der ligger
til Grund for hans Polemik mod Atomlæren.

Hvad der nødte Physikerne til at tage deres Tilflugt
til Atomlæren, det var efter hans Mening den falske
Forudsætning om det tomme Rum, hvorfra de gik
ud. Naar de nu ikke vilde miste det Reales faste Grund,
maatte det tomme Rum tænkes opfyldt af et Andet, som
var ligesaa oprindeligt, og dette i sidste Instants Reale
og Rumopfyldende tænktes saa som de udelelige Atomer,
hvilke altsaa indeholdt den Grændse for Deleligheden,
som de synlige, af Atomerne sammensatte Ting ikke
besad. Men Forudsætningen om det tomme Rum mente
han at have gjendrevet i sin Ontologi, hvor han havde
vist, at Rum og Tid ikke bestaa selvstændig og uafhængig
af Tingene, men kun ere Former for det Reale, forsaavidt
dette dels i sig selv bestaaer og vedvarer, dels hævder
sig overfor Andet ved at sætte og udfylde sin Være- og
% --- p. 172 ---
\setcounter{footnote}{0}%
\opage{172}Virkekreds. Al Kvantitet forudsætter overhovedet Kvaliteten
som det Udfyldende og Bestemmende, ligesom omvendt
enhver Kvalitet medfører sit eiendommelige kvantitative
Udtryk. Rum og Tid ere kun de i sig selv uselvstændige
Former for alt Realt; kun i Abstraktionen kunne
de faa Selvstændighed --- en saadan Abstraktion, som
efter Fichtes Mening ligger til Grund for Physikernes
Atomlære. Naar man atter ophæver denne Abstraktion
og derimod fra de rent kvantitative Former slutter tilbage
til det Kvalitative, som ligger til Grund, saa omstødes
derved ogsaa den mechaniske Atomlære. Istedetfor
de kvalitetsløse, mechanisk uopløselige Rumpunkter
eller Atomer træde saa kvalitative Atomer, som sætte
og opfylde deres Rum og ved deres indbyrdes Slægtskab
og derved betingede Vexelvirkning frembringe Phænomenet
af relativt uigjennemtrængelige Legemer.

Foruden den mechaniske Atomlære bekæmper Fichte
ogsaa den spekulativ-dynamiske Opfattelse af Materien.
Denne Opfattelse, som blev grundlagt af Kant\footnote[1]{I hans „Metaphysische Anfangsgründe der Naturwissenschaft“ (1786). Derimod stod Kant i en ældre Afhandling: „Metaphysicæ cum geometria junctæ usus in philosophia naturali“ (1756) paa det atomistiske Standpunkt. Smlgn. \emph{Lotze} i „Gött. gel. Anz.“ 1855. p. 1095---1099.}%
% CHECK p. 172: star lost in the layer: paired with the mark by order
 og senere
udvikledes videre af Schelling og Hegel, danner Modstykket
til Atomlæren. Medens denne fra først til sidst
søger at punktualisere og centralisere og først fra de
diskrete Elementer hæver sig til det Totale og Sammenhængende,
betragter den dynamiske Opfattelse Materien
som et Kontinuum, hvis sidste Grund den søger. Materiens
Grundeiendommeligheder, Uigjennemtrængeligheden
og den forskjellige Rumopfyldelse, føres tilbage til to
forskjellige Kræfter, Attraktions- og Repulsionskraften, af
hvis indbyrdes Forhold den phænomenale Materie resulterer.
Fichte indvender mod denne Konstruktion, at enhver
Kraft forudsætter et realt Subjekt, hvis Egenskab den
% --- p. 173 ---
\setcounter{footnote}{0}%
\opage{173}er, eller ved hvis Forhold til et andet Subjekt, den først
kan lægge sig for Dagen som Kraft. Men dette Kraftens
reale Subjekt lader den spekulative Konstruktion
ude af Betragtning. Den overser, at der mangler et individualiserende
Moment, hvorved Kræfterne først blive
virkelige Kræfter. Et saadant frembyder Physiken i sine
Atomer, naar vi kun berøve dem deres abstrakt kvantitative
Beskaffenhed. I sin kvalitative Atomistik mener
Fichte saaledes at have fundet et Enhedsudtryk for det
Berettigede saavel i den mekaniske Atomlære som i den
spekulative Dynamik. Men han har ved sin misforstaaede
Polemik mod den førstnævnte berøvet sig selv det reale
Grundlag, som den metaphysiske Konstruktion behøver;
og idet han først omdanner og omtyder den Virkelighed,
der skal forsones med det Ideelles Krav, kan Forsoningen
selv ogsaa kun blive ideel. Det reelle Grundlag
bliver et Phænomen uden virkelig Betydning. ---

De naturphilosophiske Undersøgelser danne Forberedelsen
til de anthropologiske. Det almindelige Forhold
mellem Sjæl og Legeme er allerede bestemt ved den
kvalitative Atomistik. Som ethvert Realt producerer sit
Rum og sin Tid, idet det opfylder dem, saaledes frembringer
Sjælen, det til Bevidsthed anlagte Reale, sit
Legeme. Herved gjøres der, efter Fichtes Mening, Ende
baade paa Dualismen og Materialismen. Den legemlige
Væren i Rum og Tid ses ikke længere som noget
for Aanden Fremmed, men fremgaaer med Nødvendighed
af dens Væsen, og altsaa kan den heller ikke have
nogen absolut Magt over Sjælen, hvis eget Produkt
den er. Men hvorledes er denne Produceren mulig?
Vanskeligheden af at indse det beroer, ifølge Fichte,
ene paa, at man har vænnet sig til kun at opfatte Sjælen
som et klart bevidst, reflekterende og overveiende
Væsen. Saaledes maatte man ved Undersøgelsen af Menneskets
Væsen standse ved en Tohed: den bevidste Tanke
og den materielle Legemlighed. Men den nærmere Over%
% --- p. 174 ---
\setcounter{footnote}{0}%
\opage{174}veielse viser snart en Mellembestemmelse. Legemet er
ikke blot Materie. Stoffet levendegjøres, formes og omdannes
af de organiske Kræfter. Vi faa altsaa en Trehed:
Aand, organisk Kraft og legemligt Stof. Stoffet
frembyder ingen Vanskelighed; dets Love gjennemskues
af den organiske Chemi. Derimod synes den organiske
Kraft at være noget Dunkelt og Gaadefuldt. Det er
den, der ved sin Omdannelse af de elementære Stoffer
frembringer den ydre, phænomenale Legemlighed, hvis
egentlig sammenholdende og forenende Baand den er.
Medens den ydre Legemlighed, som i Almindelighed
kaldes „Legeme“, aldrig er den samme, men uafbrudt
vexler i sine enkelte Dele, danner de organiske Kræfters
System derimod en fast, indre Legemlighed, en pneumatisk
Organisme, hvis Reflex eller „Effekt“ det ydre
Legeme er, som „kastes ind i den vexlende Stofverden,
omtrent som den magnetiske Kraft af Jernfilspaanerne
bereder sig et tilsyneladende fast Legeme, som dog
splittes ad til alle Sider, naar den bindende Magt unddrages
det.“

Men i hvilket Forhold staaer nu denne Kraft til
Sjælen selv? Dette Spørgsmaal er det især, hvis Besvarelse
forudsætter, at man har aflagt den ved Aarhundreders
scholastiske Dannelse nærede Fordom, at Sjælen
skulde være et rent intellektuelt og bevidst Princip. Den
organiske Kraft bærer i alle sine Virkninger Præget af
fuldkommen Fornuftighed og indre Hensigtsmæssighed.
Hvad den frembringer er ikke ringere, end om en høist
fuldkommen Intelligents med frit Overlæg havde valgt
og udført det med kunstnerisk Virtuositet. Og dog ere
alle organiske Processer ubevidste, altsaa ubevidst fornuftige,
bære Aandens Præg uden at være Aand. Den
intensive Sikkerhed i dette Princips Virken og den overlegne
Fornuft, det lægger for Dagen, gjør det til et individuelt
Forsyn for Legemet, som ikke blot virker ved
dettes første Frembringelse, naar „Urpositionen“ legem%
% --- p. 175 ---
\setcounter{footnote}{0}%
\opage{175}liggjøres, men ogsaa under det fortsatte Liv lægger sig
for Dagen i alle Ernærings-, Reproduktions- og Selvhelbredelsesprocesser.
En Analogi til en saadan Virken
have vi i Kunstnerens uvilkaarlig formende Sands og i
Dyrenes Kunstdrift. Navnlig i denne sidste lægger sig
en høist energisk, drømmeagtig Phantasivirksomhed for
Dagen. Og da nu Livsprocessen viser os en lignende
 Virken som Dyrenes instinktmæssige Arbeiden, saa ligger
det nær at antage den samme Aarsag dertil. Phantasien
er netop en saadan Mellembestemmelse mellem Bevidsthed og Ubevidsthed, som vi søgte. „Livsprocessen lader
sig i alle sine charakteristiske Phænomener kun forklare
som Sjælens mest intensive Phantasivirksomhed. Den
staaer dybest under Bevidsthedens Tærskel (dybere end
de egentlige Instinkter, der allerede føre ud over det
Ubevidste, men dog i paaviselig Kontinuitet med dem),
men netop derfor virker den rigest og skjønnest, da den
ligger længst borte fra den sandselig-reflekterende Tilstand,
som er den for Sjælens oprindelige Inderlighed
mest fremmede.“  Phantasien er altsaa den plastiske
Kraft, ved hvilken Sjælen af Stofverdenens Elementer
danner sit Legeme, en Lære, som blandt Andet allerede
Stahl (c. 1700) har banet Veien for, og som efter Fichtes
Mening kun Fordomme have trængt tilbage. „Phantasien
er aldeles ikke blot en aandelig Evne o: en Virksomhed,
der kun hører til Sjælens bevidste Sphære; den
danner ret egentlig en Mellembestemmelse, er en ligesaa
ubevidst realiserende som i sig ideel Evne, og naaer
derfor ganske paa samme Maade ned i Livsprocessernes,
den ubevidst hensigtsmæssige Legemsdannelses og Legemsopholdelses
Omraade, som den tjener de høieste Ideer
til besjælende Form.“

Sjælen som „der Herr seiner Verleiblichung“ kan
ikke ligge under for det Legemliges Forkrænkelighed.
Det ydre Legeme med Alt, som dertil hører, staaer
snarere i et uvæsentligt Forhold til Aanden, end at det
% --- p. 176 ---
\setcounter{footnote}{0}%
\opage{176}skulde være at betragte som en nødvendig Betingelse for
dens Virkeliggjørelse og dens Bevidsthed. Dette sandselige
Legeme i dets faktiske Beskaffenhed er for den en
Skranke, som virker hæmmende for alle dens Bevidsthedsfunktioner.
Der gives derfor ogsaa Tilstande (clairvoyance,
magisk actio in distans o. s. v.), i hvilke Sjælen
ideelt frigjør sig fra denne Legemligheds Skranker,
idet den skuer og virker uafhængig af de legemlige Organer.
Idet under disse ekstatiske Tilstande den sædvanlige
Bevidstheds Horizont omfattes af en friere og
mere omfattende Bevidsthed, som omvendt slet ikke existerer
for den sædvanlige, sandselig-bevidste Tilstand,
saa maa vi i den ekstatiske Tilstand have „das Vollbewusstsein
des Geistes“, „das vollgeistige Dasein“, medens
den sædvanlige Tilstand kun er en bunden og halveret
Tilværelse. Den ideelle og midlertidige Frigjørelse
fra Legemligheden bliver til en reel, naar ved Døden
det ydre Legeme igjen opløses i sine Elementer. Ved
denne Forandring gaaer Sjælen ikke over i en rent
ulegemlig Tilværelse; Døden bestaaer tvertimod kun i
den indre, pneumatiske Organismes fuldstændige Befrielse
og Udløsning. Gjennem sit indre Urlegeme, de organiske
Kræfters System kan Sjælen altsaa fremdeles sætte
sig i Forbindelse med en Omverden. Kun forudsættes
der her et Stof, som kan tjene til Medium for denne
Virken; men et saadant frembyder jo den nyere Physik
i Ætheren, der skjøndt utilgjængelig for den sandselige
Opfattelse dog er det egentlig Virkende i Naturprocesser
som Lys, Varme o. s. v.

Dette er Hovedtrækkene af Fichtes Lære om Forholdet
imellem Sjæl og Legeme. Han lægger stærkt
Eftertryk paa, at han overalt støtter sig til Erfaringskjendsgjerninger;
men disse ere rigtignok dels tagne i
meget vid Betydning dels fortolkede paa en meget fri
Maade. Kun i meget uvidenskabelig Forstand kan man
saaledes kalde Phænomener som clairvoyance o. s. v.
% --- p. 177 ---
\setcounter{footnote}{0}%
\opage{177}Erfaringskjendsgjerninger ; dertil ere de endnu for lidet
undersøgte og kritisk prøvede. I ethvert Tilfælde maa
disse til „Livets Natside“ hørende Phænomener mere
betragtes som organisk-psychiske Sygdomstilfælde end
tillægges en direkt psychologisk Betydning. Netop for
det høiere sjælelige Liv kan man ikke tilskrive dem
noget Værd, end sige at man med Fichte i dem skulde
kunne finde Sjælens „vollgeistige Dasein“. Kun Mangel
paa Kraft til at beherske den sande og virkelige Erfarings
Kjendsgjerninger kan føre en Tænker til at lægge
en saadan Vægt paa disse Phænomener.

Hvad Fichtes Anthropologi i det Hele angaaer, da
har Lotze i sit Stridsskrift træffende paavist, at han ikke
har naaet, hvad han vilde (og i temmelig høie Toner
erklærede at have naaet), overfor Materialismen og Dualismen
at paavise den sande Enhed af Sjæl og Legeme.
„\emph{Dette} Legeme,“ siger Lotze (p. 131), „paa hvis Sammenhæng
med den levende Sjæl den naturlige Følelse
ene sætter Pris, dette varme Blod, dette elastiske Kjød,
Knoglernes sikre Fasthed og Senernes Spændkraft, alt
Dette forbliver efter Deres egen Theori Sjælen ganske
fremmed. Et indre, usynligt Legeme erklærer De for
det sande; „det andet, den ydre Tilsyneladelse deraf,
som er dannet ud af en uafladelig Stofvexel, kan man
gjerne vedblive at kalde Legeme; --- det er ikke i Sandhed
ét, ikke bestaaende, men den blotte Effekt og Reflex
af hin indre Legemlighed““\dots{} Nei, det er dette, der
altid alene har heddet Legeme; kun for dette Legemes
Enhed med Sjælen kan den naturlige Følelse interessere
sig, thi det alene kjender den\dots{} At kalde dette Legeme
Billede og Virkning af et andet, er indholdsløse
Ord, saalænge indtil De har paavist, ved hvilken Kraft
overhovedet det ene er istand til at frembringe en Virkning
i det andet, eller efter hvilke Love det kan frembringe
sit Billede i de for det fuldkommen fremmede og
udvortes Masser. Det er den store Kløft, for hvilken
% --- p. 178 ---
\setcounter{footnote}{0}%
\opage{178}De stedse gyser tilbage; hvilket saa end Organismens
formende Princip monne være, det kan kun virke saa
Meget og paa de Maader, Naturens almindelige Ret tillader.
Denne den mechaniske Virkeliggjørelses Proces,
Opsøgelsen af de Udgangspunkter, fra hvilke de enkelte
samstemmende Kræfter virke,  Efterforskningen af de
specielle, i Naturløbet ved en uafbrudt Tradition vedligeholdte
og i fast Rhythme sig omformende Betingelser,
af hvilke Skabningernes Udfoldelse, Væxt og Forplantning
fremgaaer, --- alt Dette er nødvendige Opgaver
for Videnskaben, og jeg kan ikke overbevise mig om,
at de løses ved Forestillingen om en genial, organiserende
Duft“\footnote[1]{Lotzes Kritik af I. H. Fichtes Anthropologi minder i flere Henseender om Leibnitz's Kritik af Stahls Lære. Opera ed. Dutens. vol. II. pars 2. p. 131 ff.}. --- Det er mere i en uklar Tendents
end med en virkelig videnskabelig Gjerning, at Fichte
kommer ud over den spekulative Naturphilosophi og Psychologi.
Realiteten omsætter sig for ham til en ideel
„Duft“, hvis Enhed med det Aandelige det saa ikke
falder ham vanskeligt at paavise. ---

I sit sidste betydeligere Værk (Psychologie. 1ster
Theil. 1864) udtaler Fichte det som en af Hovedopgaverne
for sin Psychologi at lægge Grunden til Omdannelsen
af Religionsvidenskaben paa psychologisk Grundlag.
Hvilket dette Grundlag er, ses ved nærmere at
fremhæve den Maade, hvorpaa Fichte opfatter Forholdet
mellem Bevidstheden og Aandslivet i det Hele. --- Bevidstheden
er som saadan ikke produktiv. Den frembringer
intet Nyt, men ledsager kun med sit Lys visse
reale Tilstande og Forestillinger i Sjælen. Dens Lys
glider som en flygtig Fakkel hen over Sjælens forskjellige
Tilstande, idet den vexelvis belyser dem og lader
dem i Mørke. Men det Høieste i Aanden er, som vi
allerede have set, det bag Bevidstheden liggende reale
% --- p. 179 ---
\setcounter{footnote}{0}%
\opage{179}Indhold, som kun i enkelte overordentlige sjælelige Tilstande
træder frem for Aanden, idet denne hæves til
en høiere, mere omfattende Bevidsthed. Denne mystisk-ekstatiske
Karakter besidde nu de egentlig religiøse
Phænomener. Dette har Schelling for bestandig godtgjort
i sin Mythologiens Philosophi. Kun har han ikke
Ret i at lægge den Magt, der i hine Tilstande uimodstaaelig
virker paa og i Bevidstheden, ud over Menneskets
Væsen ved at betragte de forskjellige Religioner som
Virkninger af og Vidnesbyrd om en theogonisk Proces.
Dog skjelner Fichte ligesom Schelling mellem den blotte
Ekstase og den egentlige Aabenbaring. I hin hæves
Menneskets Aand vel ud over sin sædvanlige Begrændsning,
men forbliver dog indenfor sin egen Verden. Saaledes
den sædvanlige Sonnambulisme, Anelser o. s. v.
Men naar Aanden engang er indtraadt i denne Tilstand
af aandelig Skuen, bliver den uvilkaarlig ogsaa modtagelig
for Afspeiling af en anden Aandsverden i sin Bevidsthed;
ved denne „Ekstase i anden Potents“ træder Aanden
ind i Aabenbaringens Region og hæves ud over
det menneskelige Bevidsthedstrin. Dette er Tilfældet,
hvor der i Religionerne opgaaer et saadant Indhold, at
den mythologiske Phantasivirksomhed ikke længere er
en tilstrækkelig Forklaring. Et saadant Indhold er Guds
over Verden ophøiede, rent aandelige Enhed. Her kan
Aarsagen kun søges i en objektiv Indvirkning paa den
menneskelige Aand af den Aand, der i Aabenbaringen
giver sig selv Vidnesbyrd. Begreberne Aabenbaring og
Overnaturlighed høre ikke nødvendig sammen. Tvertimod,
Gud virker sine høieste og vældigste Gjerninger i
kontinuerlig Sammenhæng med de naturlige Phænomener,
ved Forhøielsen af den almindelige Verdensorden.

Idet Fichte nu i dette „overbevidste“ Indhold ser
Aandens sande Væsen, bliver Bevidstheden kun af formel
Betydning. Oprindelig gik han ud fra en anden Opfattelse,
forsaavidt han gjennem Erkjendelseslæren  søgte en
% --- p. 180 ---
\setcounter{footnote}{0}%
\opage{180}Begrundelse af det Absolute. Senere tog han, som vi
have set, dette tilbage og saa Erkjendelsen som en
Kjendsgjerning, hvis Grund og Aarsag man skulde søge.
I sin Psychologi endelig gaaer han endnu et Skridt
videre. Der skal ikke sluttes fra Bevidstheden til det
Absolute. Bevidsthed er ikke noget Selvstændigt og for
sig Bestaaende, men kun et endeligt Væsens Opfattelse
og Belysning af sig selv. Saamange Bevidsthedscentra
derfor, saamange Henvisninger til reale, individuelle
Aandsvæsener. Kun ved denne Tankegang skal Philosophien
kunne indledes saaledes, at Pantheisme og Idealisme
udelukkes. Bevidstheden er ikke Andet end den
mest intensive Driftvirkning, hvor Paavirkningen fører
til en indre Belysning af Driften for Aanden selv.
Sjælen er „ein instinktbehaftetes Triebwesen“. Bevidsthedens
Spire ligger i den Trang, Driften har til at søge
sin Gjenstand. Denne Søgen og Eftersporen forhøies
og opklares, naar Driften finder det, der tilfredsstiller
den; derved reflekterer den sig i sig selv og frembringer
Bevidstheden om Forholdet mellem sig og sin
Gjenstand. Driften er altsaa Mellembestemmelsen mellem
det Objektive, blot Organiske og det Subjektive,
den bevidste Fornemmelse.

Resultatet vilde altsaa blive, at Bevidstheden kun
er et Prædikat ved et forudsat Subjekt, ikke et i sig
Realt og Substantielt. Det vilde være den fuldstændige
Modsætning til den spekulative Idealisme, for hvem netop
Tanken og Selvbevidstheden var Udtrykket for Menneskets
sande Væsen. Fra theologisk Side har man med
Glæde accepteret dette Resultat, der ophæver Erkjendelsens
rationelle og nødvendige Karakter og gjør det
muligt at opbygge den kirkelige Dogmatiks transscendente
Lærdomme paa en ny Grundvold. Dog er dette
sikkert noget forhastet, selv om man ser Sagen fra Fichtes
Synspunkt. Hvor meget han end betragter det som en
Virkning af „Aarhundreders scholastiske Dannelse“ at
% --- p. 181 ---
\setcounter{footnote}{0}%
\opage{181}sætte det bevidste Liv over det ubevidste, saa er det
dog ingenlunde hans Mening derved at ville retfærdiggjøre,
hvad der ikke har nogen virkelig rationel og psychologisk
Begrundelse. Selv det „Overbevidste“ er efter
hans Mening kun et høiere Rationelt. Og ved nøiere
Eftersyn vil man finde, at han ikke konsekvent kan
hævde Driften som det oprindelige sjælelige Element.
Phantasien er jo ifølge hans Lære det oprindelige, producerende
og organiserende Sjæleprincip, og af den maa
ogsaa Driften afhænge, --- altsaa bliver alligevel det
ideelle eller intellektuelle Element det Oprindelige: Driften
forudsætter Formaalet, som sættes i den ubevidst
formende Phantasi. ---

3. \emph{Christian Hermann Weisse} (f. 1801, d. 1866)
udgik oprindelig fra Hegels Philosophi, fra hvilken han
kun successivt fjernede sig, idet han, som han udtrykte
sig, som en Sibylle, stedse fordrede flere Indrømmelser
af Hegel til Gjengjæld for ringere og ringere Anerkjendelse.
Allerede paa et tidligere Punkt (ovenfor p. 23f.)
have vi havt Anledning til at omtale de Hovedpunkter,
hvori han afveg fra Hegel. Det var ved Verdensanskuelsens
Høidepunkt, i Læren om den absolute Aand, at
Uoverensstemmelsen først opgik for ham. Hegel angiver
Forholdet mellem de høieste Ideer, som udgjøre den
absolute Aands Indhold og Væsen, saaledes, at det Absolute
først fremtræder som Skjønhedens Ide, i Kunsten,
under Anskuelsens Form; derfra skrides der frem
til den inderligere Form i Forestilling og Følelse, i Religionen,
indtil den absolute Inderlighed og den høieste
Enhed naaes i Videnskabens, i Sandhedens Ide, som for
Hegel er den logiske Ide, de rene Kategoriers Rige.
Dette var for Weisse en fuldkommen Formalisme. Enheden
af Form og Indhold blev hos Hegel tagen i den
Betydning, at Indholdet gik op i Formens rene, abstrakte
Væsen; der kunde for ham ikke være Tale om et høieste
Formaal, en sand og levende Virken i det Absolute.
% --- p. 182 ---
\setcounter{footnote}{0}%
\opage{182}Den hele Tilværelse tjener her kun til paa dialektisk
lovbestemt Vis at udvikle det Absolute i dets uendelige
Vexelspil af Former, og det Virkeliges Fuldkommenhed
beroer paa, i hvilken Grad det kan bringe det Absolutes
logiske Former til Aabenbarelse. Weisse opstiller herimod
den Fordring, at der i det Absolute, i den sidste
Grund og det høieste Formaal, maa kunne paavises et
virkeligt Værd, et i sig gyldigt og levende Gode, som
skal udvikles og virkeliggjøres i Tilværelsen paa Basis
af hine Former. Kun saa ville ogsaa Følelsen og Anskuelsen
finde Hvile i det Absolute, medens de frastødes,
ja tilintetgjøres ved den rene Abstraktion.

Det hegelske Systems Hovedsætning var, som ofte
fremhævet, Enheden af Fornuften og Virkeligheden,
men saaledes, at kun det blev anset som Virkelighed,
der var absolut Et med Fornuften, og Fornuft vil her
sige de rene logiske Kategoriers Indbegreb. Denne Sætning
giver Weisse en anden Betydning, idet han fortolker
den saaledes: „det Fornuftige er absolut Formbestemmelse,
det ubetingede Grundlag og conditio sine qva
non for det Virkelige.“ Den logiske Ide udtømmer altsaa
ikke Virkelighedens Indhold; dette er et positivt
Mere, den frie Virkelighed i Rum og Tid, som af Hegel
betragtes som et rent Negativt, et Affald fra Ideen.
Sandhedens Ide i sin rene almene Form kan da ikke
længere danne Slutningen i Læren om den absolute
Aand; den er tvertimod kun Begyndelsen, det Grundlag,
paa hvilket et Indhold af høiere Værd kan udfolde
sig. Denne Udfoldelse sker ved, at hint Almene specificeres
og individualiseres i en uendelig Mangfoldighed
af Former ved en producerende eller fødende Virksomhed,
der foregaaer i, hvad Weisse kalder den guddommelige
Natur eller det guddommelige Gemyt. Denne Proces
har Æsthetiken, som Videnskaben om Skjønhedens
Ide, at fremstille, og Weisse har forsøgt en saadan
Fremstilling i sit første større Værk „System der Aesthe%
% --- p. 183 ---
\setcounter{footnote}{0}%
\opage{183}tik (1830)\footnote[1]{Om dette Skrifts specielle Indhold se \emph{Lotze}: „Geschichte der Aesthetik in Deutshland.“}. Æsthetiken danner Overgangen til den % PRINTED AS IS (Deutshland: the image, O118)
høieste Ide i Aandens Rige, det Godes Ide, Ethikens
og Religionsphilosophiens Indhold. Det Godes Ide forudsætter
det Skjønne, fordi det bestaaer i Villien til
Meddelelsen af et Realt. Det Godes Ide bliver indholdsløs,
naar den ikke til sin Forudsætning har en levende
Fylde af Realitet, som den, det Skjønnes Ide aabenbarer
os i Guddommens Natur.

Hegels Grundfeil er, at han har hypostaseret de almene
logiske Kategorier til et positivt og realt Existerende,
gjort den logiske Ide til Gud. Det er en Feil,
han deler med hele den philosophiske Dogmatisme, der
i de nødvendige Tankebestemmelser umiddelbart ser
det virkelig Værende. Men det Store ved Hegel er, at
han først har fattet og udført Ideen om en fuldstændig
Udvikling og Fremstilling af de almene og nødvendige
Bestemmelser, der ligge til Grund for al Væren. Fordybet
i dette Arbeide ansaa han selve disse Bestemmelsers
Indbegreb, den logiske Ide for den eneste Virkelighed.
Vil man ret forstaa Hegel, da maa man integrere
denne Feil. Hans Logik maa betragtes som Fremstilling
af Virkelighedens nødvendige Tilværelsesformer, ikke af
den positive Virkelighed selv. At denne Opfattelse er
rigtig, derpaa finder Weisse et Bevis i, at den lærer
os at se Hegels Grundvildfarelse i det Samme, der
netop udgjør hans Storhed, saa at der ikke bliver noget
Tilfældigt og Uforklaret tilbage.

Men denne Grundvildfarelse med Hensynet til Forholdet
mellem Fornuft og Virkelighed fører i det Enkelte
andre Misligheder med sig. Den i Rummet og Tiden
sig udfoldende Virkelighed betragtede Hegel som et
Affald fra Ideen, fordi Rum og Tid for ham vare den
blotte Udvortesheds Former uden væsentlig Betydning
% --- p. 184 ---
\setcounter{footnote}{0}%
\opage{184}for det Absolute. Med en forandret Opfattelse af Virkeligheden
maatte der ogsaa følge en forandret Opfattelse
af Virkelighedens Former. Weisse bebreider Hegel, at
Rum, Tid og de dertil sig knyttende Bestemmelser ikke
udvikles i Logiken, da jo dog al Virkelighed er i Rum
og Tid, og disse altsaa ligesaavel som de rene Begreber
ere absolute og nødvendige Tilværelsesformer. Derved
har Hegel afskaaret sig Muligheden af en indre Overgang
fra Logiken til Naturphilosophien og bragt en
formelig Dualisme ind i sin Anskuelse.

Ved denne Kritik af Hegel var Weisse ad sin egen
Vei kommen til et lignende Resultat, som det Schelling
nogle Aar senere udtalte i sin Fortale til Cousin (se
ovenfor p. 126 f.). Til denne Schellings Udtalelse kunde
han derfor i det Hele slutte sig. Men dog forblev der
en gjennemgaaende Forskjellighed mellem de to Tænkere,
som kun ikke traadte bestemtere frem, fordi Schellings
senere Lære først langt senere blev bekjendt i sin
Helhed. Weisse optog Schellings Adskillelse mellem
negativ og positiv Philosophie. Den logiske Ide, Indbegrebet
af de rene, almene Fornuftbegreber blev da for
ham „det negative Absolute“. Men i den nærmere Bestemmelse
af Forholdet mellem det Negative og Positive
kommer Uoverensstemmelsen for Dagen. Hos Schelling 
faldt, som vi have set, den positive Philosophi sammen
med Religionsphilosophien, og det positive Udgangspunkt,
hvorfra den gaaer ud, kunde kun gribes af den, der utilfredsstillet
ved Livet i Ideen ved en religiøst bestemt
Villiesakt postulerede den „idefrie“ Gud. Vel negtede
Schelling ikke Nødvendighedens Princip, saaledes som
Nogle af hans Tilhængere (navnlig Retsphilosophen Stahl,
ved hvem Schellings nye Lære først blev bekjendt i
Nordtydskland), der ved „positiv Philosophi“ forstode
en Tilbageføren af al Lov og al Sandhed i Tilværelsen
til en ubegrændset Guddomsvillie; men i sidste Instants
kunde den negative Philosophi dog kun faa en tilfældig
% --- p. 185 ---
\setcounter{footnote}{0}%
\opage{185}og forsvindende Betydning, naar man umiddelbart, gjennem en Villiesakt kunde gaa til den positive Philosophi.
Hvad Weisse billiger i Schellings senere Lære, er
for det Første Opfattelsen af Modsætningen mellem Nødvendighed
og Frihed som en real Modsætning, der ikke
kan ophæves ved en formel Dialektik, som hos Hegel,
der i sin logiske Ide ser Enheden af disse Modsætninger.
Det er fremdeles Opfattelsen af Nødvendigheden
som det negative Grundlag eller Prius for den virkelige
Existents, efter Weisses Udtryk: som den rene Tilværelsesmulighed.
Kun det Frie er det Reale, det Existerende,
det Positive. Men dette Positive kan Philosophien
ikke gribe paa positiv Vis --- ved denne Sætning
adskiller Weisse sig fra Schelling. Den negative
Philosophi, der for Weisse falder sammen med Metaphysiken,
men hos Schelling var Et med hele den ikkereligiøse
Viden, har derfor her væsentlig Betydning for den
positive. Den logiske Ide indeholder Roden til det Positive,
og først ved det Metaphysiskes, de logiske Bestemmelsers
Iboen bliver det Positive et Frit. Friheden
udspringer saaledes af Nødvendigheden og har sin Grund
i den. Friheden eller det, at Indholdet ogsaa ikke kan
være, bestaaer altsaa ikke i et udvortes Forhold til den
rene Nødvendigheds Form eller i en Uafhængighed af
Formen; men Indholdet bliver netop kun frit ved Formen,
der i sin Fuldendelse er Frihedens Form\footnote[1]{Grundzüge der Metaphysik (1835). Das philosophische Problem der Gegenwart (1842).}.

I sin Opposition mod Hegel havde Weisse fremdeles
en Kampfælle i I. H. Fichte, med hvis Anskuelse man
derfor en Tidlang betragtede hans som væsenligt Et.
Dog ere de ikke blot uenige i Opfattelsen af Hegel,
men ogsaa i deres hele øvrige Anskuelse fremtræde saa
store Forskjelligheder, at de foranledigede Weisse til
Affattelsen af det tidligere omtalte Stridsskrift: „Das
% --- p. 186 ---
\setcounter{footnote}{0}%
\opage{186}philosophische Problem der Gegenwart.“ Fichte saa i
den hegelske Logiks Absolute Verdensaandens Begreb,
og mente at kunne acceptere dette som relativt berettiget;
hans Kritik gik saa ud paa, at Erfaringen, særlig
den indre Hensigtsmæssigheds Kjendsgjerning, nødte til
at gaa et Skridt videre, da Verdensaandens Begreb ikke
indeholdt den tilstrækkelige Forklaring. Weisse saa derimod
i Hegels logiske Ide de rene, nødvendige Tankebestemmelser,
som gjælde for al Virkelighed, men i sig
selv kun udgjøre et Mulighedens Rige. Kategorierne
vare for Fichte kun Bestemmelser ved et forudsat Virkeligt,
som han enten gjennem Erkjendelseslæren eller Erfaringen
søgte at naa Forvisningen om. Heri saa Weisse
den modsatte Ensidighed af Hegels. „Hos Dem,“ skriver
han til Fichte (p. 188), „fremtræde de samme Kategorier
som Produkter af en uselvstændig Reflexion, som
abstrakte Skygger af en hinsidig Virkelighed, hvilke vi
hos Hegel, ikke mindre imod deres Natur, se hypostaserede
til Indbegrebet af det i Sandhed Værende“. Den
rette Opfattelse er efter hans Mening den, der i Tilværelsesmulighedens
almene Bestemmelser vel endnu ikke
ser den reale Virkelighed selv, men søger at udvikle
dem til et Punkt, hvor de med indre Nødvendighed gaa
over i denne. Fichte kom, som vi have set, til at svæve
mellem Pantheismen og den vulgære Theisme, fordi han
ikke vilde vide af nogen apriorisk eller ontologisk Begrundelse
af Gudsideen; derved maatte ligesom hos
Schelling Frihedsbestemmelsen i Gud indtræde ved en
ubegreben Position.

Under denne sin Bestræbelse for en rigtig Vurdering
af den metaphysiske Erkjendelses Væsen er Weisse sig
bevidst at have en Forgænger i Kant\footnote[1]{Se, foruden den p. 119 anførte akademiske Tale, Weisses Afhandlinger i 25de og 26de Bind af Fichtes Tidsskrift: „Ueber die transscendentale Bedeutung der Urtheilsformen und Schluss%
\opage{187}figuren“ og „Ueber den letzten Grund der Gewissheit im Denken“, samt „Philosophische Dogmatik“ 1ster Band (1855).}% the note runs over from p. 186 to p. 187 (joined by draft.py)
. Kant bestred
% --- p. 187 ---
\setcounter{footnote}{0}%
\opage{187}Dogmatismen, der i de rene Fornuftbegreber troede at
have en Erkjendelse af Tingene i sig selv, navnlig en
Erkjendelse af det Uendelige og Ubetingede, det Evige
og Nødvendige i Tingene, og altsaa forvexlede Erkjendelsens
nødvendige Grundformer med de virkelige Tings
reale Væsen. Kant saa derimod Forstandens oprindelige
Eiendom som noget blot Formelt, der a priori viser os hen
til Erfaringen, ved hvilken det ene kan naa et virkeligt
Indhold. For at naa Begrebet om en Gjenstand maa
Forstanden stedse føie Noget af sit Eget til det ved
Fornemmelsen givne Materiale; af disse to Elementers
Kombination fremgaaer Gjenstandens Begreb. Hvad der
her tilføies, er i Modsætning til Erfaringens Enkelthed
det Almene, i Modsætning til det i Erfaringen givne
Stof det Formale, og i Modsætning til Erfaringsstoffets
Virkelighed det blot Mulige. Men disse Forstandselementers
Gyldighed opfattede Kant kun som en subjektiv-psychologisk
Nødvendighed for den menneskelige
Aand; han saa ikke, at den i Forstandens Prius indeholdte
Tankemulighed tillige er en Tilværelsesmulighed.
Fra denne Mislighed maa hans Anskuelse frigjøres for
at svare til Sandheden. En saadan Frigjørelse er sket
gjennem den senere philosophiske Udvikling (I. G. Fichte,
Schelling, Hegel), der førte til en Erkjendelse af Ideens
almene Væsen. Dog førte selve den Energi, med hvilken
Frigjørelsen for de subjektive Skranker gjennemførtes,
atter for en Stund tilbage til Dogmatismen, idet
man nu i den metaphysiske Ide, hvis objektive Betydning
man havde paavist, saa selve det høieste Virkelige.
Den med Kant begyndte philosophiske Bevægelse
maa saaledes, naar den rette Bevidsthed om den rene
Fornuftvidenskabs Natur og Opgave igjen trænger
igjen%
% --- p. 188 ---
\setcounter{footnote}{0}%
\opage{188}nem, atter vende tilbage til sit Udgangspunkt, men beriget
og styrket ved den vundne Indsigt.

I enhver Tankeakt, enhver Dømmen og Slutten er
hin rene Tanke- og Tilværelsesmulighed tilstede. Den
indtræder som Kopula mellem Forestillinger, som derved
ordne sig som Dommens Subjekt og Prædikat; den er
Kraften, som i Slutningen paa indre Vis sammenføier
Leddene gjennem Mellembestemmelsen. Den har sit
nærmere Indhold i de Begreber, der for Forstanden ligge
til Grund for Erkjendelsens Gjenstande og selve den erkjendende
Bevidsthed, ligesom og for Bevidstheden og
Gjenstandene i deres indbyrdes Forhold. Al Virkelighed
hviler altsaa paa disse Forudsætninger. Vel opdager den
naturlige eller umiddelbare Forstand ikke dette; men
naar Tvivlen reiser sig, søger Bevidstheden tilbage til
det almene Grundlag, den absolute Nødvendighed, der
tillige indeholder de absolute Muligheder. Vi have da
her Sandhedens Begreb i dets mest abstrakte Form, den
rent logiske Sandhed, der bestandig forudsætter sig selv,
og om hvilken Thomas Aqvinas's Ord gjælde: „Qui negat
veritatem esse, concedit veritatem esse; si enim veritas
non est, verum est, veritatem non esse.“ I denne
abstrakte Form har Sandhedens Begreb sit Udtryk i den
saakaldte Modsigelsens Grundsætning, den første negative
Form for Tilværelsesmuligheden: Alt er muligt, som
ikke indeholder en Modsigelse. Denne Sætning leder
os til at se, at allerede den rene Fornuft, uden endnu
at sætte eller forudsætte noget Virkeligt, knytter Tilværelsens
Mulighed til visse Grundbetingelser. Disse
Betingelser, der ere al Tilværelses Urformer, udgjøre
den rene Tankenødvendigheds Indhold, det rene Fornuft-Absolute,
som den philosophiske Spekulation til alle
Tider har betragtet som sin første og egentlige Gjenstand,
uden at det endnu er lykkedes den at bestemme
det i dets fulde Renhed. Det Absolute bliver altsaa den
sidste Grund til Tankens Vished, fordi i det det Objek%
% --- p. 189 ---
\setcounter{footnote}{0}%
\opage{189}tives Urmulighed er Et med Tænkningens høieste Nødvendighed.


Allerede for den umiddelbare Forstand træder dog denne
Urmulighed frem under mere konkret Form. Bliver den
umiddelbare Forstand sig end ikke som Spekulationen Fornuftkategoriernes
Nødvendighed bevidst, saa tillægger den
dog de mathematiske Bestemmelser større Evidents og fuldkomnere
Vished end den blotte Erfaring. Og dog kan den
ikke tilskrive de mathematiske Begreber en selvstændig
Væren. Den søger da at opfatte dem som knyttede til
og afhængige af det i Erfaringen Givne; men heller
ikke dette kan lykkes, thi hvorledes kan Noget, der
skylder det Givne sin Tilværelse, overgaa det i Vished
og Nødvendighed? --- Derfor er det, som allerede Platon
og Kant have indskærpet, den sikreste Vei at lede den
umiddelbare Forstand til spekulativ Erkjendelse, naar man
gaaer ud fra Undersøgelsen om Grunden til den mathematiske
Nødvendighed. I Tal, Rum og Tid ser den
spekulative Fornuft Grundforudsætninger for al Tilværelse,
hvorved denne falder ind under mathematiske
Love. Metaphysiken skal udvikle Begrebet om det Værende,
og saasnart den har frigjort sig for den idealistiske
Dogmatismes Ensidighed, maa den indse, at Værens
Begreb staaer i en oprindelig og nødvendig Sammenhæng
med Begreberne Tal, Rum og Tid. Besvarelsen af det
Spørgsmaal: hvad er Væren? fører os da gjennem de
Kategorier, ved hvilke Rums- og Tidsudfyldningens nødvendige
Grundformer betegnes, til et Grundbegreb, i
hvilket al Mulighed af et Indhold, der kan udfylde disse
Former, altsaa al Mulighed af en Væren og Vorden er
sammenfattet til en Urmulighed. Dette Grundbegreb,
den absolute Ide er Et med Guddommens Begreb, men
det vil her, hvor vi endnu staa paa det metaphysiske
Standpunkt, sige Begrebet om den mulige Guddom. Der
gives kun én for Tanken nødvendig Sandhed, nemlig at
% --- p. 190 ---
\setcounter{footnote}{0}%
\opage{190}oprindelig er kun Gud mulig, og at i hans Mulighed
alle Tings Mulighed indeholdes.

Metaphysiken følger den dialektiske Methode. Medens
Weisse snart kom til det Resultat, at den dialektiske
Methode i Hegels Forstand, som Begrebets eller
Indholdets Selvbevægelse, ikke kunde finde Anvendelse
paa det realphilosophiske Omraade, hvor Tanken Skridt
for Skridt træder i Vexelvirkning med Erfaringen, mente
han dog endnu at kunne beholde denne Methode paa
det metaphysiske Omraade. Saaledes skriver han i Fortalen
til sin Metaphysik: „Den hegelske Philosophis formale
Sandhed og materiale Usandhed, dens Methodes
gedigne Fortræffelighed og dens Resultaters trøstesløse
Tomhed paatrængte sig min Aand med lige Evidents.“
Med Rette gjorde, som vi have set (p. 24), baade Tilhængere
og Modstandere af Hegel, gjældende, at det var
en Selvmodsigelse at indskrænke den dialektiske Methode,
der er Indholdets Selvbevægelse, til blot formal
Gyldighed. Weisse har derfor nærmere præciseret sin
Opfattelse af Methoden. Ved den dialektiske Methode
forstaaer han selve den philosophiske Erkjendelses Begreb,
men ikke særlig den Form, den har antaget i Hegels
Panlogisme. Dialektiken har en forskjellig Karakter
paa det metaphysiske og paa det realphilosophiske
Omraade. Hist lægger den sig for Dagen i den Magt,
med hvilken enhver Kategori ligesom trænger sig frem
for at gjælde som Helhedsbestemmelse for Tilværelsen,
som Indbegreb af Almuligheden. Saaledes er Dialektiken
i Philosophiens Historie. Kampen og Udviklingen i denne
beroer netop paa, at de som Helhedsbestemmelser successivt
opstillede Kategorier bestandig fortrænges af
andre mere omfattende og gjennemgribende. Dette er
ogsaa Udviklingsgangen indenfor den metaphysiske Videnskab.
Maalestokken er Virkelighedens Totalitet, som
godtgjør den indre Modsigelse i de utilstrækkelige Kategorier.
Paa det realphilosophiske Omraade er Forholdet
% --- p. 191 ---
\setcounter{footnote}{0}%
\opage{191}omvendt. Her betragtes ethvert i Erfaringen givet Virkeligt,
som om det indbefatter Kategoriernes Totalitet, men
viser netop derfor med Nødvendighed ud over sig selv
til andre og høiere Udtryk for denne. Paa dette Omraade
bliver den dialektiske Methode Et med den genetiske
Methode. Thi hvad der forbinder den kategoriske Tanke
med de empiriske Phænomer, er den for disse Phænomeners
Tilblivelse og Udvikling gjældende Lov. Ved Loven,
der her ligesom i den rene Fornuftspekulation udtrykker
Tilværelsens Mulighed og Tankens Nødvendighed, beherskes
Bevægelsen paa begge Sider, Tilværelsens genetiske
Bevægelse, der ved Loven sættes som mulig, og Tankens
Bevægelse, der ved Loven sættes som nødvendig.

Paa samme Maade som det enkelte realphilosophiske
Begreb forholder sig til det existerende Reale, forholder
de metaphysiske Begrebers Totalitet sig til den konkrete
Virkelighed. Gangen i Metaphysiken falder efter Weisse
sammen med, hvad der udgjør Sandheden i det gamle
saakaldte ontologiske Bevis for Guds Tilværelse. Dette
Bevis fører i Virkeligheden ikke videre end til Begrebet
om Guddommens Mulighed, om Gud som den, der ikke
kan tænkes som ikke værende. Dog fører den metaphysiske
Betragtning i sit Slutningspunkt til en skarpere
Bestemmelse af Guddommens Væsen. Da den absolute
Ide, Urmulighedens Indbegreb er en Enhed i Uendeligheden
af de i den ophævede Forskjelle og Modsætninger,
bliver den nødvendige Form for denne Tilværelse
Tankens Form. Idebestemmelserne bestaae kun i og for
Tanken. Idet nu denne Tanke i evigt Kredsløb tænker
sig selv i alle hine Bestemmelser, er det evig Værende
selvbevidst Ursubjekt, Urpersonlighed. Vi komme altsaa
til det Resultat, at dersom den evige Mulighed ikke
forbliver afsluttet i sig selv, men gaaer over i Virkelighed,
saa maa denne Urvirkelighed være Urpersonlighed.
Man kunde tænke sig et saadant Abstraktionens Standpunkt,
hvor der ses bort fra al virkelig Væren, og hvor
% --- p. 192 ---
\setcounter{footnote}{0}%
\opage{192}derfor ogsaa Gud kan tænkes som ikke virkelig værende.
Der er ikke, som Hegel mener, en umiddelbar, dialektisk
Overgang fra Begreb til Tilværelse. Hegel opfattede det
ontologiske Bevis som det, der paaviste denne Overgang,
og betragtede derfor dette Bevis som det eneste ægte
videnskabelige. Mod ham vil man med Rette kunne
vende hin Abstraktion. Saasnart man derimod sætter
Noget som existerende, maa man ogsaa sætte Guds Tilværelse.
Man vilde ellers hilde sig i den Modsigelse
paa engang at sætte et Virkeligt og ophæve Indbegrebet
af de Betingelser, uden hvilke Intet kan blive virkeligt.
--- Dette er den paa engang dialektiske og faktiske Overgange
fra Metaphysiken til Realphilosophien eller fra
det ontologiske Bevis til det kosmologiske. Begge ere
paa det Inderligste forenede, udgjøre de to Halvdele af
samme Hele.

Men dog synes der ved denne Modsætning mellem
det negative og positive Absolute at sættes et Mulighedens
Dyb forud for Guds virkelige Existents, saa at
han først ved en Tilblivelses- eller Selvskabelsesakt arbeider
sig frem til Virkelighed. Der bestaaer efter
Weisse\footnote[1]{Se hans Artikel „Gott“ i Erschs og Grubers Encyklopædi Sekt. I. vol. 75. p. 424.} en uforsonlig Modsætning mellem den philosophiske
Spekulations Absolute og den religiøse Erfarings
Personlighedsfordring, saalænge man stiller den
Fordring, at Guds Personlighed \emph{umiddelbart} skal
sættes som det Absolute, eller det Absolute \emph{umiddelbart}
som Personlighed, istedetfor at den rene Tankes
Absolute kun indeholder den uendelige Mulighed af en
levende, personlig Tilværelse. I.H. Fichte (Verm. Schriften
I. p. 100) bemærker hertil, at en saadan Adskillelse
er utænkelig, hvis den skal gjælde objektivt, i Guds
eget Væsen og ikke være en blot theoretisk Sondring
af den menneskelige Tænkning. Lotze (Geschichte der
% --- p. 193 ---
\setcounter{footnote}{0}%
\opage{193}Aesthetik p. 202, smlgn. allerede Metaphysik p. 322 f.)
udleder Læren om det negative Absolute deraf, at idet vor
Tænkning ser bort fra Verdensindholdet og ene fordyber
sig i den logiske Sandheds Rige, opstaaer Skinnet af,
at denne logiske Sandhed er et Oprindeligt og Selvstændigt
i Forhold til den levende Virkelighed. Og vi
have allerede ovenfor (p. 120 f.) henpeget til, at dette
Punkt maatte blive det vanskeligste for hele den Retning,
til hvilken Weisse hører. At Weisse ikke vil spalte
Gudsbegrebet i to Dele, ses af hans udtrykkelige Erklæring
(Phil. Dogm. I. p. 439), at Adskillelsen mellem
Mulighed og Virkelighed i Gud kun er begrebsmæssig,
idet begge ere Momenter eller ideale Bestanddele af det
ene og udelelige Gudsbegreb. Og denne Adskillelse er
nødvendig for at indse, at Gud kun ved Ideen, ved det
negative Absolute bliver fri, bliver det absolute Væsen,
Nødvendighedens Herre. Heri ligger den reale (ikke
blot logiske) Enhed af Frihed og Nødvendighed. Uden
en saadan en indre Modsætning i Guds Væsen vil efter
Weisses Mening Begrebet om den levende Gud ikke
kunne hævdes, og han opstiller den ogsaa netop i Modsætning
til den dogmatiske Philosophis abstrakte Begreb
om Gud som actus purus.

Under en anden Form fremtræder den samme Adskillelse,
naar Weisse i Tilslutning til Schelling skjelner
mellem hvad Gud er (quid Deus sit) og at Gud er
(quod Deus est), mellem Guds Væsen og Existents, saaledes
at den ontologiske Undersøgelse skal give os Qvid,
Gudsbegrebets almindelige Indhold, og den kosmologiske
dertil føier Qvod, Existentsens Bestemmelse. At denne
Modsætning er uholdbar, have vi allerede set ved Schelling.
Men Weisse anvender den ogsaa kun for at udtrykke
det dobbeltsidige Udgangspunkt, som er nødvendigt.
Kun paa den absolute Abstraktions Standpunkt,
vilde ifølge Weisse denne Modsætning have Gyldighed.
Realphilosophien (Natur- og Aandsphilosophien) har des%
% --- p. 194 ---
\setcounter{footnote}{0}%
\opage{194}uden ikke blot den Betydning at hjælpe Erkjendelsen til
den ene Bestemmelse: Existents; tvertimod følger der
med Erkjendelsen af den existerende Virkelighed en
stedse rigere og konkretere Udvikling af Gudsideens Indhold;
der føies til det første, rent rationale Qvid et
andet, kvalitativt Qvid, et Qvale, som begrunder Guddommens
æsthetiske og ethiske Prædikater. Metaphysiken
i og for sig kan ikke give os den rette, fyldige
Erkjendelse af Gudsideen; en saadan naa vi først, naar
vi paa Basis af den metaphysiske Erkjendelse fordybe
os i Virkelighedens Indhold. Pantheismen har sin ubestridelige
videnskabelige Berettigelse i den spekulative
Erkjendelse af en den virkelige Verdens Skikkelser gjennemtrængende
og betingende metaphysisk Nødvendighed,
som den dog ikke formaaer at adskille fra Tingenes
virkelige Tilværelse, da den ikke naaer den Ide, i
hvilken denne i sig selv kun formale eller negative Nødvendighed
finder sin Afslutning. For at komme ud over
Pantheismen er det altsaa nødvendigt at føre den metaphysiske
Nødvendighed tilbage til dens rette Betydning
indenfor Gudsbegrebet og vise, at den kun er en Forudsætning
for et fyldigere Indhold. Baade overfor den
abtrakte Theisme og den ensidige Pantheisme har hin % PRINTED AS IS (abtrakte: for abstrakte; the image, s02)
Adskillelse altsaa sin store Betydning.

Overgangen fra den ontologiske til den kosmologiske
Undersøgelse, fra Metaphysisk til Realphilosophi % PRINTED AS IS (Metaphysisk til: the image, for Metaphysik; O126)
var, som vi have set, paa engang apriorisk og aposteriorisk,
en indre Overgang i Begrebet i Et med en Erfaring
af det virkelig Existerende. Det nye Element, som
herved kommer ind, kan ikke strax finde sin Retfærdiggjørelse.
Der sættes en Modsætning mellem Erfaring
og Begreb, som kun successivt kan ophæves. Dette
sker, som allerede antydet, gjennem de realphilosophiske
Videnskaber. Her, hvor vi nærmest skulle undersøge,
hvilke kvalitative Bidrag disse Videnskaber kunne yde
til Erkjendelse af Gudsideen, gjælder det dialektisk at 
% --- p. 195 ---
\setcounter{footnote}{0}%
\opage{195}udvikle Verdensindholdet, som det fremtræder i Erfaringen,
til det Punkt, hvor Verdensbegrebet gaaer over i
Gudsbegrebet. --- Naturen staaer ikke blot under mechaniske
Love, men disse Love føie sig ind i en høiere
teleologisk Sammenhæng, hvori alle dens Dele staa i
uendelig Vexelvirkning med hverandre og danne en Trinrække
ligefra den tilsyneladende døde Masse til de høiste
Fornuft- og Aandsvæsener. Dette Resultat fører Erfaringen
til, og det ligger allerede i den rene Fornuftsides
Væsen. Da denne selv er den uendelige Enhed,
maa ogsaa den Tilværelse, hvis almene Betingelser den
indeholder, have Sammenhæng og Enhed til sine Grundformer.
Derfor drager Fornuftideen sig som en ledende
Tanke fra den ontologiske Undersøgelse over i den kosmologiske
og paaviser de Steder, hvor den teleologiske
Sammenhæng ligesom tilspidser sig i Almenbegreber,
der som Prædikatsbestemmelser slutte sig udfyldende
og fuldstændiggjørende til det metaphysiske Begreb om
Ursubjektet. Den kosmologiske Undersøgelse gaaer da
paa sit Høidepunkt over i den ethikologiske, thi det er
paa det aandelige Omraade, at det sidste Formaal og
altsaa ogsaa ogsaa det høieste guddommelige Prædikatsbegreb % PRINTED AS IS (ogsaa ogsaa: the image, O129)
maa søges. Idet Tanken, nu ledet af den videnskabelige
Bevidsthed om Ideen, søger at bringe Betragtningen
af Virkeligheden i Samklang med denne Bevidstheds
Indhold, fremtræder selve det Maal, den her søger
at naa, som det, der tillige er hele Verdensudviklingens
Maal. Det høieste Princip fremstiller sig da som Erkjendelsen
af Fornuftideen i dens Enhed med Virkeligheden,
eller med andre Ord som Sandhedens Ide.
Men denne Ide har ikke blot theoretisk Betydning, den
bliver ogsaa den høieste Lov, det sædelige Princip. Sædelighedens
Princip bestaaer i, at det endelige Subjekt gjør
det til sin Handlens Formaal, hvad der er Formaalet for
den objektive Verdensorden. Hævet over den endelige
Aands empiriske Villie viser den konkrete Sandhedside
% --- p. 196 ---
\setcounter{footnote}{0}%
\opage{196}ham dog tillige de objektive Formaal, han skal gjøre til
sine. Men ved Siden af den theoretiske og praktiske
Sandhedsbestræbelse gaaer en anden Række af Virksomheder,
der ikke mindre have deres Formaal i sig selv.
Her træder Fornuftideen og dens Virkelighed i Natur
og Historie frem som Kilde til en Formernes Verden,
hvilken Aanden, befrugtet ved Berøringen med den virkeliggjorte
Ide, frembringer af sit Indre og nyder under
sin Skuen og Frembringen. Alt som hører herhen, indbefattes
under Skjønhedens Ide.

Til disse to Ideer, Sandhedens og Skjønhedens,
fører altsaa den realphilosophiske Undersøgelse. Men
de staa i en indre Modsætning til hinanden. Ifølge
Sandhedens Ide er det høieste Verdensformaal Enhed og
og absolut Idealitet; Skjønhedens Ide hævder derimod % PRINTED AS IS (og absolut Idealitet: the image, og doubled over the line turn (og|og); O130)
det selvstændige Værd af den endelige Aands æsthetiske
Tilstande og Virksomheder og sætter derved en uendelig
Mangfoldighed af reale Verdensformaal. Dette er den
sande Betydning af Kants Antinomi mellem Sædeligheden
og Lyksaligheden. Denne Modsigelse er ikke blot subjektiv,
men er ogsaa objektiv, forsaavidt man bliver
staaende ved hine to Ideer som reale Verdens- og Aandsmagter
og ikke søger en høiere Forsoning. At en saadan
maa søges, følger af selve Fornuftideens Væsen, der
ikke kan tilstede saadanne Magters Selvstændighed. De
maa finde deres Plads indenfor Gudsideen. Det ligger
i Ursubjektets Begreb, at det ikke kan sætte sig selv
som teleologisk Princip, uden at derved Muligheden sættes
af, at ogsaa den endelige, individuelle Tilværelse
bliver Formaal i og for sig selv. Baade den logiske og
den æsthetiske Virkelighed have saaledes deres Grund
i Guddommen. Denne Betragtning fører ind paa det
ethisk-religiøse Omraade, den høieste og afsluttende Erfaringssphære.
Religionsphilosophien eller den spekulative
Theologi danner det philosophiske Systems Høide og
Slutningspunkt, hvor Enheden af Ide og Virkelighed
% --- p. 197 ---
\setcounter{footnote}{0}%
\opage{197}fuldendes, idet Spekulationens høieste Begreb mødes med
den ethisk-religiøse Erfaring, som i sig indeslutter al
anden Erfaring. Det kommer da her an paa den nærmere
Opfattelse af den religiøse Erfarings og Religionsphilosophiens
Væsen\footnote[1]{„Philosophische Dogmatik oder Philosophie des Christenthums“. 1855---1862. 3 Bind.}.

Religionens Begreb gaaer ifølge Weisse ind under
det almindelige Begreb Erfaring. Erfaring er paa den
ene Side mere end en blot Modtagelighed eller rent subjektiv
Tilstand i Sjælen, men paa den anden Side mindre
end egentlig Erkjendelse eller Videnskab. Længe
førend Kant i Slutningen af forrige Aarhundrede for
første Gang gav en nøiagtig Bestemmelse af Erfaringens
Væsen, var det ofte blevet udtalt fra den religiøse Bevidstheds
Side, at Religionen væsentlig er Erfaring (παθεῖν τὰ θεῖα. Dion. Areopag.). Det er Mystikerne, der have
havt den klareste Indsigt i Religionens erfaringsmæssige
Karakter. Det Samme ligger i de gamle Kirkelæreres
Bestemmelse af Forholdet mellem πίστις og γνῶσις.
Men dog varede det, selv efterat Erfaringens Begreb
forlængst var hævdet paa andre Omraader, længe inden
en lignende Overbevisning kunde trænge igjennem med
Hensyn til Religionen. Den scholastisk-orthodoxe Retning
i Kirken ansaa den religiøse Erfarings Indhold som
Noget, der udenfra bliver meddelt Menneskeslægten som
færdige Lærdomme, og først af denne Meddelelse fremkaldes
saa efter dens Mening den subjektive Bevægelse
og Fornemmelse hos de Enkelte. Paa den anden Side
troede Mystikerne i det, de oplevede i Gemyttet, umiddelbart
at have en Erkjendelse af et Objektivt. Først
ved Kants Undersøgelser, der tilsidst ogsaa strakte sig
til Religionsphilosophien, som han egentlig har grundet,
blev det muligt at se klarere ogsaa dette Omraade. Det
rent objektive Læreindhold sætter han tilside som blot
% --- p. 198 ---
\setcounter{footnote}{0}%
\opage{198}statutarisk; men paa den anden Side falder det Religiøse
ikke for ham, som i den saakaldte Fornuftreligion, sammen
med den metaphysiske Erkjendelse; han fastholdt
et nyt Udgangspunkt for Religionsphilosophien i den
ethisk-religiøse Bevidsthed, der for ham væsenlig betinges
ved den i Menneskets Indre givne Lov, det praktiske
a priori. Han har styrtet Dogmatismen ikke blot paa
det philosophiske, men ogsaa paa det religiøse Omraade.
Dette Sidste var den vanskeligste Opgave, da det religiøse
Indhold paa Grund af sin Inderlighedskarakter er
vanskeligere at skjelne fra den rene Fornufts Indhold.

Det ethisk-religiøst Givne staaer i samme Forhold
til Erkjendelsen som alt andet i Erfaringen Givet; ogsaa
det er bundet til de Love, der give Omdannelsen af
det subjektivt Oplevede til et objektivt Erkjendelsesindhold
dens videnskabelige Form og Betydning. Al
Erfarings Love staa i organisk Sammenhæng, saa at den
aandelige Erfarings Love til deres Forudsætning have
den physiske Erfarings Love og uden denne Forudsætning
ikke kunne have nogen Betydning. En Erfaring,
der ikke hviler paa faste, uforanderlige Love,
kan ikke adskilles fra en Skuffelse. Og omendskjøndt
den religiøse Erfaring som \emph{religiøs} har sine eiendommelige
Love, saa hviler dog denne Lovmæssighed
paa de psychologiske Love i Almindelighed, og disse
igjen paa de physiske. Uden disse Forudsætninger vilde
den religiøse Erfaring ikke være \emph{Erfaring}. Forestillingen
om Underet er uethisk og ureligiøs, netop fordi
de ethisk-religiøse Begreber hvile paa det psychologiske
og physiske Grundlag. Den konsekvent gjennemførte
Supranaturalisme har ogsaa bestandig istedetfor at lede
til levende Religionserfaring begunstiget den blinde Autoritetstro,
der netop er den fuldstændigste Modsætning
til al virkelig og levende Erfaring. Naar det umiddelbare
religiøse Standpunkt alligevel er Underforestillingens
frugtbare Jordbund, da hidrører dette fra Trangen
% --- p. 199 ---
\setcounter{footnote}{0}%
\opage{199}til ved digtende Phantasi at finde et Udtryk for det
Uudsigelige i den religiøse Erfarings Kjendsgjerninger;
det Underfulde opstaaer da, fordi den naturlige („ausserreligiöse“)
Erfarings indre Sammenhæng ikke staaer som
Kjendsgjerning for Bevidstheden. Thi den religiøse Erfaring
indbefatter al anden Erfaring i sig som ideel Totalitet;
omvendt gaaer den naturlige Erfaring paa sit
Høidepunkt over i den religiøse. Derved samles alt
Erfaringsindhold i Aandsinderlighedens Midtpunkt, hvor
det gribes af den philosophiske Viden gjennem det Absolutes
Fornuftform og ikke længere forholder sig paa
udvortes Maade til denne Viden, men virkelig gaaer i
Et med den.

Spørgsmaalet bliver nu om det psychologiske Organ
for den religiøse Erfaring. Schleiermacher bestemte som
bekjendt Religionen som Fromhed og satte Fromhedens
Væsen i Følelsen. Men han indskrænkede saa Følelsens
Begreb i den Grad, at den afskjæres fra ethvert
Indhold, der gaaer ud over det rent Individuelle og Subjektive.
Og naar man pleier at anse Modsætningen
mellem Vel og Ve, Lyst og Ulyst som eiendommeligt
Karaktermærke for Følelseslivet, da skal ifølge Schleiermacher
denne Modsætning ikke i og for sig kunne tilskrives
den religiøse Følelse; først naar den lavere,
sandselige Fornemmelse træder i Forhold til den enten
hæmmende eller fremmende, faaer hin Modsætning forsaavidt
Anvendelse ogsaa for dens Vedkommende. Det
Ufuldkomne i disse Bestemmelser udleder Weisse af, at
Schleiermacher ikke har sammenlignet den religiøse Følelse
med den æsthetiske, hvis beslægtede Sidestykke
den er. Herved vilde han have set, at Modsætningen
mellem Lyst og Ulyst ogsaa kan have en aandig
Betydning og ikke behøver at skrive sig fra den sandselige
Natur. Den rent theoretiske Erfaring har fra først
til sidst en objektiv Karakter; den subjektive Fornemmelse
har her kun Værd, forsaavidt den gaaer over i
% --- p. 200 ---
\setcounter{footnote}{0}%
\opage{200}Iagttagelse og Fornemmelse, hvorfor Følelsen som saadan
her ikke afgiver noget Erkjendelsesstof. Paa det
æsthetiske og religiøse Omraade derimod spille Følelseslivets
Modsætninger en væsentlig Rolle, paa hvis nærmere
Bestemmelse igjen Forskjellen mellem det Æsthetiske
og det Religiøse beroer. Paa det æsthetiske Omraade
er Følelsen af Lyst og Ulyst den herskende Magt
og indeslutter Principet for de her virkende Natur- og
Aandskræfter. Følelsen er her, som æsthetisk Følelse,
i nøie Forbindelse med Phantasien, den skaberiske Indbildningskraft.
Paa det religiøse Omraade er Følelsen
af Lyst og Ulyst derimod ikke Formaalet, men fremtræder dels som Udgangspunkt, dels som Gjennemgangspunkt
for en Virken og en Bevægelse, der ikke igjen
ende i en Følelse, men gaa op i indre eller ydre Handlen
eller i en blivende Form for den aandelige Personligheds
Væsen. Den religiøse Erfaring hører altsaa ind
under den praktiske eller sædelige Erfarings Omraade,
der danner en Modsætning til det Æsthetiske paa den
ene Side, det Theoretiske paa den anden Side. Forskjellen
mellem det Æsthetiske og det Religiøse træder især
slaaende frem derved, at Følelsesmodsætningerne paa
det ene Omraade kunne faa den omvendte Betydning
af den, de have paa det andet. Paa det æsthetiske Omraade
er Lystfølelsen saaledes altid af positivt Værd,
Ulystfølelsen af negativt Værd; men religiøst set kan Lystfølelsen
blive af negativ, hindrende Betydning og omvendt
Følelsen af Ulyst virke fremmende.

Følelsen har altsaa ikke absolut Betydning i religiøs
Henseende, som Schleiermacher mente. Den er ikke
Aandslivets Inderste, men Aabenbarelse af et bag den
skjult og fra dens bevægelige Omskiften aldeles forskjelligt
Indre. Kun dette Indre, den ethiske Substants,
er den religiøse Erfarings egentlige Element; Følelsen
er kun det Kjendemærke, paa hvilket den forskende
Tanke erkjender den sædelige Empiris Stof. Dette leder
% --- p. 201 ---
\setcounter{footnote}{0}%
\opage{201}til det Spørgsmaal, hvad der i sidste Instants er det
objektive Moment i den religiøse Erfaring. Schleiermacher
kunde kun inkonsekvent antage et saadant objektivt
Moment. Men naar den religiøse Erfaring hører
ind under den almindelige sædelige Erfarings Omraade,
maa den ogsaa stilles under det Begreb, der behersker
Ethiken, det Godes Begreb. Paa dette Begreb har saavel
den religiøse som den philosophiske Bevidsthed bestandig
bygget Gudsideen. Den philosophiske Ethik forstaaer
ved et Gode et Livssamfund af fornuftige, personlige
Aander, der ved organiske Love ere indbyrdes
forbundne til et levende Fælleslegemes Enhed. Den
indre, sædelige Erfaring er da Iagttagelsen af et sædeligt
Samfunds Livsfunktioner gjennem Følelsestilstandene
hos det Subjekt, der hører til Samfundet og virker og
lider med i dettes Funktioner. En rent subjektiv religiøs
Følelse vilde umuliggjøre et saadant Samfund. Som de
ethiske Goder maa ogsaa det høieste Gode, den religiøse
Erfarings Indhold og Gjenstand, bestaa i et saadant
Livssamfund. Men medens de ethiske Samfund
ogsaa kunne blive Gjenstand for en ydre (sandselig eller
historisk) Erfaring, er den indre Erfaring paa det religiøse
Omraade den eneste mulige Forvisning om den
tilsvarende Objektivitet. Men paa dette Høidepunkt
forenes saa det Religiøse og det Spekulative, den religiøse
Forestilling om Guds Rige og det philosophiske
Begreb om det høieste Gode.

Den religiøse Erfarings subjektive Moment er Følelsen,
dens objektive Moment det Godes Ide. De Love,
som gjælde med Hensyn til Forholdet mellem disse to
Momenter, maa være de samme, som gjælde for det
aandelige Liv i det Hele. Da Aandslivet ikke indskrænker
sig til de enkelte Individer hvert for sig, men disse
for at udtale og udtrykke, udvikle og berige deres eiendommelige
aandelige Liv slutte sig sammen og danne Samfund,
træde de Ideer, der bevæge Aandslivet, ud af den
% --- p. 202 ---
\setcounter{footnote}{0}%
\opage{202}personlige Inderlighed frem paa Historiens Skueplads. Det
bliver derfor ogsaa for den religiøse Erfarings Vedkommende
den paa Kritik byggede historiske Betragtningsmaade,
der afgiver Organet for Opfattelsen af Forholdet
mellem det Subjektive og det objektive Moment. Dog
maa man med en saadan historisk Betragtning, med den
historiske Skole, der er den moderne Dannelses og Videnskabs
eiendommelige Frembringelse, ikke forvexle den
theologiske Brug af Begrebet „historisk Opfattelse“, hvor
der vel ligger en Anelse til Grund om de positive Religioners
historiske Karakter, men hvor man ikke gjør Alvor
deraf, tvertimod troer at hylde det Historiske ved at
løsrive det fra alle paa andre Omraader gjældende Erfaringslove.

Intet kan blive til religiøs Erfaring ad ydre Meddelelses
Vei, naar det ikke allerede var det i dem, fra
hvem Meddelelsen oprindelig udgik. Naar nu den religiøse
Erfaring udformes gjennem Meddelelse og Vexelvirkning
i et Samfund, finder den sit Udtryk i en bestemt,
positiv Form. De positive Former opstaa altsaa
historisk og maa forklares historisk. Det Positive opstaaer
paa et vist Punkt i den historiske Udvikling og
har sin hele Bestaaen og Virken indenfor denne. De
historiske Religioner ere forskjellige Forsøg paa at gribe
det høieste Gode, at knytte et Baand, der organisk kan
binde den menneskelige Aand til Guddommen og det
Oversandselige. Men foruden det egentlig religiøse
Moment spille ogsaa andre Momenter ind med. For det
Første øver den politisk-sociale Drift sin Virkning og
gaaer umiddelbart i Et med den religiøse Erfaring, saa
at hins Skranker ogsaa blive dennes. Derved blive
de historiske Religioner Folkereligioner. Dernæst trænger
ofte den skaberiske Indbildningskraft med sine yppige
Frembringelser det egentlig Religiøse næsten ganske
tilbage. Derved blive de historiske Religioner Mythologier,
hvilket dog ikke udelukker, at de kunne være
% --- p. 203 ---
\setcounter{footnote}{0}%
\opage{203}Vidnesbyrd om og Kilder til religiøs Erfaring. Men det
politiske og æsthetiske Dække maa afkastes, for at den
religiøse Erfaring kan fremtræde i sin Renhed.

Enhver Religion hviler paa Forudsætningen af et
Livssamfund mellem den guddommelige og den menneskelige
Aand. Derfor gjør enhver Religion Fordring paa
at være guddommelig Aabenbaring og maa erkjendes som
saadan, forsaavidt den virkelig besidder religiøs Gehalt.
Aabenbaringens Begreb bliver i denne Almindelighed Et
med den religiøse Erfaring, men da denne ikke fremtræder
rent i Folkereligionerne, indføres under Aabenbaringens
Navn ogsaa æsthetiske, spekulative og nationale
Elementer. En snevrere, mere bestemt Betydning
faaer Aabenbaringen, naar dens Begreb udskilles fra de
ikke umiddelbart religiøse Elementer. I Christendommen
fremtræder Aabenbaringens Begreb allerede hos
dens Stifter og første Tilhængere som Udtryk for en
Befrielse, en Modsætning til en tidligere bunden og formørket
Tilstand. Dette kommer af, at det egentlig religiøse
Erfaringsindhold opgaaer i sin fulde Klarhed uden
fremmede Former. Den historiske Christus fattede først
klart og udtrykte i Lære, Liv og Død den sande Ide om
det høieste Gode, Guds Rige eller Himmeriges Rige, til
hvilket Mennesket paa aandelig Vis maa lade sig føde.
Frelsen, Adgangen til dette Rige beroer kun paa de Betingelser,
der ligge i det høieste Godes eget Væsen, og
hindres ikke af Skranker som de, der i Folkereligionerne
og Mythologierne reistes dels af national Egoisme
dels af digtende Phantasi. Forsaavidt kan den Form,
den religiøse Erfaring har vundet ved den historiske
Christus, betegnes som den ved ham bevirkede Forløsning.
Til denne Christendommens Grundide slutter Religionsphilosophien sig, idet den, efter ved den historiske Kritik
at have sigtet den religiøse Erfarings Indhold, specielt
bliver en Christendommens Philosophi. Religionsphilosophien
knytter sig til Ethiken ved sit Grundbegreb om
% --- p. 204 ---
\setcounter{footnote}{0}%
\opage{204}det høieste Gode. Men først paa det Punkt i den historiske
Udvikling, hvor dette Grundbegreb træder klart
frem for den religiøse Bevidsthed, kan der knyttes en
sand og inderlig Enhed af det philosophiske Begreb og
den religiøse Erfaring. Religionsphilosophiens Karakter
beroer paa dens Forhold til den philosophiske Ethik paa
den ene Side, og til den religiøse Erfaring paa den
anden Side\footnote[1]{Weisses Afhandling „Ueber Eintheilung und Gliederung des Systemes der Philosophie“ i 27de Bind af „Zeitschr. für Phil. und phil. Kritik.“ p. 88 ff.}. Derfor bliver en Religionsphilosophi først
mulig ved Christendommen. Men ogsaa kun mulig; thi
dels forudsætter den en Bevidsthed om den philosophiske
Erkjendelses og om Erfaringens Væsen, som først den
nyere Tids philosophiske og naturvidenskabelige Udvikling
kunde frembringe, dels kunde Christendommens store
Grundide ikke i ren og uforandret Form gjennemtrænge
Verden og Folkenes Bevidsthed.

Vel var den historiske Christus, der erklærede, at
han ikke vilde give andet Tegn end Propheten Jonas's
Tegn (ɔ: selve Personlighedens og Sandhedens Under),
for sin Person hævet over alle mythologiske og supranaturalistiske
Elementer. Men i sine Omgivelser kunde
og turde han ikke afværge Driften til digterisk religiøs
Sagndannelse. Den mægtige Aandsrystelse, den nye
Tros Fødselsveer vakte paany denne Drift, som nu gjennemvirkede
hans Liv og Personlighed med et Væv af
betydningsfulde Underbilleder, hvis Karakter væsentlig
er mythologisk. Allerede for de første Troendes Bevidsthed
maatte saadanne Elementer af den almene
Frelsessubstants, som ligge i enhver ægte Religion, sammenblandes
med det historiske Begreb om Christi Person
og Liv. Derved antog Forestillingen om Frelsen en exklusiv
Form, der strider mod det Universelle i Frelsens
Væsen, ifølge hvilken denne er knyttet til den religiøse
% --- p. 205 ---
\setcounter{footnote}{0}%
\opage{205}Erfaring i det Hele og ikke blot til den specifik christelige.
Allerede de første Disciple overførte Logosideen,
dette store Resultat af den førchristelige Spekulation,
Ideen om Guds evige, med ham selv væsenslige Tanke
eller Ord, paa den historiske Christus, og ved denne
Anvendelse, der som et Lynglimt slog ned i Disciplenes
Sjæl, stiftedes det kirkelige Samfund bygget paa det
derved grundlagte kirkelige Lærebegreb\footnote[1]{Dette er et af de Punkter, hvor Weisse i sin historiske Kritik afviger fra Tübingerskolen. Ifølge denne er Ideen om Christus som Logos, som den menneskeblevne Guds Søn, Frugten af en lang Udviklingsgang, der først fuldendtes mod Slutningen af 2det Aarhundrede efter Christus, til hvilken Tid derfor Affattelsen af Johannesevangeliet, der indeholder dette Resultat, henlægges. Weisse gaaer derimod ud fra, at ogsaa de ubestridte paulinske Breve forudsætte Menneskeblivelsens Ide, hvorfra denne maa have udviklet sig hos Apostlene selv, idet de søgte at klare sig Indholdet og Betydningen af, hvad de havde oplevet. Hvad Johannesevangeliet angaaer, da anser Weisse kun den didaktiske Del (Prologen og Talerne) for Apostelen Johannes's Værk, medens Indføielsen heraf i en fortællende Fremstillings Form kun er „et lidet lykket, om gjennemgaaende Ubekjendtskab med Sagens sande Sammenhæng vidneude % PRINTED AS IS (vidneude: for vidnende; the reading, tesseract and the Minnesota copy agree)
Forsøg fra en fremmed, ubekjendt Haand.“}. Hvad der
ifølge Christendommens Grundide maa ske med Enhver,
der skal naa Frelsen, at han „fødes af Gud“, det anskuede
Kirken i Inkarnationsdogmet som ydre historisk
Kjendsgjerning, fuldbyrdet i sin Stifters Person. Der
blev ikke længer nogen Forskjel mellem den ideale og
og den historiske Christus. Paa anden Maade kunde % PRINTED AS IS (og den historiske: the image, og doubled over the line turn (og|og); O133)
hin Grundide ikke opbevares og bringes ind i Slægtens
Bevidsthed. Trangen til positivt at afstikke sit Omraade,
navnlig overfor Gnostikernes dristige Spekulationer, førte
til en bestemt Troesbekjendelse og en fastslaaet Samling
af hellige Skrifter. Paa samme Tid benyttede Kirkelærerne
sig under deres dogmatiske Undersøgelser af den
græske Philosophi, som ikke havde nogen klar Bevidsthed
om Forholdet mellem Spekulation og Erfaring, men
% --- p. 206 ---
\setcounter{footnote}{0}%
\opage{206}i sit Væsen var Dogmatisme. Alt Dette førte til en i
dogmatiske Læreformer stivnet Kirkeautoritet, indenfor
hvis snevre Grændser kun Scholastikernes utrættelige
Tankearbeide og Mystikernes begeistrede Inderlighed viste
hen til de oprindelige Kilder. End ikke Luthers naturfriske og dybsindige Aand kunde unddrage sig Dogmatismens
og Supranaturalismens Magt; ogsaa han, men
endnu mere de, som have kaldt sig efter ham, hildedes
i Bogstavtroens Baand\footnote[1]{Smlgn. Weisses dogmehistoriske Afhandling: „Die Christologie Luthers.“. (1852).}. Den derpaa følgende Adsplittelse i flere mindre Kirkesamfund maatte medføre en
større Vægt paa det Statutariske, de fastslaaede dogmatiske
Bestemmelser. Theologien blev nu kun en thetisk
Fremstilling af disse, som tilmed manglede det Storartede
i den middelalderlige Scholastik. Den religiøse
Erfarings Væsen blev mere og mere ubekjendt, jo mindre
Theologerne brøde sig om at følge med de store Fremskridt
paa den naturvidenskabelige Erfarings Omraade.
Først Schleiermacher hævdede energisk den religiøse
Erfarings Betydning, og samtidig har den historiske Kritik
skaffet det Redskab tilveie, ved hvilket hin Erfarings
Indhold videnskabelig kan sigtes og lægges til Rette for
Religionsphilosophien. Ved dennes Udvikling ville de
dogmatiske Former og Skranker falde, og den frie Erkjendelse
af den religiøse Sandhed træde istedetfor den
kirkelige Theologi. Ogsaa hvad de kirkelige Troslærdomme
angaaer, vil da det evangeliske Ord bekræfte
sig, at Sædekornet maa dø for at kunne bære Frugt.

Den historiske Kritik fører Weisse til en skarp
Sondring mellem den ideale og den historiske Christus.
--- Han finder i Christi Lære tre Hovedbegreber, udtrykte
under prægnante Former: Himmeriges Rige, ---
den himmelske Fader, --- Menneskens Søn.

Grundtanken er, som allerede berørt, at Menneskets
% --- p. 207 ---
\setcounter{footnote}{0}%
\opage{207}Frelse bestaaer i det levende Samfund med Gud og alle
de til samme høieste Maal stræbende Aander. Denne
store Tanke kan ikke være fremgaaet af den blotte Omdannelse
af en i Folketroen faststaaende, halvt overovertroisk % PRINTED AS IS (overovertroisk: over-|overtroisk, for overtroisk; the image, s03)
Forestilling, hvorimod man vel kan tænke
sig, at den i deres Gemytter, der endnu ikke formaaede
at bære den, kunde blive Kilden til en Række phantastiske
Forestillinger om hint Riges Fortid og Fremtid.
Det er den sande Apologetiks Opgave at paavise, at
Christus har været saa overlegen over sine Omgivelser
og en lang Række af følgende Generationer, at disse for
at fastholde hans Grundtanke maatte indklæde den i
mythiske Former, der snart stivnede til dogmatiske Former;
enhver anden Art af Apologetik stammer selv fra
disse Vildfarelser og arbeider paa at forevige dem.

Andre Betingelser for Adgangen til dette Rige end
de, der ligge i Sagens Natur, stilles ikke. Christus har
sammenfattet disse Betingelser i den prægnante Fordring
om at lade sig føde paa Ny, at lade sig føde af Gud,
blive Guds Barn. Herved forstaaes den dybe Betydning
af Betegnelsen af Gud som den himmelske Fader.

I ligesaa nær Forbindelse med Gjenfødelsens Ide
staaer Udtrykket Menneskens Søn. Det betegner Indbegrebet
af de ud af den naturlige Menneskehed gjenfødte
Mennesker. Det er altsaa et idealt Kollektivbegreb.
Men det blev tillige en Betegnelse af den historiske Christus
som den personlige Repræsentant for dette ideale
Begreb. At dette Begreb gaaer ud over denne individuelle
Anvendelse, ses bl. A. af de Steder, hvor Christus
adskiller sig fra „Menneskenes Søn“ ved at lade
denne fremtræde som tredie Person (især i Antydningerne
om Lidelsen og Verdensdommen). Først i de Visioner,
i hvilke Disciplene troede at se deres korsfæstede Mester
som opstanden, smeltede den ideale og den historiske
Menneskesøn saaledes sammen, at Mesterens Person nu
blev Bærer af al ideal og religiøs Gehalt og Troen paa
% --- p. 208 ---
\setcounter{footnote}{0}%
\opage{208}ham som historisk Person blev derved sat som nødvendig
Betingelse for Frelsen. Man tabte snart den Erkjendelse,
at Menneskeblivelsens Mysterium i Alle, som
gjenfødes, er væsentlig det samme,  som i den Ene, i
hvem Kirken anskuer Guddommens Fylde i legemlig
Iboen; kun Mystikerne søgte atter tilbage til den religiøse
Grundide.

Hermed fulgte tillige den Forandring, at Frelsesbevidstheden
udelukkende byggedes paa Syndsbevidstheden.
Med den exklusive Opfattelse af Frelsens Begreb
fulgte den ensidige Opfattelse deraf. Hos Christus selv
var det negative Moment, Modsætningen til „denne Verden“,
ganske underordnet og indesluttet i den store positive
Anskuelse af det nye Livs Kræfter. Allerede Disciplene
anskuede dette Negative som en positiv Grundfordærvelse.
Religionsphilosophien maa ogsaa her føre
tilbage til den sande Grundtanke. Den ser ikke den
guddommelige Naades Virken blot som en Ophævelse af
Synden, men som en fremskridende Livsproces, hvorved
det guddommelige Ideal virkeliggjøres i Menneskeheden
i det Store som i det Enkelte, en Menneskebliven af
det Guddommelige, Avling og Fødsel af en „Sønmenneskelighed“.
Det sande Indhold af den kirkelige Dogmatiks
Lære om de to Naturers Forening i Christus er
netop denne Livsproces, hvis Resultat er den med Guddommen
forenede Menneskenatur. Denne i Menneskeslægten
vordende Sønmenneskelighed er den ideale Christus,
fra hvem Frelsen udgaaer.

I den historiske Christus personificerede denne ideale
Christus sig ved en Akt, der har fundet sit mythiske
Udtryk i Fortællingerne om hans Undfangelse og Daab,
og som nærmest maa sammenlignes med den Maade,
hvorpaa hos Oldtidsfolkene de mythologiske Gudeforestillinger
og den eiendommelige Folkekarakter bleve til
ved én og samme Akt. En saadan Akt er fremdeles
% --- p. 209 ---
\setcounter{footnote}{0}%
\opage{209}ogsaa Troen, ved hvilken det enkelte Menneske i sin
hele Personlighed tilegner sig det høieste Gode. Troen
er paa det religiøse Omraade, hvad Selvbevidstheden er
paa det logiske og psychologiske Omraade. Troens
Gjenstand er selve det høieste Gode, i hvis subjektive
Tilegnelse Frelsen bestaaer, og er i sig selv uafhængig
af alle historiske og dogmatiske Forudsætninger.
Eller skulde en Antigones Tro, hun som véd sig født
„ikke til at hade med, men til at elske med“, som véd,
at „den Tid er længere, i hvilken hun skal søge at behage
Magterne i det Hinsides, end den, i hvilken hun
skal behage Magter her paa Jorden“, --- skulde denne Tro
være af ringere Værd end Horen Rachabs\footnote[1]{Hebr. XI, 31.} Tro? Skulde
en Sokrates's og en Platons spekulative, paa det Godes Ide
og Bevidstheden om dens Evighed grundede Forvisning
være af ringere Værd end en Baraks og Gideons, en
Jephtas og Samsons Tro\footnote[2]{Hebr. XI, 32.}, der er blandet med saa
mørke Elementer af vild og grusom Overtro?

Hvad Kirken har sammenblandet med den universelle,
frelsende Tro, er de særegne Bestemmelser, der
have Gyldighed indenfor et snevrere Troessamfund. Den
christne Kirke fører sin Bevidsthed om Frelsen tilbage
til den historiske Christus, der har udtalt Ideen om det
høieste Gode i en Renhed og Klarhed som ingen Anden.
Troen paa denne Christus bliver derfor her efter Sagens
Natur en Betingelse for personlig Deltagelse i det ydre
Kirkesamfund. Dog kun denne Tro overhovedet, som Henvisning
til det historiske Tilknytningspunkt for Frelsesbevidstheden;
en nærmere dogmatisk Bestemmelse vilde
krænke Troens Universalitet. Den Vei, Kirken slog ind
paa, da den søgte at give saadanne nærmere Bestemmelser,
har ført til en Flerhed af særskilte Kirker, hver
med sin dogmatiske Bekjendelse. Den sande evangeliske
% --- p. 210 ---
\setcounter{footnote}{0}%
\opage{210}Kirke vil i Modsætning til dem være det Samfund, der
udtrykker Christendommens ideale Indhold frit for al
dogmatisk Indskrænkning. Indtil den udvikler sig, er
Staten som Bærer af den almindelige humane og ethiske
Fornuftreligion en nødvendig integrerende Magt overfor
Særkirkernes dogmatiske Sneverhed,  hvor hint humane
Element er trængt tilbage og fordunklet. Som saadan
repræsenterer Staten tillige Tolerancens Princip, der ikke
kan have sit Hjem i nogen af disse Kirker. Derved
bliver Staten et Middel til den sande Kirkes Udvikling
og Selvbefrielse ud af de hæmmende Skranker. ---

Weisses Lære, som vi nu have givet et almindeligt
Overblik over, men saaledes, at vi navnlig dvælede ved
Begyndelses- og Slutningspunktet, søger i sine forskjellige
Dele at give et fremskridende Bevis for Guddommens
Realitet. Metaphysiken paaviser det almene Grundlag,
der paa indre Vis oplukker og specificerer sig i
konkrete Former, hvis Udviklingsstadier de realphilosophiske
Videnskaber forfølge, indtil den hele Verdensanskuelse
naaer sit høieste Punkt i Religionsphilosophien,
den inderligste Enhed af Spekulation og Erfaring. Som
i vor Oversigt, saaledes er det ogsaa i vor Bedømmelse
Begyndelses- og Slutningspunktet, vi især ville dvæle
ved; det er dem, der give Weisses Lære dens Eiendommelighed.


Den metaphysiske Erkjendelse gaaer ifølge denne
Lære ikke ud paa et Virkeligt, men paa al Virkeligheds
almene Grundbetingelser, den rene Tilværelsesmulighed.
Derved er Befrielsen fra Begrebsdogmatismens
Herredømme gjort mulig, samtidig med, at den
spekulative Begrebsudviklings Betydning fastholdes. Derved
bliver det fremdeles muligt at tilskrive Erfaringen
en principiel Betydning for Verdensanskuelsen uden at
henfalde til en umiddelbar Empirisme. Men dog har
Weisse selv, som vi have set, ikke undgaaet den Mislighed,
som for ham maatte ligge nær, nemlig at give
% --- p. 211 ---
\setcounter{footnote}{0}%
\opage{211}hint Indbegreb af Kategorier, „Tilværelsesmuligheden“,
en afsluttet Væren for sig i Modsætning til den konkrete
Gudside og den virkelige Verden, som det dunkle Fatum,
der i den græske Mythologi dannede Baggrunden for det
guddommelige og det menneskelige Liv. Tildels beroer
vel dette paa en Ufuldkommenhed i Fremstillingen, der
staaer i Forbindelse med Weisses hele theosophiske Retning,
som ofte fører ham til at lade Begreberne krystallisere
sig for Phantasien; og hint Fatum optages
ogsaa, som ovenfor vist, under den videre Udvikling i de
konkrete Former. Men Weisse er dog paa dette Punkt
endnu stærkt paavirket af Hegel. Det viser sig navnlig
i, at han udvikler Metaphysiken rent apriorisk og stiller
den i Spidsen for det hele System. Er den metaphysiske
Erkjendelse end i sit Princip uafhængig af Erfaringen,
saa forudsætter den dog i sin virkelige Udvikling
denne, forsaavidt den menneskelige Aand kun gjennem
den kritiske Gjennemarbeiden af Erfaringens Indhold
formaaer at blive sig de almene og nødvendige
Sandheder bevidst. Paa denne Maade faaer selve Virkelighedserkjendelsen
methaphysisk Betydning. Weisse har % PRINTED AS IS (methaphysisk: the image, O136)
forsaavidt selv gjort Begyndelsen hertil, som han i Modsætning til Kant og Hegel har hævdet Rummets og
Tidens metaphysiske Betydning; i dette Punkt er han
Forløber for Trendelenburgs skarpsindige logiske Undersøgelser.


I Religionsphilosophien har Weisse den Fortjeneste
at have ført ud over den spekulative Opfattelses ensidige
Fremhæven af det Intellektuelle, uden at han derfor
som Feuerbach opløser det Religiøse i pathologiske Illusioner
eller som Schelling fører os ud i phantastisk Transscendents.
Han knytter Religionens Væsen nøie til det
Psychologiske (Følelsen) og det Ethiske (det høieste
Gode), og kan derved paa engang hævde det specifik
Religiøse i dets Eiendommelighed og bevare dets rationelle,
humane Karakter. Det vanskelige Punkt  ligger
% --- p. 212 ---
\setcounter{footnote}{0}%
\opage{212}her i selve Forholdet mellem det subjektive (psychologiske)
og det objektive (ethiske) Moment. Weisse viser
os kun ad Analysens Vei, at disse to Momenter, den
personlige Erfaring og det objektive Ideal,  indeholdes i
det Religiøse, men det indre Forhold mellem dem, det
Baand, der knytter det psychologisk Individuelle til det
høieste, Alt omfattende Formaal, paavises ikke. Han
har her efterladt en Opgave, hvis Løsning kun kan
vindes gjennem dybtgaaende psychologiske, ethiske og
religionshistoriske Undersøgelser, men som han dog selv
paa mange Maader har givet Forarbeider til.

Weisse søger ved sin Behandling af religionsphilosophiske
Problemer bestandig Tilknytningspunkter indenfor
den tidligere Kirkelære. Denne Bestræbelse for organisk
Udvikling har --- saameget Belærende og Interessant
den og paa Grund af Weisses ualmindelige Lærdom
fremdrager --- ført ham til at overvurdere Betydningen
af de kirkelige Symboler, efterat deres dogmatiske
Gyldighed er omstyrtet. At den kirkelige Theologi
har udspillet sin Rolle, og at kun Religionsphilosophien
kan indtage den ledige Plads, derom nærer han ingen
Tvivl. Men som han i sin Opfattelse af Christendommen
i det Hele lægger Vægten paa den humane og immanente
Side i dens Væsen, paa dens ideale Indhold,
ikke paa den Transscendentsens Form, under hvilken
det fremtræder, en Form, han skriver paa den kirkelige
Mythologis og Dogmatiks Regning, saaledes mener han,
at Religionsphilosophien vil blive Troeslære i en Fremtidskirke,
hvor Christendommens Ide er kommen til sin
Ret overfor Særkirkernes Forvanskninger. Han gaaer
ud fra, at den historiske Christi Bevidsthed i sig selv
var hævet over det mythologiske og supranaturalistiske
Standpunkt, og vil ligesom Nyrationalismen føre tilbage
til Christi Religion i Betydning af den Religion, som
Christus selv havde. Men dette er i hvert Tilfælde
Noget, som kun historiske Undersøgelser kunne godt%
% --- p. 213 ---
\setcounter{footnote}{0}%
\opage{213}gjøre, og det bliver, naar det opstilles som Faktum forud
for en saadan Undersøgelse, ligesaavel et blot Dogme som
de kirkelige Lærdomme. Og hvad Muligheden af et paa
en saadan ideal Christendom bygget kirkeligt Samfund
angaaer, da er det Centraldogme, der har konstitueret
og hidtil sammenholdt Kirken, netop opstaaet ved Sammensmeltningen
af „den ideale“ og „den historiske Christus“,
og om Muligheden af et kirkeligt Samfund, hvor
med dette Dogme alle andre Dogmer skulle udelukkes
af Bekjendelsen, kan der sikkert tvivles ligesaa meget
som om den platoniske Stat. Alligevel er det maaske
ad denne Vei, at Menneskeaanden efterhaanden vil arbeide
sig ud af Autoritetstroens og den dogmatiske Intolerances
Herredømme, og Weisses Arbeider ville ogsaa
i denne Henseende bevare deres Betydning. ---

Førend vi forlade Weisse, maa vi endnu til den almindelige
Oversigt over hans Lære føie en nøiere Fremstilling
af hans Gudside\footnote[1]{„Das philosophische Problem der Gegenwart“. --- „Philosophische Dogmatik.“ Erster Band.}%
% CHECK p. 213: star lost in the layer: paired with the mark by order
, hvormed atter Hovedpunkterne
af hans Naturphilosophi\footnote[2]{„Philosophische Dogmatik. Zweiter Band.“} staa i indre Forbindelse.

Grundlaget for hans Gudside dannes, som vi have
set, af den absolute Ide, Indbegrebet af Virkelighedens
evige Betingelser og Muligheder. Denne Ide er absolut
iboende i Gud, fra Begyndelsen af nærværende for hans
Bevidsthed, ja den er denne Bevidstheds nødvendige Betingelse.
Det er ikke Ideen selv i dens mod al Vexel
og Forandring i Tiden absolut ligegyldige Ro og ubetingede
Lighed med sig selv, der er Gud; men den
Tanke, der til sit eneste Indhold har denne Ide og selv
udspringer af dens Skjød, denne Tanke er det ligesaa
uforanderlig hvilende som rastløst og ustandselig bevægede
Grundlag for den guddommelige Selvbevidsthed. I
Almindelighed anser man Selvbevidsthed og Uendelighed
% --- p. 214 ---
\setcounter{footnote}{0}%
\opage{214}for uforenelige. Men Forholdet er tvertimod dette, at
det er Uendeligheden, der først gjør Selvbevidstheden
mulig. Baade sig selv og Tingene opfatter Bevidstheden
kun i Kraft af den uendelige Tilværelsesmulighed, der
forsaavidt bestandig er dens egentlige og umiddelbare
Gjenstand (smlgn. ovenfor p. 188—192). Vel betinges
denne Bevidsthed i det endelige Jeg ved et sandseligt
Forestillingsliv; men den existerer først fra det Øieblik
af, da den uendelige Fornuftide bliver Gjenstand for
den, om end nærmest kun i almindelige og ubestemte
Omrids, som kun ved en udtrykkelig derpaa rettet Tænkning
udvikles til Bestemthed, til den udtrykkelige Viden,
om det Muliges Grændser, som udgjør Mathematikens og
Methaphysikens Indhold. I Guddommen er det derimod % PRINTED AS IS (Methaphysikens: so twice in the book (also p. 154); the image, s03)
Ideens uendelige Tankeindhold selv i dets absolute Inderlighed,
der danner Bevidstheden, og ikke som i det endelige
Jeg Tanker, der afspeile et ydre og fremmed Indhold.
Men den guddommelige Selvbevidsthed kan ikke
være uendelig, udenat den uendelige Tid og det uendelige
Rum ere den iboende. Enhver guddommelig Bevidsthedsakt
er tillige en Akt af Rums- og Tidsopfyldelse;
den guddommelige Tankes Liv er en evig Avlen og Føden,
hvorved det Almene i Urenhedens Væsen specificeres
i det Uendelige. Det nye Indhold er tilstede i den almene
Ide paa tilsvarende Maade, som efter den logiske
Tænknings Love det Særskilte og Enkelte i det almene
Begreb, der indeslutter en Mangfoldighed af Muligheder,
som ved at virkeliggjøres fremtræde dels ved Siden af
dels efter hverandre.

Det Princip, hvorefter Specifikationen af det almene
Indhold finder Sted, er Materien i Gud. Allerede Scholastikerne
have kaldet Materien principium individuationis
formarum. Dette Princip kan theosophisk kaldes Guds
Gemyt, hans skabende Phantasi, det expansive eller
centrifugale Princip i hans Væsen. Anskuelsen af et
saadant centrifugalt Princip laa til Grund for Heden%
% --- p. 215 ---
\setcounter{footnote}{0}%
\opage{215}skabets pantheistiske Rørelser, der mere eller mindre
holdtes nede af de sædelige og intellektuelle Magter, men
dog hist og her kom til Gjennembrud i den løsladte
Phantasis vilde Fostre, navnlig i Forestillingen om en
blindt avlende og ødelæggende Verdensmagt. I sin inderlige
Enhed med Ursubjektet, med den absolute Ide fremtræder
dette Princip i den christelige Lære om den evige
Logos. At Gud som Fader fra Evighed af avler sig en
væsenslig Søn, vil altsaa sige, at Gud ikke blot fører et
uendeligt, mod al Grændse og Tidsbestemmelse ligegyldigt
Tankeliv i den rene Fornuft, men ogsaa et producerende
Naturliv, saaledes som bl. A. Giordano Bruno
med saa stor Begeistring har skildret. I denne indre,
uendelige Produceren virker Ursubjektet som teleologisk
Princip, der ordner og samler Alt.

Men hvis vi fastholdt Gudsideen paa dette Trin,
vilde vi ikke komme ud over et spontant Naturliv, ledet
af et metaphysisk Princip, ikke naa den frie Virkelighed.
Til Tanken og Naturen i Gud maa der som tredie,
afsluttende Element føies Villien. Villien er væsentlig
en begrændsende og ved Begrændsning formende Magt,
ikke umiddelbart producerende. I dens Begreb ligger
det derfor, at Stof og Indhold maa være givet eller rettere
stedse paa Ny gives ved de levende Tanker, der
udspringe af samme Kilde som Villien selv. Den guddommelige
Natur er saa at sige Villiens Legeme; de
evige Naturformer blive til som den guddommelige
Villies Udtryk. Den supranaturalistiske Opfattelse feiler
ved at betragte den guddommelige Villie som naturløs;
derved ophæver den selve Villiens Begreb og kan kun
ved den absolute Tankeløsheds salto mortale bringe en
Overgang tilveie fra Gudsbegrebet til Verdensbegrebet.
Men ligeledes feiler man, naar man med Platon, Leibnitz
og Hegel umiddelbart tillægger Guds rene absolute
Væsen det Godes Prædikat, Dette Prædikat kan saa % PRINTED AS IS (Prædikat, Dette: the image, a comma for a point; O138)
ikke faa noget virkeligt Indhold, forskjelligt fra de rent
% --- p. 216 ---
\setcounter{footnote}{0}%
\opage{216}metaphysiske Attributer. Det Gode bestaaer væsentlig
i Villien til Meddelelse af et Realt, hvis Besiddelse altsaa
forudsættes hos det villende Subjekt. Det Ethiske
har det Æsthetiske, det Gode det Skjønne, Villien det
rige Gemyt og den skabende Phantasi til sin reale Forudsætning.
De ethiske Begreber ligge vel som Muligheder
i den guddommelige Forstands ideale Forudsætninger,
der som saadanne ogsaa ere Villiens Forudsætninger;
men ogsaa kun som Muligheder, der først virkeliggjøres
ikke blot ved, men i Villien. Ved den Retning, der i Villien
gives de reale, producerende Kræfter, fremkaldes først
det Indhold, som de ethiske Begreber udtrykke, nemlig
dels Villien selv, der først nu er en virkelig Villie, dels
de guddommelige Naturbestemmelser, der nu gaa op i
og udformes ved Villien. Denne gode ɔ: meddelende
Villie er det, der i den kirkelige Lære fremtræder som
den hellige Aand. Hvad Kirkelæren fremstiller mythologisk
som en Trehed af Personer, er altsaa nødvendige
Momenter i Gudsideen, der kun begrebsmæssig, for at
ses i deres Klarhed og Eiendommelighed, kunne adskilles
fra hverandre.

Guds Væsen staaer ikke afsluttet og færdigt i sig.
Da Tidsuendeligheden er det Element, i hvilket Guds
Liv evig udfolder sig, saa maa der ogsaa finde en evig
Fremgang og Udvikling Sted af hans Væsen. Alle guddommelige
Elementer ere i en uendelig Forhøielsesproces.
Ene herved bliver ogsaa en overfor Gud relativ selvstændig
Tilværelse mulig. Den guddommelige Villie er
meddelende; den sætter sig som Basis og Mulighed, for
at den evige Naturs Indhold kan gaa over i Virkelighed.
Den bliver derved til Verdensmaterien, som er Mellembestemmelsen
mellem Naturen i Gud og Naturen udenfor
Gud. Ved Naturbegrebet knyttes Verdensbevidstheden
til Gudsbevidstheden. Naar vi betegne Verdensindholdet
som Natur, betragte vi det som det, der bevæger
sig i en stadig Proces, hvor det vexelvis føder og
% --- p. 217 ---
\setcounter{footnote}{0}%
\opage{217}fødes, frembringer og frembringes. Denne Proces har
sit Princip i den evige Natur eller Materie i Gud, hvorfra
den udspringer gjennem Verdensmaterien, den som
Basis og Substants satte guddommelige Villie, som Mellembestemmelse.
Den Forandring, som herved bevirkes i
det evige Naturindhold, har sit Udtryk i Modstandskraften,
i Inertiens Lov, som derfor er Verdensmateriens
Karaktermærke\footnote[1]{Naar Weisse bekæmper den moderne Atom- og Monadelære (som vi i næste Kapitel ville lære nærmere at kjende), da grunder dette sig dels paa en lignende Misforstaaelse som den, vi allerede have paavist hos I. H. \emph{Fichte} (p. 170 f.), dels paa faktiske Feiltagelser (som naar han antager, at med Tyngden hører ogsaa Inertiens Lov op). I sin almindelige Opfattelse staaer han egentlig ikke den moderne Monadelære saa fjernt, da han ligesom den ser Materiens Princip i den spontane, individualiserende Kraft.}.

Den ydre, materielle Tilværelse, som saaledes udvikler
sig, er ikke ide- eller aandsforladt. Vel hersker
Inderligheden ikke her umiddelbart som i Guddommens
indre Liv. Det Teleologiske har nu paa ethvert Punkt
det Mechaniske til sit nødvendige Grundlag og maa arbeide
sig frem af det. Men dog er der en iboende Idealitetens
Strøm i det Materielle. Mellem disse Kræfter,
det aandig-teleologiske Princip og det mechanisk-materielle
føres der en Kamp for Tilværelsen, hvoraf Materiens
successive Betvingelse fremgaaer. Saalænge den
formende Kraft endnu umiddelbart har med Masserne at
gjøre, kan intet Bevidsthedsliv opstaa. Enden falder ikke
umiddelbart sammen med Begyndelsen. Udviklingen gaaer
ikke engang i lige Linie. Thi i de materielle Kræfter rører
sig en spontan Selvstændighed, hvori det Ondes Mulighed
ligger; denne Spontaneitet kan nemlig forfeile sit Maal,
bøie af fra Veien. Det er en inhuman og barbarisk Anskuelse,
naar den christelige Orthodoxi udleder det physiske
Onde af det moralske Onde, af Synden, og atter opfatter
denne som Menneskets bevidste Skyld. Tvertimod, som
% --- p. 218 ---
\setcounter{footnote}{0}%
\opage{218}den kontrære Modsætning stedse forudsætter den kontradiktoriske\footnote[1]{Smlgn. S. \emph{Heegaards} formelle Logik p. 13 f.},
saaledes forudsætter det moralske Onde det
physiske. Den guddommelige Naturs Salighed lader sig
ikke umiddelbart overføre paa den endelige Natur. En
saadan Overførelse maa jo bestandig forudsætte og modificeres
ved det materielle Grundlag som Mellembestemmelse.
Derved slaaer Lystfølelsen stedse over i Usalighed,
og det animalske Liv er en stadig Vexlen af Lyst
og Ulyst. I Følelsen af Ulyst og Smerte udtrykkes den
Kamp, ved hvilken det Levende ved sin Tilværelses Begyndelse
saavel som i ethvert Øieblik af sin Tilværelse
maa hævde sin Selvhed overfor de døde Stoffer. Heri
ligger det physiske Ondes metaphysiske Nødvendighed.

Dyrets Sjæl er ikke en særskilt Ting, som udenfra
føies til Materien, men er selve den materielle Substants's
oplukkede Kjerne. Dyrets Sjæleliv udspringer af den
Naturaand, der fra først af iboer Verdensmaterien og
lægger sig for Dagen i enhver enkelt Akt af den kosmogoniske
Proces. Ved de organiske Processers Kredsløb
fastholdes denne Aand i Dyrerigets Individer i en
vedvarende Rækkefølge af indre Livsfunktioner og Livstilstande.
Fra de laveste Organisationstrin til de høieste
strækker sig en fælleds Grundstamme, hvis Udløbere ere de
til forskjellige Tider uddøde eller bestaaende Dyreslægter.
Antagelsen af en saadan Grundstamme, hvortil Lamarques
og Geoffroy St. Hilaires Undersøgelser synes at føre,
staaer ligesaalidt i Strid med det rigtig opfattede Skabelsesbegreb
som f. Ex. den naturlige Sammenhæng i
Menneskeslægtens Forplantning.

Allerede hos Dyrene forekommer en sporadisk Tænken.
Men de egentlige Fornuftvæsener opstaa først,
naar Bevidstheden kan hæve sig over det Sporadiske og
Spontane og ved en fri Akt tage sig selv i Besiddelse.
Selvbevidstheden erhverves, er ikke oprindelig given.
% --- p. 219 ---
\setcounter{footnote}{0}%
\opage{219}Først gjennem Sproget vindes en saadan Fasthed i Bevidstheden,
at det Sporadiske kan trænges tilbage og Selvbevidstheden
udvikle sig. Først ved Sproget bliver derfor
Fornuften Artskarakter for Mennesket. Den selvbevidste
Fornuft er dog endnu kun det formelle Udtryk
for Menneskets Væsen. Det høieste Maal, Tilegnelsen
af det reale „Guds Billede“, naaes, som det har vist sig
i det Foregaaende, kun ved en fri Gjenfødelse, hvorved
Individet indtræder i, hvad den nyere Philosophi kalder
den absolute Aands Sphære, hvad Christendommen kalder
Guds Rige. Hermed ere vi da atter komne tilbage til
Religionsphilosophien, som vi tidligere have fremstillet. ---

Mod denne hele Lære har man indvendt\footnote[1]{\emph{Pfleiderer}: Das Wesen der Religion. 1869. p. 236.}, at den
gjør Pantheismen og Materialismen den betænkelige Indrømmelse,
at Aanden ikke skal kunne være det Oprindelige,
men maa være Resultat af det Ikke-Aandeliges Udvikling.
Dog kan denne Indvending kun fremføres mod
Weisse fra et Standpunkt, der betragter den guddomme-Aand % PRINTED AS IS (guddomme: guddomme-|Aand, for guddommelige Aand: the line ends guddomme- in both copies; the image, s02)
som afsluttet og isoleret i sig. Det er netop
Weisses Fortjeneste, og det er det Punkt, hvorved han
især udmærker sig fremfor den spekulative Theismes
øvrige Repræsentanter, at han har ført den guddommelige
Aands Begreb ud over hin dualistiske Forudsætning.
Derved er det blevet ham muligt at gaa ind paa de storartede
Anskuelser, der trænge sig frem i vor Tids Naturvidenskab,
og vise, hvorledes det Ideelle og Ethisk-Religiøse,
rigtig og universelt opfattet, ikke krænkes,
men netop kommer til sin Ret derigjennem. Som det
synes, uden at kjende Darwin selv, er han ved Benyttelse
af dennes Forløbere kommen til de Problemer, som
Darwinismen nødvendigvis maa sætte i Bevægelse. Og
hvor Meget der end i det Hele og i det Enkelte vil
kunne indvendes mod hans Behandling af disse Problemer,
saa viser den dog utvivlsomt i en rigtig Retning.

% --- p. 220 ---
\setcounter{footnote}{0}%
\chaphere{FEMTE KAPITEL.}
\opage{220}De Forskere, vi i det foregaaende Kapitel have beskjæftiget
os med, havde alle udviklet sig under det
overvældende Indtryk af den spekulative Idealisme og
maatte trods al deres Modstræben vedblive at bevare
dens Præg. Derved vedligeholdtes vel Kontinuiteten i
Tænkningens historiske Udvikling; men paa den anden
Side hindredes derved den rene og fulde Udvikling af
de nye Elementer, der trænge sig frem, og disse faa
endnu kun en underordnet Indflydelse paa den hele Anskuelse.
Med større Energi vil man kunne hengive sig
til disse nye, tidligere tilbagetrængte Elementer, naar
den historiske Kontinuitet ikke umiddelbart spiller nogen
Rolle, men Forskeren fra først af staaer friere og mere
kritisk overfor den foregaaende Tidsalders Tænkning.

Den spekulative Idealisme var en Helhedsanskuelse.
I det Absolutes Ide saa den Tilværelsens Væsen og Indhold,
og ved denne Ides Udvikling gjennem den rene,
dialektiske Tankebevægelse mente den at naa Tilværelsens
hele Fylde. Det Enkelte fik kun forsaavidt Berettigelse,
som det kunde finde sin Plads paa et eller
andet Punkt af denne dialektiske Udvikling. I Modsætning
hertil lægges der nu Vægt paa Enkeltundersøgelser;
gjennem Analysen af det Enkelte søger man at
nærme sig de Høider, hvorfra de store, omfattende Syntheser
bliver mulige, hvorimod man betragter det som
Illusion, naar Tanken ud af sig selv vil udspinde en
Erkjendelse af Virkeligheden. Derfor betones Nødvendigheden
af et inderligt Forhold mellem Philosophien og
den exakte Erfaringsvidenskab. „Videnskabens Liv,“
siger Trendelenburg allerede i Forordet til første Udgave
af hans „Logische Untersuchungen“ (1840), „bestaaer som
alt Liv i Kamp, Kamp saavel med Meninger, der stille
% --- p. 221 ---
\setcounter{footnote}{0}%
\opage{221}sig imod, som med Kjendsgjerninger, der ikke ville bøie
sig for Tanken\dots{} Kampen med Kjendsgjerninger er
i det Hele vanskeligst; thi de staa, naar de ere rigtig
opfattede, ubøielige, og Tanken maa føie sig efter dem
for at kunne underkaste sig dem.“ I denne sidste Sætning
ligger den eiendommelige Modsætning til Spekulationen
udtrykt. Spekulationen opfattede Tankens Seir
som en umiddelbar og saa ikke, at Veien til denne Seir
ofte gaaer gjennem tilsyneladende Underkastelse og Resignation.
Med Betoningen af Kjendsgjerningerne mister
det Ideelle ikke sit Værd og og sin Realitet. Tvertimod % PRINTED AS IS (og og: the image, O142)
fastholder man netop Overbevisningen om, at i det Ideelle
indesluttes alt Værd, og at overfor det alt Faktisk,
al blot Væren forholder sig tjenende. „Enhver Tings
sande og levende Substants bestaaer i Virkeligheden kun
i den Ide, som den udtrykker; det Reale er overhovedet
ikke Andet end det Ideale.“ (Lotze. Streitschriften 1857.
p. 45). Den exakte, punktuelle Opfattelse af Virkeligheden
fører af sig selv til en ideal Anskuelse, der har
det Fortrin frem for den spekulative Idealisme, at den
ikke som denne har den reale Tilværelse udenfor sig og
imod sig, men netop vindes ved at gjennemtrænge denne.
Selve det mechanisk virkende Atom bliver fra dette
Synspunkt et ideelt Moment, der er optaget i det Heles
indre Sammenhæng. Derfor betegne vi denne hele Retning
som \emph{Virkelighedsidealismen}.

Denne Retnings Fortjeneste bestaaer fra den anden
Side set i, at det Ideelle, der i Spekulationen ligesom
flagrede løst for alle Vinde, punktualiseres og bestemmes
ved at bringes i sand og inderlig Forbindelse med Virkeligheden.
Men netop i denne Forbindelse maa efter
Sagens Natur den største Vanskelighed ligge. I sidste
Forstand maa enhver Philosophi hvile i et høieste Princip,
hvoraf al Realitet udspringer ifølge en indre Udviklingslov.
Denne sidste og høieste Opgave træder her
i Skygge for Analysen og Induktionen og paa Grund
% --- p. 222 ---
\setcounter{footnote}{0}%
\opage{222}af den Overmættelse, den forudgaaende Tidsalders dristige Deduktioner havde efterladt. Dog anerkjendes 
udtrykkelig Opgavens Betydning. „Den dialektiske Udvikling af den Ide, som enhver Skabning og ethvert Phænomen er bestemt til at udtrykke i Verdens fornuftige
Hele, gjælder ogsaa for mig som den eneste sande Angivelse
af deres Begreb, og kun deraf venter jeg en udtømmende
Indsigt i alle underordnede Tilværelsesformer
og i Grunden til disses Lovmæssighed. Men jo mere
levende jeg er gjennemtrængt af denne Overbevisning,
des mindre tilbøielig er jeg til at tage den gode Villie
for Gjerningen og beundre den Munterhed, hvormed jeg
saa mangen En bestandig med godt Mod ile til denne
uløselige Opgave. Tro derfor ikke, at dette videnskabe-Ideal % PRINTED AS IS (videnskabe-Ideal: the image, videnskabe-|Ideal, for videnskabelige Ideal; recrop rc2)
skulde være ubekjendt for mig og for Mange, der
i denne Henseende ytre sig ligesom jeg; det bevæger
ogsaa vort Hjerte, men for dybt til, at vi skulde ville
hjælpe os ud over Problemernes Alvor med saa kummerlige
Svar, som de, vi fordetmeste se afgivne.“ (Lotze.
Streitschriften p. 98). I Modsætning til den dialektiske
Færdighed, med hvilken Spekulationen medierede alle
Tilværelsens Modsætninger, støde vi her ofte paa Problemer,
der erklæres for uløselige. Men det Absolute
fastholdes som videnskabeligt og praktisk Ideal, med en
Begeistring, som er mere stille, men ikke mindre dyb
end hos Spekulationens Heroer.

1. \emph{Adolf Trendelenburg} (f. 1802, d. 1872). --- Vi
begynde med en Tænker, der endnu bestemt viser hen
til historisk Kontinuitet, om end ikke specielt med den
nærmest forudgaaende Periode. Trendelenburg lægger
endog en særegen Vægt paa Philosophiens Historie, ligesom
han selv er en af de aandfuldeste Bearbeidere af
denne Videnskab\footnote[1]{„Historische Beiträge zur Philosophie“. 3 Bind. 1856---1867.}. Det er en tydsk Fordom, siger han
% --- p. 223 ---
\setcounter{footnote}{0}%
\opage{223}(Fortale til anden Udgave af „Logische Untersuchungen“
1862), at enhver Philosoph skal begynde paa egen Haand,
have sit ganske eiendommelige Princip, et Speil, der er
slebet efter en særegen Formel for deri at opfange Verden.
Philosophien vil ikke opnaa sin gamle Magt, førend
den faaer en fast Bestaaen ved at voxe paa samme
Maade som de andre Videnskaber og udvikle sig kontinuerlig,
uden at begynde paa Ny i ethvert Hoved, men
saaledes, at den historisk optager Problemerne og fører
dem videre. Principet for den philosophiske Erkjendelse
er fundet og behøver ikke først at søges: det ligger i
den af Platon og Aristoteles grundlagte organiske Verdensanskuelse,
der efter dem er bleven videre udviklet
og efterhaanden skal fuldendes gjennem dybere Undersøgelser
af Grundbegreberne og i Vexelvirkning med de
reale Videnskaber. --- Det er den antike, særlig den aristoteliske
Opfattelse, Trendelenburg vil gjennemføre med
de Omdannelser, som de forandrede Forudsætninger og
Erkjendelsens hele mellemliggende Historie gjøre nødvendige.
Ikke uden Grund har man anket over, at han,
der er saa ubønhørlig en Kritiker af Hegel, hos Aristoteles
gaaer let hen over aabenbare Vildfarelser. Men
denne Tilknytning til det Antike og den dermed beslægtede
nøie Forbindelse mellem Tanken og Anskuelsen
hos ham give hans Fremstillinger deres ædle plastiske
Form; det bliver ham let at lade Tanken bevæge sig i
Virkelighedens Verden uden at tabe sin oprindelige
Kraft. Han er, efter Søren Kierkegaards Udtryk om
ham\footnote[1]{„Uvidenskabelig Efterskrift“. p. 79.}, „en Mand, der tænker sundt og er lykkelig dannet
ved Grækerne“. Dog kan man maaske heraf ogsaa
udlede en vis Mangel paa gjennemført Tankebevægelse,
en Hvilen i Anskuelsen uden tilsvarende Begrebsgrundlag,
navnlig ved Behandlingen af de sidste, afsluttende
Problemer.
% --- p. 224 ---
\setcounter{footnote}{0}%
\opage{224}Medens de fleste philosophiske Systemer gaa ud fra
Helheden for at naa en Erkjendelse af Delene, som
indeholdes deri, vil Trendelenburg slaa ind paa den modsatte
Vei. Erkjendelsen stræber at tyde den guddommelige
Skabelses Under ved at skabe det paa Ny i
Tanken; men naar man begynder paa denne Opgave i
det Enkelte, saa fører den af sig selv videre. Med
samme Magt, hvormed Tingene ere fremgaaede af den
evige Grund, vise de igjen tilbage til Grunden. Naar
da Undersøgelsen af det Enkelte fører til faste Punkter,
naaes derved et Grundlag, hvorfra der kan søges en
fuldkomnere Forstaaelse af det Hele.

Naar Erkjendelsen har naaet sin Fuldendelse, idet
Universet fuldstændig gjenfrembringes af den erkjendende
Aand, ophæves Modsætningen mellem Philosophi og Fagvidenskab,
idet alle Videnskaber blive Lemmer paa
samme Organisme. Men dette er et Ideal. I Virkeligheden
repræsenterer hver Videnskab Erkjendelsen af en
særegen, begrændset Sphære, og Philosophien søger Helhedens
Ide i Delene, det Almene i det Særskilte. De
enkelte Videnskaber vise nemlig ud over sig selv. De
angaa hver kun én særegen Væren, og Spørgsmaalet
bliver da om Væren i det Hele og om Grunden til de
særegne Former for Væren. Fra denne Side set vise
Fagvidenskaberne hen til Metaphysiken, der beskjæftiger
sig med det Værende som saadant, den almene Grund
til de særskilte Gjenstande. Enhver Videnskab har dernæst
sin eiendommelige Methode, sin særegne Fremgangsmaade,
hvormed den omfatter sin Gjenstand og
gjennemskuer dens Grund. Men i alle Methoder ligger
den samme Tanke til Grund, og det Spørgsmaal opstaaer,
hvorledes de forskjellige Methoder udspringe af
Tankens Væsen. Fra denne Side set vise Videnskaberne
tilbage til Logiken. Grundvidenskaben, den almindelige
Videnskabslære bliver paa engang Logik og Methaphysik % PRINTED AS IS (Methaphysik: the image, O145)
(naar den betegnes som Logik, tages dette Navn altsaa
% --- p. 225 ---
\setcounter{footnote}{0}%
\opage{225}i udvidet Betydning). I denne Grundvidenskab paavises
Enheden af Tænken og Væren i Nødvendigheden.
Uden Tanken vilde der kun gives et dødt Virkeligt,
intet Nødvendigt. Men uden Væren vilde der ligesaalidt
være Nødvendighed. Tanken maa for at naa Nødvendigheden
paa ethvert Punkt lade sig bestemme og binde
ved Sagens Natur, gjøre sig til Et med Sagen og ud af
den paavise Loven. I Nødvendigheden, der gjør Videnskaben
til Videnskab, er Væren paa eiendommelig Maade
optaget i Tanken, Tanken i Væren.

Derfor blive saadanne Standpunkter at afvise, hvor
Tanken enten ganske afsondres fra Væren eller absolut
gjøres identisk med den. Det Første er Tilfældet i
den formelle Logik og hos Herbart, det Andet hos
Hegel. --- Den formelle Logik skjelner mellem Tanken
og dens Gjenstand og vil undersøge Tankens Former,
afset fra dens Indhold. Men Tankens Væsen fremtræder
først i Vexelvirkning med Gjenstandene, ligesom
Sandseorganerne kun kunne opfattes efter deres
Væsen ved at ses i Funktion. Og den formelle Logik,
der selv definerer Sandhed som Tankens Overensstemmelse
med Gjenstanden, udelukker sig selv fra
Sandheden ved at isolere Tanken fra Gjenstanden. ---
Herbarts Metaphysik er et supplerende Sidestykke til
den formelle Logik. Han betragter Indholdet som
givet ved en absolut Position, og Tanken har kun den
Opgave, støttet til den formelle Logiks Grundsætninger. % PRINTED AS IS (Grundsætninger.: the image, a point for a comma; O146)
at skaffe de Modsigelser bort, hvormed det Givne
umiddelbart er beheftet. Men det viser sig, at hvad Herbart
kalder Modsigelsen i Erfaringsbegreberne, netop kun
fremkommer ved hans skarpe Distinktion mellem Tænken
og Væren, og at han altsaa ved denne Distinktion baade
forsynder sig mod Tanken og mod Erfaringens Gjenstand.

Hegel fastholdt den absolute Enhed af Tænken og
Væren og mente i sin dialektiske Methode at have en
Selvbevægelse af den rene Tanke, der tillige var Sagens
% --- p. 226 ---
\setcounter{footnote}{0}%
\opage{226}egen Selvfrembringelse; den sig selv udfoldende Tanke
skulde udfolde Tingenes inderste Væsen. Men ved nærmere
Undersøgelse viser det sig, at den rene Tanke er
en Illusion. a) Hegel forudsætter fra først af og helt
igjennem Bevægelsen. Dette ses allerede ved de første
Kategorier i Logiken: „Væren“ og „Intet“ ere begge i
sig ubevægelige og kunne ikke ved deres Enhed frembringe
den bevægede Vorden. „Vorden“ kunde ikke
\emph{vorde} af Væren og Ikkeværen, naar Forestillingen om
Vorden ikke gik forud. Den „forudsætningsløse“ Dialektik
har altsaa idetmindste én Forudsætning, Bevægelsen.
b) I Bevægelsen indeholdes Rum og Tid, og Begrebet
har derfor bestandig Anskuelsen til sit Grundlag.
Uden den vilde den dialektiske Proces være umulig.
Ifølge Hegel er det den logiske Negation, der skal være
den fremaddrivende Magt i denne Proces. Men den logiske
Negation kan ikke selv frembringe noget Nyt; den
kan kun ophæve det Foregaaende. Kun ved Hjælp af
Anskuelsen, der sætter en real Modsætning, kan Bevægelsen
føres videre, men saa afbrydes Begrebsudviklingens
Immanents jo ved enhver Overgang. c) Hvis den
rene Tænken var Et med Sagens eget Væsen, maatte
den dialektiske Methode være en genetisk Methode, Et
med en Paavisning af Sagens egen Oprindelse. Men dette
er ikke saa; den dialektiske Tanke udvikler sig ifølge
Hegel apriorisk, uden at støtte sig til Erfaringen, som
dog ene kan føre til Indsigt i Tingenes Tilblivelse. Vi
komme her til det Dilemna: enten er den dialektiske % PRINTED AS IS (Dilemna: for Dilemma; both copies)
Erkjendelse selvstændig og apriorisk, men saa faaer
den Erfaringen udenfor sig og bliver ingen Sagerkjendelse,
eller ogsaa forudsætter den Erfaringsvidenskaben,
men saa er den ikke selvstændig og apriorisk.

Ved denne Kritik naaes som Resultat, at Grundproblemet
om Videnskabens og Nødvendighedens Væsen
bliver uløseligt, baade naar man absolut adskiller og absolut
identificerer Tænken og Væren. Vi maa altsaa
% --- p. 227 ---
\setcounter{footnote}{0}%
\opage{227}fastholde, at Tænken og Væren i deres Forskjellighed
maa staa i en indre Forbindelse, og Spørgsmaalet bliver
da om denne Forbindelses Mulighed. Forsaavidt
begyndes der med en Modsætning, en Dualisme. Men
den menneskelige Aand er endelig og begrændset og
lever af de Paavirkninger, den modtager. Naar Menneskeaanden
blot var fri, blot selvstændig, saa at den ikke
modtog Noget, men selv dannede Alt, saa vilde den
vistnok være sin egen Herre; men dette ensomme Herredømme vilde være som Fuglens paa Sneregionens øde
Marker, uden Forbindelse med den levende Verden. Heller ikke kan man begynde med det Absolute for deraf
at udlede Erkjendelsen af det Endelige og Relative. Vel
indeholder det Ubetingede den sidste Grund for alt Endeligt; men for de endelige Væsener ligger det Endelige
og Betingede nærmest, og de maa derfra hæve sig til
Erkjendelsen af det Uendelige. Saa lidet Mennesket
kan springe over sin egen Skygge, saa lidet kan han
overspringe den endelige Tænken.

Vi maa derfor søge den elementære Mellembestemmelse,
som i den endelige Erkjendelse forener Tænken
og Væren. Det maa være en Virksomhed; thi en hvilende
Egenskab vilde ikke kunne forene, men vilde forblive
afsluttet i sig selv. Denne Virksomhed maa være
den almindeligste og oprindeligste, og den maa være enkelt,
ikke kunne opløses i Elementer. En saadan Virksomhed
er \emph{Bevægelsen}. At denne ligger til Grund for
Væren, lader sig let vise. Alt, hvad der synes at hvile,
er, dybere betragtet, dog stedt i Bevægelse. Astronomien
har forgjæves søgt et fast Punkt i Universet. Ligeledes
er Tanken i sig selv Bevægelse, en Væren opfattende,
konstruktiv Bevægelse, der lægger sig for Dagen
baade i Anskuelsen og i Forstandsvirksomheden (Adskillen
og Forbinden, Slutten o. s. v.). Bevægelsen kan
ikke forklares af noget Tidligere, men maa forudsættes
i det Uendelige. Heller ikke er den sammensat.  Rum
% --- p. 228 ---
\setcounter{footnote}{0}%
\opage{228}og Tid, som Nogle (Zenon, Kant og Herbart) have adskilt
fra Bevægelsen, ere dens Elementer, som vise
den fra to Sider: Tiden er Bevægelsens indre Maal, Rummet
dens ydre Phænomen. --- Da Bevægelsen saaledes
er den almindeligste, oprindeligste og absolut enkelte
Virksomhed, kan der ikke gives nogen Definition af den,
thi en Definition maatte jo gaa ud fra et mere almindeligt
og oprindeligt Begreb.

Den menneskelige Aand besidder altsaa en oprindelig
Virksomhed, som er Modbilledet til den ydre Bevægelse,
og som betinger al Opfattelse. Bevægelsen er forud
for Erfaringen, idet den først gjør denne mulig. Herpaa
beroer Muligheden af en apriorisk Erkjendelse, og af
Aandens frie, spontane Liv, der udvikler sig med
indre Nødvendighed og udspringer af Aanden selv.
Uden en saadan Selvstændighed vilde Aanden heller ikke
kunne forholde sig receptiv. Vor Sandseanskuelse forudsætter
bestandig en geometrisk Anskuelse, efter hvilken
vi ordne Tingene perspektivisk og i Afstande. Indtrykkene
tjene som Motiver for en Konstruktion, hvorved
de forvandles til Billeder af ydre Gjenstande.

Men denne frie, konstruktive Bevægelse er ikke absolut.
Der forudsættes betandig et Substrat, en Bærer % PRINTED AS IS (betandig: for bestandig; both copies)
bag Virksomheden. Dette Substrat er Materien. Den
kan ikke udledes af Bevægelsen, saaledes som den dynamiske
Opfattelse vil. Bevægelsens Kræfter bæres af
en ubekjendt Ting, som ikke mere er Bevægelse. Der
bliver da her et Hul i Begrebsudviklingen, hvor et i Erfaringen
Givet skyder sig ind. Vel trænger Bevægelsen
ogsaa ind her og opløser det Faste, men Forestillingen
kan ikke undvære Substratet, sætter det bestandig igjen.
Saaledes komme vi i en Kredsgang: Bevægelsen betinger
alt Værende, men for at der kan være Bevægelse, maa
der være Noget, som bevæger sig. Paa den anden Side
kan den atomistiske Anskuelse heller ikke gjennemføres,
men viser tilbage til den dynamiske: Atomerne kunne
% --- p. 229 ---
\setcounter{footnote}{0}%
\opage{229}ikke tænkes absolut træge og hvilende, thi saa vilde
deres indbyrdes Paavirkning blive umulig. Atomerne maa
altsaa tænkes som Kraftpunkter, deres egentlige Væsen
bestaaer i Bevægelsen, og denne er følgelig det Oprindelige.
Ogsaa i Atomlæren kan altsaa Materien kun tænkes
formedelst Bevægelsen. --- Vi maa lade os nøie med
det almindelige Resultat, at forsaavidt Materiens Væsen
oplukker sig for Aanden, sker det kun gjennem Bevægelsen.

Ved den konstruktive Bevægelse bliver Mathematiken
som ren, apriorisk Videnskab mulig. Af Bevægelsen
i og for sig udspringe de mathematiske Kategorier
(Figur og Tal). Ved Bevægelsen trænge fremdeles Konstruktion
og Kalkule ind i de materielle Elementer,
hvorved den anvendte Mathematik opstaaer. Idet Aanden
paa denne Maade fordyber sig i Virkeligheden,
kunne vi af Bevægelsen, der er fælles for Tanken og
Væren, udlede de reale Kategorier; den ydre Væren er
blind og døv og mærker ikke sin egen Virken og Liden,
men Tanken kan som Selvbevidsthed betragte sig selv
under Bevægelsen, og ved Iagttagelsen af Bevægelsens
forskjellige Former og Frembringelser opstaa Kategorierne.

I den rene mathematiske Anskuelse frembringes en
Figur eller et Tal. Under Forudsætning af det materielle
Substrat træder dette Aarsagsforhold tydeligere frem.
Kausaliteten, den virkende Aarsag er den første reale
Kategori. Og ligesom Figuren og Tallet opstaa ved
den ideelle, konstruktive Frembringen, saaledes den materielle
Skikkelse og Størrelse ved den reale Bevægelse.
Ligesom i Sproget Substantiver dannes af Verber, saaledes
udspringer der af Bevægelsen en Ting, en afsluttet
Substants. Dog døer Bevægelsen ikke hen i sit
Produkt: Substantsen er selv Kausalitet, virker gjennem
sine Kvaliteter paa eiendommelig Maade. I den levende
Sammenhæng mellem de virkende Egenskaber lægger
Vexelvirkningens Begreb sig for Dagen, og med det er
% --- p. 230 ---
\setcounter{footnote}{0}%
\opage{230}atter Kraftens Begreb nøie forbundet: Alt hvad der i
Vexelvirkningen bidrager til Resultatet, maa betegnes
som Kraft. --- Alle disse Kategorier tjene paa den
virkende Aarsags Omraade til fast Grundlag for vore
Tanker og kunne sammenlignes med Legemets Benbygning.
Forestillingens Kategorier og Tingenes Principer
ere ikke to forskjellige Verdener, men ere Et i
deres Oprindelse. Saalænge Bevægelsen som ren Konstruktion
holdes indenfor det mathematiske Omraade,
kan man forud bestemme Virkning og Modvirkning, og
Erkjendelsen er apriorisk. Men saasnart Materien optages
som Forudsætning, maa Erfaringen afgjøre, hvorledes
Tingene virke og Bevægelserne møde hverandre.

Men naar vi nu, udrustede med de ved Betragtningen
af Bevægelsen i og for sig og Materien vundne Kategorier
(de mathematisk-physiske), der have deres typiske
Udtryk i den virkende Aarsag (causa efficiens), træde ud
i Erfaringens Verden, viser det sig, at disse Kategorier
ikke strække til for at forklare Alt, som her møder os.
Naturen frembyder Kjendsgjerninger, som det spirende
Frøkorn, Øiet, der aabner sig for Lyset, Dyrenes til
deres Element svarende Bevægelsesorganer, hvor en Flerhed
af Momenter ere saaledes forbundne, at det ene viser
hen til det andet, saa at Forklaringen af de enkelte
Dele kun kan søges i den hele Sammenhæng Ved den % PRINTED AS IS (Sammenhæng Ved: the image, no point before Ved; O150)
virkende Aarsag forklare omvendt Delene Helheden som
deres Virkning; her derimod bliver Helheden (Virkningen)
Formaal og som saadant Aarsag, Enden leder Begyndelsen,
Fremtiden Nutiden. Dette er kun muligt ved,
at Virkningen er forud tilstede som Tanke, og at denne
Tanke er Et med den virkende Aarsag. Ved den virkende
Aarsag naaer Formaalet (Tanken) sin Virkelighed;
i Formaalet (Tanken) har den virkende Aarsag sin
Sandhed.

Under Forudsætning af Formaalets Realitet i Naturen
viser der sig et nyt Forhold mellem Tænken og
% --- p. 231 ---
\setcounter{footnote}{0}%
\opage{231}Væren. Hidtil have vi set Væren virke paa Tanken og
derved blive Erkjendelsens Aarsag (causa cognoscendi).
I Formaalet virker Tanken derimod paa Væren (som
causa finalis),  idet den fordrer en virkende Aarsag til
at realisere dens Indhold. Den nødvendige Forudsætning
er da her, at Tanken selv er Et med det Værendes
Væsen; kun derved kan den have Magt over det Reale.
Ogsaa herfor ligger Muligheden i Bevægelsen. Ligesom
vi kun erkjende den ydre Bevægelse ved vor Aands egen
Bevægelse, saaledes erkjende vi ogsaa kun de i Naturen
virkeliggjorte Formaal, fordi vor egen Aand formaaer at
sætte sig Formaal. Bevægelsen var det første a priori,
Formaalet er det andet. Men Formaalet forudsætter Bevægelsen
og den materielle Væren og er derfor kun et
relativt a priori. Hvad der nøder til at opgive den virkende
Aarsags blinde Eneherredømme, er dens egen Afmagt.
Overfor Phænomener som det organiske Liv maa
der slaaes ind paa en anden Vei. Vel bliver her i det
Enkelte Feilgreb mulige, men i det Hele staaer det fast,
at Lyset ikke kan forklares af Mørket. Formaalets Begreb
gjør en Forandring mulig, hvor den virkende Aarsag ikke
slaaer til. Men Formaalet kunne vi kun begribe fra den
ideelle Side, som den Helhedstanke, der bestemmer
Delene, idet den med Konsekvents sætter Midlerne til
til sin Fuldbyrdelse. Derimod kunne vi ikke iagttage % PRINTED AS IS (til sin Fuldbyrdelse: the image, til doubled over the line turn (til|til); O152)
det Punkt, hvor Tanken griber Kraften og leder den
mod Maalet. Ligesom de Gamle fremstillede Helios
staaende paa sin Vogn og styrende Solhestene med sin
Haand, men ikke gave Haanden Tøiler, disse den menneskelige
Krafts Redskaber, saaledes styrer Formaalets
Tanke med usynlige Tøiler de virkende Kræfter.

Det indre Formaal er det egentlig individualiserende
Princip i Verden. Naar Kraft og Formaal falde sammen
i det samme Subjekt, opstaaer Selvet, og først Selvet
er et Individ i egentlig Forstand. I den blotte Maskine
er der ikke en saadan Enhed. I Planten, Dyret og
% --- p. 232 ---
\setcounter{footnote}{0}%
\opage{232}Mennesket viser denne Enhed sig paa forskjellige Trin.
Sjælen er en sig virkeliggjørende Formaalstanke. De
sjælelige Virksomheder ere reflexive; de udgaa fra
Selvet og have Selvet til deres Øiemed. I Følelsen
af Lyst og Ulyst fremtræder det individuelle Liv som
fremmet eller hæmmet; Væsenets Virken er ikke længer
fremmed for det selv. I Dyret er Formaalstanken endnu
skjult for sig selv; den blinde Drift retter sig mod det
til Grund liggende Formaal, og Lyst eller Ulyst følge
af dens Tilfredsstillelse eller Skuffelse. Men Mennesket
bliver sig Formaalet bevidst og hæver sig i Erkjendelse
af det Almene over den umiddelbare Drift og Fornemmelse.
I det Menneskelige gaaer det Organiske over i
det Ethiske. Den ethiske Ide er det indre Formaal, der
af Tanken gjøres til Opgave for Handling og Liv.
Formaalet, som i Naturen er rent objektivt, bliver personligt
i Mennesket, og derpaa beroer Muligheden af, at
Fornuftens Herredømme kan voxe i Verden. Den ethiske
Verden har sin Historie; den virker organisk og udvikler
sig selv som en Organisme.

Formaalets Begreb fører os saaledes ind i to nye
Sphærer, den organiske og den ethiske. Men derfor
opgives ikke de tidligere vundne Kategorier (de mathematiske
og physiske). Disse vedblive at ligge til
Grund, men gjennemtrænges af Formaalet og samles
under et nyt Centrum, hvorved de vinde en nærmere
Bestemmelse. Saaledes bliver den virkende Aarsag nu
Middel, Substantsen Maskine eller, med en dybere Enhed
af Stof og Form, Organisme. Vexelvirkningen er ikke
længere et Spil af blinde Kræfter, men et indre Forhold
mellem det Hele og Delene. Materien er ikke længere
den døde Substants, der ligger bag ved Alt, med fordres % PRINTED AS IS (med fordres: the image, med for men; O155)
med Nødvendighed af Tanken, der behøver den til sin
Tjenerinde for ikke selv at blive ensom og uden Liv.
Uden den sondrende og bærende Materie vilde der ikke
gives noget Hold for Tanken, ikke noget individuelt Liv.
% --- p. 233 ---
\setcounter{footnote}{0}%
\opage{233}Paa sit højeste Punkt fører Formaalets Begreb os til
det Sandes, det Skjønnes og det Godes Ideer. Den virkende
Aarsag giver os kun det Virkelige; Sandheden
aabenbarer sig, hvor den reale Helhed svarer til sit Formaal.
I Skjønheden fremtræder det indre Formaal tillige
i en harmonisk Form for Anskuelsen. I det Ethiske
optages Formaalet i Erkjendelse og Sindelag.

Det er altsaa de samme Grundbegreber, der frit
konstrueres i det Mathematiske, udfyldes i det Physiske,
fordybe sig i det Organiske og forhøies i det Ethiske.
Saaledes saa vi f. Ex. Begrebet om Helheden først i
abstrakt Form som Figur eller Tal, dernæst som materiel
og konkret Substants i Naturen, som Organisme i
det Levende og endelig som Personlighed og sædelig
Organisme i det Ethiske. Formaalet optager tilsidst
som et høiere Forsyn alle andre Kategorier i sig; for
dets Skyld ere de alle. Paa det høieste Standpunkt
vilde man ud fra det sidste Formaal ved at gaa tilbage
kunne komme til alle Kategorierne. Men vi, der ikke staa
paa det Absolutes Standpunkt, maa gaa ud fra Bevægelsen
(den virkende Aarsag) og igjennem den arbeide
os frem til Erkjendelsen af det i sig Første og Oprindelige,
Formaalet.

De Grundbegreber, der hidtil ere udviklede, angaa
alle Sagens indre Natur. Men under Tankens Arbeide
paa Erkjendelsens Værk opstaa nye Grundbegreber,
der have deres Rod i den tænkende Erkjendelse og
angaa Sagens Forhold til den. Det er de saakaldte
modale Kategorier, Mulighed, Virkelighed og Nødvendighed.
Overfor Sandseanskuelsen staa Tingene som et
umiddelbart Virkeligt, som Phænomen. Den begribende
Forstand vender sig nu til Phænomenet og spørger om
dets Grund. Grunden er ligesaalidt som Aarsagen en
absolut Enhed, men et Indbegreb af sammenhørende Betingelser.
Er nu kun en eller nogle af disse Betingelser
tilstede, opstaaer Mulighedens Begreb, idet Tanken sup%
% --- p. 234 ---
\setcounter{footnote}{0}%
\opage{234}plerer det Manglende. Fra Muligheden stræber Tanken
videre for at finde den hele Grund, det hele System af
Betingelser. Herved sker Overgangen fra Mulighed til
Nødvendighed. Den indre, fuldt udviklede Mulighed gaaer
i Retning af Væren over i Virkelighed og fremstiller sig
for Tanken som bindende Nødvendighed. I sin umiddelbare
Form er Nødvendigheden vel kun det Uundgaaelige,
men ved Tankens Anerkjendelse bliver dette Negative
til en positiv Begrundelse. Vel er Tanken ikke
noget Led i Tildragelsernes Sammenhæng, men Begrebet
om deres Nødvendighed opstaaer først, naar de maales
af Tanken. Indtil da ere de kun Begivenheder; nu blive
de Nødvendighed. Den Anerkjendelse, Aanden her yder,
er ikke Svaghed, som om den hellere vilde noget Andet,
men maatte ligge under; men dens Væsen og dens
Styrke, idet den paatrykker den blinde Tilværelse Nødvendighedens
høiere Præg.

Spørgsmaalet bliver nu, i hvilke eiendommelige Former
Tanken løser den Opgave, hvis Mulighed saaledes
er paavist. Tankens Former og Forbindelser maa svare
til Værens Former, da Erkjendelsens Mulighed beroer
paa Bevægelsen, som er en for Tænken og Væren fælles
Virksomhed. I Formernes Parallelisme speiler den oprindelige
Enhed sig. --- Ligesom Substantsen fremgaaer
som Virksomhedens Resultat og atter er virksom i sine
Egenskaber, saaledes er Begrebet Resultatet af en oprindelig
Dømmen og udvikler sit Væsen gjennem fortsat
Dømmen. Begrebet er Substantsen opfattet som almen.
Som Opfattelse af Substantsen har Begrebet sit bestemte
Indhold, som alment sit bestemte Omfang. Indhold og
Omfang forholde sig til hinanden som Loven til Phænomenerne.
--- Til disse Begrebets Momenter slutte Dommens
Hovedarter sig: den kategoriske (og hypothetiske)
Dom er væsentlig en Indholdets Dom og gaaer ud paa
en Almindeliggjørelse af Subjektet; den disjunktive
Dom er en Omfangets Dom og gaaer ud paa en Specifi%
% --- p. 235 ---
\setcounter{footnote}{0}%
\opage{235}kation af Subjektet. --- Ogsaa Slutningens Væsen betinges
ved Begrebets to Momenter. Den er Tankens Bevægelse
enten fra Loven (Begrebets Indhold) til Phænomenet
(Begrebets Omfang) --- den kategoriske Slutning;
eller omvendt --- den disjunktive Slutning. --- Ved Hjælp
af Slutningen skrider Erkjendelsen synthetisk og deducerende
frem. Men den analytiske Undersøgelse af Erfaringen
er den bestandige Forudsætning, og ikke engang
i Mathematiken kan Erkjendelsen stedse skride frem ad
Deduktionens Vei; den behøver apagogiske Beviser og
Hjælpelinier. Istedetfor det genetiske Bevis, der paaviser
Nødvendigheden i Sagens egen Tilblivelse, træder
det indirekte Bevis, der afgrændser den ved at paavise
det Modsattes Umulighed. Kun i Forbindelse med en
disjunktiv Dom, der opstiller alle Tilfælde, som overhovedet
ere mulige, kan det indirekte Bevis føre til et
positivt Resultat. Navnlig kunne de høieste Principer
ikke bevises direkte eller genetisk; i saa Tilfælde vare
de jo ikke Principer, men havde en fremmed Oprindelse.
De kunne kun godtgjøres ved en Erkjendelsesgrund:
Principernes Nødvendighed bestaaer i det Modsattes Umulighed.
Men i dette negative Udtryk ligger et positivt
Princip skjult, og jo mere omfattende Synspunktet er,
des mere tvingende Kraft har det. For det ubetingede
Princip, for Gud maa det indirekte Bevis hentes ikke
fra noget Enkelt, men fra den hele Tilværelse.

Men det er kun den endelige Erkjendelse, hvis Mulighed
og hvis Former vi nu have undersøgt. Naar vi
efterat have gjennemløbet Endelighedsformerne søge det
Absolute, strække de tidligere Midler ikke længere til.
Man kunde mene, at ligesom de mathematisk-physiske
Kategorier undergik en Omdannelse paa det organiske
og ethiske Omraade, saaledes kunde Endelighedens Kategorier
ved en lignende Metamorphose udvikles til at gjælde
for det Absolute. Men en saadan Uendeliggjørelse af
Kategorierne er umulig, og det er en Illusion, naar man
% --- p. 236 ---
\setcounter{footnote}{0}%
\opage{236}ad denne Vei troer at kunne naa en konstruktiv Erkjendelse
af Guds Væsen. Indholdet bliver et helt andet,
naar det, der er Endeligheden eiendommeligt, skal negeres.
Logisk set mangler enhver Mellembestemmelse,
enhver Analogi mellem det Betingede og Ubetingede.
Tanken kan i sin Sphære ligesaalidt, som Øiet i sin, opfatte
det rene og klare Lys; dens Gjerning er at skue
de dæmpede og i Farvespil brudte Straaler --- og
 naar
den gjør det, negter den jo ikke, hvorfra de stamme.
Om Gud gjælder det: nesciendo scitur. Kun den theosophiske
Phantasi sætter sig ud over disse kritiske Betænkeligheder.


Men skjøndt man ikke kan konstruere det Ubetingede,
saa maa dog enhver Anskuelse, der ikke som Kants
betragter Kategorierne som rent subjektive Former, i
det Endelige finde væsentlige Fingerpeg med Hensyn
til, hvorledes Forestillingen om det Uendelige maa dannes.
Hvad der i det Endelige er sandt ved sin Nødvendighed,
kan ikke være usandt i det Uendelige. Det
Nødvendige i den betingede Erkjendelse er det faste
Punkt, hvortil Forudsætningen om det Ubetingede kan
knytte sig. Den Nødvendighed, som den endelige Aand
har indset, kan maaske, set fra det Uendeliges Synspunkt,
indtræde i en større Sammenhæng og selv forstørres,
men forbliver dog i sig selv, hvad den er. Den
guddommelige Aand er som en Digter, af hvem vi kun
kjende et enkelt Værk. Vor Verden er vel kun et Brudstykke,
men dog et Hele for sig, ligesom et enkelt Drama
i en Trilogi. Den Maade, hvorpaa det Absolute opfattes
beroer paa den hele Verdensanskuelse. Her staaer navnlig
den organiske Verdensanskuelse overfor den blot physiske
eller mechaniske\footnote[1]{Smlgn. Afhandlingen „Ueber den letzten Unterschied philosophischer Systeme“ i „Historische Beiträge zur Philosophie“. 2ter Band.}. Denne ser Verden under Syns
% --- p. 237 ---
\setcounter{footnote}{0}%
\opage{237}punktet af den virkende Aarsag: Bevægelse og Materie
ere det Første og det Sidste, og Formaalet er kun Skin.
Den organiske Verdensanskuelse ser derimod Verden
under Formaalets Synspunkt og Tanken som den sidste
Grund til Alt. Afgjørelsen kan ifølge det Foregaaende
ikke være tvivlsom; thi det har vist sig, at den virkende
Aarsag var utilstrækkelig til at forklare Erfaringens
Kjendsgjerninger, og at den opgik som Element i det
omfattende Formaal. Den organiske Verdensanskuelse
naaer sin Fuldendelse i Ideen. Begrebet bliver Ide,
naar det ses som bestemt ved det høiere Formaal eller
viser sig i det Ubetingedes Lys. Ideen er indesluttet
i Guds Væsen, thi uden Gud er der intet Guddommeligt.


--- Trendelensburgs „Logische Untersuchungen“, hvis % PRINTED AS IS (Trendelensburgs: for Trendelenburgs; the crop reader and the Minnesota layer)
Grundtanker vi nu have søgt at fremstille, er ubestridelig
det vigtigste logisk-metaphysiske Værk siden Hegels
Logik. Men hvilken Modsætning er der ikke mellem
disse to Værker! Hegel gaaer ad den rene Tankes Vei,
siger sig løs fra Anskuelsen og Erfaringen og mener
paa denne Maade, ved Begrebernes rhythmiske Selvbevægelse,
at naa den høieste Sandhed og Virkelighed. Hos
Trendelenburg er det netop Tankens første Gjerning at
iføre sig Anskuelsen som sit Klædebon for saaledes at
foretage sig den møisommelige Vandring gjennem Virkeligheden.
Hegels rene Begreber ere efter hans Mening
kun svækkede og fortyndede Anskuelser. Den menneskelige
Tænkning lever af Anskuelsen og døer Hungersdøden,
naar den skal nære sig af sine egne Indvolde.
Philosophien forudsætter derfor Erfaringsvidenskaberne,
og dens eiendommelige Problemer opstaa, naar den erkjendende
Aand søger at blive sig sin Virksomhed indenfor
disse bevidst og derfor opkaster det Spørgsmaal:
hvorledes er Videnskab mulig? hvorpaa grunder Nødvendigheden
sig?

Trendelenburgs logiske Undersøgelser betegne den
% --- p. 238 ---
\setcounter{footnote}{0}%
\opage{238}første ædruelige Selvbesindelse efter Spekulationens Begeistringsrus,
i hvilken Tankens Vinger bleve saa lette,
at den troede at kunne svinge sig op i den rene Æther.
Allerede Kant havde lagt Grunden til Læren om
den rene Tanke ved sin skarpe Distinktion mellem Rum
og Tid som Anskuelsesformer og de rene Forstandsbegreber.
Hegel udviklede saa i sin Logik Ideens Væsen i den
rene Tankes Form og stødte først i Naturphilosophien
paa Rum, Tid og Bevægelse. Den spekulative Idealisme
blev derved egentlig, som Weisse paaviste, en fuldstændig
Dualisme. Paa denne Ensidighed søgte Weisse og
I. H. Fichte at raade Bod, men de vare endnu selv for
stærkt paavirkede af den spekulative Dialektik til fuldstændig
at gjennembryde dens Skranker. Den metaphysiske
Erkjendelse blev endnu ikke sat i nøie Forbindelse
med en analytisk Kritik af det Givne. Paa dette Punkt
ligger Trendelenburgs Fortjeneste.

I Modsætning til Spekulationen erklærer han kun
at ville beskjæftige sig med det Endelige. Han negter
ikke det Absolute, men anser kun en indirekte Begrundelse
af det for mulig: Philosophien fører til Erkjendelsen af
Ideen; Ideen er guddommelig; uden Gud er der intet
Guddommeligt. Gjennem denne indirekte og ubestemte
Opfattelse af det Absolute er der hos Trendelenburg en
Tilknytning til den religiøse Tradition, „den Overlevering
fra Slægt til Slægt, som Gudsbegrebet skylder sit egentlige
Liv“. Kun overfor Endeligheden, i hvilken hans
Tanke bevæger sig med græsk Livsglæde og Harmoni,
føler han Trang til bestemt og afsluttet Opfattelse; det
Absolute er ham den utilgjængelige, skjøndt Alt bærende
Uendelighed. Dog træde i hans egen Fremstilling to
Synspunkter her i Modsætning til hinanden. Paa den
ene Side hævder han, at Forskjellen mellem det Endelige
og det Uendelige er for stor til, at de Kategorier, der
gjælde for det Endelige skulde kunne undergaa en saadan
Forandring, at de kunne blive Udtryk for det Uendeliges
% --- p. 239 ---
\setcounter{footnote}{0}%
\opage{239}Væsen. Men ligesaa bestemt hævder han paa den anden
Side, at Nødvendighedens Væsen maa være det samme,
hvad enten det ses i Endelighedens Verden eller indtræder
i det Absolutes uendelige Sammenhæng. Naar
nu Nødvendighedens Begreb danner den faste Kjerne i
den menneskelige Erkjendelse, saa have alle andre Kategorier
kun Betydning, forsaavidt de staa i indre Forbindelse
med den og ere forskjellige Former for den. Og kan
Nødvendighedens Begreb ikke tabe sin Betydning ved at
overføres fra det Endelige til det Uendelige, saa maa
ogsaa de særskilte Former, under hvilke den fremtræder,
i sig have mere end blot endelig Betydning. Saaledes
korrigerer Trendelenburg selv i Principet den Modsætning,
han har sat mellem det Endelige og det Uendelige.

Ikke blot overfor det Absolute, men ogsaa overfor
Kjendsgjerningens Verden er der hos Trendelenburg en
Standsning i Begrebsudviklingen. Han begynder med
den mathematiske Konstruktion, men maa for at komme
ud over denne Sphære empirisk optage Materiens Begreb.
Efterat nu de physiske Kategorier ere udviklede,
bliver en ny Henvendelse til Erfaringen nødvendig for
at begrunde Formaalets Begreb. Trendelenburg er her
aabenbart paavirket af Kants kritiske Sondring mellem
det Aprioriske og det Aposterioriske, Tænkningen og
Tingen i sig selv. Han indrømmer selv, at Materien,
forsaavidt vi kunne opfatte og kjende den, bestemmes
ved og udtrykker sig i Bevægelse. Med hvad Grund
tillader han da Forestillingen bestandig paa Ny at sætte
et dødt Substrat bag det levende Begrebs Indhold? Anskuelse
og Begreb, der efter hans egen Lære skulle være
i Harmoni, ere her i uforsonet Splid. Paa lignende
Maade gaaer det med Formaalets Begreb. Det optræder
i Modsætning til de foregaaende Kategorier, en Modsætning,
som hverken Philosophien eller den nyere Erfaringsvidenskab
kan tilstede. En virkelig Gjennemførelse
af Begrebsudviklingen vilde her have tilfredsstillet begge,
% --- p. 240 ---
\setcounter{footnote}{0}%
\opage{240}idet Formaalets Begreb paa indre Vis bragtes i Forbindelse
med de tidligere Kategorier i Kraft af Trendelenburgs
egen Grundtanke, at det er Helhedens Ide,
der betinger de enkelte Dele. Kapitlet om Formaalet
er fra Fremstillingens Side et af de smukkeste i Trendelenburgs „logiske Undersøgelser“, men han har her i
flere Henseender ikke hævet sig over den populære
Opfattelse. ---

Der staaer endnu tilbage at betragte Trendelenburgs
praktiske Philosophi, som han navnlig har udviklet i sin
„Naturrecht auf dem Grunde der Ethik“ (1860). Til
Grund ligger den ved de „logiske Undersøgelser“ vundne
organiske Verdensanskuelse, hvor Tanken ses som den
sidste Aarsag, ud fra hvis Centrum Kraften bestemmes.
Den i en oprindelig Tanke grundede Helhed fuldender
sig i sig selv og i sine Dele. Til dette Princip fører
med Nødvendighed det Ethiske tilbage. De Anskuelser,
der finde Udgangspunktet for det Sædelige i Lysten,
Selvkjærligheden eller Selvopholdelsen, gaa ud fra de
isolerede Enkelte og kunne kun naa en mechanisk Sammenhæng
mellem disse. Selve det Enkelte kommer først
til sin Ret ved at knyttes til Idealet, ved at ses som
det, der skal udvikle sig efter et Formaal. Herved vises
vi hen til et andet Udgangspunkt, nemlig Formaalet, det
Almene. Dog vilde dette være et tomt Ideal, naar
det ikke i indre Artikulation fæster Rod i det Enkelte.
Begge Udgangspunkter komme til deres Ret, naar vi
bestemme det ethiske Princip som Menneskets Ide. Det
eiendommelig Menneskelige beroer, som allerede de
„logiske Undersøgelser“ paaviste, paa, at Tanken, det
indre Formaal, der i de lavere organiske Væsener er
skjult for sig selv eller kun fornemmes som iblinde, i
Mennesket kommer til Selvbevidsthed. Idet Mennesket
erkjender og vil sit Væsens skaberiske Tanke, bliver
det Ethiske det frigjorte Organiske. Men det enkelte
Menneske for sig vilde ikke kunne hæve sig til dette
% --- p. 241 ---
\setcounter{footnote}{0}%
\opage{241}Punkt; først i Samfundet lægger Menneskets Væsen sig
for Dagen i sit fulde Omfang, først der bliver ogsaa de
Enkeltes Selvbefrielse mulig. Samfundet udvikler sig atter
gjennem Historien, hvis indre bevægende Kraft er Menneskets
Ide. I denne store historiske Sammenhæng, hvori
det enkelte Menneske saaledes staaer, dreier det sig fra
hans Side om Bestyrkelse og Befæstelse, om at udvide
sine Kræfters Raaderum. Hertil kommer fra Helhedens
Side Trangen til Artikulation. Begge Retninger maa
falde sammen, for at det Ethiske skal være tilstede;
Ideens Virkeliggjørelse og Individets Selvopholdelse gaa
i det Ethiske Haand i Haand.

I den personlige Hengivelse til den ideale Fornuft
med Opgivelse af den egoistiske Drift viser det Ethiskes
inderligste Element sig, det Sindelag, ved hvilket Villien
bliver den gode Villie. Gjennem Sindelagets Inderlighed
har det Ethiske sin Grund i det Religiøse. Naar Menneskets
Ide sættes som det Bevægende i den Enkelte og
i Historien, saa viser Ideens Begreb, den skabende Tanke
tilbage til Gud. Det Menneskelige staaer ikke i Modsætning
til det Guddommelige. Den positive Religion
griber den guddommelige Ide i visse historiske Fakta,
men Philosophien har en universellere Karakter og hævder
det Almenmenneskelige, af hvilket det Christelige
kun kan være en eiendommelig bestemt Retning, som
ikke kan fornegte det almene Grundlag. Religionens
Element er Sindelaget; dens Væsen ligger paa det Punkt,
hvor Forestillingen om det Guddommelige mødes med
den menneskelige Affekt. Den har sit Tilknytningspunkt
i Menneskets Frygten og Haaben, forsaavidt disse
bestemmes ved Forestillingen om et overmenneskeligt
Væsen. Medens det i den sande Tro er guddommelige
Tanker, der fremkalde Affekten, hersker i Overtroen ---
og Mennesket er af Naturen overtroisk --- den tilfældige
Drift gjennem Ideassociationerne. I Aandsreligionen befries
Mennesket fra Overtroen. Christendommen er for%
% --- p. 242 ---
\setcounter{footnote}{0}%
\opage{242}saavidt Verdenshistoriens Midtpunkt, som i den Ideen
om det Guddommeliges og det Menneskeliges Enhed fremtræder
for Bevidstheden. Den forudsætter den organiske
Verdensanskuelse (ἐν ἀρχῆ ἦν ὁ λόγος), og Menneskets Ide
har fundet sit konkrete Udtryk i Forestillingen om det
gjennem Slægterne voxende Christi Legeme og dets Lemmer.
Men Religionen fremhæver kun Sindelaget, Inderligheden. Philosophien er ogsaa her universellere; den
retter tillige sin Opmærksomhed mod det Godes Gjennemførelse i Verden. Her kommer det da ikke blot an
paa det gode Sindelag, men ogsaa paa Handlingens Sandhed,
at den nemlig svarer til Menneskets ideale Formaal,
og paa Udførelsens Harmoni og Skjønhed. Det
Godes Ide i videre Forstand indeholder disse tre Elementer:
det Gode i Sindelaget, det Sande i Begrebet,
det Skjønne i Fremstillingen.

Det høieste Gode er den virkeliggjorte Ide som sædelig
Organisme, hvoraf de enkelte Goder, Pligter og Dyder
udspringe. Heri ligger ogsaa Rettens Kilde. Med Urette
har man opstillet en skarp Modsætning mellem det Moralske
og det Juridiske. Af samme Aand, af hvilken
Pligten opstaaer, udspringer ogsaa Retten, som med Helhedens
Magt bevarer de ydre Betingelser for det Ethiskes
Virkeliggjørelse. Retten skal ikke blot værge og
bevare; som i det organiske Liv, saaledes er ogsaa her
Opholdelsen en videre Udvikling, hvis Mulighed altsaa
maa ligge i Rettens Væsen. Pligten er det sædelig
Nødvendige, Retten det sædelig Mulige; Pligten er derfor
af snevrere, Retten af videre Omfang. Men deraf
følger ikke, at Retten, den mere omfattende Mulighed,
ikke selv skulde være sædelig betinget og bære et sædeligt
Formaal i sig. Dette Formaal er Bevarelsen af Betingelserne
for den ethiske Ides Liv og Udvikling i Samfundet.
Heraf forklares den med Retten forbundne Tvang;
den grunder sig paa Samfundets Magt som Helhed. Men
Retten naaer ikke sit Maal, naar de Enkeltes Tænke-
% --- p. 243 ---
\setcounter{footnote}{0}%
\opage{243}og Handlemaade ikke kommer den imøde, og der ikke
i dem hersker den samme sædelige Aand, som Retten
skal bevare mod al Modstand og Hindring. Retten bortskjærer
Udvæxter, men kan ikke selv helbrede en indre
Sygdom. Saaledes blev Romerretten staaende, men Romerfolket
forfaldt.

Den Enkelte har ingen Ret udenfor Samfundet, og
indenfor Samfundet bliver Pligten til sammen med Retten.
Den moderne Lære om medfødte Rettigheder er derfor
ensidig, navnlig naar den bestemtere formuleres saaledes,
at den Enkelte overfor Samfundet gjør Fordring paa
Livsunderhold. De Gamle betragtede, fulde af Pietet,
som en Velgjerning, hvad de Moderne opstille som en
Retsfordring. De Gamle stillede sig ikke isolerede overfor
Staten. Derimod taler man i nyere Tid om en Modsætning
mellem Stat og Kirke, Stat og Videnskab o. s. v.
som om Staten var en særegen Faktor og ikke netop
den Magt, der omfatter alle Menneskelivets Elementer.
Staten er Virkeliggiørelsen af det universelle Menneske, % PRINTED AS IS (Virkeliggiørelsen: the image, O161)
af Menneskets Ide i et Folks individuelle Form. Ethvert
Folk arbeider paa at virkeliggjøre Statens Ide i
Forhold til sine aandelige og materielle Kræfter. Medens
den menneskelige Virksomhed i den Enkelte kun
er et ringe Brudstykke, fremtræder den i Staten som et
mangesidigt Hele. Hvad i Mennesket den ledende Fornuft
er, det er i Staten Regeringen; til Menneskets enkelte
Kundskaber svare i Staten Videnskabernes forgrenede
System, til hans enkelte Færdigheder Kunster
og Haandværker o. s. v. ---

Trendelenburgs ethiske Anskuelse minder ved Læren
om den Enkeltes Opgaaen i Staten om den antike
Opfattelse. Dog er denne ikke blot hos ham ført tilbage
til en dybere og alsidigere Begrundelse, men navnlig
i to væsentlige Punkter stiller han sig i Modsætning
til den. --- I den antike Stat blev det humane Liv kun
Faa tildel paa den større Masses Bekostning, som  maatte
% --- p. 244 ---
\setcounter{footnote}{0}%
\opage{244}arbeide for hines Fornødenheder. Denne Mislighed har
Trendelenburg ikke blot anerkjendt, men han gjenfinder den
til alle Tider. Hvad Slaverne vare i Oldtiden, vare de
Livegne i Middelalderen og ere Fabrikarbeiderne i den
nyere Tid. Deres Vilkaar staa i skjærende Modsætning
til den voxende Kultur, der stedse skaffer større Mulighed
for menneskelig Tilfredsstillelse. Dog frembringer
Kulturudviklingen selv sin Kræftskade. Den gjør
Fornødenhederne mangfoldigere, udvider og forøger Trangen
og Begjæret, forøger Folkemængden og dermed de Begjærendes
Antal. De Begjærende begjære Mere og blive
flere. Det er lettere at gjendrive de kommunistiske og
socialistiske Theorier end at helbrede de Onder, der give
dem deres Magt. Thi denne Magt hidrører fra Lidenskaberne,
som ere i Bevægelse, ikke fra det Element
af Sandhed, de indeholde. Lovene og Retten formaa
her kun lidet; ved direkte at modsætte sig Ondet træffe
de kun Virkningen, ikke Aarsagen. Den Opgave staaer
bestandig tilbage, aandelig og sædelig at beseire den
physiske Kraft, som bryder frem ved Overbefolkningen,
ved at Menneskekræfterne tabe i Værd (vor Tids Fatum),
og som truer med at opløse Samfundet i en Kamp af
Alle mod Alle. --- Den antike Sands for Harmoni og
for den plastisk udformede Helhed har altsaa ikke blændet
Trendelenburgs Blik overfor Tilværelsens Disharmonier
eller ført ham over i en haardnakket Optimisme.

Den antike Stat lod Individet saaledes gaa op i
Staten, at det ikke beholdt nogen personlig Inderlighed
tilbage, ingen indre Uendelighed; Staten beherskede
umiddelbart Alt og Alle som Ideens adækvate Form.
Trendelenburg lægger derimod, som vi have set, Vægt
paa Sindelaget, paa det indre Villiesforhold til Ideen og
finder heri en Tilknytning af det Ethiske til det Religiøse.
„Mennesket tænker i det Guddommelige et Begreb,
der har sit Værd i sig selv. Hvor han derfor,
som i de ethiske Religioner, træder i et væsentligt For%
% --- p. 245 ---
\setcounter{footnote}{0}%
\opage{245}hold til det Guddommelige, erhverver han sig derved
selv et indre Værd.. I Religionen hører han op at % PRINTED AS IS (Værd..: the image, two points, spaced; O163)
have en blot Markedspris\dots{} I Verdenshistoriens almindelige
Strøm og under Folkenes almindelige Skjæbne
arbeider Christendommen paa at lade Sjælen finde sin
Grund i sig selv og ikke at slippe den, uden at den
gjennem Gud griber sig selv. Saaledes forbliver Mennesket,
der ellers taber sig i det Store, Menneske i sig
selv, og kun ved Mennesket i det Enkelte er Staten,
Mennesket i det Store, Menneske.“

Men Trendelenburg vender atter tilbage til en antik
Tankegang, idet han opfatter det Religiøse paa objektiv
Vis og som en direkte Gjenstand for Statens Omsorg;
idetmindste ligger det hos ham nærmere at udlede dette
af det Antikes end af det Middelalderliges Eftervirkninger.
Staten, siger han, kan ikke undvære Kirken og dens
aandelige Hjælp, Muligheden af den mest indgribende
Virkning paa Sindelaget til at hæve Individerne over
Egoismen og Nydelsessygen. Det er, som allerede anført,
efter hans Opfattelse urigtigt at sætte Staten i
Modsætning til Kirken; den behøver Kirken, eller rettere
Statens verdslige Side behøver den geistlige. Deraf
følger ikke, at den skal slutte sig til en bestemt Religion
eller Konfession og have en Statskirke; den skal
frit lade de forskjellige Retninger udfolde sig og repræsenterer
overfor dem den humane Ethiks og Tolerancens
Ide. Men Trendelenburg hævder, at Staten maa sørge
for, at Børnene opdrages i en bestemt Religion, selv om
Forældrene sige sig løs fra alle anerkjendte Kirker.
„Den Forudsætning, at det skulde være dem, der ville
det, muligt at leve i Staten uden Religion [ɔ: efter
Sammenhængen: uden positiv Religion eller Konfession]
tør ligesaalidt tænkes, som de gamle Lovgivere vilde
tænke sig Muligheden af Fadermord.“ Som Sokrates
blev anklaget for ikke at antage de Guder, Staten antog,
saaledes vilde man altsaa nu kunne tænke sig en Anklage
% --- p. 246 ---
\setcounter{footnote}{0}%
\opage{246}for ikke at høre til nogen af de Religioner, Staten anerkjender.
Trendelenburg har ikke bestemt nok fremhævet
Modsætningen mellem det positive og det rationelt-humane
Element i Religionen. Hans Begrundelse
af det Religiøses Betydning for Staten støtter sig egentlig
til det sidste Element; først i Konklusionen indskydes
det første. Og dog maatte det ligge ham nær at
fremdrage det store Problem, han her har skudt tilside,
eftersom han selv har fremhævet Ideens Begreb som
Principet baade for Erkjendelsen af det Guddommelige
og for Erkjendelsen af det Menneskelige. En
skarpere Opfattelse af Ideens Væsen, støttende sig til Nødvendighedstankens
absolute Gyldighed, som Trendelenburg
selv anerkjender, vilde ogsaa have ført ham til en
skarpere Opfattelse af det religionsphilosophiske Problem,
der ikke er af mindre Betydning end det sociale Problem,
som Trendelenburg med saa ædel Humanitet henpeger paa,
og som sikkert har mange Berøringspunkter med hint.

2. \emph{Gustav Theodor Fechner} (f. 1801) var en
Tidlang Professor i Physik i Leipzig, men udgav foruden
Arbeider i denne Videnskab ogsaa forskjellige poetiske
og humoristiske Skrifter. Tankens og Phantasiens frie
Løb blev hos ham ikke lammet ved nogen Fagensidighed.
Snart søgte han ogsaa philosophisk at udvikle
sin Verdensanskuelse, saaledes som den havde udviklet
sig for ham i Vexelvirkning med hans exakte
Studier. Fra den forudgaaende Philosophi har han efter
sin egen Erklæring modtaget vigtige Impulser. Det
var et Skrift, der var udsprunget af Schellings Anskuelser,
nemlig Okens Naturphilosophi, som ved sin titaniske
Dristighed først førte ham ud over den almindelige
Naturopfattelse. „Den bedste Frugt“ har han „plukket
af en rigtignok langt til Siden bøiet Gren af Hegels
Philosophi“, idet han skylder Weisse sin religionsphilosophiske
Opfattelse af Christendommen. Endelig har
Herbarts mathematiske Psychologi været ham et For%
% --- p. 247 ---
\setcounter{footnote}{0}%
\opage{247}billede, skjøndt nærmest ved dens Misligheder: „jeg har
brændt et Stykke Kul fra Herbarts Aske paa min Arne“\footnote[1]{Ueber die physikalische und die philosophische Atomenlehre. 2te Aufl. p. XIV.}.
Men dog stiller han sin Philosopheren i en afgjort Modsætning
til den tidligere. Dennes Grundvildfarelse bestaaer
for Fechner i Bestræbelsen for at finde det,
som ligger bag ved Phænomenet, i den Forudsætning,
at Alt, hvad vi se, høre, føle, ja endog hvad vi tænke,
skulde være et subjektivt Skin, som maa hæves, for at
man kan naa det virkelig Reale. Hverken Erfaringsvidenskab
eller Philosophi kan gaa bagved det Sandselige,
Phænomenet og hvad man maa forestille sig i Sammenhæng
med det. I denne lyse Verden skal Tanken forblive
og ikke mene at vinde kostbare Skatte ved at fordybe
sig i Underverdenen, hvor den kun vil finde Skygger.
Hvorledes kan da den ideelle Interesse tilfredsstilles,
den Længsel, som ikke slaaer sig til Ro med det blot
Givne? Aldrig vil man, siger Fechner, kunne indespærre
den menneskelige Aand og det menneskelige Gemyt saaledes
i Phænomenet, at den finder sig tilfreds i dette
Fangenskab og ikke stræber at naa et Dybere, et Oversandseligt.
Men denne Trang skal ikke tilfredsstilles
ved at gaa tilbage til et Væsen bag Phænomenet, men
ved at forfølge Phænomenernes hele Sammenhæng i de
mindste Enkeltheder og de fineste Nuancer saavelsom i
i den totale, afsluttende Enhed. Vel naaer Erfaringen % PRINTED AS IS (i den totale: the image, i doubled over the line turn (i|i); O167)
ikke det Uendelige, ja ikke engang langt i Endelighedens
Verden; men den Traad, der leder os indenfor Erfaringens
Omraade, kunne vi følge ud over dette. De Synspunkter,
der gjælde indenfor Erfaringen, behøve kun at
udvides, forhøies, potentseres, for at man ved Hjælp
af dem skal kunne naa de høiere Sphærer. Naturforskeren,
der Skridt for Skridt støtter sig til Erfaring og
Mathematik, har kun at gjøre med Naturens ydre Side
% --- p. 248 ---
\setcounter{footnote}{0}%
\opage{248}ɔ: med Naturen som Gjenstand for ydre Iagttagelse. Men
ved Induktion, Analogi og fornuftig Kombination kan
en videregaaende Naturbetragtning saaledes udvide og
uendeliggjøre det umiddelbart Givne, at dette fører ud
over sig selv. Dette er den eneste theoretiske Vei, ad
hvilken man kan naa en Viden om, hvad der ligger ud
over det i Erfaringen Givne. Det er en Gang gjennem
den synlige Verden for at finde den usynlige. At ville
gaa den omvendte Vei, saaledes som den spekulative
Philosophi, det er at ville begynde Bygningen af et
Taarn fra oven af istedetfor først at sørge for en solid
Grundvold.

Det Første, man altsaa maa sikre sig, er Erfaringsvidenskabens
Methode og Resultater, for paa dette Grundlag
at bygge videre. Det kommer da her dels an paa
de Grændsebegreber, med hvilke den exakte Erfaringsvidenskab
afslutter, for at Metaphysiken saa paa sin Vis
kan arbeide videre i det samme Spor, dels paa saadanne
Synspunkter, ved hvilke en exakt Opfattelse af
de Modsætningsforhold, som spille en væsentlig Rolle i
metaphysisk Henseende, bliver mulig. I begge Henseender
har Fechner ydet overordentlig fortjenstlige Arbeider.
I sit Skrift: „Ueber die physikalische und die
philosophische Atomenlehre“ (1855. Andet Oplag 1864)
har han præciseret Materiens Begreb, saaledes som det
stiller sig fra Naturvidenskabens Synspunkt, og har dernæst
søgt at vise, hvorledes dette Begreb kan gaa ind
i en philosophisk Anskuelse. I sine „Elemente der Psychophysik“
(2 Bind. 1860) har han givet en exakt Lære
om Forholdet mellem Sjæl og Legeme, der vel er uafhængig
af videregaaende metaphysiske Forudsætninger,
men dog med Nødvendighed viser hen til saadanne. I
sine øvrige Skrifter\footnote[1]{Navnlig: Zendavesta (3 Bind 1851). Ueber die Seelenfrage (1861). Die drei Motive des Glaubens (1863).} søger han at føre det store Analogi%
% --- p. 249 ---
\setcounter{footnote}{0}%
\opage{249}og Induktionsbevis, hvorpaa efter hans Overbevisning en
philosophisk Verdensanskuelse ene kan bygges, og han
skyer ingen Konsekventser, selv om de synes mystiske
eller phantastiske. Ofte gaaer Prosastilen (især i Zendavesta)
over til dithyrambisk Digt. --- Det er ikke Fechners
exakte Undersøgelser, der have ført ham til hans
philosophiske Verdensanskuelse; men omvendt førte hans
philosophiske Forsken ham til at stille den Opgave, hvorved
Grunden blev lagt til „Psychophysiken“\footnote[1]{Elemente der Psychophysik, II. p. 553.}. ---

Spekulationen, siger Fechner i sin Atomlære, stiger
ned i Naturen, ligesom Bjørnen i et Bistade. Den roser
den sammenhængende Honning, men ændser ikke, at de
enkelte smaa Væsener have bragt den sammen og have en
Ret til den; den tager Honningen med ét Greb og omstyrter
hele Bistadet. Naturvidenskaben er derimod som
Bivogteren, der véd, at Bygningen i al sin Enhed er
utallige Enkeltes Værk, udført under en stille og skjult
virkende Lovs Herredømme, og som derfor pleier dem,
følger deres Flugt og gjør deres hele Leven og Færden
til Gjenstand for sin opmærksomme Iagttagelse. Som
der hersker uforsonlig Strid mellem Bivogteren og Bjørnen,
saaledes ogsaa mellem Naturvidenskaben og Spekulationen.
Denne vil ad apriorisk-dialektisk Vei udgrunde
Naturens Væsen og ender derfor i en dynamisk Anskuelse,
som udleder Materien af almene Kræfters Vexelvirkning,
hvorved ethvert fast Grundlag ophæves. I
Modsætning hertil hævder Naturvidenskaben Materiens
exakte Diskretion. Ved en Række af Kjendsgjerninger
fra Physikens og Chemiens forskjellige Egne (Undulationstheorien,
Varmelæren o. s. v.) søger Fechner at
godtgjøre, at Naturvidenskaben ikke kan opfatte Materien
som et Kontinuum, men maa betragte det som et
System af diskrete Dele, der for den ere uopløselige.
Alt individualiserer sig saaledes indtil det mest Enkelte,
% --- p. 250 ---
\setcounter{footnote}{0}%
\opage{250}men forbindes netop derved  til det fasteste Hele. Som
Himmellegemerne bevæge sig i Verdensrummet, saaledes
Atomerne i det uendelig lille Rum, og her som hist
hersker en almen Lov. Den moderne Naturvidenskab
er derfor atomistisk i den Forstand, at den fastholder
Materiens Diskretion. Men den indlader sig ikke som
den gamle Atomistik paa at udgrunde det absolut Sidste.
Naar man ved Atom forstaaer det absolut Udelelige, da
passer Benævnelsen Atomistik ikke paa den moderne
Naturvidenskab, som kun hævder, at Diskretionen maa
gaa videre end Øiet og Mikroskopet naa, men ingenlunde
negter en Delelighed in indefinitum. Den moderne Atomlære
er altsaa relativ. Physikeren paastaaer ikke, at
Rummet mellem Atomerne er tomt, og at der ikke mellem
dem kan findes et fint Stof, der ikke har nogen Indflydelse
paa de Atombevægelser, han kan kontrollere,
ligesom Ætheren opfylder Verdensrummet mellem Himmellegemerne.
Mellem dette Stofs Atomer kunde der da
igjen tænkes en endnu finere Æther o. s.v. Vi bevæge
os i det Hele her paa det Relatives Omraade. Det er
kun af Vigtighed at fastholde saadanne Punkter, hvor
den exakte Analyse kan finde sin Tilknytning. Men
dette er umuligt for den spekulativ-dynamiske Anskuelse,
der forklarer Materien som Resultatet af en kontinuerlig
virkende, rumopfyldende Kraft. Dette Begreb er et Absurdum
for Physiken. For denne er Kraften inkommensurabel
med Materien og kan ikke frembringe den eller
dialektisk blive Et med den. Alle Kræfter ere Forholdsbegreber
og forudsætte derfor Noget, som staaer i disse
Forhold. Hvad man tilskriver et Legeme som dets bestemte
Kraft, er den Andel, hvormed det virker til Lovens
Opfyldelse. Det Stof, som forudsættes, er Atomerne,
der ikke maales (som Kraft, Rum og Tid), men
tælles. Atomernes Bevægelse udfylder Tid og Rum.
Atomets Masse er ikke Udstrækning; den er physisk set
Et med det konstante Forhold, i hvilket ved dette Atom
% --- p. 251 ---
\setcounter{footnote}{0}%
\opage{251}Kraften stedse staaer til Hastighedens Forøgelse. Massens
Begreb er altsaa ikke materialistisk. Tvertimod er
den Grundopfattelse, mod hvilken Physiken i sidste Instants
peger, absolut idealistisk. Materiens Udstrækning
er = o, og det materielle Atom er et Enkeltvæsen, en
punktuel Intensitet, et Kraftcentrum.

Uagtet Fechner indrømmer den moderne Atomlæres
relative og skeptiske Karakter, gaaer han dog selv ud
over dens Grændser og gjennemfører Atomistiken absolut.
Og han gjør dette udtrykkelig af philosophiske
eller metaphysiske Grunde. Metaphysikens Opgave bestaaer
for ham i at finde og udforske Erfaringens
Grændsebegreber i deres almindelige Forhold og Sammenhæng,
og dens Methode i at almindeliggjøre og gjennemføre,
hvad den exakte Erfaringsvidenskab har fundet og
godtgjort, til en saadan Afslutning, som Tænkningen
fordrer, idet selve Betingelserne for at være et Sidste
og Høieste danne Formen for de derved fundne Begreber.
Metaphysiken er altsaa ikke apriorisk, men svarer
til sit Navn, idet den kommer efter Physiken. Den metaphysiske
Erkjendelse kan ikke ligefrem udvide den
exakte Viden; den bliver, physisk set, bestandig en Hypothese,
men er dog Physikens begrebsmæssige Slutning,
til hvilken man kommer, naar man fører den Vei, som
den exakte Viden gaaer med Sikkerhed, til Ende i
Ideen. Ikke physisk, men metaphysisk kommer Fechner
altsaa til den absolute Atomistik, der opfatter de
materielle Atomer som Kraftcentra. Men han kommer
til den ved umiddelbart at fastholde de exakte Grændsebegreber,
til hvilke Physiken leder. Han vægrer sig ved
at gaa bag ved Phænomenet for at finde Materiens Væsen.
Det Materielle er det Haandgribelige, det almindelige
Grundlag for Naturphænomenerne; bag det kan Intet
erkjendes, men i Sammenhæng dermed og ud derover er
der Plads for Abstraktioner og høiere Synspunkter. Substantsens
Begreb er Et med Begrebet om den hele phæ%
% --- p. 252 ---
\setcounter{footnote}{0}%
\opage{252}nomenale Sammenhæng, i hvis Lovmæssighed det Objektive
i Materien bestaaer. Fra dette Udgangspunkt er
det saa, at Physiken kommer til Begrebet om Materiens
Diskretion, og et andet Begreb om Materien kan heller
ikke Metaphysiken lægge til Grund. Men naar man
gjennemfører dette Begreb, forsvinder al Byrde og Tyngde;
enhver uigjennemtrængelig Skranke, som staaer Aanden
imod, forvandler sig til Kraft.

Bag Phænomenet kan Intet erkjendes. Men det
Phænomenale selv indeslutter i sig en vigtig Forskjel,
hvorved det deler sig i to store Grupper. Det Legemlige
er det, der kun lader sig tilsyne for Andet; det Sjælelige
beroer paa en indre Tilsyneladelse for sig selv. Vi
faa altsaa en Modsætning mellem indre og ydre Phænomener.
I Sagen selv er der dog ingen Modsætning, men
den fremkommer ved de to forskjellige Standpunkter,
der kunne indtages. Hvad der paa det indre Standpunkt
fremtræder som Aand, fremtræder paa det ydre Standpunkt
som Legeme, ligesom den samme Cirkel indenfra
set viser sig konkav, udenfra set konvex. Den, der derfor
kun kjender det ydre Standpunkt, vil betragte Alt
som legemligt og negte Aanden; saaledes Materialisterne,
der kun se sig om i Peripherien af Cirkelen, og da de
her ikke opdage noget Centrum, negte, at der er
noget. Omvendt ville Idealisterne ud fra Centrum bestemme
Peripherien; men om samme Centrum kan der
jo beskrives mangfoldige Cirkler. Dog bliver Modsætningen
mellem disse Standpunkter ikke staaende: Erfaringen
viser os et bestemt og lovmæssigt Forhold mellem
det Sjælelige og det Legemlige, og vi søge i Psychophysiken
en exakt Lære om dette Forhold. For at
kunne begrunde en saadan Lære behøves der et Maal
for de sjælelige og de legemlige Bevægelser. Vanskeligheden
er her at finde en Maalestok for det Sjælelige, da
dette ikke umiddelbart kan maales, skjøndt der ikke kan
være nogen Tvivl om, at det er undergivet kvantitative
% --- p. 253 ---
\setcounter{footnote}{0}%
\opage{253}Bestemmelser, eftersom der jo gives en større eller ringere
Styrke af Fornemmelser og Drifter, høiere og lavere
Grader af Opmærksomhed, af Bevidsthedens Klarhed
o. s. v. Den eneste Maade, hvorpaa disse aandelige
Bevægelser kunne maales, er at se dem i Forhold til de
legemlige Bevægelser, til hvilke de med Nødvendighed
ere knyttede. Holde vi os nu til Fornemmelsen, den
mest elementære sjælelige Bestemmelse, da har denne
til nödvendig Betingelse dels det ydre Indtryk, dels de % PRINTED AS IS (nödvendig: with ö, a wrong sort; the crop f:C175 and the Minnesota layer)
organiske Funktioner, hvorved Indtrykket omsættes til
indre Fornemmelse. Men disse organiske Funktioner
(det Psychophysiske i snevrere Forstand) ere os ubekjendte
i deres indre Væsen, og det er dog dem, der
ere Mellemledet mellem det Physiske og det Psychiske.
Den indre Psychophysik, der undersøger Forholdet mellem
det Sjælelige og det Organiske (mellem det Psychiske
og det Psychophysiske), maa derfor bygges paa den
ydre Psychophysik, der umiddelbart undersøger Forholdet
mellem Fornemmelse og Indtryk (mellem det Psychiske
og det Physiske). Fra denne ville vi da kunne slutte
tilbage til hin.

Ingen Fornemmelse staaer pludselig paa sin fulde
Høide, men enhver gjennemløber en hel Skala ligefra
det Punkt, hvor den endnu er umærkelig, og det ofte i
saa kort Tid, at den synes pludselig at være tilstede i
sin hele Kraft. Denne successive Væxt gjør det muligt
at maale den. Thi naar man iagttager Fornemmelsens
Tilvæxter ved de successive Forøgelser af Indtrykkets
Styrke eller Udstrækning, vil man have de Elementer
givne, paa hvis Forhold Maalet beroer. Vel synes det
ikke let at forvisse sig om Ligestorheden af Fornemmelsens
Tilvæxter. Men dog kan man jo bedømme to Toners
eller Toneintervallers Lighed, og desuden angiver
Fechner tre Methoder, ved hvilke idetmindste en exakt
Tilnærmelse vil kunne naaes. Den første er „de netop
mærkelige Forskjelles Methoder“, som ved Forsøg be%
% --- p. 254 ---
\setcounter{footnote}{0}%
\opage{254}stemmer den Størrelse af Forskjel (i Vægt o. s. v.),
som er nødvendig for, at Forskjellen overhovedet skal
kunne fornemmes; de to andre støtte sig til Sandsynlighedsberegningen
og søge ved Hjælp af denne gjennem
gjentagne Forsøg at eliminere de Feil, der nødvendig
maa indsnige sig, hvor som her den subjektive Fornemmelse
spiller en saa vigtig Rolle.

Det simpleste Forhold vilde være det, at Tilvæxten
i Fornemmelse var ligefrem proportional med Indtrykkets
Forøgelse. Dog er dette ikke Tilfældet. Physiologen
E. H. Webers Forsøg have derimod godtgjort, at
det næst simpleste Forhold finder Sted, at nemlig de \emph{relativt}
lige store Tilvæxter i Indtryk svare til ligestore
Tilvæxter i Fornemmelse --- naturligvis forudsat, at Modtageligheden
er den samme og Forholdene i det Hele
normale. Dette er Indholdet af, hvad Fechner kalder
den Weberske Lov, der kun er det exakte Udtryk for,
hvad der er Enhver bekjendt fra den daglige Erfaring.
Man hører tydelig sin Nabos Stemme, naar det er stille,
men under stærk Larm kan man ikke høre, hvad man
selv siger, ɔ: ikke mærke den ved Ens egne Ord frembragte
Indtryksforøgelse, fordi denne er forsvindende i
Forhold til det allerede forud tilstedeværende Indtryk.
Naar et Værelse er stærkt oplyst, mærker man ikke, at
der kommer et Lys til o. s. v.

Ethvert Indtryk eller Forskjel mellem Indtryk maa
allerede have naaet en vis endelig Størrelse, før det kan
mærkes. Fornemmelsens Nulpunkt ligger altsaa over
Indtrykkets Nulpunkt. Fornemmelsens Nulpunkt, hvor
den netop begynder eller forsvinder, kalder Fechner
Tærskelen. Saalænge Indtrykket ikke er stærkt nok
til at hæve Fornemmelsen over dette Punkt, er den ubevidst,
og jo dybere den ligger under Tærskelen, des
høiere er Ubevidsthedens Grad. En i sig selv umærkelig
Tilvæxt kan her gjøre Udslaget og ved at slutte sig til
% --- p. 255 ---
\setcounter{footnote}{0}%
\opage{255}allerede forhaandenværende Indtryk hæve Fornemmelsen
over Tærskelen ɔ: lade den blive bevidst.

Et mathematisk Udtryk for den Weberske Lov haves
i Formler, der allerede opstilledes af Dan. Bernouilli i
Begyndelsen af forrige Aarhundrede for at udtrykke Forholdet
mellem relativ og absolut Værdi (ex conditione personæ
--- ex re ipsa). Den første af disse Formler, som
Fechner kalder den psychophysiske Fundamentalformel,
er: dγ = k dβ/β, hvor γ betegner Fornemmelsen, β Indtrykket og k en af de for β og γ valgte Enheder afhængig
Konstant. Denne Formel udsiger altsaa, at Fornemmelsestilvæxten (dγ) beroer paa Forholdet mellem Indtrykstilvæxten (dβ) og det forud tilstedeværende Indtryk (β).
Hertil slutter sig dernæst, idet Forholdet mellem
Fornemmelsens og Indtrykkets Tilvæxter er det
samme som mellem Logarithmens og Tallets Tilvæxter,
og den saakaldte Maalformel, hvis simpleste Form er : γ =
log β. Denne Formel udsiger, at Fornemmelsen stiger
i logarithmisk Forhold til Indtrykket. Og dersom den
fører til negative Værdier angive disse Ubevidsthedens
Dybde. De positive Værdier angive Bevidsthedens Styrke.

Den Weberske Lov viser sig dog ikke at være gyldig
i alle Tilfælde. Dersom Maalformelen havde ubegrændset
Gyldighed, maatte Fornemmelsen voxe i det
Uendelige ved Forøgelse af Indtrykket. Men paa Grund
af vor Organismes, navnlig Sandseorganernes Indretning
ophører Loven at gjælde ved meget stærke Indtryk.
Derimod synes den at maatte have uindskrænket Gyldighed
paa den indre Psychophysiks Omraade, da der her,
paa Grund af det umiddelbare Forhold mellem det Psychiske
og det Psychophysiske, ikke er saadanne Grunde
til dens Begrændsning som ved Forholdet mellem Fornemmelsen
(det Psychiske) og Indtrykket (det Physiske).
Men mere end den almindelige Overbevisning om den
Weberske Lovs Gyldighed i den indre Psychophysik
% --- p. 256 ---
\setcounter{footnote}{0}%
\opage{256}kan ikke naaes. Vi kjende nemlig endnu for lidet Forholdet
mellem det Psychiske og det Psychophysiske, til
at vi her kunne faa et sikkert Tilknytningspunkt for
Analysen. Medens vi i den ydre Psychophysik altid
kunne bestemme Værdien af Maalformelens β, der betegner
Indtrykkets Størrelse (Vægt o. s. v.), saa er det i
den indre Psychophysik, hvor β betegner den organiske
(psychophysiske) Bevægelse, der betinger Bevidsthedslivet,
umuligt at bestemme denne Størrelses Værdi; og
da γ jo er den ubekjendte Størrelse, der skal findes, bliver
Regningen umulig\footnote[1]{Fechner havde fattet Ideen til sin Psychophysik med Herbarts mathematiske Psychologi for Øie, men saaledes, at han vilde undgaa den Ensidighed i Udgangspunktet, der hindrede dens Gjennemførelse. Han saa nemlig Herbarts Mangel deri, at denne ikke sikrede sig paalideligt Maal i det materielle Grundlag for de sjælelige Phænomener. Derfor gaaer han ud fra Fornemmelsen, medens Herbart gik ud fra Forestillingen. Men dette sikre Grundlag forlader ham jo netop, naar han fra den ydre Psychophysik gaaer over til den indre, og denne ligger derfor under for den samme Mislighed som Herbarts mathematiske Psychologi.}. Derimod kan den indre Psychophysik,
naar det almindelige lovmæssige Grundlag staaer
fast, føre os til almindelige theoretiske Resultater angaaende
Sjælelivet og dets Forhold til det Legemlige.

Hvad for det Første den almindelige Opfattelse af
Sjælens Væsen angaaer, da erklærer Fechner sig bestemt
mod den monadologiske Opfattelse, der navnlig
repræsenteres af Lotze, og som i Sjælen ser en punktuel
Enhed. I Modsætning hertil betegner Fechner sin
egen Opfattelse som den synechologiske og ser Sjælen
som det sammenfattende Princip for den legemlige Flerhed
af Dele og de legemlige Funktioners Succession.
Kun i sin indre Fremtræden for sig selv er Aanden en
Enhed; i den ydre Tilsyneladelse udfolder den sig i en
Flerhed, ligesom det, der i det Indre er en enkelt Fornemmelse,
i det Ydre fremtræder som en sammensat og
% --- p. 257 ---
\setcounter{footnote}{0}%
\opage{257}successiv Nerveproces. Efter den synechologiske Opfattelse
ere altsaa Bevægelserne i det Sjælelige og i det
Legemlige samtidige, parallele Udtryk for samme Faktum.

Tærskelens Kjendsgjerning er af Betydning ogsaa
for den indre Psychophysik ved at afgive et fast Fundament
for Ubevidsthedens Begreb. Saaledes ved Forholdet
mellem Søvn og Vaagen. Søvnen er ikke et sjæleligt
Intet. Sjælens Livsbevægelse finder ikke blot Sted
i den vaagne Tilstand. Denne er kun Sjælens Liv over
Tærskelen, medens Søvnen er Sjælelivet under Bevidsthedens
Tærskel; betegne vi Tærskelen ved 0, saa betegne
de positive Værdier den vaagne Bevidstheds Høide og
de negative Værdier Ubevidsthedens Dybde under Søvnen.
Vi have her en Skala, som viser en kontinuerlig
Overgang i Subjektet mellem den vaagne og den
sovende Tilstand. Og ligesom i den ydre Psychophysik
Indtrykket kunde have en positiv Værdi paa det Punkt,
hvor Fornemmelsen er 0 (Tærskelen), saaledes indtræder
Grændsepunktet mellem Søvn og Vaagen ikke ved en
Nulværdi men ved en vis positiv Værdi af den til Grund
liggende legemlige Virksomhed.

Dette Forhold kan nu faa sin videre Anvendelse.
Mellem de forskjellige Organismer viser den Bevidsthedens
Kontinuitet sig afbrudt, hvorved en Flerhed af
Phænomener sammenknyttes i samme Sjæleenhed. Og dog
ere alle Organismer ved den almene Natur forenede til
et eneste System. Den psychophysiske Virksomhed strækker
sig altsaa ikke videre \emph{mellem} Organismerne paa
samme Maade som i dem. Men en Ophøren af den psychophysiske
Virksomhed vil ikke sige Andet end den størst
mulige Dybde under Tærskelen. Hvad der er under Tærskelen
adskiller som Ubevidsthed de bevidste Væsener, idet
det tillige danner Grundlaget for dem og underholder
den physiske Forbindelse mellem dem. Den synechologiske
Opfattelse maa da her tilsidst føre os til Antagelsen
af en almindelig Besjæling af Naturen paa forskjellige
% --- p. 258 ---
\setcounter{footnote}{0}%
\opage{258}Trin, idet Alt, som Thales sagde, er fuldt af Guder, ---
til „Anskuelsen af en i Naturen allestedsnærværende
Gud, i hvem alle Aander leve, røres og ere, som han i
dem, med Mellemtrin mellem ham og os i Himmellegemernes
Sjæle, saa at de skabte Aander bæres ligesaa
nøie sammenknyttede i dem, som disse selv i den
guddommelige Aand, og som de skabte Aander igjen
bære deres Sandser og disse igjen deres særlige Fornemmelser
i sig.“

Her staa vi da ved det Grændsepunkt, hvor Fechner
fra den exakte Videns begrændsede Sphære med den
dristigste, ingen Paradoxier frygtende Tankeflugt begiver
sig ud paa sin poetisk-spekulative Verdensanskuelses
Baner. Der er, som allerede omtalt, ingen blot tilfældig
eller udvortes Forbindelse mellem hans exakte og hans
spekulative Anskuelse. Ikke blot fremhæver han selv
gjentagne Gange, at det var hans spekulative Ideer om
Forholdet mellem det Sjælelige og Legemlige, der først
førte ham til Tanken om Muligheden af en Psychophysik;
ikke blot føre omvendt de psychophysiske Undersøgelser
ham tilbage til hans spekulative Grundanskuelse;
--- men Methoden og de ledende Principer
ere de samme: han søger at blive paa Induktionens og
Analogiens Vei, selv efterat det Omraade er forladt,
hvor en exakt Erfaring er mulig.

I Almindelighed anser man, siger Feehner, kun % PRINTED AS IS (Feehner: the image, O173)
Mennesker og Dyr for besjælede, ellers Intet, idetmindste
paa vor Jord. Grunden til, at man saaledes indskrænker
Sjælelivets Sphære og lader Livet være Snyltegjæst
i Dødens Rige, er navnlig den, at man troer at
maatte hæve Gud saa høit op over Verden, at denne
ikke tør anse sig som hans Legeme. Saaledes er man
da kommen til den ulykkelige, fredløse Tro paa en naturløs
Gud og en aandløs Natur. Man troer, at den exakte
Erfaring forbyder at udstrække Besjælingen videre. Men
det hele Sjælespørgsmaal er fra første Færd en Troes%
% --- p. 259 ---
\setcounter{footnote}{0}%
\opage{259}sag. Et exakt Bevis hviler paa Erfaring og Mathematik;
men kun om sin egen Sjæl har man direkte Erfaring,
og Mathematiken mangler ethvert Grundlag til at bevise
en anden Sjæls Virkelighed. Selv min Faders, min Broders
Sjæl maa jeg tro paa, idet jeg ad Analogiens og
Induktionens Vei slutter mig til den. Vi begynde i
Virkeligheden som Cartesius med et cogito ergo sum, og
en ubevidst Slutning om de andre Menneske- og Dyresjæle
bliver en dogmatisk Forudsætning, som man kun
i Blinde kan antage for en bevist Kjendsgjerning. Men
naar dette forholder sig saaledes, hvad kan da hindre os
i at udstrække den samme Fremgangsmaade, vi ubevidst
anvende ved Mennesker og Dyr, videre --- navnlig til
Planteverdenen og til de himmelske Kloder, forudsat, at
der viser sig gyldige Grunde dertil?

Man indvender, at Planterne ikke kunne have nogen
Sjæl, da de ingen Nerver have; men have da Polyperne
Nerver? --- at Planten ikke er et i sig afsluttet Individ;
men det Individuelle er kun noget Relativt, da jo intet
levende Væsen staaer absolut afsondret overfor Oververdenen.
Hvad er der til Hinder for, at der til det af
Planten, vi kunne opfatte med vore Sandser, svarer en
indre sig selv fornemmende Enhed? Vel er der en saa
karakteristisk Forskjel mellem Planten og Dyret, at der
ogsaa maa være en eiendommelig Forskjellighed i det
indre Væsen; men paa den anden Side viser Planten
saa stor Analogi med den dyriske Organisme og er der
saa kontinuerlig en Overgang fra Dyre- til Planteverdenen,
at der ikke er Grund til at fastslaa en Modsætning
mellem dem som mellem Besjælet og Sjælløst.
Der maa i Planten findes et eiendommeligt Sjæleliv, om
det end maaske staaer ligesaa langt under Dyrets,
som dette staaer under Menneskets. --- Hvad der gjælder
om Planterne, gjælder og om Himmellegemerne. Hvorfor
skulde Jorden ikke være besjælet? Man maa blot
ikke begynde med at sætte en skarp Grændse  mellem
% --- p. 260 ---
\setcounter{footnote}{0}%
\opage{260}det Organiske og det Uorganiske. Det Uorganiske er
kun noget Lavere, forsaavidt det, som i Physiken og
Chemien, tænkes løsrevet fra sin levende og naturlige
Forbindelse med det Organiske; begge høre sammen i
deres Bestaaen som i deres Oprindelse. Mennesker og
Dyr ere endnu inderligere sammenvoxede med Jorden
end Klipper og Stene. Vore Legemer ere Jordens Dele
og Organer, vore Sjæle Dele af Jordens Sjæl. Saa ulig
Jorden end er os, Uligheden betyder ei Andet end større
Høide og Fylde, ikke Mangelen af, hvad en Sjæl behøver
til Udtryk for sit Væsen. Med et aandigt Baand
ere alle jordiske Sjæle forenede i én, som en Mængde
mindre Kredse indeni den større, mere omfattende. Man
har ofte talt om Menneskeaanden, om Ideen i Menneskehedens
Historie. Hvad her forstaaes ved Jordens Sjæl
er egentlig ikke Andet, kun at det Menneskelige ikke
isoleres fra den store Sammenhæng, hvortil det hører,
og navnlig at vi ikke opfatte Ideen som bevidst i det
Enkelte, men ubevidst i det Hele. De lavere Aander
kunne kun finde deres fælles Enhed i en høiere bevidst
Aand, hvis enkelte Anskuelsesbilleder og Motiver de ere.

Som de lavere Aander ere Dele af de høiere, saaledes
disse tilsidst af den høieste. Beviset for dennes
Tilværelse er derfor af samme Natur, som Beviset for
hines. Hvad Menneskehedens historiske Troesformer vise
hen til, hvad den praktiske Trang fordrer, det følger
ogsaa af det eneste theoretiske Bevis, som her er muligt.
Vi finde ved at undersøge Tilværelsen i dens hele
Omfang en høieste, almindelig Verdenslov, Principet for
alle enkelte Love, den nemlig: naar og hvorsomhelst
de samme Omstændigheder vende tilbage, vende ogsaa
de samme Følger tilbage; men under andre Omstændigheder
indtræffe andre Følger. Denne Lov viser, at ét
Væsen gjennemtrænger alle Steder og Tider, alt Aandeligt
og Legemligt. I det er den høieste Lov Et med
den høieste Frihed. Det staaer nu tilbage at erkjende,
% --- p. 261 ---
\setcounter{footnote}{0}%
\opage{261}hvorledes Guddommens Væsen er, at erkjende dette som
Aand, som uendelig Bevidsthed. Herom gjælder det,
ligesom ved de endelige Aander, at kun den egne Bevidsthed
er vis paa sig selv. Heller ikke overfor den
guddommelige Aand have vi andre Midler til Erkjendelse
end Analogi og Induktion. Fra Betragtningen af
vor Aand hæve vi os til Tanken om den absolute Aand,
fra Betragtningen af vort Legeme som Afspeiling og
Bærer af vor Aand hæve vi os til Anskuelsen af den
hele Verden som Afspeiling og Bærer af den uendelige
Aand. Vel maa vi, som vi have set, tænke os som
Mellemled en Kjede af høiere Aander; men den absolute
Aand er dog umiddelbart tilstede selv i de allerlaveste,
ligesom en lille Kreds umiddelbart er i en større,
selv om den indenfor denne omsluttes af en tredie Kreds,
der i Størrelse ligger imellem dem begge. Idet vi tænke,
føle og anskue, tænker, føler og anskuer den uendelige
Aand umiddelbart ved og i os. Her er ingen Udvorteshed,
men absolut Inderlighed. I Almindelighed hævder
man rigtignok Modsætningen mellem den uendelige og
den endelige Aand som mellem endelige Aander. Man
stiller det Uendelige overfor, over, hinsides, udenfor det
Endelige, ja anbringer maaske endog en uhyre Skillevæg,
for at de ikke skulle komme hinanden for nær.
Men de ligge slet ikke udenfor hinanden. Det Endelige
er det Uendeliges Indhold. Derfor er det Uendelige
heller ikke ufatteligt; det kan fattes i sit Indhold; kun
omfattes kan det ikke. Derfor behøver man ikke at befrygte
en Sammenblanding af det Guddommelige og det
Menneskelige; er du Gud, fordi du er i Gud? er det
mindste Moment det Samme som det Hele? --- Ikke
blot de endelige Aander ere i Gud, ogsaa den materielle
Verden. Verden er ikke noget Lavere end Gud, men
noget Lavere i Gud. En Yderverden gives der kun for
den endelige, ikke for den uendelige Bevidsthed. Gud
ser med nogle af sine Skabninger som med sit Legemes
% --- p. 262 ---
\setcounter{footnote}{0}%
\opage{262}Dele og Organer andre overfor dem stillede Dele og Organer
og griber med sin høiere Bevidsthed ud derover,
ligesom vor Aand over alle enkelte sandselige Iagttagelser.
De Naturkræfter, der synes os ubevidste, ere i sig
selv Elementer i en Villie, der rager ud over vor Bevidsthed
og virker gjennem hele Verden. Derfor behøver
ikke enhver enkelt Virken af en Naturkraft at
betyde en enkelt Villiesakt. Villien er rettet mod Maalet
for alle Kraftvirkningerne, og disse afgive ubevidst
deres Bidrag til det totale Øiemed. Saaledes tænker
Spindersken kun paa sit Garn, medens Foden mechanisk-ubevidst
træder Hjulet.

Idet Verdensbegrebet saaledes i sidste Instants falder
sammen med Gudsbegrebet, bliver Anskuelsen Pantheisme
i Ordets fulde Betydning. Den videste Opfattelse af
Gudsbegrebet er den mest begrundede og den snevrere
Opfattelse kun en Abstraktion. Dog har den almindelige
Pantheisme (den hegelske) Uret i kun at betragte
Enkeltvæsenerne som bevidste; al Bevidsthed maa tvertimod
tilsidst samle sig i et høieste bevidst Væsen\footnote[1]{Ligeledes adskiller Fechner sin Anskuelse fra Ørsteds i „Aanden i Naturen“ (der nylig var udkommen, da Fechner skrev sin „Zendavesta“). Han vil ikke som Ørsted kalde Naturlovene Naturtanker, da Lovene ere Tankens Gjenstand og Indhold, ikke selv Tanker; at gjøre Love og Tanker til Et kan efter hans Mening let føre til, at man glemmer at spørge efter de virkelige Tanker og den uendelige, Alt omfattende Bevidsthed, hvorfra de udgaa. Derfor kommer --- siger han fremdeles --- en bevidst Naturens Aand ikke til Gjennembrud hos Ørsted uden i det blotte \emph{Navn} af Gud.}.

Da Verdens Liv er Guds Liv, saa kan dette ikke
være afsluttet i sig selv, men udvikler og udfolder sig
med og ved Verdensudviklingen. Guds Fornuft er ikke
fuldstændig udviklet og bevidst forud for hans Virken;
han beskuer ligesom den geniale Kunstner først sit Værk,
naar det er fuldendt („Og Gud saa Alt, hvad han havde
gjort, og se, det var saare godt“). Vel herskede den
% --- p. 263 ---
\setcounter{footnote}{0}%
\opage{263}guddommelige Urfornuft fra Evighed af, men den hæver
sig stedse høiere over det materielle Grundlag ved stedse
mere at indlemme dette i sit Aandsrige. Gud er ikke
ufuldkommen, fordi han udvikler sig saaledes, thi hans
Fuldkommenhed bestaaer ikke i en færdig Afslutning,
men i en ubegrændset Fremadskriden. Heller ikke krænker
dette inderlige Forhold til Verden Guds Hellighed.
Thi den endelige Villies Synd kan ikke tilskrives Gud,
saalidt som en enkelt Tanke eller Tilskyndelse, der opstaaer
i vor Aand, kan gjøre denne ond. Ligesom
Mennesket kaldes godt eller ondt efter den hele Retning,
den høiere, omfattende Villie, de herskende Synspunkter
i hans Aand, saaledes maa vi maale Gud efter det
Totale og Evige i hans Væsen. Guds evige Villie fører
ud over det Onde, saa at dette tilsidst opløses i den
seirrige, Alt forsonende Harmoni. Spørges der, hvorfor
og hvorledes det Onde i det Hele taget da er til, kan
der ikke svares Andet end, at det Onde er med metaphysisk Nødvendighed, vel ikke uafhængigt af Guds Væsen,
men dog af hans Villie, der netop stedse arbeider paa
dets Forsoning og Ophævelse. Gud føler det Onde med
og i den Lidende; men det hindrer ikke hans Salighed,
da han i sin absolute Bevidsthed er hævet over alt Enkelt.

Denne Lære er ikke Andet end en virkelig og alvorlig
Gjennemførelse af Christendommens Udsagn, at
Gud er Aand og maa dyrkes i Aand og Sandhed, ---
at i ham ere, leve og røres vi, --- Udsagn, der i Almindelighed
staa hen som blotte Ord. Ingen har fremstillet
Gud og det høieste Bud saa rent og helligt som
Christus. Det Evige i Christendommen, det, hvorved
den gik ud over alt Tidligere, dens Ide er, „at det Høieste
sættes som det, der forener Alle, at det mest Omfattende
sættes som det, der skal forenes, og at det Bedste
sættes som det Høieste. Ideen om Alles Forening fra det
Synspunkt, fra hvilket en Enhed af Alle ene er mulig,
er først af Christus med Bevidsthed bleven bragt den
% --- p. 264 ---
\setcounter{footnote}{0}%
\opage{264}jordiske Verden til Bevidsthed, og ved Lære og Liv har
han gjort en levende Begyndelse til denne Ides Udbredelse
og Virkeliggjørelse.“ --- Derimod er den fremstillede
Lære ganske vist ikke christelig, naar man til Christendommen
regner som noget Væsentligt „Troen paa
Æblebiddet i Paradis med dets mystiske Følger, paa de
ikke Udvalgtes uigjenkaldelige Fordømmelse, paa Underne
mod Naturens Love, paa Guds Fjernhed fra hans
Verden og paa alt det lidet Opbyggelige, hvormed Theologerne
i Almindelighed udstyre Christendommen, ja
hvoraf de opbygge den.“ Christendommens Ide maa
derimod fattes saaledes, at den falder sammen med selve
Ideen om Religionen som det mest universelle, humane
Foreningspunkt. Christendom er det, fordi Christus først
har bragt denne Ide til Bevidsthed. Men ved sine
ideale ethiske Bud har han indirekte begrundet Spiritualismen
og Dualismen, skjøndt han selv ikke løser Gud
fra Verden. Denne Spiritualisme er det, der atter
har ført til at sætte det Uendelige udenfor det Endelige.
Naar denne den nuværende Verdensanskuelses
Ensidighed ophæves, vil Urgrunden, hvoraf saavel Hedenskab
som Christendom have udviklet sig, træde frem i
al sin Glands, idet Hedenskabet forklares ved Christendommen
og Christendommen forynges ved Hedenskabet.
Det er en Fornyelse af den gamle Naturreligion, men
saaledes at Aandsreligionen langt fra derved at krænkes
i sit sande og evige Væsen, netop beriges og udvikles,
ja først nu kommer til sin fulde Ret.

Fechner erklærer om sin Lære, at den er materialistisk,
forsaavidt den ser Alt som legemligt, men ligesaa
vel idealistisk, fordi den ser Alt som Aand; den er dualistisk,
idet den ser Aand og Legemlighed som parallele
Rækker, men ligesaavel Identitetslære, idet den med Spinoza
ser begge disse Rækker som Former for et og
samme Grundvæsen. Dog adskiller den sig fra Spinozismen
ved at hævde en Vexelvirkning mellem de to
% --- p. 265 ---
\setcounter{footnote}{0}%
\opage{265}paralelle Rækker, hvorved ikke blot et virkeligt Kausalitetsforhold % PRINTED AS IS (paralelle: for parallele; the crop reader and the Minnesota layer)
og dermed et exakt Funktionsforhold bliver
muligt, men ogsaa, hvad Spinoza benegter, en teleologisk
Sammenhæng, idet den sidste Grund til begge Rækkers
Uadskillelighed søges i den guddommelige Bevidstheds
Enhed. Men om det Grundvæsen, der ligger bag ved
hine to Rækker, kan der ikke spørges, naar man ikke
vil gaa bag ved Phænomenet, ud i det Transscendente. ---

Vel har efter Fechners eget Udsagn et Studium af
den nyere Philosophi grebet dybt ind i hans Udviklingsgang,
men dog førte det ham tillige, som vi have set,
i en eiendommelig Modsætning til den. Udgaaende fra
Erfaringsvidenskaben og besjælet af en barnlig Glæde
over alt Godt og Skjønt i Verden, kunde han kun tilfredsstilles
ved en Verdensanskuelse, der vel udvidede
og forhøiede Phænomenernes Liv (som han saa sympathetisk
fordybede sig i, at han endog kunde mærke Sjælelivets
Pulse der, hvor alle Andre negtede dem) men
ikke forklarede det ved at vise tilbage til en dunkel
Baggrund, hvorfra det skulde udspringe. Denne umiddelbare
Hengivelse til det Phænomenale er det, der giver
hans Anskuelse dens tiltrækkende Friskhed; men det er
tillige den, der philosophisk set er den væsentlige Anke
imod den. Den dybere Begrundelse kommer til at mangle.
Han skyer den, idet han anser den som et Mørke, der
fordunkler det klare Lys. Og dog fører en virkelig Begrundelse
selv en høiere, ideal Klarhed med sig, ved
hvilken Tanken føler sig lutret og frigjort, selv om der
for den umiddelbare Anskuelse Intet synes at være vundet.
Her gjælde Stenos Ord: pulchra sunt qvæ videntur,
pulchriora qvæ sciuntur. Denne ideelle Begrundelse
søger Philosophien dels i Metaphysiken dels i Erkjendelseslæren;
en saadan Klaring af Grundbegreberne og
de aandelige Grundlove er nødvendig, for at det gjennem
Erfaringen og Beskjæftigelsen med Phænomenerne
Vundne skal kunne faa philosophisk Betydning. Fechner
% --- p. 266 ---
\setcounter{footnote}{0}%
\opage{266}fordrer den umiddelbare Anskuelse fuldt tilfredsstillet
ogsaa paa Begrebsudviklingens Omraade, hvilket er en
ligefrem Modsigelse, da Abstraktionen her netop gaaer
ud over det umiddelbart Anskuelige. Philosophisk Tænkning
er, som allerede Platon lærte, en Afdøen fra det
ydre Liv for at vinde Ideen; men denne Afdøen er det,
Fechner frygter. Den umiddelbare Hengivelse til det
Phænomenale fører ham, da han overspringer hine philosophiske
Forundersøgelser, til en phantastisk Naturreligion,
der i mange af sine Træk er af høi poetisk
Skjønhed og udvikles med elskværdig Naivitet, og som
tillige paa sin Vis begrundes med en logisk Skarpsindighed,
der gjør det vanskeligere at gjendrive den, end det
ved første Øiekast synes. Modstanderne af Hypothesen
om Plantesjælen ville saaledes gjøre vel i at læse Fechners
Bevisførelse i Skriftet „Ueber die Seelenfrage“.
Men dels kommer Fechner herved til at lægge for stor
Vægt paa Noget, som ikke i sidste Instants er det Afgjørende,
dels føres han ogsaa ind paa en Slags Scholastik,
idet han med et forstandsmæssigt Apparat af Analogi-
og Induktionsslutninger vil bevise Noget, som
dog tilsidst beroer paa en sympathetisk Phantasitilfredsstillelse.


Jo mere Fechner søger at føre Alt tilbage til Bevidstheden
som sidste Holdepunkt, hvorfra ogsaa al Erfaring
udgaaer, des nærmere synes det og at maatte
ligge ham at fordybe sig i Bevidsthedens Væsen og at
søge dens Grundlove og dens Betydning i theoretisk og
praktisk Henseende. Mangelen af en saadan Analyse
viser sig i den Modsætning, der hos Fechner stedse bliver
tilbage mellem Aanden og Phænomenerne. Han opfatter
Aanden (synechologisk) som den, der sammenfatter
og sammenknytter Phænomenerne; som disses høieste
Enhed er Aanden den sidste Grund til Alt. Men Phænomenerne
staa uforklarede af Aandens Væsen, dykke
op deri som spredte Punkter, hvis Oprindelse ikke under%
% --- p. 267 ---
\setcounter{footnote}{0}%
\opage{267}søges. Derfor kommer Fechner til at vakle mellem en
umiddelbar Identitet og en ligesaa umiddelbar Dualisme:
Alt er én Substants --- men denne kan sés fra to aldeles
forskjellige Standpunkter. Navnlig staaer Atomlæren,
ved hvilken Fechner har indlagt sig saa megen
og almindelig anerkjendt Fortjeneste af Naturphilosophien,
i et aldeles tilfældigt Forhold til hans øvrige Anskuelser.
Atomerne skulle kun være de faste Punkter
for Anskuelsen paa den exakte Erfarings Omraade. Men
dette er endnu ingen metaphysisk Forklaring; disse faste
Punkter blive, saa længe de staa uopløste, til mørke
Punkter, der træde saa meget des stærkere frem i det
Lyshav, i hvilket Fechner svælger. Hans Antipathi mod
den monadologiske Anskuelse og hans hele pantheistiske
Retning have hindret ham i selv at gjøre sin Atomlære
philosophisk frugtbringende. Hos Lotze og Hartmann
ville vi finde Forsøg i denne Retning.

Fechner vil undgaa den „transscendente“ Philosophis % PRINTED AS IS (Philosophis: the image, for Philosophies; O177)
Søgen efter den dunkle „Ting i og for sig“ bag
Phænomenet. Og dog spøger denne „Ding an sich“ i
hans egen Anskuelse, især i Antydningen om et Grundvæsen,
hvortil den dobbelte (psychiske og physiske) Række
af Phænomener viser tilbage. Vel maatte efter hans
Betragtningsmaade Aanden, den høieste Enhed og Sammenhæng,
ogsaa være det Iogforsigværende. Men paa Grund
af Manglen paa metaphysisk Afklaring dukker for hans
af Anskuelsen umiddelbart beherskede Tanke stedse paa
Ny den Mulighed frem, at der dog kunde ligge Noget
bag Bevidsthed og Phænomener. Deraf hans stedse
gjentagne Polemik mod den nyere Philosophi, en Polemik,
der altsaa i Grunden gjælder hans egen af Tanken
endnu ikke gjennemtrængte Anskuelse.

Dog antyder Fechner selv det Punkt, hvor det egentlig
philosophiske Grundlag, vi savne hos ham, maatte
have sit Sted. Han udtaler (Atomenlehre. 2te Ausg.
p. 149), at naar man ad Analogiens og Induktionens Vei
% --- p. 268 ---
\setcounter{footnote}{0}%
\opage{268}gaaer ud over Erfaringen for at finde den Afslutning,
Tanken søger, bestemme selve Betingelserne for at vinde
et Sidste og Høieste de fundne Begrebers Form, medens
Erfaringen maa lede til at bestemme deres Indhold.
Hvad han her kalder „Betingelsen for at finde et Sidste
og Høieste, er de Væsensbestemmelser, uden hvilke
det Absolute ikke vilde være absolut (kfr. Weisse). Uden
en forudgaaende Undersøgelse af disse Væsensbestemmelser,
er det umuligt at anvende Analogien og Induktionen
paa Philosophiens Omraade; man vilde jo ellers,
hvilket ligger i Fechners egne Ord, ikke kunne vide,
om det virkelig var det Absolute, det Sidste og Høieste,
man naaede til. ---

Fechners positive og direkte Fortjeneste ligger i
hans grundige og skarpsindige Analyse af Grændsebegreberne
mellem Erfaringsvidenskab og Philosophi (navnlig
i hans Atomlære og Psychophysik). Men ogsaa indirekte
have hans aandfulde og begeistrede Skrifter deres store
Betydning ved at hævde det Anskueliges Værd i Modsætning
til den abstrakte Spekulation og ved at opponere
mod Spiritualismen og Materialismen. ---

3. \emph{Eduard von Hartmann} (f. 1842). --- For Fechner
var Alt Sjæl og Bevidsthed; en evig Klarhed og Harmoni
lyste for ham gjennem hele Tilværelsen. Vi synes da i
Forfatteren af „Philosophie des Unbewussten“, et nyt
Værk, der har vakt usædvanlig megen Opsigt (1869 udkom
første, 1872 fjerde Oplag), at have en fuldstændig
Antipode til Fechner. Thi Hartmann gjør, som allerede
Titelen paa hans Værk antyder, det Ubevidste til Principet
for sin Verdensanskuelse. Og dog have begge ikke
faa vigtige Berøringspunkter. Ikke blot lægge begge
Erfaringen og Induktionen til Grund for Spekulationen
og finde i de psychologiske Grundbestemmelser Nøglen
til Forstaaelsen af Tilværelsens Væsen; men netop ved
at begge opfatte deres tilsyneladende ganske modsatte
% --- p. 269 ---
\setcounter{footnote}{0}%
\opage{269}Principer (Bevidsthed --- det Ubevidste) paa extrem og
abstrakt Vis, mødes de tilsidst i samme Knudepunkt.

Fechner stødte under sine psychophysiske Undersøgelser
paa det Ubevidste, idet han paaviste, at naar
Indtrykket ikke overskred en vis Grændse, kom Fornemmelsen
ikke til Bevidsthed. Men en saadan ubevidst
Fornemmelse er det ikke, Hartmann mener, naar han
philosopherer over det Ubevidste. En ubevidst Fornemmelse
i Fechners Betydning er efter hans Mening ingen
virkelig Fornemmelse, er noget rent Negativt. Derimod
er det Ubevidste for ham noget Positivt, en virkelig Forestilling,
det vil sige: et idealt Indhold, og en virkelig
Villie, det vil sige: en Stræben efter at omsætte dette
ideale Indhold i Virkelighed. --- Motivet til at opstille det
Ubevidste (Enheden af ubevidst Forestilling og ubevidst
Villie) som Princip har Hartmann fundet i den tidligere Philosophi.
Denne var væsenlig en Gjennemarbeiden af Bevidstheden
og dens Indhold. Dette Vinbjerg er nu saa
gjennemarbeidet, at Tiden er kommen til at grave ned
i den underjordiske Verden for at finde de Skatte, den
skjuler. Mange Gaader, der tidligere forbleve uløste,
ville da maaske ses i et klarere Lys. Allerede de tidligere
Philosopher have paa mange Punkter forberedt
Udviklingen af det Ubevidstes Begreb. Først og fremmest
Leibnitz, der udtrykkelig skjelnede mellem ubevidst
og bevidst Forestilling (perception og apperception),
hævdede, at Sjælen godt kunde tænke uden at være sig
det bevidst, og paaviste de ubevidste Forestillingers Betydning
for Erkjendelse og Handling. Kant viste nøiere,
hvorledes Menneskeaanden anvender Tankeformerne (Kategorierne)
længe førend den bliver sig bevidst, at det
er paa dem, al Erfaring og Erkjendelse beroer. Fichte
lader Universet fremgaa af Subjektets ubevidste Produceren,
der gaaer forud for enhver bevidst Reflexion, og
Schelling gjør denne Produceren til en Virken af det
Absolute, „det evig Ubevidste, den evige Sol i Aander%
% --- p. 270 ---
\setcounter{footnote}{0}%
\opage{270}nes Rige, der skjuler sig i sit eget uformørkede Lys.“
Hegels logiske Ide, hvis Liv pulserer i Alt, er den ubevidste
Tanke, der kun i det endelige Subjekt fremtræder
under Bevidsthedens Form, Den har sit Modstykke i % PRINTED AS IS (Form, Den: the image, a comma for a point; O181)
Schopenhauers ubevidste Villie, som for ham er det eneste
metaphysiske Princip, medens Forestilling og Tanke
kun ere forsvindende Hjernephænomener. --- Men ikke
blot Philosopher, ogsaa Naturforskere ere ved deres
Undersøgelser blevne førte til at fremhæve det Ubevidstes
Begreb. Saaledes erklærer Physiologen Carus, at
Nøglen til Erkjendelsen af det bevidste Sjælelivs Væsen
ligger i det Ubevidstes Region. Og den berømte Helmholtz
har som væsentlig Faktor i den sandselige Iagttagelse
fundet en ubevidst logisk Slutning, der i sit
Væsen er beslægtet med den bevidste Slutning, forsaavidt
ogsaa denne foregaaer paa umiddelbar og spontan Vis.

Naar nu dette Princip, der saaledes synes at trænge
sig frem, skal udvikles og retfærdiggjøres, kan dette kun
ske ved en Forening af den induktiv-empiriske og den
deduktiv-spekulative Methode. Hin vilde i og for sig
give en god Grundvold, men naaer ikke til de sidste
Principer; denne gaaer ud fra Principerne, men naaer
ikke til den faktiske Virkelighed. Derfor maa man gaa
ud fra Erfaringen og vise, hvorledes denne ved at fortolkes
ret mødes med Spekulationen. Hartmann gjør
derfor til sit Motto: spekulative Resultater ifølge induktiv-videnskabelig
Methode.

Først maa det fremhæves, at det Ubevidste er Enhed
af Forestilling og Villie. Hegel og Herbart hævde
kun Forestillingen (Tanken, Ideen) og betragte Villien
kun som Forestillingernes eller Idemomenternes indbyrdes
Spænding. Paa den anden Side hævder Schopenhauer
Villien som det eneste Reale og betragter Forestillingen
som Illusion. Men uden Forestilling er der ingen Villie;
thi enhver Villen maa ville Noget, have et Indhold, og
dette Indhold kan kun være givet som Forestilling. Om%
% --- p. 271 ---
\setcounter{footnote}{0}%
\opage{271}vendt behøver Forestillingen Villien for at være virkelig.
Villien er selve Overgangen fra Mulighed til Virkelighed,
i sig selv en tom Form, der finder sin Udfyldning i
Forestillingen, som paa sin Side er aldeles indifferent,
en mundus intelligibilis, der kun ved Villien føres over
i virkelig Existents. Denne virkelige Existents er det,
Erfaringsvidenskaben godtgjør.

--- I Almindelighed tager man Villie i snevrere Betydning,
om den bevidste Selvbestemmelse. Men om Selvbestemmelsen
bliver bevidst eller ikke, er en Omstændighed,
der ikke vedrører Villiens Væsen som saadant,
men kun beroer paa, om den gaaer gjennem Hjernen
eller ikke. Ogsaa fra de lavere Nervecentra (Rygmarv
og Ganglier) kan der udgaa virkelige Villiesyttringer,
som ere ubevidste, forsaavidt de ikke opfattes af det
Centralorgan (Hjernen), hvorpaa det hele Individs Bevidsthed
beroer. De lavere Nervecentres Virksomheder ere
Mere end blotte Reflexvirkninger. Erfaringen godtgjør
Tilstedeværelsen af en virkelig Affekt og en gjennemført
Konsekvents, der vise os ud over den rent mechaniske
Sphære. Saaledes kæmpe f. Ex. de to Halvdele af en
australisk Myreart med Raseri mod hinanden; og den
hovedløse Frø viser tydelige Spor paa Beregning og
Frygt. Den sammenlignende Anatomi fører os desuden
til at betragte Hjernen som en Samling af Ganglier, der
forbindes ved Nervetraade, og altsaa i sin væsentlige Indretning
er ensartet med Rygmarven og Gangliesystemet.
Ja, ogsaa i aldeles nerveløse Dyr som Polyperne viser
der sig tydelige Spor af Affekt og Villie. Naar vi altsaa
hos de høiere Dyr have baade lavere og høiere
Nervecentra, da er der Intet til Hinder for, at Noget
kan være bevidst i et lavere Centrum, uden at være
det i det høiere, der betinger hele Individets Bevidsthed,
og det viser sig altsaa nødvendigt at udstrække Villiens
Begreb ud over det bevidste Livs Omraade.

Men herved have vi endnu kun naaet det relativt
% --- p. 272 ---
\setcounter{footnote}{0}%
\opage{272}Ubevidste. At der gives en absolut ubevidst Villie, indses
derimod, naar man undersøger, hvorledes den
vilkaarlige Bevægelse sker. Forat den bevidste Villie skal
føre til virkelig Bevægelse, maa der ske en Impuls i et
vist Punkt i Hjernen; men hvilket Punkt dette er, véd
Bevidstheden selv ikke. Den eneste mulige Sammenhæng
er den, at den bevidste Villie, f. Ex. til at bevæge
Fingeren, fremkalder en ubevidst Villie til at bevæge den
Nerveende i Hjernen, som bevirker Fingerens Bevægelser;
og denne ubevidste Villie forudsætter da igjen en
ubevidst Forestilling om hin Nerveendes Beliggenhed.
Da denne Forestilling og denne Villie, der danne de
de nødvendige Mellemled mellem den bevidste Villie og % PRINTED AS IS (de nødvendige Mellemled: the image, de doubled over the line turn (de|de); O184)
den organiske Bevægelse, ikke komme til Bevidsthed i
Hjernen, hvor Processen foregaaer, kunne de heller ikke
komme til Bevidsthed i andre Centra, og ere altsaa absolut
ubevidste.

Ogsaa i Instinktet lægger en ubevidst Villie sig for
Dagen. Instinktet kan ikke være en Følge af den legemlige
Organisation; thi det kan være forskjelligt ved
samme Organisation og omvendt. Ligesaa lidt beroer
det paa en engang for alle indplantet Hjerne- eller Aandsmekanisme;
thi saa kunde ikke Midlerne variere, medens
Formaalet bliver det samme. Heller ikke kunne Instinkthandlingerne
fremgaa af bevidst Overlæg; ikke blot forekomme
de hos de laveste Dyr, men de foretages med
umiddelbar Sikkerhed og forudsætte Kundskab om det
Fremtidige og i det Hele om Ting, Dyret ikke kan
vide, en Slags clairvoyance. Saaledes danne f. Ex. Hanlarverne
af visse Insekter sig ved Forpupningen en Hule,
der er dobbelt saa stor som de selv, for at der kan
være Plads til det udviklede Insekts lange Følehorn.
Instinktet kan kun opfattes som en bestemt Villen af
et Middel til et ubevidst villet Øiemed. Kun ved denne
Opfattelse forstaa vi, hvorledes Instinktet saa inderlig
kan være Et med Individets hele Væsen. Dette ses
% --- p. 273 ---
\setcounter{footnote}{0}%
\opage{273}især ved Driften til Selvopholdelse og til Artens Forplantning.
Naar man sønderriver Larvens Spind, udbedrer
den det igjen og bliver ved dermed, indtil den
døer af Udmattelse. En Fugl, som man bestandig berøvede
et Æg, vedblev at parre sig og lægge et nyt Æg,
indtil man tilsidst (ved det 29de) fandt den død paa
Reden. En saadan absolut Selvopofrelse er kun forstaaelig,
naar en almen Magt driver Individet til at opbyde alle
sine intellektuelle og physiske Kræfter for at tjene den.
Men selve dette Formaal bliver Individet sig ikke bevidst,
saalænge det beherskes af Instinktet; ellers vilde dets
Egoisme sætte sig op imod det Almenes Herredømme.

Reflexbevægelserne ere de underordnede Nervecentras
Instinkthandlinger. Ligesaa kunne ogsaa den naturlige
Lægekraft og den organiske Nydannelse føres tilbage
til en ubevidst Virken. Der kan kun være en
Gradsforskjel mellem Aarsagen til, at Fuglen istandsætter
sin ødelagte Rede eller Edderkoppen sit sønderrevne
Væv, og Aarsagen til at Sneglens Hus og Fuglens Fjederdragt
igjen fornyes efter en Beskadigelse. Vi have i
alle disse Phænomener kun Midler til et Øiemed, som
ubevidst forestilles og villes, og Forskjellen mellem dem
beroer kun paa, at de forskjelligartede ydre Impulser
fremkalde forskjelligartede Reaktioner.

Instinktet godtgjør ogsaa sin Magt i den menneskelige
Aand. Saaledes i Dødsfrygten, Undseelsen, Medfølelsen,
Elskovsdriften o. s. v. En i sig selv bevidst
Følelse kan grunde sig paa en ubevidst Villie og ledsages
af ubevidste Forestillinger; herpaa beroer det Uklare
og Uudsigelige i Følelsernes Væsen. Ja, selve de Forestillinger,
der fremkalde Følelsen af Lyst eller Ulyst,
kunne være os ubevidste, idet de kun komme til Bevidsthed
i lavere Nervecentra, hvor de vække en Følelse,
der ledes til Hjernen, uden at den bevirkende Aarsag
naaer til denne. Tydelig viser dette sig ved Hypochondri
og lignende Tilstande. Vi kjende i det  Hele ved
% --- p. 274 ---
\setcounter{footnote}{0}%
\opage{274}vor Handlen kun det bevidste Motiv, men ikke den
Maade, hvorpaa Villien reagerer mod Motivet og derved
gaaer over i virkelig Villen. Deraf kommer det, at vi
kun gjennem praktisk Erfaring kunne lære vor egen
Karakter at kjende og ofte overraskes ved de Indblik,
vi derved gjøre i vort eget Væsen. Villiens Værksted
ligger i det Skjulte, og kun indirekte naaer vort Blik
derned. Naar man mener at være sig Villien umiddelbart
bevidst, saa forvexler man den enten med dens
Motiv eller dens Følge (Handlingen) eller med de Følelser,
der ledsage den. Ved denne den egentlige Villies
absolut ubevidste Karakter bliver det forstaaeligt, hvorledes
Tænkere som Spinoza, Hegel og Herbart have kunnet
opløse dens Væsen i Tanke og Forestilling, medens
paa den anden Side den mest dilettantiske af de nyere
Philosopher, Schopenhauer, ligesom den almindelige
Menneskeforstand, mente at være sig Villien umiddelbart
bevidst.

I den sandselige Iagttagelse foregaaer der en ubevidst
psychisk Proces, der leder til Anskuelsen af Gjenstanden
som en ydre; uden en saadan vilde vi ikke
komme ud over den subjektive Fornemmelse, da denne
i sig selv ikke indeholder nogen Grund til at gaa ud
over Subjektets Sphære. Det er denne Proces, der, naar
Subjektet bliver sig philosophisk bevidst, viser sig at
bero paa et Kausalitetsforhold mellem det og Gjenstanden.
Det Ubevidste virker ligeledes i Ideassociationen,
idet det som Reaktion mod den bevidste Interesse lader
netop den Forestilling dukke op, som behøves, men
som Bevidstheden ved udtrykkelig Søgen maaske netop
vilde hindre sig i at finde. Alt i Sandhed Apriorisk er
sat af det Ubevidste og er kun som Resultat Gjenstand
for Bevidstheden; idet Bevidstheden nemlig reflekterer
over det forefundne Indhold, erkjender den a posteriori
det ubevidst virksomme Aprioriske. Al Mystik bestaaer
% --- p. 275 ---
\setcounter{footnote}{0}%
\opage{275}i Bevidsthedens Opfyldelse med et Indhold derved, at
dette uvilkaarlig dukker op af det Ubevidste.

Vi træffe fremdeles det Ubevidstes Virken overalt,
hvor en Flerhed virker sammen til et fælles Maal, i en
og samme Retning, og hvor saavel bevidst Aftale som
umiddelbar Mechanisme er umulig. Saaledes ved Sprogets
Dannelse, der kun er tænkelig under Forudsætning
af et Masseinstinkt af den Art, vi træffe det i en Bisværm,
en Myre- eller Thermitstat. Saaledes i Historien, hvor
de Enkelte gjennem deres bevidste, endelige Formaal
ubevidst arbeide i det hele og store Verdensformaals
Tjeneste. Jeg vil ubevidst noget ganske Andet end hvad
min Bevidsthed troer udelukkende at ville. Hvad er
Skjebne eller Forsyn Andet end det Ubevidstes, det historiske
Instinkts Herredømme i Menneskenes Handlinger,
saalænge deres bevidste Forstand ikke er moden
nok til at gjøre Historiens Formaal til sine egne?

Resultatet er altsaa, at det er det Ubevidste, der
danner og opholder Organismen, leder dens Bevægelser
paa hensigtsmæssig Vis og danner Mellemleddet mellem
den og den bevidste Villie. Det giver i Instinktet ethvert
Væsen hvad det behøver til Selvopholdelse, og
hvortil den bevidste Tænkning ikke er tilstrækkelig.
Saaledes giver det Mennesket Instinkt til Forstaaelse af
den sandselige Fornemmelse, til Sprog- og Statsdannelse
o. s. v. Det opholder Arterne ved Forplantningsdriften
og Moderkjærligheden, forædler dem ved Parrings- og
Kvalitetsvalget og fører gjennem den historiske Udvikling
Menneskeslægten fremad mod den høieste Fuldkommenhed.


Men endnu have vi kun beskjæftiget os med organiske
og aandelige Phænomener. Herved kunne vi ikke
blive staaende. Den moderne Naturvidenskab lærer os
(kfr. Fechner) at opfatte Materien som et System af
Atomkræfter. Hvert enkelt Atom kan ingen Masse
have, ligesaalidt som man kan tale om Talstørrelsen  af
% --- p. 276 ---
\setcounter{footnote}{0}%
\opage{276}en Ener. Holde vi os altsaa til Atomkraften som det
sidste Element, hvortil Analysen af det Virkelige fører
os, da er det klart, at i Atomkraftens Stræben kan der
kun findes bestemt Sammenhæng og Konsekvents, naar
denne Stræbens Indhold og Maal paa ideal Vis er forud
givet som en logisk Nødvendighed, hvoraf den ydre faktiske
Nødvendighed i Atomkraftens Virken under de
forskjellige Omstændigheder udspringer. Hvad er altsaa
Atomkraftens Stræben Andet end en Villie, hvis Indhold
er den ubevidste Forestilling om hvad der tilstræbes?
Kraftens Begreb gaaer saaledes med Nødvendighed
over i Villiens Begreb. Kraften i sin sædvanlige Betydning
(om de elementære Kræfter) bliver først forstaaelig
ved Villien. Saaledes er da Materien opløst i Forestilling
og Villie, og den radikale Modsætning mellem Aand
og Materie er ophævet. Forskjellen mellem dem beroer
kun paa en høiere eller lavere Form for det samme
Væsens, det evig Ubevidstes Aabenbaring. Materie, organisk
Dannelse eller Instinkt, og Bevidsthed ere tre forskjellige
Virkeformer for det Ubevidste, som udgjør Verdens
inderste Væsen.

Det Problem fremstiller sig da, hvorledes Bevidsthedens
Oprindelse af det Ubevidste er mulig. I det
Ubevidste ere Forestilling og Villie Et; Intet villes
uden at forestilles, Intet forestilles uden at villes. Bevidstheden
derimod indeholder Muligheden af Tankens
Emancipation fra Villien. Den ubevidste Aandsvirksomhed
er den oprindelige; først i og ved den normale Hjernefunktion
bliver den Bevidsthed. Den materielle Indvirkning
paa den ubevidste Aand bestemmer Forestillingens
Indhold; men dette Forhold kan i sig selv ogsaa tænkes
ubevidst. Bevidsthedens Form opstaaer ved Villiens
Opposition mod den ved den organiske Funktion producerede
Forestilling. Uden Modsætning ingen Bevidsthed.
Men Bevidstheden hører ikke nødvendig til Forestillingens
Væsen; den er kun en særegen Form for dette.

% --- p. 277 ---
\setcounter{footnote}{0}%
\opage{277}Bevidsthedens Opstaaen er Høidepunktet af det
Ubevidstes Virksomhed. Men denne er overalt den samme;
det Ubevidste reagerer mod de materielle Betingelser og
Virkninger, og disse forskjellige Reaktioner betegne de
forskjellige Trin i Tilværelsesformernes Udvikling. Det
Ubevidste griber Livet overalt, hvor det kan gribe
det. Saaledes virker det ved Avlingen paa den organiske
Spire og udvikler derved Sjælelivet som ved en Slags
Krystallisationsproces. I Analogi hermed maa ogsaa
Uravlingen, det organiske Livs Oprindelse tænkes: da
alle Betingelser vare tilstede, var den første organiske
Celle det Ubevidstes nødvendige og lovmæssige Reaktion.
Denne Faktor overser Darwin i sin geniale Theori,
der ene holder sig til Betingelserne, men stedse maa
forudsætte den indre producerende Kraft. Dette er anerkjendt
af Mænd som Wallace og Nägeli. Navnlig
sammenblander Darwin Varieren med Overgang i en
høiere Orden. Hvis Kvalitetsvalget virkelig havde den
uhyre Betydning, Darwin tillægger det, maatte alle høiere
Arter have større Livskraft end de lavere; men dette
er langtfra Tilfældet. Desuden forklarer Kampen for
Livet vel de physiologiske, men ikke de morphologiske
Ændringer. Kvalitetsvalget angiver altsaa et af de
mechaniske Midler, det Ubevidste benytter, men ikke
den tilstrækkelige Forklaring af alle de Former, til hvilke
det arbeider sig frem.

Alt Existerende er individuelt. Vi have allerede i
Atomet fundet det absolute Enkeltvæsen. Den levende
Organisme er en Verden af Individer. Indenfor det organiske
Individ træffe vi nemlig som relative Individer
først Cellerne, dernæst de forskjellige organiske Systemer
og de underordnede Nervecentra. Ud over disse
rager saa Individets almene Bevidsthed. Det høiere Individ
har bestandig de lavere til sin Forudsætning. Individualitetens
Begreb er altsaa relativt. Det have allerede
Aander som Spinoza, Leibnitz og Goethe indset,
% --- p. 278 ---
\setcounter{footnote}{0}%
\opage{278}og det bekræftes i fuldt Maal af den moderne Naturvidenskab.
--- Tilsidst gaaer Alt op i det absolute, Alt
omfattende Individ, det Ubevidste. Erkjendelsen af denne
det Ubevidstes Al-Enhed, vinder man kun ved at frigjøre
sig for det praktiske Instinkt, som fastholder Jeget
og tilspidser sig i den rene Selvbevidsthed, og hæve
sig til den Overbevisning, at Bevidstheden ikke hører
til Væsenet, men kun til Phænomenet: „Was an mir
\emph{Wesen} ist, bin ich nicht.“ Det Ubevidstes Philosophi
angiver den sande Forsoning af Monisme og Individualisme
ved at henføre hin til det metaphysiske, denne til
det physisk-reale Omraade. Flerheden har kun phænomenal,
ikke metaphysisk Betydning. Vel ere Rum og
Tid ikke, som Schopenhauer mente, blotte Former for
den subjektive Hjerneanskuelse; de ere ligesaavel Former
for den ydre Virkelighed. Men de gjælde kun for
det Ubevidstes Virksomhed ikke for dets Væsen i og
for sig. Virksomheden, ikke Væsenet er i Rum og Tid.

Da de ubevidste Funktioner ere tilstrækkelige til
Forklaring af Phænomenerne, er det en ren Hypothese,
at det Absolute skulde blive sig disse bevidst, medens
vi ikke gjøre det. Om det Absolutes Forestillen kunne
vi ingen Forestilling gjøre os. Mod hin Hypothese strider,
at Bevidstheden, forsaavidt vi kjende den, kun opstaaer
i og ved Hjerne og Ganglier. Desuden er det
Absolutes Tænken tidløs, hævet over al Succession; det
tænker Alt i Et og paa engang, er absolut fordybet i
Intuitionen og kan derfor ikke vende sig reflekterende
mod sig selv. Dog ligger den egentlig afgjørende Grund
mod Theismen i det Ondes Virkelighed, som under Forudsætning
af den bliver aldeles uforklarlig. Pessimismen
har afgjort Ret i, at denne Verden ikke kan være et
Værk af en bevidst Tanke, men kun kan være udsprungen
af en blind Villie.

Dog skal hermed ikke siges, at det Ubevidste i sig
selv er blindt. Tvertimod overgaaer dets Visdom al mu%
% --- p. 279 ---
\setcounter{footnote}{0}%
\opage{279}lig Bevidsthed (det er forsaavidt snarere at kalde overbevidst
end ubevidst), og det er den, der bestemmer
Skabelsens og Verdensprocessens Indhold. Det Ubevidstes
Philosophi er derfor i Principet enig med den
spekulative Theisme, der tilskriver Gud en over den
endelige, individuelle Bevidstheds Skranker absolut ophøiet
Bevidsthed; det Absolute er absolut Individualitet
med Villie og Forestilling. Men dog foretrækker
den af de anførte Grunde at betegne det Absolute ved
det negative Udtryk: det Ubevidste. Vel vilde dette
have mere Ret til Benævnelsen „Gud“ end Spinozas
Substants; men Gud er et Udtryk af særlig religiøs Oprindelse
og bør derfor saameget som muligt undgaaes i
Philosophien. Det negative Udtryk har vel sine Ulemper,
men det virker lutrende, og naar de anthropomorphistiske
Forestillinger engang have tabt deres Magt, saa
vil det sikkert allerede være afløst af et positivt.

Det Villende og Forestillende er én Stubstants. Det % PRINTED AS IS (Stubstants: for Substants; the reading of the Minnesota image (5x zoom))
er ikke en Blind, der bærer en Lam, men en Sund og
Hel, der rigtignok ikke kan se med Benene eller gaa
med Øinene. Denne Substants er den absolute Aand.
Men dennes oprindelige absolute Hvile er paa en grundløs
Maade ophævet ved den i sig selv blinde Villie, der
har taget Initiativet til at forlade den potentielle Væren.
Paa Grund af sin nødvendige Enhed med Forestillingen
(Ideen) river den denne med sig og fører den
over i Væren. Kun saaledes, paa blind, tilfældig Vis,
kan Verdens Tilværelse begribes. Thi skjøndt Schopenhauer
ikke har Ret i sin absolute Pessimisme og kun
paa sophistisk Vis kunde bevise, at denne Verden er
den slettest mulige, saa er det dog sikkert, at dens Tilværelse
er værre end dens Ikkeværen. Selv om en fuldkommen
Tilfredshed var mulig, vilde Livet kun være
ligesaameget værdt som Ikkeværen; men der gives ingen
virkelig Tilfredshed. Hvad man anpriser som Goder, er
kun privative Goder, det vil sige, man savner dem, naar
% --- p. 280 ---
\setcounter{footnote}{0}%
\opage{280}man ikke besidder dem, men deres Besiddelse skaffer
ingen positiv Lyst. Derfor opdagede Menneskeheden
ogsaa ved den gamle Verdens Slutning, at det var en
Illusion at vente Lykken her i Verden. Selv for den
gunstigst Stillede har Smerten Overvægten over Nydelsen.
Men fra Illusionens første Stadium gik man
over til det andet, idet man med Christendommen forsagede
al Lykke i denne Verden for at gjenfinde den i
overvættes Grad i et tilkommende, overnaturligt Liv.
Denne Illusions Magt knækkedes ved en ny Tilbagevenden
til det jordiske Liv og en ny Tillid til dets Gyldighed.
Men belært af Erfaringen venter man nu ikke
længer en umiddelbar Nydelse af Lykken; man haaber
en positiv Lyksalighedstilstand for Menneskeheden engang
i Fremtiden og anspænder alle sine Kræfter for
at bane Veien for den. Dette er Illusionens tredie og
sidste Stadium, i hvilket Menneskeheden endnu maaske
skal forblive i lange Tider. At ogsaa dette Stadium er
under Illusionens Magt, ses af, at det størst mulige
Fremskridt jo dog kun kan skaffe forøgede Midler og
Betingelser, uden at derfor en positiv Lykke bliver mere
mulig end den altid har været. Ja, den pessimistiske
Betragtning vil netop vinde i Kraft ved dette stigende
Misforhold. Men dette er ogsaa det Ubevidstes Hensigt
med den hele historiske Proces, og allerede med Bevidsthedens
Opstaaen. Verdensvirkeligheden oprandt ved den
blinde Villies Initiativ, som Forestillingen paa Grund af
sin Enhed med den var nødt til at følge. Kan det nu
lykkes at emancipere Forestillingen fra Villien, der i
sig selv er alogisk, men ved Følgen af hint Initiativ
er bleven antilogisk, saa vil Forløsningen være mulig.
Allerede Bevidstheden opstaaer ved en Seir over den
blinde Villie, der kommer i Opposition mod sig selv.
Og naar Illusionerne falde, den ene efter den anden,
kommer tilsidst det Punkt, hvor Mennesket gjør det
Ubevidstes Øiemed til sin Bevidstheds Øiemed og hen%
% --- p. 281 ---
\setcounter{footnote}{0}%
\opage{281}synsløst hengiver sig til dets Tjeneste for Verdensforløsningens
Skyld. Naar denne Udbrænding af Egoismen,
som gaaer i Et med Pessimismens Gjennemførelse, er
trængt saaledes igjennem, at Størstedelen af alle aandige
Væsener i Universet giver sig hen til den, da slynges
den aktuelle Villen igjen tilbage i det Intet, hvorfra
den kom, og Verdensprocessen ophører. Det Ubevidste
er da, hvad det altid har været; det kan Intet lære og
Intet erindre, og er efter Verdensprocessens Afslutning,
hvad det var før dens Begyndelse. ---

Det vil have sin Interesse til dette Omrids af Grundtankerne
i „det Ubevidstes Philosophi“ at føie en kort
Angivelse af Hartmanns Forhold til den tidligere Tænker,
han i sin hele Verdensanskuelse staaer nærmest,
nemlig Schopenhauer. Han korrigerer denne i hans
Hovedtanke ved at hævde den nødvendige Sammenhæng
mellem Villie og Forestilling. Schopenhauer betragtede
Tingenes Væsen som Villie, men ansaa Forestillingen
kun for et tilfældigt Skin. Men naar han tager en saa
afgjørende Bestemmelse bort af Villiens Begreb, saaledes
som vi ene kjende det, mister han Retten til at udstrække
det ud over den psychologiske Sphære. En
saadan Berettigelse kan Hartmann derimod godtgjøre,
da han viser, hvorledes ogsaa Forestillingens Begreb kan
udstrækkes til at gjælde for Tingenes Væsen. Derved
korrigerer han tillige Schopenhauers Opfattelse af Tanken
som blot Hjernephænomen; Tanken har ikke blot
psychologisk, men ogsaa metaphysisk Betydning. Derved
medføres igjen en anden Opfattelse af den historiske
Udvikling. Schopenhauer betragtede denne som noget
Tilfældigt og Betydningsløst, som Figurerne i Skyerne
eller Isblomsterne paa Ruden. Forløsningen fra Tilværelsens
Usalighed, der ogsaa for ham er det sidste
Formaal, skal enhver Enkelt gjennemføre ved Askese og
Ophævelse af Villien til Livet. Herimod indvender Hartmann,
at denne individuelle Befrielse ikke udretter Mere
% --- p. 282 ---
\setcounter{footnote}{0}%
\opage{282}end hvad der vilde ske ved Individets naturlige Død;
kun naar Slægten som saadan gjennem den historiske
Udvikling føres til Ophævelse af Villien til Livet, opnaaes
en virkelig Forløsning, en saadan som ogsaa ifølge
Christendommen Alskabningen længes efter og forventer
af dem, der have Aandens Førstegrøde. --- Selve den
pessimistiske Grundtanke har Hartmann altsaa beholdt,
og hans Udviklinger i denne Retning bære Præget af
at være udførte paa anden Haand.

I andre Dele af sin Lære stemmer Hartmann med
Fechner og Lotze. Saaledes i hans Atomlære og Monadologi.
Om Lotze minder særlig hans Lære om Atomerne
som Monader og deres Iboen i det Absolute. Med
Fechner faaer han ikke blot, som allerede antydet, et
Berøringspunkt, forsaavidt den absolute Bevidsthed, der
hos Fechner er Et og Alt, i Forhold til den endelige
Bevidsthed er Et med den absolute Ubevidsthed. Men
han har ham ogsaa til Forgænger i at føre Kraftens Begreb
tilbage til Villiens og dog stille sig i Modsætning
til Schopenhauer (Fechner Atomenlehre, 2te Ausg. p. 132
--- 135; kfr. ovenfor p. 262). Saa meget des mærkeligere
er det, at han selv aldeles ikke omtaler sin Anskuelses
Forhold til disse Tænkere, hvis Skrifter (Fechners
Atomlære og 1ste Del af Lotzes Mikrokosmus) forelaa
13 Aar før Hartmanns Bog udkom.

Det Tiltrækkende ved Hartmanns Værk beroer især
paa den klare og sindrige Benyttelse af det naturvidenskabelige
Stof. Kun ufuldkomment have vi i vor Fremstilling
kunnet antyde den Rigdom af specielle Analyser,
hvorpaa han grunder sine metaphysiske Principer. Hans
philosophiske Talent ligger aabenbart paa Naturphilosophiens
Omraade, hvorimod han i psychologisk Finhed
ikke kan maale sig med Schopenhauer og i Metaphysiken
ender i en eventyrlig Potentslære, der minder om
Schellings. Selve det Princip, han vil gjøre gjældende,
og hvorefter han har opkaldt sit Værk, har aabenbart
% --- p. 283 ---
\setcounter{footnote}{0}%
\opage{283}sin store Berettigelse, hvilken han selv har angivet ved
at vise, hvorledes det dukker frem paa engang i Philosophien
og i Naturvidenskaben. Men dels ved den Overvurdering,
et nyt Princip altid er udsat for, dels ved
Eftervirkninger fra Schopenhauer er der i hans Anvendelse
af dette Princip afgjørende Misligheder, som Kritiken
maa paapege.

Hartmann synes selv at glemme sin Kritik af Schopenhauer.
Paa det egentlig afgjørende Punkt i hans
Verdensanskuelse, ved Problemet om det Onde, ved hvis
Behandling han, som vi have set, i det Væsentlige slutter
sig til Schopenhauer, beholder han ogsaa dennes Modsætning
mellem Villie og Forestilling. Han finder, ligesom
Schopenhauer, Principet for Verdensudviklingen i
en blind og grundløs Villie til Livet, der ved at tage
Initiativet har revet Forestillingen ind med i Bevægelsen.
Her ligger altsaa den Forudsætning til Grund, at
den fra Forestillingen isolerede Villie er selvstændig og
kan bevæge sig uden Impuls fra Forestillingen. Men
dette strider baade mod Hartmanns egen Lære og
mod den psychologiske Iagttagelse, at enhver Drift
bestemmes ved en Forestilling, selv om denne er
nok saa dunkel eller ubevidst. Ved nærmere Eftersyn forudsætter ogsaa selv hin „grundløse“ Villie en
Forestilling som sit Motiv; den vil nemlig sig selv, er
efter Hartmanns Udtryk „velle volens“, og har altsaa
sit eget Væsen til ubevidst Indhold. Udvikles denne
Tankespire, da vil den Meningsløshed forsvinde, at det
blinde Tilfælde skulde være Virkelighedens Princip,
og det Mørke spredes, i hvilket den haardnakkede Pessimisme
lader Alts Oprindelse gaa tilbage. Vi slutte nu
ikke med Schopenhauer og Hartmann, at paa Grund af
det Ondes og Smertens Virkelighed og overvældende
Magt kan Tanken ikke være det sidste Princip for Alt;
men vi slutte, at da Tanken nødvendigvis er det Oprindelige,
den sidste Forudsætning, kan det Onde kun
% --- p. 284 ---
\setcounter{footnote}{0}%
\opage{284}være en metaphysisk Nødvendighed (kfr. Weisse og Fechner),
paa hvis Grundlag Livsudviklingen gaaer fremad
mod sit Maal.

Vi have set, at Hartmann ikke opfatter det Ubevidste
rent negativt, men erklærer sig enig i Principet
med den spekulative Theisme, der hævder Guddommens
Bevidsthed, naar man kun opfatter denne Bevidsthed
som aldeles hævet over det endelige Jegs Skranker.
Men denne Erklæring (som iøvrig først findes i Værkets
tredie Udgave) synes ikke at stemme med, at det Ubevidste
alligevel staaer i et rent negativt Forhold til Naturens
og Historiens Liv. Dette Liv har kun Betydning,
forsaavidt det Ubevidste gjennem det befrier sig
fra sig selv. Naar Udviklingen er fuldendt, har det Intet
vundet. Det kommer derved ogsaa i et negativt Forhold
til Bevidstheden og staaer lavere end denne, istedetfor
at være det „Overbevidste“, der naaer ud over
Bevidsthed og Personlighed og optager disse i sig efter
deres Væsens indre Kjerne. Thi paa sin Udviklings
Høidepunkt besidder Bevidsthedslivet en Rigdom og
en Fylde, der staaer i en grel Modsætning til det rene
Intet, det Ubevidste i sidste Instants er og vil. Hartmann
har i sin Opfattelse af det Absolute som det Ubevidste
ikke vist dialektisk Varsomhed nok\footnote[1]{Smlgn. derimod en Forfatter, med hvis Lære Hartmauns % PRINTED AS IS (Hartmauns: an inverted n, for Hartmanns; the note reader)
har Meget tilfælles, K. G. \emph{Carus}: Psyche. (Pforzheim 1837). p. 402: „Den høieste guddommelige Bevidsthed, Guds Aands Bevidsthed i og for sig, er af os kun at tænke som et saa Umaaleligt, Uendeligt og Altomfattende, at den for en saa aldeles betinget og til Endelighed bunden Bevidsthed, som den menneskelige, altid tilsidst vil falde fuldkommen sammen med selve det Ubevidstes Mysterium; men omvendt ligger netop derfor ogsaa det, vi have kaldt det Ubevidstes Guddommelighed kun i den høieste guddommelige Bevidstheds Umaalelighed og Ubegribelighed.“ Denne sidste Sætnings Indhold har Hartmann ikke ladet komme til sin Ret; derfor er hans „Ubevidste“ i Virkeligheden heller ikke det Absolute.}. Der bliver
bestandig en skarp Modsætning tilbage mellem Væsen
% --- p. 285 ---
\setcounter{footnote}{0}%
\opage{285}og Virksomhed, Væsen og Bevidsthed. Virksomheden
skal kun høre Phænomenet til; men da Phænomenets
Former slet ikke vedkomme Væsenet, forbliver dette
ganske uberørt af sin egen Virksomhed. Saaledes angaaer
Bevidstheden, der opfattes som blot Form, heller
ikke den egentlige Kjerne: „was an mir Wesen ist, bin
ich nicht.“ --- Vi ende her i en Dualisme, der gjennem
Schopenhauer viser os tilbage til Kant og hans Distinktion
mellem Phænomenet og Tingen i og for sig.

Det er derfor intet Under, at Hartmanns Forsøg
paa at forklare Bevidsthedens Oprindelse mislykkes.
Gjennem de billedlige Udtryk, som her ganske vist ere
uundgaaelige („Villiens Opposition mod den organisk
producerede Forestilling“), skinner det igjennem som
Forudsætning, som netop skulde bevises. Thi hvad skal
man tænke sig ved den organisk producerede Forestilling,
der fremkalder Villiens Opposition, Andet end den
første Fornemmelse, der overskrider Bevidsthedens Tærskel?
Før den har overskredet denne Tærskel, har den
jo ingen Grund til at vække Villiens Modstræben; og
naar den har overskredet den, kan denne Modstræben
vel udvikle og potentsere den, men er unødvendig til at
frembringe den.

Bevidsthedslivet har ifølge Hartmann kun den Betydning
successivt at emancipere Forestillingen fra Villiens
Herredømme. Menneskeheden skal tilsidst ende i
Pessimismens tusindaarige Rige, der danner Overgangen
til den fuldstændige Nirvana. Men søger man den sidste
Aarsag til Hartmanns Pessimisme, da ligger den
ligesom hos Schopenhauer i Sondringen mellem den individuelle
Lykke og den ideale Opgave. Skal Lykken
absolut findes udenfor den Opgave, der er stillet Individet
og Slægten, da har sikkert den pessimistiske Betragtning
Ret. Men den adskiller, hvad der ikke kan
og ikke skal adskilles. Individets og Slægtens Berettigelse
beroer kun paa deres eiendommelige Formaal og
% --- p. 286 ---
\setcounter{footnote}{0}%
\opage{286}Opgave, gjennem hvis Løsning deres eget Væsen udvikles,
og kun i selve denne Udvikling kan Lykken og Tilfredsstillelsen
søges. Hvad der arbeides og udvikles ad
denne Vei, det kan saa heller ikke være betydningsløst
og negativt; det maa, hvorledes man saa nærmere vil tænke
sig Sagen, bevare sin Realitet, fordi det bevarer og udvikler
selve det Væsentlige i Virkeligheden. Det Eudæmonologiske
og det Ethiske staa i en indre Sammenhæng,
og Pessimismens Betydning ligger i indirekte at
at godtgjøre Nødvendigheden af denne Enhed. --- % PRINTED AS IS (at godtgjøre: the image, at doubled over the line turn (at|at); O188)

4. \emph{Rudolph Hermann Lotze} (f. 1817). ---
Ligesom Fechner er ogsaa Lotze fra naturvidenskabelige
Undersøgelser gaaet over til philosophiske. Men heller
ikke hos ham staa disse to Studieretninger i ydre og
tilfældigt Forhold til hinanden; de slutte sig hos ham
endnu inderligere sammen end hos Fechner. Det var
en levende Interesse for Kunst og Poesi, der først førte
ham til Philosophien. Han følte sig tiltrukken af den
store Tankekreds, som Fichte, Schelling og Hegel havde
aabnet Indgangen til, men tilegnede sig den mere som
en charakteristisk Retning i den almindelige Dannelse
end som afsluttet Lærebygning. En bestemtere Indflydelse
paa hans Anskuelser fik Weisse, under hvem han
studerede i Leipzig, og som han endnu i sine seneste
Skrifter mindes med Taknemlighed og Pietet. I sine
„Streitschriften“ (1857), hvor han giver disse Oplysninger
om sin Udviklingsgang, erklærer han, at han ikke
blot i Almindelighed skylder Weisse mange Paavirkninger,
men særlig en saadan Belæring angaaende en snevrere
Kreds af Tanker, at han senere ikke har havt Grund
til at opgive den. Denne „Kreds af Tanker“ angaaer
aabenbart den almindelige metaphysiske og ethiske Grundanskuelse,
til hvilken vi kunne finde Spirerne i Weisses
Lære, men saaledes, at Lotzes klare og selvstændige
Aand og hans dybtgaaende naturvidenskabelige og kulturhistoriske
Indsigt har udviklet og omdannet disse Spirer
% --- p. 287 ---
\setcounter{footnote}{0}%
\opage{287}paa eiendommelig Maade. Navnlig henledte det medicinske
Studium, som han først havde valgt sig til Livsopgave,
tidlig hans Opmærksomhed paa den nødvendige
reale Sammenhæng, uden hvilken selv den høieste og
herligste Ide kun er en tom Phantasi. I disse to Yderpunkter,
Idealet og de faktiske Betingelsers System, satte
hans Tanke sig fast for ud fra dem at opsøge den Enhed,
som Spekulationen havde troet at kunne bemægtige
sig paa engang ved et dristigt Greb. 22 Aar gammel
kunde han optræde som Docent baade i det medicinske
og i det philosophiske Fakultet (senere blev han Professor
i Philosophi i Göttingen). Hans Skrifter gaa i begge
Retninger. I Aarene 40 udgav han paa den ene Side
en Metaphysik og en Logik, paa den anden Side en
Pathologi og en Physiologi, foruden nogle æsthetiske Afhandlinger.
Hans medicinske Skrifter vakte stor Opmærksomhed,
navnlig ved hans Stræben efter ad exakt
mechanisk Vei at søge Forklaring for de Phænomener,
man hidtil paa en uklar og mystisk Maade havde udledet
af en saakaldet Livskraft. Han søger at vise, at
de organiske Funktioner kun adskille sig fra de uorganiske
ved de elementære Kræfters eiendommelige Ordning,
men ikke ved en principiel Forskjellighed i de
virkende Kræfters Væsen. Hans gjennemtrængende Skarpsindighed
lod ham indse det Ubevislige i Meget, der
hidtil havde staaet som afgjort, og Mange betragtede
ham derfor som Skeptiker. I philosophisk Henseende
fik disse Arbeider navnlig Betydning, fordi Lotze gjennem
dem videre udviklede sin Overbevisning om Nødvendigheden
af at antage en Mangfoldighed af Elementer
i mechanisk Vexelvirkning som det reale Grundlag
for Virkeligheden. Den Maade, hvorpaa han udtalte
denne Overbevisning, lod Mange anse ham for Herbartianer;
men i det nys anførte Stridsskrift udleder han
den udtrykkelig af sine naturvidenskabelige Studier. Derimod
henviser han til Leibnitz som det philosophiske
% --- p. 288 ---
\setcounter{footnote}{0}%
\opage{288}Forbillede, der stod ham nærmest med Hensyn til den
videre Betragtningsmaade af det reale Grundlags Natur.
--- Om sine almindelige philosophiske Anskuelser erklærer
han sammesteds, at de nærmest kunne sammenstilles
med den ældre Fichtes Grundoverbevisning, at
kun det Godes Ide og dens Indhold er den tilstrækkelige
Grund til al Væren og Vorden, at de værdifulde
Formaals Verden (die Welt der Werthe) indeholder Nøglen
til Formernes Verden. Kun at man ikke som Fichte
maa indskrænke det Værdifulde til det strengt Moralske;
ogsaa det Skjønnes rolige Salighed, den hvilende og
affektløse Stemnings Hellighed og det Sandes indre
Konsekvents med al den Fred og Harmoni, den fremkalder,
hører med til Idealverdenen. Ei heller maa man
mene stykkevis at kunne fordre eller paavise Naturens
forskjellige Dele som Grundlag for tilsvarende ideale
Momenter; Ideen lever ikke saaledes fra Haanden i Munden,
men Naturen er som en stor Husholdning, der følger
sine egne eiendommelige, af de enkelte Formaal uafhængige
Love og først som Totalitet viser sig knyttet
til Ideen.

Da Lotzes physiologiske Skrifter hævdede den mechaniske
Naturopfattelse, toge Materialisterne (der naturligvis
ikke læste hans philosophiske Skrifter) ham til
Indtægt. Herpaa gjorde hans „Medicinische Psychologie
oder Physiologie der Seele“ (1852) Ende. Vel gaaer
han heller ikke her tilbage til de sidste Forudsætninger,
men han viser dog, at Materiens Begreb ikke er en Forudsætning,
hvorved Tanken kan blive staaende, og at
Mechanismen, hvor nødvendig dens Begreb end er for
Naturvidenskaben, dog kun faaer sin Betydning ved at
være Form og Grundlag for en indre Livsfylde. Endelig
giver Lotzes Hovedværk „Mikrokosmus. Ideen zur Naturgeschichte
und Geschichte der Menschheit. Versuch einer
Anthropologie“ (3. Bind 1856—1864) en Udsigt over
hans hele Verdensanskuelse. Da han i en philosophisk
% --- p. 289 ---
\setcounter{footnote}{0}%
\opage{289}Psychologi vilde fremstille sin Lære om det menneskelige
Liv, stod det klart for ham, at vor Naturs dybeste
Kræfter kun kunne iagttages tydelig, hvor de fremtræde
i den menneskelige Udviklings store Sammenhæng. Derfor
blev Historiens Philosophi den nødvendige Fuldstændiggjørelse
af Psychologien; Betragtningen af det individuelle
Liv og af Menneskeslægtens Kulturhistorie forene
sig i Anthropologien. Støttet til dette rige Grundlag
søger han saa ved Værkets Slutning at besvare de afgjørende
Problemer om Tingenes Sammenhæng, om Gud
og Verden. --- Ogsaa i hans seneste Skrift Geschichte der
Aesthetik in Deutschland“ (1868) findes mange Bidrag
til hans Verdensanskuelse.

Som Forfatter indtager Lotze en fremragende Plads.
Klar og præcis Opfattelse af Kjendsgjerningerne forenes
hos ham med en digterisk Sympathi for det Virkelige
og med en energisk Overbevisning om det Ideales Gyldighed.
Det exakte, det poetiske og det philosophiske
Element sammensmelte paa eiendommelig Maade i hans
Fremstilling. Og denne er heri en tro Afbildning af
hans Anskuelse. Idealet er for ham ikke den i abstrakte
Former sig udviklende Ide, som i Spekulationen, ei heller
den ydre mechaniske Sammenhæng, som i Materialismen;
men hine Former saavel som denne Sammenhæng
ere kun de ydre Skaller for Livets evige Kjerne.
Og dog have de afgjørende Betydning. Man kan ikke
umiddelbart gribe Kjernen; Forsøgene derpaa føre kun
til at gribe Skallerne i Kjernens Sted.

„Sandheden,“ siger Lotze i Slutningen af sin Metaphysik,
„er ikke det Første, men betinges ved, at det
Godes Rige frembringer den som sin nødvendige Forudsætning.“
Dette oversaa Spekulationen. Thi Videnskaben
anser let sit eget Omraade, den tænkende Erkjendelse,
for det aandelige Livs Helhed eller dog for dets Høidepunkt.
Men Tanken forudsætter stedse et realt Indhold.
Tankens Væsen er i sig Kritik af det Givne. Ved den
% --- p. 290 ---
\setcounter{footnote}{0}%
\opage{290}kritiske. sammenlignende Virksomhed, hvorved Erkjendelsen % PRINTED AS IS (kritiske: the point after it, for a comma (kritiske. sammenlignende); the crop f:C449 and the Minnesota layer)
af Loven opstaaer, adskiller den menneskelige
Aand sig fra de lavere Sjæleformer, men ikke ved at
gaa absolut op i den rene Tanke. De oprindelige aandelige
Kjendsgjerninger, som Fornemmelsen af Farver og
Toner, Rumsanskuelsen, Følelsen af Lyst og Ulyst og
al æsthetisk og ethisk Vurdering af det Erkjendte, ere
ikke blot foreløbige eller mislykkede Forsøg paa at tænke,
men udgjøre det Stof, som Tanken ikke selv kan frembringe,
men kun sammenligne og bearbeide. Baade i
den antike og den moderne Philosophi savner Lotze
Anerkjendelsen heraf. Grækerne have grundlagt Logiken,
et storartet Værk, der vel tør stilles ved Siden af
den moderne Naturvidenskab. Men den græske Tanke
var tilbøielig til at sammenblande det Logiske og det
Metaphysiske, Sandheder som gjælde, og Ting som ere.
Fra Mythologien arvede Tænkningen derhos Troen paa
symmetriske og rhythmiske Former i Verdensløbet, og
denne Tro var stærkere end Overbevisningen om, at der
dog maatte være bestemte og nødvendige Love for de
enkelte Elementers Samvirken indenfor det hele Verdensløb.
Derfor sammenblandede Grækerne den ideale Tydning
og den kausale Forklaring, og denne sidste kom ei
til sin Ret. Endnu i den nyere Philosophi virker Oldtidens
Overvurdering af Logos. Realismen lærer os kun,
hvad der er, Idealismen kun, hvad der skal være; hverken
i det Reale eller i det Ideale fremdrages det, hvorved
de ere Mere end al Fornuft. I sidste Instants bestaaer
enhver Tings Væsen i den Ide, den skal virkeliggjøre,
og den sande Angivelse af Tingens Begreb kan
derfor kun vindes ved en dialektisk Udvikling af denne
Ide i dens Begrundelse ved de høieste Ideer. Men man
kan ikke ved enhver Enkelthed gaa tilbage til de første
Principer; de ere for vægtfulde til umiddelbart at kunne
anvendes.

% --- p. 291 ---
\setcounter{footnote}{0}%
\opage{291}Det er altsaa to Hovedindvendinger, Lotze reiser
mod den tidligere Philosophi. Den har ikke lagt tilbørligt
Vægt paa den mechaniske Sammenhæng, hvorved
Ideens Virkeliggjørelse ene er mulig, og den har ikke
opfattet Ideen i dens hele Fylde, som det absolute Formaal,
det høieste Gode. Erkjendelsens høieste Tendents
er teleologisk og indbefatter i sig en Tendents til genetisk
Forklaring og en Tendents til ideal Udtydning,
som ofte med Urette træde stridende op mod hinanden.
--- Af Lotzes eget metaphysiske Ungdomsarbeide ville
vi her kun fremdrage en Hovedtanke, til hvilken han
stedse kommer tilbage i sine senere Undersøgelser. ---
Efterat have undersøgt Værens Former (Grund, Aarsag,
Formaal) gaaer han over til Læren om Phænomenet.
Phænomenets Begreb fordrer ikke blot et Værende, som
ligger til Grund for det, men ogsaa et Værende, for
hvilket det lader sig tilsyne. Ikke blot det, der gaaer
for sig \emph{mellem} Væsenerne, men ogsaa det, som gaaer
for sig \emph{inden} i dem, er en sand og virkelig Tildragelse.
At Tingenes Virkelighed bestaaer i, at de aabenbare sig
i Phænomenet, vil altsaa sige, at de ere for hverandre.
Virkeligheden er ikke blot den ydre Gjenstand; ogsaa
vor Opfattelse af den hører med til det Virkelige, ligesom
Lyden først er virkelig, naar Luftsvingningerne ved at
paavirke Hørenerverne forvandles til Fornemmelser. Det
Aandige er ikke en overflødig Tilgift, der i det Høieste
skulde have den Opgave at afspeile den i sig selv fuldt
færdige Virkelighed. Den aandelige Verden er tvertimod
Universets væsentligste Bestanddel; den Maade, hvorpaa
Virkeligheden opfattes af Aanden, er den væsentligste
Del af Alt, som sker, uden hvilken Verden ikke vilde
være færdig eller afsluttet. --- Saaledes siger Lotze ogsaa
senere i sin Æsthetikens Historie: „Der gives ingen Skjønhed
uden for den Aands Følelse, der nyder og beundrer
den; men Tingenes Sammenhæng er saaledes  ordnet, at
% --- p. 292 ---
\setcounter{footnote}{0}%
\opage{292}den hos Aanden kan vække Bevægelsesformer, gjennem
hvilke hin Nydelse bliver den tildel. --- At Virkeligheden
i det Store er anlagt til at gjøre et saadant Sammentræf
muligt, --- at den virkelige Verdens Sammenhæng passer
til Aandens Modtagelighed, --- at Tingenes Forbindelser
ske og kunne ske i Former, hvis Indtryk vække Sjælens
Funktioner til harmonisk Virken, --- det at Verden
og Aanden saaledes ere for hinanden, er den store
Kjendsgjerning, som vi nyde i Følelsen af Skjønhed.“
--- Denne Harmoni mellem Tilværelsens Former har sin
sidste Grund i det høieste Formaal, og derfor har Metaphysiken
heller ikke sin Begyndelse i sig selv, men i
Ethiken. ---

Efterat have skildret Lotzes almindelige Opfattelse
af den philosophiske Erkjendelses Væsen og Betydning
ville vi nu med ham fordybe os i „Mikrokosmus“; dog
saaledes at vi foruden Skriftet af dette Navn ogsaa benytte
de vidtløftigere Undersøgelser i „Medicinische Psychologie“.
Gjennem de naturphilosophiske, psychologiske
og historiske Problemer vil han da tilsidst føre os
til den afsluttende metaphysiske og ethisk-religiøse Anskuelse.


Den oprindelige Opfattelse af Naturen saa alle dens
Phænomener besjælede af guddommelige Væsener, af
hvilke alt Liv og al Virken udsprang. Men selv efterat
denne barnlige Betragtningsmaade var bleven afløst af
den videnskabelige Forsken efter Tingenes Aarsager, vedligeholdt
dens Spor sig i Forestillingen om plastisk virkende
Drifter som Livets Grund i Naturen. Det var
navnlig de organiske Phænomener, til hvilke denne idealistiske
Opfattelse knyttede sig. Man gjorde sig ikke
ret Rede for, hvortil en saadan ubevidst virkende Drift
skulde nytte. Thi den vilde kun da indeholde Forklaringen
af Hensigtsmæssigheden i det Organiske, naar
den virkede ifølge en Fornuft, der overveiede Omstændighederne;
men dette er netop udelukket ved, at den
% --- p. 293 ---
\setcounter{footnote}{0}%
\opage{293}opfattes som ubevidst\footnote[1]{Denne Bemærkning vilde ikke, som det maaske kunde synes, ramme Hartmann, da dennes Anvendelse af „det Ubevidste“, idetmindste i Principet, netop har den mechaniske Sammenhæng til Forudsætning. Derimod rammer den saadanne Theorier som den yngre Fichtes.}. Videnskaben kan ikke søge
Forklaringen af det organiske Liv andensteds end i en
saadan lovmæssig Virken, som vi træffe i hele Naturen.
Den levendegjørende Magt, hvorpaa den spekulative
Naturphilosophi troede, opløses da her i en Mængde
Elementarkræfter. Den mechaniske Naturopfattelse bliver
nu herskende. Dens Grundtanke er, at Evnen til
at frembringe en bestemt Virkning aldrig fuldt og færdig
indeholdes i et enkelt Atom eller i et enkelt Legemes
Natur, men at ligesom Nødvendigheden af en Virken
overhovedet kun fremgaaer af det gjensidige Forhold
mellem flere Elementer, saaledes beroer ogsaa hvert Elements
Virkemaade og Størrelsen af dets Indflydelse paa
de andre Elementers Natur. Det Levende adskiller sig
ikke fra det Livløse ved en høiere, eiendommelig Kraft,
ved hvilken det skulde være unddraget den almindelige
mechaniske Sammenhængs Lov. Men det Levendes indre
Ordning har en saadan Form, at Elementernes naturlige
Kræfter under de ydre Betingelsers Indflydelse maa udvikle
en sammenhængende Række af Virkninger efter
de samme almindelige Love, efter hvilke ogsaa ellers
overalt den ene Tilstand pleier at følge paa den anden.
Saaledes er Legemet ikke længer en Klædning, Sjælen
væver sig selv, men et mechanisk System af ydre Betingelser
for Sjælens Liv, --- det Stykke af Omverdenen,
som staaer Sjælen nærmest.

Denne mechaniske Opfattelse faaer ofte en saadan
Magt over Tanken, at denne ikke kan frigjøre sig fra
den paa Omraader, hvor dens Gyldighed ikke længer er
absolut. Vi vænne os ved Betragtningen af Naturen i
den Grad til de middelbare Virkninger, som vi forklare
% --- p. 294 ---
\setcounter{footnote}{0}%
\opage{294}ved Sammensætning af enkelte Bidrag, vænne os til at føre
betydningsfulde Forskjelligheder tilbage til ubetydelige
Forandringer i ensartede Elementers Størrelse og Forbindelsesmaade,
at vi tilsidst tabe Forstaaelsen af alt
Umiddelbart og uvilkaarlig forudsætte et indviklet Maskineri
for Alts Opstaaen og Tilværelse. Al vor Viden
synes da at være en cognitio circa rem, en Kundskab
om de Betingelser, under hvilke et Phænomen træder
frem for os og forandres i sin Vexelvirkning med andre
Phænomener. Og dog besidde vi ogsaa paa visse Punkter
en cognitio rei, en Kundskab om Gjenstandens egen
væsentlige Natur. En saadan er mulig, hvor vi ikke
blot kunne iagttage en Gjenstand i dens ydre Forhold,
men have den i en saa umiddelbar Anskuelighed, at vi
baade i Følelse og Forestilling kunne gjøre os til Et
med den, sætte os ind i dens Tilstand og fornemme,
hvorledes en saadan Existents er tilmode ifølge sit eiendommelige
Væsen. En saadan positiv og umiddelbar
Anskuelse besidde vi kun af det Levende og Virksomme;
kun dette forstaa vi og gjennemtrænge vi. Materien som
saadan er en mørk Kjerne for os.

Den Kjendsgjerning, at vi kunne blive os selv bevidste,
forudsætter en saadan Enhed i vort Væsen, at
det bliver nødvendigt at antage et eiendommeligt sjæleligt
Princip, der træder i Vexelvirkning med den legemlige
Mechanisme. Ellers vilde man ende i den uklare
og selvmodsigende Antagelse af en Virksomhed uden noget
virkende Princip. Men Sjælen er ikke virksom uden
Paavirkning og Medvirkning af legemlige Funktioner. I
denne Vexelvirkning mellem Sjæl og Legeme ligger der
ingen større Gaade end i et hvilketsomhelst andet Aarsagsforhold.
Kun er man ofte tilbøielig til at opfatte
Kausalitetsprincipet ensidig, idet man ene dvæler ved
Virkningens Afhængighed af Aarsagen. Sætningen kan
ogsaa udsiges under den Form, at enhver Aarsag har
sin Virkning, hvilket betyder, at ethvert Element i
% --- p. 295 ---
\setcounter{footnote}{0}%
\opage{295}Virkeligheden paa sin eiendommelige Maade griber ind
i den hele Sammenhæng, i Overensstemmelse med de i
denne gjældende Love. Vor Naturkundskab bestemmer
Intet om det indre Baand mellem de virkende Elementer.
Standpunktet er ved en saadan cognitio circa rem
egentlig okkasionalistisk: vi søge kun Anledninger og
Betingelser. Dette Standpunkt maa vi ogsaa fastholde,
naar vi ville fremstille de første Grundtræk af Psychologien
og undersøge Forholdet mellem Sjæl og Legeme.
Først naar disse Grundtræk ere vundne, ville vi kunne
forandre Standpunktet og overvinde den Modsætning,
som har vist sig for os mellem den ydre Mechanisme
og det indre Liv, mellem cognitio circa rem og cognitio
rei.

Vi karakterisere først nærmere Sjælens Væsen for
derpaa at se den i dens Forhold til Legemet. --- Aandens
Enhed lægger sig for Dagen i, at den ikke blot
samler og sammenfatter sine forskjellige Tilstande til
en indre Helhed, men ogsaa ved en kritisk og sammenlignende
Virksomhed stræber at tyde det Enkelte og
anvise det dets bestemte Plads indenfor denne Helhed,
som derved bliver en ideel Verden. I denne finder da
Aanden et Gjenskin af sin egen Enhed. Selv i den rent
sandselige Opfattelse viser denne Eiendommelighed sig.
Hvad vi opfatte ved Sandserne, optage vi ikke blot som
et ligegyldigt Indhold, ei heller blot som det, der volder
Lyst eller Smerte, men vi tillægge det et eget Værd,
hvorved det udfylder sin Plads i Phænomenernes betydningsfulde
Sammenhæng. Allerede Sandsningen hæver
sig op over den rent mechaniske Nødvendighed og søger
at indordne sin Virken under en Livsorden, der ikke
umiddelbart er given i Naturen.

Den Enhed, der saaledes er gjennemgaaende i Sjælelivet, forklares ikke i de Anskuelser, der enten som Materialismen
udlede det Sjælelige som Produkt af det
Legemlige eller dog betragte Bevidsthedens Enhed som
% --- p. 296 ---
\setcounter{footnote}{0}%
\opage{296}Resultat af legemlige og sjælelige Kræfters Vexelvirkning.
Forsaavidt man her beraaber sig paa Analogien
med Kræfternes Parallelogram, da viser denne netop i
modsat Retning. Den diagonale Bevægelse fremkommer
nemlig kun, naar to Kræfter virke paa et og samme
Punkt; det er da dette Punkt, der bevæger sig efter
Diagonalen. Altsaa forudsættes her udtrykkelig en Enhed,
hvorfra den nye Kraft kan udgaa. --- Men, vil man
spørge, hvor finde vi denne Enhed i Sjælen? vi kjende
den jo kun som erkjendende, følende, villende o. s. v.
Lotze svarer hertil, at enhver Tings Væren bestaaer i
dens eiendommelige Indhold; udenfor eller bagved dette
Indhold har det ingen Væren. Vilde man fordre Oplysning
om, hvad der gjør dette Indhold til et Værende,
da vilde det være at spørge, hvorledes Væren bliver til.
Som det er Vandets Væsen at være flydende under nogle
Omstændigheder, dampformigt eller fast under andre,
saaledes er det ogsaa Sjælens Væsen under visse Betingelser
at fremtræde som forestillende, under andre som
følende, under andre igjen som villende, og dens eiendommelige
Enhed lider ikke herved.

Den nyere Psychologi har ofte ment at maatte hævde
Sjælens Enhed ved at føre dens forskjellige Funktioner
tilbage til en eneste. Saaledes betragtede Hegel alle
aandelige Phænomener som forskjellige, lavere og høiere
Former for den begribende Erkjendelse. Herbart saa
Sjælens Virken som Forestillen; Følelse og Villie vare
for ham kun Resultater af særegne Modsætningsforhold
mellem Forestillinger. Lotze indrømmer, at Forestillinger
fordetmeste ere Tilknytningspunkterne for Følelserne,
og at Drifterne igjen udvikle sig af disse. Men
deraf følger ikke, tilføier han, at Forestillingerne ogsaa
ere den hele og fulde Aarsag til de Følelser og Drifter,
de fremkalde. I Forestillingen selv ligger ingen Grund
til den eiendommelige Funktion, der fremtræder i Følelsen
af Lyst og Ulyst. Evnen til at frembringe denne
% --- p. 297 ---
\setcounter{footnote}{0}%
\opage{297}maa søges i Sjælen selv og træder i Virksomhed, naar
Sjælen ved Forestillingslivets Gang vækkes til at ytre
sig paa denne Maade. Paa samme Maade frembringes
Villiesytringen ikke af Følelsen, men udløses ved denne
af Sjælen selv. De tre Grundevner ere altsaa trinvise
Anlæg, der gjensidig vække hverandre til Virksomhed.
Medens ved en blot mechanisk Enhed Resultatet af
Vexelvirkningen mellem to Elementer er fuldkommen
bestemt ved Naturløbets almindelige Love og de givne
Omstændigheder, saa træder i det aandelige Liv Sjælens
Natur til som virkende Faktor, der medbetinger
Bevægelsens Resultat.

Naar Sjælen er et saadant oprindeligt Princip, saa
følger det af sig selv, at saa gjennemgribende den end
er betinget ved den legemlige Organisme, kan den dog
ikke umiddelbart være modtagelig for Indvirkninger fra
denne, men maa bestandig omforme disse paa eiendommelig
Maade. De legemlige Indvirkninger tilføre ikke Sjælen
færdige Fornemmelser eller Forestillinger, men kun Signaler,
som den skal oversætte i sit eget Sprog, ligesom
omvendt Sjælens indre Tilstande ikke som saadanne
overføres i de legemlige Organer, men blive Anledning
til, at disse fremtræde i Funktion. Villien véd Intet
om, ved hvilke Midler Armen sættes i Bevægelse, og
hvad Fornemmelsen angaaer, da har den nyere Videnskab
klarlig paavist Uensartetheden mellem den og den
ydre Paavirkning. Dette gjælder ogsaa med Hensyn til
vor Opfattelse af Tingene som indtagende hver sin Plads
i Rummet. Hvorledes opstaaer denne Rumsanskuelse?
Et er Gjenstandenes ydre Udstrækning, et Andet vort
indre Billede af denne Udstrækning. Dette indre Billede
er selv ikke udstrakt. Vi gjenfrembringe paa ideal
Maade Rumsforholdet i vor Anskuelse, det Extensive
omsættes til et Intensivt. De modtagne Indtryk udbrede
sig i Sjælen til en Rummets Verden, ikke fordi der
findes et virkeligt Rumsforhold indenfor Bevidstheden,
% --- p. 298 ---
\setcounter{footnote}{0}%
\opage{298}men fordi der i de af Indtrykkene fremkaldte sjælelige
Bevægelser bevares Forhold, der svare til Forholdene
mellem de ydre Gjenstande. Enhver Paavirkning faaer
paa Grund af det Punkt i Nervesystemet, hvor den finder
Sted, en eiendommelig Farve, en charakteristisk Bibestemmelse,
som kan kaldes dens Lokaltegn, fordi det
er den, der leder Sjælen til at henføre Gjenstanden,
hvorfra Indtrykket skriver sig, til et vist Sted i Rummet.
Lokaltegnene kunne ikke være af passiv Natur, men
maa tænkes som en Art Reflexbevægelse; tillige maa
de, for at en klar, geometrisk Rumsforestilling kan blive
mulig, udgjøre et System. Men Lokaltegnene frembringe
ikke selve Rumsanskuelsen; de tjene kun til at specificere
den. De føre Sjælen, der ifølge sin egen Natur
tragter efter Udfoldelse i Rummet af sit intensive Indhold,
til at anvende sin Rumsforestilling i Overensstemmelse
med Gjenstandenes Natur og gjensidige Forhold.
Oprindelig ubevidst og mechanisk kan den rumlige Lokalisering
ved Øvelse udvikles til høiere Grader. ---
Denne Lotzes Lære om Lokaltegnene, der paa naturlig
Maade udspringer af hans psychologiske Grundanskuelse,
har vundet megen Anerkjendelse, ogsaa hos Naturforskere
(f. Ex. Hemholtz). % PRINTED AS IS (Hemholtz: for Helmholtz; both copies)

Der staaer endnu to Spørgsmaal tilbage; det ene
angaaer Sjælens Sæde, det andet Sjælelivets Organer. ---
Hvad det første angaaer, gaaer Lotze ud fra den Sætning,
at enhver Ting er, hvor den virker. Sjælens Sæde
maa søges i Hjernen, som er den Del af Legemet, med
hvilken den staaer i umiddelbar Forbindelse. Vel tale
Mange om Sjælens Allestedsnærværelse; men de kunne
da ikke forklare, til hvad Nytte Nervesystemet er. En
bevidst Fornemmelse kommer først istand, naar Indtrykket
ledes fra Sandseorganet til Hjernen. Vel kan en
immateriel Substants, som mangler al Udstrækning, ikke
udfylde noget Rum; men der er intet til Hinder for, at
den har et bestemt Sted i Rummet, ud fra hvilket den
% --- p. 299 ---
\setcounter{footnote}{0}%
\opage{299}virker som fra et Kraftcentrum, og til hvilket alle fra
den ydre Natur stammende Paavirkninger maa forplante
sig for overhovedet at være til for den. --- Forestillingen
om et Sjæleorgan anser Lotze for tankeløs. Grunden
til Forestillingens, Følelsens og Driftens eiendommelige
Kvalitet maa søges i Sjælens egen Natur, og ingen legemlig
Virken formaaer at meddele den, skjøndt en saadan
er den nødvendige Betingelse for hine aandelige Bevægelsers
virkelige Indtræden. De legemlige Organer
arbeide i de høiere Aandsvirksomheders Tjeneste ved
at overlevere og forberede det Materiale, paa hvilket
Sjælen skal øve sine Kræfter; men deri bestaaer ogsaa
hele deres Betydning. Reflexbevægelsernes Betydning
for Rumsanskuelse og Bevægelse nøder os til Antagelsen
af Centralorganer for den rumlige Anskuelse og for Reguleringen
af Bevægelserne, hvorimod det er meningsløst
at antage Centralorganer for Forstandsvirksomheden,
den sædelige og æsthetiske Følelse o. s. v., da disse
Virksomheder bestaa i et ideelt Forhold mellem givne
Elementer, et Forhold, der ikke selv kan forklares mechanisk.
Derved bliver dog Sjælens Afhængighed af
det Legemlige ikke mindre. Den Inderlighed og Høide,
dens Virksomhed skal naa, betinges ved, at Sandseorganerne
og Centralorganerne nøiagtig og med tilstrækkelig
Intensitet tilføre den Indtryk udenfra. Jo rigere og
ædlere Midler, des ædlere Virksomhed. Den høieste
aandelige Virken er som Plantens Blomst; men Ingen
vil dog forlange, at Blomsten skal have sin særegne
Rod, forskjellig fra den, der bærer den hele Plante.

Hvis den methodiske Okkasionalisme, ved hvis Hjælp
denne Opfattelse af Forholdet mellem Sjæl og Legeme
er vunden, skulde være en virkelig Theori, vilde den
føre til en Dualisme mellem to kun ved ydre, mechanisk
Vexelvirkning forbundne Verdener, den sjælelige
og den legemlige, hver med sine eiendommelige Love.
Denne Dualisme gjælder det nu at komme ud over ved
% --- p. 300 ---
\setcounter{footnote}{0}%
\opage{300}at ophæve de methodologiske Forudsætninger, der laa
til Grund i det Foregaaende.

Al mechanisk Sammenhæng forudsætter en Flerhed
af Elementer, mellem hvilke den finder Sted. Derfor
er den moderne Physiks sidste Forudsætning en uendelig
Mængde af Atomer, som ikke kunne opløses ved de
os bekjendte Kræfter. Men ved denne Forudsætning kan
den metaphysiske Tænkning ikke blive staaende. Den
kan ikke, som Physiken, der kun har til Opgave at
undersøge selve Mechanismen, blive staaende ved Begrebet
om Materien som en mørk Kjerne, der bevæger
sig i et lyst Net af Relationer. Forestillingen om et Værende,
der ikke er til for sig selv, er selvmodsigende.
Et saadant Værende vilde da kun være for den Bevidsthed,
der opfatter det, og den ene Del af Verden vilde
altsaa være det blinde og livløse Middel for den andens
Øiemed. Kun levende Væsener tilkommer der en virkelig
Væren; alle andre Tilværelsesformer maa forklares
ud af det aandige Liv. Vi maa lægge en cognitio rei ind
i ethvert af de Elementer, som cognitio circa rem nøder
os til at antage. Ethvert af de Atomer\footnote[1]{Da Fechners „Atomenlehre“ udkom, kunde Lotze erklære (Göttinger gelehrte Anzeiger 1855 p. 1093 ff.), at han længe havde næret en saadan Anskuelse som den, Fechner i det nævnte Skrift udviklede. Hvad han savner i den moderne saavel som i den antike Atomlære, er en tilstrækkelig Grund til den absolute Fasthed, med hvilken Atomerne under enhver tænkelig Vexelvirkning dog skulle beholde deres indre Konstitution. Dette Spørgsmaal troer han ikke at kunne løse uden at antage kvalitativt forskjellige Atomer, en Antagelse, Oldtidens Atomistik benegtede paa Grund af sin antinominalistiske Retning, og som Fechner kun har forkastet af en uberettiget Kjærlighed til Simpelhed.}, Physiken
leder os til at betragte som faktisk uopløselige, maa
tænkes at bestaa af væsentlig forskjellige Urbestanddele,
mellem hvilke der finder et saadant Valgslægtskab Sted,
at det ikke overbydes af noget andet. Ved disse Ur%
% --- p. 301 ---
\setcounter{footnote}{0}%
\opage{301}bestanddele kan man ganske se bort fra al rumlig Udstrækning
og betragte dem som oversandselige Væsener,
der ud fra bestemte Punkter i Rummet ved deres Kræfter
beherske et vist Maal af Udstrækningen uden i egentlig
Forstand at udfylde det. Deres indbyrdes Stilling
betinges ved deres Vexelvirkning, og en Rumfigur vilde
herved beskrives ligesaa sikkert og bestemt, som hvis
de kontinuerlig opfyldte dens Indre. Til disse enkelte
Punkter maa vi tænke os de tiltrækkende og frastødende
Kræfter knyttede. I disse Centra, fra hvilke Kræfterne
skulle udspringe, maa man tænke sig et indre Liv, som
i sidste Instants betinger Kraftens Eiendommelighed.
Vel kunne vi, der ikke staa i Verdens Centrum, ikke
erkjende, hvorledes Kræfternes mechaniske Vexelvirkning
betinges ved Monadernes indre Væsen; men vi føres
med Nødvendighed til at opløse Materiens Begreb og
se den som Materialitet o: som Form for et oversandseligt
Reales Aabenbaring. Materiens Liv er saaledes et
dobbelt, et ydre og et indre. Den i det Uendelige delelige
Udstrækning er kun et Skin. Udstrækning kan ikke
være et enkelt Væsens Prædikat, men betegner et Forhold
mellem mange Væsener, der ikke selv ere udstrakte.
Ethvert af disse Væseners udelelige Enkelthed gjør det
muligt i dem at antage en Sammenfatten af de ydre
Indtryk i Fornemmelse og Nydelse. Kun herved bliver
hvert Enkeltvæsen delagtigt i Formaalet for det hele
Verdensliv, som bestaaer i Virkeliggjørelsen af Goder,
der kunne føles og nydes af Sjælen. Vi behøve ikke
med Leibnitz at tillægge enhver Monade en Forestilling
om sin Plads i Verden eller en Stræben efter at udvikle
sine Kræfter; men vi give ham dog Ret i, at ethvert
Enkeltvæsen er et Speil af Universet, idet det
som et eiendommeligt Midtpunkt i sin indre Fornemmelse
optager det, som den hele Verdenssammenhæng
virker paa dette bestemte Sted. Ingen Del af det Værende
er nu uden Liv og Sjæl; kun en Del af det, der
% --- p. 302 ---
\setcounter{footnote}{0}%
\opage{302}tildrager sig, nemlig de Bevægelser, der sætte det ene
Værendes Tilstand i Forbindelse med det andets, slynger
sig som en ydre Mechanisme gjennem det Besjæledes
Fylde og tilfører det Leiligheder og Opfordringer til
vexlende Udfoldelse af det indre Liv.

Dualismen mellem Sjæl og Legeme er nu forsvunden,
ikke ved paa materialistisk Vis at gjøre Sjælen til
et Stof, men omvendt ved at gjøre Stoffet til en med
det Sjælelige beslægtet Substants. Sjælen er kun den
ved sin Stilling begunstigede Monade, primus inter pares.
Udviklingen af den aandige Tilværelse i et monadisk
Element afhænger af gunstige Betingelsers Indvirkning,
der fremstille sig under Form af legemlig Organisation.
Kun det Element, der er stillet paa en saadan Organisations
gunstigste Punkt, modtager alle Indtryk saaledes
ordnede, og kan saaledes virke tilbage paa Legemets
modtagelige Dele, at det kan udvikle sig til en egentlig
Sjæl. Sjælen staaer nu ikke i Vexelvirkning med Materien
som Materie, men virker som oversandseligt Reale
nærmest paa de andre oversandselige Realer, der i sig
ere ensartede med Sjælen, men for os frembyde Skinnet
af den materielle Tilværelse, og paa samme Maade virker
det Legemlige tilbage paa Sjælen. Der bliver saaledes
en Vexelvirkning mellem ligeartede Led. Men
Sjælen smelter ikke sammen med de Atomer, der ligge
til Grund for den legemlige Organisation. Det levende
Væsen er et Samfund af mange Enkeltvæsener, der beherskes
af en fælles, i deres Indre begrundet Lov.

Det var ikke særlig psychologiske, men metaphysiske
Motiver, der førte os til denne spiritualistiske og
monadologiske Opfattelse. Men opløses herved ikke Alt
i en Uendelighed af Enkeltvæsener, der kun paa udvortes
Vis staa i Forbindelse med hverandre? Istedetfor
Dualismen mellem Sjæl og Legeme synes der nu at være
kommen en uforsonet Modsætning mellem en uendelig
Mængde Elementer. Med andre Ord, Sammenhængen
% --- p. 303 ---
\setcounter{footnote}{0}%
\opage{303}mellem Enkeltvæsenerne staaer endnu uforklaret. Vi
begyndte med et mørkt Punkt i et lyst Net af Relationer;
nu synes vi at være naaet til en Mængde lyse
Punkter i et mørkt Net. --- Hverken Enkeltvæsenerne
eller den mechaniske Sammenhæng er det Absolute.
Alle Enkeltvæsener ere kun indre Dele, særlige Indholdsbestemmelser
af et uendeligt Væsen. Ligesom vor Sjæls
indre Enhed afspeiler sig i forskjellige Former af aandig Virken, der ikke kunne reduceres til hverandre, men
hver for sig og dog i nødvendig Sammenhæng med hverandre
udspringe af Sjælens Natur, saaledes forholder den
uendelige Substants sig til Enkeltvæsenerne. Kun det
er, som har sin Plads i den evige Ides fornuftige Sammenhæng.
Den mechaniske Vexelvirkning er den Form,
ved hvilken det Uendelige sikrer ethvert Element dets
eiendommelige Virken i Forhold til de andre. Uden det
Absolute vilde Mechanismen være uforklarlig; kun ved
det kan den Kløft overvindes, der evig vilde skille de
enkelte selvstændige Elementer fra hverandre. Ligesom
i Sjælen Forestillingernes indbyrdes Forhold kun kunne
fremkalde eiendommelige Følelsestilstande ved at virke
tilbage paa selve Sjælens Natur, saaledes er i den uendelige
Substants enhver Paavirkning af et enkelt Element
tillige en Paavirkning af det hele Uendelige og
kun derved er en tilsvarende Virkning i et andet Element
mulig.

Saaledes er Mechanismens Ide paavist i det Absolute.
Hvad der for den spekulative Idealisme faldt sammen
med den rene Endelighed og Udvorteshed, og derfor
i sidste Instants stod som et Negativt, et Ikkeværende,
det ses her som den nødvendige Betingelse for det Uendeliges
Udvikling. Medens Idealismen satte Ideen i Modsætning
til Mechanismen, hævder Lotze, at hvad enten
en Ide styrer Tingen eller ikke, er enhver Væren eller
Vorden kun mulig i Kraft af mechaniske Aarsager. Om
Ideens Bud udføres, kan kun erkjendes ved Indsigt i
% --- p. 304 ---
\setcounter{footnote}{0}%
\opage{304}den virkelige Natursammenhæng. De to Tendentser i
Erkjendelsen, den kausale Forklaring og den ideale Udtydning,
kunne derfor bestaae ved Siden af hinanden.
Men den kausale eller mechaniske Forklaring fører, som
vi have set, ved at underkastes en metaphysisk Kritik
tilbage til en Mangfoldighed af aandige Elementer som
Kræfternes Centrer. Den Identitet af det Ideale og det
Reale, som Spekulationen hævdede, gjælder altsaa punktuelt
for Verdens enkelte Elementer. Ethvert af disse
er ideelt, fordi dets Væsen bestaaer i den eiendommelige
Betydning, det har som særegen Indholdsbestemmelse
i det Uendelige; det er realt, thi Realitet er den
Existentsform, hvorved Noget fremstiller sig som selvstændigt
Udgangspunkt for Virkninger, som det udøver
eller lider. Aanden er da ikke længere, som for Spekulationen,
et svævende Lysvæsen, men opløses i en
utallig Mængde af skarpt begrændsede, straalende Punkter.
Intet enkelt Element fremgaaer af den absolute
Grund, uden at der af denne tillige fremgaaer en bestemt
Mængde andre, som sammen med hint udfylde
Virkeligheden til at være den fuldstændige Aabenbaring
af det Absolute. I Enkeltvæsenernes aandige Indre har
det Absolute det aabne Sted, hvor det, uden at det mechaniske
Verdensløbs Nødvendighed ophæves, kan øve sin
Indvirkning og fremkalde saadanne Livsytringer, ved
hvilke en almindelig Fremskriden bliver mulig gjennem
hele den naturlige og historiske Udvikling. Denne Udvikling
er, da Enkeltvæserne ere det Absolutes Indholdsbestemmelser, % PRINTED AS IS (Enkeltvæserne: the image, O193)
Et med det Absolutes Selvudvikling gjennem
en Række af Verdensaldre. Men dette er en Tanke,
som den menneskelige Erkjendelse ikke kan gjennemføre.
Kun fordi Spekulationen betragtede de jordiske
Tilværelsesformer som de absolute, kunde den mene at
have udmaalt det Uendeliges Livsgang. Denne Forudsætning,
at den skabende Urgrund kun aabenbarer sig paa
% --- p. 305 ---
\setcounter{footnote}{0}%
\opage{305}den jordiske Naturs snevre Vei, maatte føre til saadanne
dialektiske Idyller som de spekulative Konstruktioner.

Dog er dette endnu ikke den væsentligste Indsigelse,
der maa reises mod den spekulative Betragtning af Naturen
og Historien. Det er ikke nok, for hver Skabning
og hvert Phænomen at finde en Ide, hvis Virkeliggjørelse
de tjene. Det Spørgsmaal kan ikke afvises, hvormeget
af denne Ides betydningsfulde Indhold der er tilstede
i Væsenet selv som Motiv for dets Virken eller
Gjenstand for dets Nydelse. Ellers faa vi kun at vide,
hvad et Phænomen er eller betyder for os, men ikke,
hvad det er i sig selv. Saaledes er det ved Betragtningen
af Historien ikke nok at opregne Rækkefølgen
af de Ideer, som hver for sig udtrykke en Tidsalders
Karakter og almindelige Retning; det kommer an paa at
vide, hvorledes enhver Tidsalder selv og den Enkelte i
den var tilmode, hvormeget af den Ide, som vi betragte
i Afstand, der virkelig levede i Gemytterne. De Enkeltes
Tanker, Følelser og Handlinger have for den
spekulative Opfattelse kun Værd og Betydning, forsaavidt
de tilfældigvis stemme overens med den almene Ides
enkelte Udviklingsmomenter. Historiens Gang bliver den
store Slagterbænk, paa hvilken al individuel Lykke, alt
individuelt Liv ofres, for at Menneskehedens almene
Ide kan skride frem i sin Udvikling. Og dog er det en
Selvmodsigelse, at der skulde gives Formaal, hvis Indhold
og Virkeliggjørelse ikke blev Nogen bevidst. Hvad
der skal være et Gode, maa have sin Tilværelse i
et aandigt Væsens levende Følelse; Alt hvad der kun
er Kjendsgjerning, Ting, Egenskab, Forhold eller Begivenhed,
er kun Middel, ikke selv et Gode.

Den spekulative Opfattelse søger klassifikatorisk at
paavise en Trinrække i Naturen fra det Ufuldkomnere
til det Fuldkomnere. Men der er ingen fornuftig Grund
til, at det Ufuldkomne skulde vedblive at bestaa ved
Siden af det Fuldkomne. Derimod kunde man tænke
% --- p. 306 ---
\setcounter{footnote}{0}%
\opage{306}sig, at det underordnede Trin opnaaer en eiendommelig
Fuldendelse, som det ikke vilde naa, naar det gik over
som Led i høiere Organismer. Da vilde den samtidige
Bestaaen af det Ufuldkomnere og det Fuldkomnere bevirke
en større kvalitativ Rigdom i Tilværelsen. Mange
physiske Processer, som vi ikke kunne fornemme, træffe
maaske i lavere Dyr modtagelige Organer. Saaledes
navnlig mange chemiske, elektriske og meteorologiske
Processer. Man maa antage, at Alt som sker, ogsaa
paa et eller andet Punkt fornemmes eller inderliggjøres
af en psychisk Substants.

Sjælelivet i de enkelte materielle Atomer er kun
Øieblikkets Værk. Naturløbets Elementer blive til aandige
Begivenheder i Følelsen, men der dannes heraf intet
sammenhængende Hele, og de opbevares ikke i nogen
Erindring. Denne de psychiske Bevægelsers regelløse
Virken bliver i Planten til en bestemt uforanderlig Drøm,
en musikalsk Udvikling uden Frihed og Bevægelighed.
Først i Dyrene er den organiske Fasthed saaledes forbunden
med bevæget Mangfoldighed, at et Sjæleliv i
snevrere Forstand kan opstaa. Forskjellen mellem Mennesket
og Dyret bestaaer ikke i særegne Former eller Elementer,
der komme til hos Mennesket. Man kan ikke
stille Menneskeaanden som et eget Væsen i Modsætning
til Dyresjælen; det vilde være ligesaa umuligt at fastsætte
en Grændse mellem disse to Væsener som at paavise,
hvorledes de kunne blive til Et i Mennesket. Eiendommeligt
for Mennesket er det derimod, at den sjælelige
Virksomhed hæver sig til en saadan Intensitet, at
den sammenlignende og kombinerende kan træde i et
frit Forhold til de umiddelbart modtagne Indtryk. Dyret
fornemmer vel ogsaa en fast Sammenhæng mellem Phænomenerne,
men det opfatter ikke denne Sammenhængs
Nødvendighed; det lider under Loven, men bliver sig
ikke Loven bevidst. Denne Forskjel rokkes ikke ved,
at man maaske vil kunne godtgjøre Menneskets Ned%
% --- p. 307 ---
\setcounter{footnote}{0}%
\opage{307}stammen fra undergaaede Dyrearter; thi denne Kjendsgjerning
vil Intet afgjøre med Hensyn til Bedømmelsen
af de nu faktisk bestaaende Tilværelsesformers Værd.

Erkjendelsen af Nødvendigheden er det eiendommelig
Menneskeliges Karaktermærke. Den Anelse om en nødvendig
Sandhed og Sammenhæng, som ligger til Grund for
al menneskelig Tænken og Virken, kan vel lede feil ved
forhastet Afslutning, men i sig selv er den rigtig, uagtet
den ikke ligefrem anviser Midlerne til at naa det
Maal, den stiller. Forsøgene paa i Sammenhæng at
fremstille Indbegrebet af de nødvendige Sandheder mislykkes
gjerne, fordi det er Elementer af forskjellig Oprindelse
og Gyldighed, der skulle forenes. Nogle Sandheder
gjælde abstrakt, for al Virkelighed (Identitets- og Kausalitetssætningen),
andre støtte sig kun til vor Erfaring
(Forestillingen om Stoffets Uforgængelighed), atter andre
grunde sig paa vort Gemyts æsthetiske og ethiske Idealer
(Fordringen om en høieste Enhed i Tilværelsen).

Følelseslivet spiller en stor Rolle ved Udviklingen
af de høiere Fornuftsideer, hvor Spørgsmaalet om Værd
og Betydning træder istedetfor Spørgsmaalet om virkende
Aarsager. Derfor er det saa vanskeligt at omsætte disse
Ideer til almindelig Forstaaelse. Videnskaben maa være
tilfreds, naar det lykkes den at eftervise, at Forstandens
klare og uafviselige Grundsætninger ikke ere Andet end
Formerne for den ideale Fornufts eget Øiemed: Virkelighedens
Forbindelse til en levende Totalitets Enhed.
--- Kun Følelsens Indgriben kan fremdeles forklare Selvbevidsthedens
Væsen. Thi naar man i Almindelighed
lader sig nøie med at bestemme Selvbevidstheden som
Enhed af det Tænkende og det Tænkte, da er dette en
rent formel Bestemmelse, der ikke forklarer den Inderlighed
og det uendelige Værd, som ligger i dette indre
Forhold til sig selv. Selvbevidstheden er kun at opfatte
som Udtydningen af en Selvfølelse, hvis oprindelige Liv
ikke umiddelbart forhøies ved vor Erkjendelses  Udvik%
% --- p. 308 ---
\setcounter{footnote}{0}%
\opage{308}ling. Begrebet Værd er uadskillelig knyttet til Følelsens
Begreb, som langt eiendommeligere end Erkjendelsen
betegner Aandens sande Natur. Følelsen er ikke,
som man ofte mener, uimodtagelig for kvalitative Forskjelle,
som bero paa dets Værd, der fremkalder Følelsen.
Af mindre Værd er det, der kun svarer til en øieblikkelig
og tilfældig Tilstand eller en individuel Egenhed;
af større Værd hvad der harmonerer med de almindelige
og normale Grundtræk i den Organisation, ved
hvis Hjælp Aanden naaer sit Formaal; det Høieste vilde
være, hvad der tilfredsstiller et idealt Gemyts bestandige
Stemning, hvor enhver Vildfarelse fra Udviklingens
Maal er umulig. Men udover dette kan man ikke komme:
Tanken om et uendeligt Værd, der ikke godtgjør sin
Realitet ved at fremkalde Lyst, indeholder en Selvmodsigelse.
Den moralske Rigorisme, der i sin Iver for at
bekæmpe Egoismen vilde udelukke Lystfølelsen fra Handlingens
Motiver, formaaede derfor heller ikke at give
Andet end rent formelle Maximer; Morallovens forpligtende
Værdighed formaaede den ei at forklare. Det Ethiske
viser her hen til det Religiøse, de sædelige Forpligtelser
følge af den religiøse Verdensanskuelse: kun hvor
Mennesket ser sig som Medarbeider i en oversandselig
Verdensorden, faaer han positive Formaal, der komme
ham selv til Gode og hæve ham op over hans endelige
og sandselige Væsen. Det Religiøses Væsen bestaaer i
den bevidste Anerkjendelse af det Sædeliges og det Helliges
ubetingede Gyldighed overfor Alt, hvad der er
eller synes at være Kjendsgjerning. I de positive Religioner
fremtræder det specifik Religiøse ikke rent; det
fordunkles af kosmologiske Spekulationer, ved hvis Sammensmeltning
med hin Kjerne den religiøse Tro kan
blive en inhuman og fanatisk Magt. Ogsaa i Christendommen
er denne Fordunkling kommen ind. Grundtanken
i den oprindelige Christendom, som ikke maa
sammenblandes med den Tro paa historiske Kjendsgjer%
% --- p. 309 ---
\setcounter{footnote}{0}%
\opage{309}ninger, Kirken fordrer, gik netop ud paa at hæve Mennesket
op over enhver given Tilstand og udrydde Dyrkelsen
af blind, faktisk Væren. Men hertil staaer Kirkens
Fordring om Tro paa en hellig Historie i Modsætning,
især naar denne Tro ved uvis Tydning af uvisse Skrifter
udformes til et kosmologisk System, der maaske
endog (i Treenighedslæren) søger en Støtte i den Trehed
af Grundbegreber, hvori den tænkende Erkjendelse
ender, en Trehed, der, som det vil vise sig, ikke er et
Bevis paa Tankens Styrke, men netop paa dens Svaghed.

I Historien træffe vi Menneskeaandens Udvikling og
Virkeliggjørelse af sit Indhold. Undersøgelsen af Aandens
Væsen har godtgjort Uholdbarheden saavel af den
materialistiske som af den spekulativ-idealistiske Anskuelse.
Vilde man vende tilbage til den Opfattelse,
der ser Historien som Menneskeslægtens Opdragelse, da
vilde denne, der ellers synes at komme Sandheden nærmest,
nødvendigvis fordre en nærmere Bestemmelse.
Thi hvorledes kan man tale om Opdragelse, naar Udviklingens
Resultater komme Slægter til Gode, der ikke
have været med at frembringe dem? Menneskeslægten
er jo ikke en fast Enhed, men en Succession af Slægter.
De Goder, der opnaaes, blive desuden neppe Gjenstand
for Bevidsthed og Nydelse, før de gaa over til at
blive et ligegyldigt, tilvant Grundlag for videre gaaende
Ønsken og Stræben. Desuden bliver under Udviklingens
Gang Kulturens Indhold saa mangfoldigt, spalter sig i
saamange Grene, at ingen Enkelt kan overse, end sige
tilegne sig det Hele. Det viser sig endelig ved nærmere
Betragtning, at Kulturens Goder kun komme et ringe
Antal til Gode; under disse Begunstigede finde vi gjennem
hele den historiske Udvikling et Proletariat, over
hvis Hoveder Kulturen gaaer sin Gang uden væsentlig
at gribe ind i dets Vilkaar. De talløse Fremskridt i det
Enkelte have endnu ingenlunde samlet sig til et almin%
% --- p. 310 ---
\setcounter{footnote}{0}%
\opage{310}deligt Fremskridt i Livslykke; tvertimod synes Livet
mere og mere at blive til en Kamp for Tilværelsen.

Men derfor kunne vi dog ikke opgive Tanken om
et virkeligt Fremskridt og et virkeligt Formaal i den
historiske Udvikling. Kun derved finde vi Forklaringen
af al begeistret Selvopofrelse i Ideens Tjeneste.
Staaer det fast, at ingen Kjendsgjerning og ingen Form
har Værd i sig selv, men at kun det er et Gode, der
levende tilegnes af et aandigt Væsen, saa maa vi ogsaa
ved den Menneskehed, hvis Væsen udvikles gjennem Historien,
ikke forstaa en Sum af Enkelte, paa hvem en
vis Artskarakter passer, men et levende og realt Samfund,
i hvilket de menneskelige Aander, der ellers adskilles
ved Tidsforskjellene, sluttes sammen til at være
for hverandre saaledes, at Enhver indtager sin eiendommelige
Stilling, og Alt kommer Alle til Gode. Kun
under denne Forudsætning faaer Forestillingen om Menneskeslægtens
Opdragelse, ligesom ogsaa den spekulative
Opfattelse af Historien som Ideens Udvikling, sin Betydning;
thi kun da bliver det af noget Værd, at Ideen
udvikles. Ikke for Individernes Skyld maa denne Fordring
stilles, men for Verdenssammenhængens Skyld,
der vilde blive meningsløs uden dens Opfyldelse. Kun
Indsigten i det, som skal være, kan aabne os Indsigten
i det, som er. Deri har Idealismen Ret, at der kun
existerer Saadant og Saameget, som har sin nødvendige
Plads i den høieste Ide. ---

Til Slutning samle vi, ligesom Lotze selv i sidste
Afsnit af „Mikrokosmus“, de i det Foregaaende givne
Motiver til en afsluttende Verdensanskuelse. --- Lotze
gaaer ligesom Herbart ud fra det Givne og slutter tilbage
derfra. Men allerede ved den første Slutning skiller
han sig fra Herbart. Denne bestemmer Væren som
absolut Position o: absolut, enkelt Kvalitet. Men dette
er en Feilslutning. Væren vil tvertimod sige, at et Indhold
er sat i visse bestemte Forhold. Enhver Tings
% --- p. 311 ---
\setcounter{footnote}{0}%
\opage{311}eiendommelige Natur kan kun begribes som et bestemt
Led i det realt mulige Verdensindhold. Vexelvirkningen
mellem disse enkelte Led er kun forklarlig, naar de ere
Dele af en uendelig Substants, der omfatter Alt. For
os træde Tingene frem hver paa sit Sted i et uendeligt
Rum; men ifølge Læren om Lokaltegnene er Rumsanskuelsen
kun den Form, hvorunder Sjælen omsætter
Indtrykkene i sit eget Sprog. Dog er Rummet ikke en
blot subjektiv og tilfældig Form. Til Ordningen i Rummet
for vor Anskuelse svarer den Maade, hvorpaa Tingene
ere ordnede i det Uendeliges intellektuelle Sammenhæng,
og denne beroer paa Betydningen og Værdien af
den Ide, hver Ting udtrykker. Paa lignende Maade forholder
det sig med Tiden. Heller ikke denne har umiddelbar
objektiv Betydning. Men hvad der for os fremtræder
som Tildragelse i Tiden, det er kun Aabenbaringen
af det indre Betingelsesforhold mellem det Uendeliges
enkelte Elementer. Det Enkeltes Tidspunkt og
Varighed beroer paa dets intellektuelle Betydning o: paa
Betydningen af den Ide, det udtrykker.

Men naar vi saaledes føre alt Existerende tilbage
til den Ide, det udtrykker, og se dets Væsen udtrykt
deri, synes vi derved at hildes i adskillige Vanskeligheder.
Tingen kan ikke være blot Ide; skal den have
Virkelighed, maa den være en virkelig Ide, men Ideen kan
kun virke ved at tænkes. Vi tillægge Tingen Enhed;
men vi kjende kun Enheden fra vort eget Jeg; kun
fordi vi blive os bevidste, ere vi Et. Vi tillægge Tingen
en eiendommelig Virken og Liden efter de Forhold, i
hvilke den er sat; men kun det kan i Sandhed virke
og lide, som kan føle Vel og Ve.

Disse Vanskeligheder ved Opfattelsen af Tingenes
Væsen have foranlediget forskjellige Løsningsforsøg. ---
Man har ladet det Reale i Tingene staa hen som en
dunkel, uforstaaet Kjerne, og betragtet vor Forestilling
om Ting overhovedet som en subjektiv Erkjendelsesform,
til hvilke Virkeligheden ikke behøvede at svare. Men
% --- p. 312 ---
\setcounter{footnote}{0}%
\opage{312}man nødes dog stedse til at tillægge Tingene i sig
selv en Væren og en Vexelvirkning med os; og man
bevæger sig da i en Cirkel, idet man paa engang anvender
og negter Anvendelsen af vore Begrebsformer paa
Tingenes Væsen. --- Der staaer da kun tilbage enten
aldeles at negte Tingenes Realitet eller at søge at udvikle
Tingens Begreb saaledes, at Vanskelighederne
forsvinde. Den første af disse Udveie fører nu med Nødvendighed
til den anden. Thi Idealismens Sætning, at
kun det Aandige er virkeligt, maa, naar Virkeligheden
ikke skal blive et rent Skin, vendes om, saa at den
kommer til at lyde: alt Virkeligt er Aand. Denne Sætning
har i det Foregaaende vist sig at indeholde den
sidste Forklaring af, hvad Erfaringen fremstillede for os.
Den forflygtiger ikke Tingenes Begreb, som Idealismen;
Tingen faaer Selvhed, Forsigværen, det almindeligste
Udtryk for den aandelige Natur, der i de selvbevidste
Jeger naaer sit høieste Trin, uden derfor at mangle der,
hvor Tilværelsen ikke opfattes med klar Bevisthed, men % PRINTED AS IS (Bevisthed: the image, O196)
under en dunklere Følelses Form. Forskjellen mellem
denne Opfattelse, den spiritualistiske Realisme, og Idealismen
bestaaer ikke, som man maaske kunde mene,
deri, at hin skulde tillægge Tingene en Væren uden for
Gud, medens denne opfatter dem som aldeles iboende
det Absolute. Thi det beroer paa en uklar Tankegang
at mene, at det, som er udenfor Gud, skulde være mere
virkeligt end det, som er i Gud. Ja, hvad Mening vil
der i det Hele kunne forbindes med det Udtryk „at
være udenfor Gud?“ Forskjellen bestaaer derimod i, at
Idealismen negter Tingenes Selvhed, medens den spiritualistiske
Realisme, der gaaer ud fra, at vi i Aanden,
i Selvet idetmindste have \emph{ét} Reale, hvis Natur vi kjende,
mener at maatte tillægge alle Realer denne aandige
Natur i forskjellige Grader.

De tidligere Undersøgelser førte til Antagelsen af
en uendelig Substants, hvis Indholdsbestemmelse alle
% --- p. 313 ---
\setcounter{footnote}{0}%
\opage{313}Ting ere. Men dog kan dette Uendelige, Pantheismens
Gudsidé, ikke være det afsluttende Begreb. Vi vilde da
ende med J. G. Fichte i Forestillingen om en uendelig
Verdensorden. Men Orden kan ikke tænkes uden et
Ordnende. Og den hele Opfattelse af det Reales Væsen,
som fremgaaer af det Foregaaende, at alt Realt maa
være for sig selv, have et aandigt Midtpunkt, maa dog
ogsaa finde sin Anvendelse her. Indsigelserne mod at
anvende Personlighedens Begreb paa Gud tabe deres
Vægt, da vi have vist, at Selvbevidsthedens Væsen ikke
blot beroer paa Modsætningen mellem et Jeg og et Ikke-Jeg,
men kun er det reflekterede Udtryk for Selvets
oprindelige Væsen. Sagen stiller sig tvertimod ved nærmere
Betragtning saaledes, at Personlighedens Begreb
kun ufuldkomment kan anvendes paa de \emph{endelige} Væsener,
hvis indre Oprindelighed ikke er absolut; det er et
Ideal, der kun i det absolute Væsen kan tænkes fuldstændig
virkeliggjort.

Det høieste Virkelige, det Absolute maa være absolut
Enhed. Men her viser den menneskelige Tankes
Utilstrækkelighed sig. Den formaaer ikke at samle Erkjendelsens
forskjellige Traade til en fuldstændig Enhed.
De tidligere Undersøgelser have paa ethvert Punkt vist
os en Trehed af Kræfter, som ikke kunne føres tilbage
til Enhed: 1) de almindelige og abstrakte (metaphysiske
og mathematisk-mechaniske) Love; --- 2) selve de reale
Væsener, i hvilke Kræfterne koncentrere sig, og ud fra
hvilke de virke ifølge hine Love; --- 3) Formaalene, den
omfattende Plan, efter hvilken Verdenslivet udvikler sig.
Hvad der hindrer Tanken i at lade Loven og Kraften
smelte sammen med Formaalet til ét Væsen, det er i
sidste Instants det store, uløste Problem, der ligger i
i det physiske og moralske Ondes Virkelighed. Kun % PRINTED AS IS (i det physiske: the image, i doubled over the line turn (i|i); O198)
saameget maa Tanken fastholde, at Lovens og Kraftens
Grund maa ligge i selve Formaalet; det høieste Virke%
% --- p. 314 ---
\setcounter{footnote}{0}%
\opage{314}lige udfolder sig i én Bevægelse, der for os spalter sig
i tre Kræfter. ---

Det vil have sin Interesse at sammenligne Lotzes
Lære med andre, tidligere og samtidige Philosophers,
med hvilke den har Berøringspunkter. Blandt ældre
Philosopher komme her især Leibnitz, Kant og Herbart,
blandt senere Comte, Mill og Fechner i Betragtning\footnote[1]{Vi have ovenfor (p. 288) hørt, i hvilket Forhold Lotze stiller sin Anskuelse til den ældre Fichtes. Hans Forhold til den Schelling-Hegelske Spekulation er blevet fremhævet paa flere Punkter i Fremstillingen.}.

Leibnitz's Monadelære har strax ved første Øiekast
stor Lighed med Lotzes Anskuelse. Men afset fra den
allerede tidligere omtalte Forskjel i Opfattelsen af Monadens
indre Væsen, som Lotze selv har fremhævet,
bestaaer den afgjørende Uoverensstemmelse i, at Lotze
identificerer Monaden med Atomet, saa at en virkelig
Vexelvirkning bliver mulig, kun at man ikke kan opfatte
Materien som et Kontinuum. Leibnitz tillægger
derimod ikke umiddelbart sine Monader nogen physisk
Betydning; en egentlig Vexelvirkning mellem dem er
umulig, og den tilsyneladende Vexelvirkning forklares
ved en forudbestemt Harmoni imellem Monadernes Selvudfoldelse.
Medens Leibnitz derfor er Idealist, er Lotze
Realist. --- Derved faaer han et Berøringspunkt med
Herbart. Men ikke blot adskiller denne udtrykkelig
sine Realer fra Physikens Atomer; han betragter dem
ogsaa som absolut enkelte i den Forstand, at ingen
indre Udvikling, kun Selvopholdelse overfor ydre Paavirkning
bliver mulig. Lotzes Reale er ikke uforanderligt,
ligesaa lidt som Leibnitz's Monade; Lotze er Realist,
men spiritualistisk Realist. I Psychologien adskiller
han sig fra Herbart ved at hævde Følelsens og Villiens
Oprindelighed, medens Herbart udleder dem af bestemte
Forhold mellem Forestillingerne. Og medens
% --- p. 315 ---
\setcounter{footnote}{0}%
\opage{315}Herbart endelig stiller det Theoretiske og det Praktiske
uafhængige og ganske sondrede overfor hinanden, finder
Lotze i det Ethisk-Religiøse Metaphysikens Udgangspunkt
og i det, som skal være, den sidste Forklaring
for det, som er. --- Hvad Kant angaaer, da maa det her
være nok at henvise til Læren om Rum og Tid, om
Tingen i og for sig, om det Gode. Ja, endog for Atomlærens
Vedkommende har Lotze, som allerede berørt,
fundet et Tilknytningspunkt, nemlig i en ældre Afhandling
af Kant, hvor denne opstillede en atomistisk Theori,
som han senere forlod for i sine „Metaphysische Anfangsgründe
der Naturwissenschaft“ at grundlægge den
dynamiske Opfattelse, som blev herskende i den spekulative
Naturphilosophi. Lotzes hele kritisk-skeptiske, antidogmatiske
Aandsretning minder om Kants Kritiker.

Med den moderne Positivisme i Frankrig og England
stemmer Lotze (ligesom tildels ogsaa Fechner) overens
i en dobbelt Henseende. Han gaaer ligesom den
ud fra Phænomenet, og han lægger Vægten paa en exakt
fattelse af den phænomenale Sammenhæng. I denne % PRINTED AS IS (fattelse af den phænomenale: the image, Op lost at the line turn (exakt|fattelse); O200)
Del af sin Lære er han ganske Positivist. Men han
gaaer videre og søger at løse Problemet om Phænomenernes
Væsen og Grund. Dette Problem betragte Comte
og Mill som uopløseligt; men da Ingen af dem formaaer
at fastholde denne Skepticisme, besvare de det alligevel
dels i materialistisk, dels i idealistisk Retning. --- Fra
Fechner adskiller Lotze sig, som allerede tidligere bemærket,
ved at opfatte Atomerne som Monader med et
indre Liv, et Analogon til det bevidste Aandsliv. Saavidt
udstrækker Fechner ikke sin Besjælingshypothese,
hvorfor Atomerne hos ham staa som den sidste Forudsætning
for Erfaringsanalysen, men uden Sammenhæng
med hans øvrige Anskuelse. I Opfattelsen af selve
Sjælelivet ere de paa en Maade Antipoder, idet Fechner
ser Bevidsthedens Væsen i, at den sammenfatter de enkelte
Fornemmelser og Forestillinger (den synechologiske
% --- p. 316 ---
\setcounter{footnote}{0}%
\opage{316}Opfattelse), medens Lotze lægger Vægten paa den Koncentration
i et Midtpunkt, som Bevidstheden forudsætter
(den monadologiske Opfattelse). Forholdet mellem dem
kan sammenstilles med det mellem Spinoza og Leibnitz.

Betragte vi endelig Lotzes Lære i Forhold til den
efterhegelske Philosophi i det Hele, da finde vi hos ham
optaget som væsentlige Elementer, hvad der i de andre
Retninger fremtræder isoleret og eneraadende. Nødvendigheden
af en kausal Forklaring ved elementære
Kræfters mechaniske Vexelvirkning betoner han ligesaa
stærkt som Materialismen og paa mere videnskabelig
Vis. Han opponerer trods den negative Kritik mod al
philosophisk og religiøs Dogmatisme og lukker ikke
Øinene for Erkjendelsens psychologiske Betingethed. Og
han bestrider endelig ligesom den spekulative Theisme
Forgudelsen af det abstrakte Almene og hævder det Individuelles
og Personliges metaphysiske og religionsphilosophiske
Betydning. Ikke uden Grund slutte vi
derfor vor Fremstilling med denne aandfulde Tænker.

Lotzes Blik for Tingenes og Kræfternes Eiendommelighed,
hans Sands for det fyldige Livsindhold
og hans Iver for at lade det komme til sin Ret over for
en død Formalisme og Mechanisme, gjør Behandlingen
af de konkrete Problemer, hvor Tanken endnu færdes i
den individuelle Virkelighed, til hans Philosophis stærke
Side. Hans Ærlighed og skarpe kritiske Syn gjør desuden,
at han hellere aabent erklærer et Problem for
uløseligt end beroliger sig ved en tilsyneladende Besvarelse
deraf (sml. ovenfor p. 222). Derfor aabner han kun
sin Læser Udsigten til de spekulative Høider, men fører
ham ikke derop. Tankens spredte Traade kan efter
hans Overbevisning den menneskelige Erkjendelse ikke
samle til én. Den høieste ethiske Ide staaer for ham
som det, der ogsaa maa være det høieste Virkelige.
Men han savner her den bestemte Tilknytning for Loven
og den virkende Kraft. Dog har han for en Del selv
% --- p. 317 ---
\setcounter{footnote}{0}%
\opage{317}fremkaldt dette Savn ved sin Opfattelse af Tankens Væsen
som væsentlig formelt reflekterende; derved kommer
Tanken ikke til at gjenfinde sig selv i sit Indhold.
Lotze hævder, at de ethiske Ideer maa gives Tanken
fra Følelsen og først derpaa kunne udvikles til Erkjendelse.
Men naar man ikke gaaer ud fra en staaende
Modsætning mellem Tanken og dens Indhold, er det
under Erkjendelsens Udvikling ligesaameget Indholdet,
der bestemmes ved Tanken, som omvendt. Ligeledes
er det ensidigt, naar Lotze udtaler, at det, som skal
være aabner os Indsigten i det, som er; ogsaa det omvendte
Synspunkt er af største Vigtighed, at nemlig
Erkjendelsen af det virkelig Værende belærer os om
Formaalets Væsen. Ingen har mere udøvet denne Grundsætning
i Praxis end Lotze selv, der slaaende har godtgjort,
hvor urigtig man opfatter Ideen, naar man overser
de virkelige Betingelser for dens Gyldighed og Realitet.
Metaphysiken faaer da ikke blot, som Lotze siger, sit Princip
i Ethiken, men Ethiken maa omvendt have sin Begrundelse
i Metaphysiken; eller bestemtere udtrykt, den
Sætning, at Metaphysiken har sin Begyndelse i Ethiken,
kan ikke betyde Andet end at den, der skal udvikle
Væsensbestemmelserne for alt Værende, da ogsaa maa
paavise det almene Grundlag for det Ethiske. Dette stiller
altsaa en Fordring til Metaphysiken, men raader ikke
umiddelbart for, hvorledes denne Fordring sker Fyldest.

Naar Lotze udtaler, at det Absolute, der i sig
selv evig er én Bevægelse, for os spalter sig i tre Kræfter,
da synes det at følge af sig selv, at man ikke kan
udlede den ene af disse Kræfter af nogen af de
andre (f. Ex. Loven og Kraften af Formaalet). De maa
derimod saameget som muligt føres tilbage til deres
fælles Udspring, og dette tilstræber vor Erkjendelse ved
at se dem i deres gjensidige Forbindelser. Udspringe
de virkelig af ét Princip, da maa en indre Vexelvirkning
mellem dem ogsaa være mulig og idetmindste til%
% --- p. 318 ---
\setcounter{footnote}{0}%
\opage{318}nærmelsesvis kunne begribes af os. Den ethisk-religiøse,
den metaphysiske og den physiske Erkjendelse (Erkjendelsen
af Formaalet, af Loven og af Kraften) staa saaledes
ikke i Modsætning til hverandre, men samle sig
fra forskjellige Udgangspunkter til et Centrum.

Uagtet det, som er, ifølge Lotze, skal forklares af
det, som skal være, er hans egen Lære dog ikke opstaaet
paa denne Maade, men er væsentlig en Hypothese,
der er udsprungen af Bestræbelsen for at forklare
den givne Sammenhæng. Med Rette har man her savnet
den dybere philosophiske Begrundelse\footnote[1]{\emph{Weisse} i „Zeitschr. für Philos. und philos. Kritik.“ 47de Bd. p. 279 f.}. Ligesaalidt
som Ethiken faaer Naturphilosophien sit metaphysiske
Grundlag. Derved er det ogsaa sket, at Lotze ikke overvinder
den Vanskelighed, der møder alle monadologiske
Systemer, nemlig at vise, hvorledes de selvstændige
Enkeltvæsener kunne forenes med det Absolute. Monaderne
fremtræde snart som fuldkommen selvstændige,
snart gaa de ganske op i Ideen, af hvis Naade, de ene
ere til. Det sidste Alternativ maa forsaavidt stemme
mest med Lotzes Lære, som jo hævder, at det som er,
maa underordnes det, som skal være. Derfor vil Lotze
heller ikke indlade sig paa Spørgsmaalet om den individuelle
Udødelighed. Men gribes dette Alternativ, hvad
bliver der da af Atomlæren? --- Naar fremdeles det
Absolute skal være et personligt Selv, maa det Spørgsmaal
opstaa, om det da ikke maa være en Monade, ligesom
de endelige Væsener, og hvorledes dette vil kunne
forenes med det Absolutes Væsen.

Men Lotze har ikke lovet os noget Svar paa disse
Spørgsmaal. Og det maa fremhæves, at han, der væsentlig
har fornyet Leibnitz's Anskuelse, har korrigeret denne
paa afgjørende Punkter, dels ved at opfatte Monaden
% --- p. 319 ---
\setcounter{footnote}{0}%
\opage{319}som Atom, hvorved han har sikret sig de exakte og
reale Tilknytningspunkter, dels ved fra første Færd at
hævde, at Monaderne kun ere Indholdsbestemmelser i den
uendelige Substants.

\chaphere{SJETTE KAPITEL.}


Der staaer nu tilbage i korte Træk at sammenfatte
den efterhegelske Philosophis Væsen og Betydning. Den
Fremstilling vi nu have sluttet, har paa mange Maader
vist os Modsætningen til den foregaaende Periodes Tænkning,
samtidig med, at det blev klart, hvorledes de afgjørende
Problemer her som altid gaa over fra den
ene Slægt til den anden. Forholdet stiller sig, naar vi
se tilbage paa de Phænomener, vi have lært at kjende,
ved nærmere Betragtning saaledes, at den spekulative
Philosophi i sin første storartede Udvikling førte til Anskuelsen
af et høit videnskabeligt Ideal, der først successive,
ved trofast og udholdende Bearbeidelse af den
givne Virkelighed, kan modtage sit faste og urokkelige
Grundlag. Dette Arbeide har den efterhegelske Philosophi
indledet, og heraf udspringer baade dens Polemik
mod Spekulationen og dens Sympathi for denne.

Den spekulative Philosophi var væsentlig dogmatisk
og idealistisk. Den var dogmatisk, idet den hvilede paa
Forudsætningen om Tankens og Virkelighedens Enhed
og alle dens Bestræbelser kun gik ud paa at udvikle
Alt, hvad der indeholdes i denne oprindelige Enhed.
Naturligvis var denne Forudsætning ikke vilkaarlig; den
var igjen Resultatet af den kritiske Philosophis dybtgaaende
Undersøgelser af Erkjendelsens Væsen. Men
% --- p. 320 ---
\setcounter{footnote}{0}%
\opage{320}dette Resultat hævedes dog over enhver endelig og
subjektiv Begrundelse; man mente engang for alle at
have naaet det Punkt, hvorfra en apriorisk og synthestisk % PRINTED AS IS (synthestisk: for synthetisk; crop g:C005 and the Minnesota layer)
Udvikling af den philosophiske Verdensanskuelse
var mulig. Dog er det klart, at hvad man ad denne Vei
kan naa, kun kan være, hvad man allerede har i Forudsætningen,
hvorfra man gaaer ud. Da nu Forudsætningen
vel var Enheden af Tanken og Virkeligheden, men
Virkelighed her kun opfattedes som ren abstrakt Væren,
det vil sige som tænkt Væren, saa maatte ogsaa Verdensanskuelsen
blive en absolut Idealisme, der lod den virkelige
Tilværelse staa som uløst Problem.

Derfor er Virkelighed, Existents ogsaa det Slagord,
med hvilket der fra alle Sider rykkes i Marken
mod Spekulationen. Vi have fundet det i alle de
fire Grupper af philosophiske Phænomener, vi have
fremdraget. Man kommer altsaa til Erkjendelse af
at Spekulationens Forudsætning maa underkastes en
ny Prøvelse. Forsaavidt gjenoptager den efterhegelske
Philosophi Kants Arbeider og faaer en kritisk Karakter
i Modsætning til Spekulationens Dogmatisme. Dog
opgiver den derved ikke dennes væsentlige Udbytte.
Den falder ikke tilbage til Kriticismens skarpe Modsætning
mellem Tanken og Existentsen, men fastholder Enheden
som den sidste Forudsætning, kun at denne ikke
umiddelbart kan gjøres gjældende. Derfor hævder den,
at overfor den konkrete Virkelighed med dens Forskjelle
ere særegne Methoder nødvendige. Philosophien kan
ikke umiddelbart spekulere sig ind i Virkelighedens
Væsen, men behøver Erfaringsvidenskaben som Mellemled.
Hvad den ad denne Vei kan tilegne sig, maa den
saa selv give den særegne Begrundelse, hvorved en videnskabelig
Verdensanskuelse bliver mulig. Ideen mister
Intet af sit Glands, fordi dens Virken foregaaer efter
bestemte Love og under faste Betingelser. Dette er den
efterhegelske Philosophis realistiske Karakter.

% --- p. 321 ---
\setcounter{footnote}{0}%
\opage{321}Men det er ikke nok at bestemme denne Philosophis
Væsen i Forhold til den forudgaaende. For at vinde
et fyldigere Indblik i, hvad den har at sige, maa vi
præcisere de væsentlige Resultater, hvortil den har ført.
Det følger af sig selv, at intet videnskabeligt Resultat
er afsluttet i den Forstand, at det nu kunde staa fast
herefter som Dogme. Hvad her menes med Resultater,
kan kun være den almindelige Retning, i hvilken den
philosophiske Tænkning peger under sin Stræben efter
at vinde en afsluttende Verdensanskuelse.

De philosophiske Phænomer, vi have beskjæftiget % PRINTED AS IS (Phænomer: so printed here and on p. 191 (Phænomener elsewhere); the blind page read)
os med, dele sig strax, som vi allerede i Begyndelsen
(ovenfor p. 16) antydede, i to større Grupper. Den ene
af disse har en væsenlig negativ Karakter. Den kritiske
Bevidsthed opløste Spekulationen ved at hengive sig til
den formelle Konsekventses Magt, der tilsidst førte den
ud over al Sandhed og Virkelighed. Tilsidst maatte her
ikke blot al Philosophi, men ogsaa al Videnskab i det
Hele vise sig som Illusion. Selv Feuerbach, denne Proteusskikkelse,
der kun ved sin Aands egen ideale Rigdom
fik Kraft til at kæmpe mod Ideen, og som bestandig
viste fremad til en positiv Anskuelse, som han
dog aldrig naaede, selv han ender med at sige: „meine
Philosophie ist keine Philosophie“. Men den blot formelle
Konsekvents er usand. Den griber et enkelt Moment
af Tilværelsen i dets Isolation og afviser saa med
sophistisk Dialektik ethvert Forsøg paa at ophæve denne
Isolation. Det første Skridt til Sandhed er at se, at al
Virkelighed er Sammenhæng, Flerhed i Enhed. --- Til
den negative Gruppe maa ogsaa den moderne Materialisme
henføres. Vel mener den selv at have et positivt
Princip, nemlig Stoffet, Materien; men dels er det Væsentlige
ved Materialismen dens Polemik mod ethvert
aandeligt Princips Realitet, dels har det vist sig umuligt
paa positiv Vis at bestemme Materiens Væsen, hvorfor
Materialisterne selv faktisk opgive dette Princip for at
% --- p. 322 ---
\setcounter{footnote}{0}%
\opage{322}søge en Begrundelse ad Postulatets Vei --- og derved
komme de da til at appellere til ideelle Motiver! ---

Men uagtet den kritiske og materialistiske Retning
væsentlig er negativ, saa har den derfor alligevel en
stor philosophisk og kulturhistorisk Betydning. Den er
en Forskole, hvor de dogmatiske og idealistiske Illusioner
underkastes en Ildprøve. Kun Autoritetstroen kan
uden Videre afvise disse Retninger. For Philosophien
ere de en nødvendig Lutringsproces, gjennem hvilken
der vindes Kræfter til fornyet Stræben efter en positiv
Verdensanskuelse. Desuden tjener den negative Kritik
og Materialismen til at bryde Supranaturalismens og Dualismens
Herredømme i den almindelige Bevidsthed. Det
maa engang for alle staa fast, at kun det har Realitet,
der kan godtgjøre sig som sandt for den prøvende Kritik,
og som ikke trodser den faktiske Virkeligheds Betingelser.


Vende vi os nu til de positive philosophiske Phænomener,
som vi have behandlet i 4de og 5te Kapitel,
da ville vi sammenfatte deres væsentlige Betydning ved
at betragte dem fra tre Hovedsynspunkter, nemlig først
med Hensyn til Forholdet mellem Spekulation og Erfaring,
dernæst med Hensyn til Forholdet mellem Aand
og Materie, og endelig med Hensyn til Forholdet mellem
Philosophi og Religion.

Det første Forhold er allerede omtalt i sin Almindelighed,
idet vi have paavist den efterhegelske Philosophis
realistiske Karakter. Hvad vi her skulle henlede
Opmærksomheden paa, er, at Erfaringen ikke blot faaer
negativ Betydning, som en Instants, Spekulationen ikke
maa overse, men at der tilstræbes en positiv Enhed og
Sammensmeltning af begge. Vor Periode begynder med,
at man ud over det Omraade, den logiske Tankebevægelse
kan udmaale, troer at øine en fyldig Virkelighed,
hvis Tilegnelse skal danne et Supplement til den rent
immanente Spekulation. Saaledes den Retning, man en
% --- p. 323 ---
\setcounter{footnote}{0}%
\opage{323}Tidlang kaldte Frihedslæren eller den positive Philosophi,
og som vi have behandlet under Navn af den spekulative
Theisme. Allerede indenfor denne Retning fandt vi
en Bestræbelse for at overvinde Modsætningen mellem
de to Sphærer og stille dem i et indre Forhold til hinanden.
De Anskuelser, vi have sammenstillet under
Virkelighedsidealismens Navn, føre dette videre ved at
godtgjøre, at selve Erfaringen viser tilbage til spekulative
Forudsætninger. Her er det Punkt, hvor den franske
og engelske Realismes Utilstrækkelighed viser sig.
Den gjennemfører ligesaalidt som Materialismen sine
egne Forudsætninger og ender derfor principløst eller
skeptisk. „Virkelighedsidealismen“ viser os derimod,
hvorledes netop en fuldstændig Anerkjendelse af Erfaringens
Ret og en exakt Opfattelse af dens Resultater
philosophisk set fører ud over den blotte Empirisme.

Den spekulative Philosophi saa al Virkelighed i
Ideen. Men den kunde ikke udlede den faktiske, konkrete
Virkelighed af Ideens Væsen. Virkelighedsidealismen
gaaer nu den omvendte Vei og søger at vise, at al Virkelighed
i sidste Instants er af ideal Natur. Og den gjør
dette uden derved at opgive det exakte Fodfæste i den
virkelige Verden. Under en mere beskeden Form peger
den dog hen mod en Anskuelse, der i Dristighed og
Storartethed ikke viger for de spekulative Systemer.
Her er da Veien banet for den Forening, Philosophien
og i den Menneskeaanden bestandig søger, Foreningen
mellem Idealet og Virkeligheden. Denne Forening skal
paavises i selve Virkelighedens Kjerne, ligesom Spekulationen
paaviste den i Ideen. At Materien i sidste Instants
er Aand, har Philosophien ofte udtalt; men først
nu er det lykkedes at indlede en saadan Gjennemførelse
af denne Sætning, at det Materielle og Virkelige ikke
forflygtiges men bevares i sin væsentlige Betydning.
Først nu er Materialismen optaget som Element i den
philosophiske Verdensanskuelse, og netop  derved vil
% --- p. 324 ---
\setcounter{footnote}{0}%
\opage{324}denne kunne gjennemføre den Idealisme, den ikke kan
opgive uden at opgive sig selv.

Heraf følger da igjen, at Philosophien mindre end
nogensinde kan føle Trang til at bøie sig for en Autoritetstro,
der negter Muligheden af, at den menneskelige
Tanke ad sin egen Vei kan naa Erkjendelsen af
det Guddommelige. Philosophien bøier sig for enhver
Kjendsgjerning paa det aandelige og legemlige Omraade,
men kun fordi enhver Kjendsgjerning i sit Væsen og sin
Grund godtgjør det samme rationelle Princip, hvorfra
den videnskabelige Forsken selv udspringer. Vi have
set, hvor nøie en Forbindelse der er mellem den historiske
og den philosophiske Kritik i vor Tid; begge have
udviklet sig Haand i Haand. Philosophien har derved
faaet et paalideligt Grundlag for sin Opfattelse af det
Religiøse. Den kan nu stille sig i et frit og af Antipathi
og Sympathi lige uhildet Forhold til de historiske
Religioner. Og det specifik Religiøse selv er ikke et
for Philosophien fremmed Element. Thi den philosophiske
Tanke er ikke længere formel og abstrakt; den
nærer sig med selve Virkelighedens Substants og er sig
sit inderlige Forhold til Menneskets psychologiske Væsen
bevidst. Derfor kan den rumme alle virkelige Livsrørelser
og give dem deres faste og bevidste Form. Philosophien
er selv religiøs, idet den fører Mennesket til
Erkjendelsen af Tilværelsens Væsen og af hans Forhold
til dens ideale Magter.

Den spekulative Philosophis Gudside var den rene
Aand, der i evig Selvudfoldelse udvikler Alt, som er.
Den var væsentlig abstrakt Pantheisme. Fra forskjellige
Synspunkter have vi set den efterhegelske Philosophi
stræbe hen til en konkretere, mere levende Gudside.
Naar Virkeligheden bliver mere end et Skin, en blot
negativ Form for det Absolute, da kan dette heller ikke
være den rene, abstrakte Ide. Og naar alt Virkeligt er
et Enkelt, et Individuelt, da maa dette ogsaa gjælde for
% --- p. 325 ---
\setcounter{footnote}{0}%
\opage{325}det høieste Virkeliges Vedkommende. Derfor søger den
spekulative Theisme at paavise den guddommelige Aand
som den absolute Personlighed: som absolut er Gud Et
med Ideen og staaer i et immanent Forhold til Verden;
men indenfor denne Immanents gjør sig en Modsætning
gjældende, idet den guddommelige Tanke samler
sig i et indre Midtpunkt, hvorfra Verdenslivet udgaaer,
og som er det høieste Forbillede for de endelige Personligheder.
Fra pantheistisk Side kommer man denne
Tankegang imøde ved at hævde Forskjellens og Modsætningens
væsentlige Betydning. Hvad enten man nu vil kalde
en saadan Anskuelse „immanent Theisme“ eller „transscendent
Pantheisme“, saameget er klart, at her i almindelige
Træk er tegnet et rationelt Gudsbegreb, som paa engang
er Udgangspunktet og Slutningspunktet for al philosophisk
og videnskabelig Forsken og for al virkelig human
Stræben i det Hele.

\end{document}
